\documentclass[11pt,twoside,openany]{book}
\usepackage[paperwidth=7.5in,paperheight=10in,inner=.82in,outer=.72in,top=.78in,bottom=.8in,headheight=15pt,headsep=18pt,footskip=28pt]{geometry}
\usepackage[T1]{fontenc}
\usepackage{lmodern,microtype}
\usepackage{amsmath,amssymb,amsthm,mathtools}
\usepackage{enumitem,booktabs,longtable,array,multicol}
\usepackage{tikz}
\usetikzlibrary{calc,angles,quotes,intersections,arrows.meta,patterns,decorations.pathreplacing}
\usepackage{xcolor,fancyhdr,needspace,etoolbox}
\definecolor{ink}{RGB}{22,42,62}
\definecolor{accent}{RGB}{23,91,122}
\definecolor{soft}{RGB}{241,246,249}
\usepackage[colorlinks=true,linkcolor=accent,urlcolor=accent,citecolor=accent,bookmarksdepth=2]{hyperref}
\hypersetup{pdftitle={Olympiad Mathematics},pdfauthor={},pdfsubject={A four-subject course from first proofs to advanced problem solving}}
\pagestyle{fancy}
\fancyhf{}
\fancyhead[L]{\small Olympiad Mathematics}
\fancyhead[R]{\small\nouppercase{\rightmark}}
\fancyfoot[C]{\thepage}
\renewcommand{\chaptermark}[1]{\markboth{\thechapter.\ #1}{}}
\renewcommand{\sectionmark}[1]{\markright{\thesection}}
\fancypagestyle{plain}{\fancyhf{}\fancyfoot[C]{\thepage}\renewcommand{\headrulewidth}{0pt}}
\setlength{\parindent}{1em}
\setlength{\parskip}{3pt}
\setlength{\emergencystretch}{2em}
\raggedbottom
\clubpenalty=10000
\widowpenalty=10000
\linespread{1.06}
\setlist{itemsep=3pt,topsep=6pt}
\setlist[enumerate]{leftmargin=*}
\setcounter{tocdepth}{3}
\setcounter{secnumdepth}{2}
\allowdisplaybreaks[2]
\newcommand{\NN}{\mathbb N}
\newcommand{\NNz}{\mathbb N_0}
\newcommand{\ZZ}{\mathbb Z}
\newcommand{\QQ}{\mathbb Q}
\newcommand{\RR}{\mathbb R}
\newcommand{\CC}{\mathbb C}
\newcommand{}{\mathbb F}
\newcommand{\eps}{\varepsilon}
\newcommand{\dang}{\measuredangle}
\newcommand{\abs}[1]{\left|#1\right|}
\newcommand{\norm}[1]{\left\lVert#1\right\rVert}
\newcommand{\floor}[1]{\left\lfloor#1\right\rfloor}
\newcommand{\ceil}[1]{\left\lceil#1\right\rceil}
\DeclareMathOperator{\lcm}{lcm}
\DeclareMathOperator{\ord}{ord}
\DeclareMathOperator{\rad}{rad}
\DeclareMathOperator{\sgn}{sgn}
\DeclareMathOperator{\conv}{conv}
\newcommand{\vp}[2]{v_{#1}(#2)}
\newcommand{\cyc}{\mathrm{cyc}}
\newcommand{\sym}{\mathrm{sym}}
\newcommand{\idea}{\par\medskip\noindent\textbf{\color{accent}The idea.}\ }
\newcommand{\intuition}{\par\medskip\noindent\textbf{\color{accent}Build the intuition.}\ }
\newcommand{\strategy}{\par\medskip\noindent\textbf{\color{accent}Choosing the method.}\ }
\newcommand{\recognition}{\par\medskip\noindent\textbf{\color{accent}Recognition cues.}\ }
\newcommand{\pitfall}{\par\medskip\noindent\textbf{Common mistake.}\ }
\newcommand{\checkpoint}{\par\medskip\noindent\textbf{Checkpoint.}\ }
\newcommand{\extension}{\par\medskip\noindent\textbf{Extension.}\ }
\newcommand{}[1]{\textnormal{\small#1}}
\theoremstyle{plain}
\newtheorem{theorem}{Theorem}[chapter]
\newtheorem{lemma}[theorem]{Lemma}
\newtheorem{proposition}[theorem]{Proposition}
\newtheorem{corollary}[theorem]{Corollary}
\theoremstyle{definition}
\newtheorem{definition}[theorem]{Definition}
\newtheorem{example}{Worked example}[chapter]
\newtheorem{exercise}{Problem}[chapter]
\newtheorem{remark}[theorem]{Remark}
\newenvironment{solution}{\begin{proof}[Solution]}{\end{proof}}
\newenvironment{exploration}{\begin{quote}\small\textbf{Try before reading.}\ }{\end{quote}}
\AtBeginEnvironment{theorem}{\needspace{5\baselineskip}}
\AtBeginEnvironment{example}{\needspace{4\baselineskip}}
\AtBeginEnvironment{exercise}{\needspace{4\baselineskip}}
\tikzset{every picture/.style={line width=.65pt,scale=1.35},point/.style={circle,fill=ink,inner sep=1.3pt},aux/.style={dashed,gray},main/.style={draw=accent,thick}}
\begin{document}
\frontmatter
\begin{titlepage}
\vspace*{1.4in}
{\Huge\bfseries Olympiad Mathematics\par}
\vspace{.3in}
{\Large From First Proofs to Advanced Problem Solving\par}
\vspace{.6in}
{\large Algebra \(\cdot\) Combinatorics \(\cdot\) Geometry \(\cdot\) Number Theory\par}
\vfill
{\large 3 November 2026\par}
\end{titlepage}

% ---- front/notation ----
\chapter*{Symbols and Conventions}
\pdfbookmark[0]{Symbols and Conventions}{notation-guide}
\markboth{Symbols and Conventions}{}
This guide comes before the contents so that the notation is available
from the first chapter. You do not need to memorize it. Return here when
a symbol is unfamiliar. Every specialized symbol is also introduced in
the chapter where it first matters.

\section*{Numbers, sets, and logical language}
\renewcommand{\arraystretch}{1.12}
\begin{longtable}{@{}>{\raggedright\arraybackslash}p{.23\textwidth}>{\raggedright\arraybackslash}p{.71\textwidth}@{}}
\toprule Symbol & Meaning\\ \midrule\endhead
$\NN,\NNz$ & Positive integers $\{1,2,3,\ldots\}$ and nonnegative
integers $\{0,1,2,\ldots\}$. Some books include $0$ in $\NN$; this one does not.\\
$\ZZ,\QQ,\RR,\CC$ & Integers, rational numbers, real numbers, and complex numbers.\\
$a\in S,\ a\notin S$ & $a$ belongs to the set $S$, or does not belong to it.\\
$\varnothing,\ \{a,b\}$ & The empty set, and the set whose elements are $a,b$.\\
$A\subseteq B$ & Every element of $A$ belongs to $B$; equality is allowed.\\
$A\cup B,\ A\cap B$ & Union: elements in either set. Intersection: elements in both sets.\\
$A\setminus B$ & Elements in $A$ that are not in $B$.\\
$A\times B$ & Ordered pairs $(a,b)$ with $a\in A,b\in B$.\\
$|S|$ & Number of elements of a finite set $S$. On a number, vertical bars mean absolute value instead.\\
$[a,b],(a,b)$ & Real intervals with endpoints included or excluded. Parentheses may also denote an ordered pair; context distinguishes them.\\
$P\Rightarrow Q$ & If statement $P$ is true, then statement $Q$ is true.\\
$P\Leftrightarrow Q$ & Both implications hold: $P$ is true exactly when $Q$ is true.\\
$\forall,\ \exists$ & For every; there exists. Uniqueness also requires at most one such object.\\
$:,\ \text{such that}$ & A condition: $\{n\in\ZZ:n>0\}=\NN$.\\
$:=$ & A definition, rather than an equality requiring proof.\\
$\square$ & End of a proof or solution.\\
\bottomrule
\end{longtable}

\section*{Algebra and functions}
\begin{longtable}{@{}>{\raggedright\arraybackslash}p{.23\textwidth}>{\raggedright\arraybackslash}p{.71\textwidth}@{}}
\toprule Symbol & Meaning\\ \midrule\endhead
$|x|$ & Distance of a real number $x$ from $0$; modulus for a complex number.\\
$\sqrt{x}$ & Nonnegative square root, defined for real $x\ge0$.\\
$x^{1/n}$ & Nonnegative $n$th root for $x\ge0$. Odd roots of negative reals are specified when used.\\
$\sum_{k=1}^n a_k$ & Sum $a_1+\cdots+a_n$. The index $k$ is a temporary label.\\
$\prod_{k=1}^n a_k$ & Product $a_1\cdots a_n$. An empty sum is $0$; an empty product is $1$.\\
$n!,\binom nk$ & Factorial $1\cdot2\cdots n$, with $0!=1$; number of $k$-element subsets of an $n$-element set. For $n\ge0$, $\binom nk=0$ if $k<0$ or $k>n$.\\
$\floor{x},\ceil{x}$ & Greatest integer at most $x$; least integer at least $x$.\\
$\{x\}$ & Fractional part $x-\floor{x}$ when $x$ is a real number, not a set expression.\\
$f:A\to B$ & A function with domain $A$ and codomain $B$. Its image is $\{f(x):x\in A\}$.\\
$f\circ g,\ f^n$ & Composition $(f\circ g)(x)=f(g(x))$; $f^n$ denotes $n$ iterations only when explicitly stated. A power $f(x)^n$ is different.\\
$f^{-1}$ & Inverse function when $f$ is bijective, or an inverse image of a set when specified. It does not mean $1/f$.\\
$P(x),\deg P$ & A polynomial; its degree when $P$ is nonzero. The zero polynomial has no ordinary finite degree.\\
$\ZZ[x],\QQ[x]$ & Polynomials with integer or rational coefficients. Integer values do not imply integer coefficients.\\
$\sum_{\mathrm{cyc}}$ & A cyclic sum: for three variables, use $(a,b,c),(b,c,a),(c,a,b)$. A symmetric sum uses all distinct permutations.\\
$\log_a x$ & Exponent $t$ satisfying $a^t=x$, for $a>0,a\ne1,x>0$.\\
$i,\overline z$ & $i^2=-1$; complex conjugate. If $z=x+iy$, then $\overline z=x-iy$.\\
$a_n\to L$ & A sequence approaches a limit $L$, defined precisely when used.\\
\bottomrule
\end{longtable}

\section*{Number theory}
\begin{longtable}{@{}>{\raggedright\arraybackslash}p{.23\textwidth}>{\raggedright\arraybackslash}p{.71\textwidth}@{}}
\toprule Symbol & Meaning\\ \midrule\endhead
$a\mid b,\ a\nmid b$ & $a$ divides $b$: $b=ak$ for an integer $k$; or does not divide it. This bar is not absolute value.\\
$\gcd(a,b)$ & Greatest common positive divisor; $\gcd(a,0)=|a|$ for $a\ne0$.\\
$\lcm(a,b)$ & Least common positive multiple of positive integers $a,b$.\\
$a\equiv b\pmod m$ & $m\mid(a-b)$, for a positive modulus $m$.\\
$a^{-1}\pmod m$ & A residue $u$ with $au\equiv1\pmod m$; it exists exactly when $\gcd(a,m)=1$.\\
$v_p(n)$ & Exponent of prime $p$ in nonzero integer $n$. We use $v_p(0)=+\infty$ only when explicitly stated.\\
$\varphi(n)$ & Number of integers in $\{1,\ldots,n\}$ coprime to $n$.\\
$\tau(n),\sigma(n)$ & Number of positive divisors of $n$; their sum.\\
$\ord_m(a)$ & Least positive $k$ with $a^k\equiv1\pmod m$, for $m\ge2$ and $\gcd(a,m)=1$.\\
$\left(\frac ap\right)$ & Legendre symbol for odd prime $p$: $0$ if $p\mid a$, $1$ if $a$ is a nonzero square modulo $p$, and $-1$ otherwise.\\
$\FF_p$ & Residue classes modulo prime $p$, with addition and multiplication. Nonzero classes have inverses.\\
$\rad(n)$ & Product of distinct prime divisors of $n>0$; $\rad(1)=1$.\\
\bottomrule
\end{longtable}

\section*{Geometry}
\begin{longtable}{@{}>{\raggedright\arraybackslash}p{.23\textwidth}>{\raggedright\arraybackslash}p{.71\textwidth}@{}}
\toprule Symbol & Meaning\\ \midrule\endhead
$AB,\ |AB|$ & Length of segment $AB$ unless a directed ratio is specified.\\
$AB\parallel CD$ & Lines $AB$ and $CD$ are parallel.\\
$AB\perp CD$ & The lines are perpendicular.\\
$\angle ABC$ & Ordinary angle with vertex $B$, between rays $BA$ and $BC$.\\
$\dang(AB,CD)$ & Directed angle between the indicated lines, modulo $180^\circ$ unless specified otherwise.\\
$(ABC)$ & Circle through noncollinear points $A,B,C$. A general circle is often $\omega$.\\
$[ABC]$ & Triangle area, nonnegative unless signed area is expressly used.\\
$a,b,c$ & In triangle $ABC$, sides $BC,CA,AB$, respectively. Its angles are $A,B,C$.\\
$s,\Delta$ & Semiperimeter $(a+b+c)/2$; area of the reference triangle.\\
$O,R;\ I,r$ & Circumcenter and circumradius; incenter and inradius.\\
$G,H$ & Centroid and orthocenter of the current reference triangle.\\
$\operatorname{Pow}_\omega(P)$ & Power to circle $\omega$: $PO^2-R^2$, for center $O$ and radius $R$.\\
$\vec u\cdot\vec v$ & Dot product. $\|\vec u\|$ is vector length.\\
$\det(\vec u,\vec v)$ & Signed parallelogram area of two plane vectors.\\
$(x:y:z)$ & Homogeneous barycentric coordinates relative to a specified triangle; common nonzero scaling changes no point.\\
\bottomrule
\end{longtable}

\section*{Combinatorics and probability}
\begin{longtable}{@{}>{\raggedright\arraybackslash}p{.23\textwidth}>{\raggedright\arraybackslash}p{.71\textwidth}@{}}
\toprule Symbol & Meaning\\ \midrule\endhead
$G=(V,E)$ & Finite graph with vertex set $V$ and edge set $E$. Unless stated otherwise, graphs are undirected, with no loops or repeated edges.\\
$d(v),N(v)$ & Degree; set of neighbors of $v$. For a vertex set $S$, its neighborhood $N(S)$ is defined in context.\\
$K_n,\ K_{r,s}$ & Complete graph on $n$ vertices; complete bipartite graph with parts of sizes $r,s$.\\
$\Pr(A)$ & Probability of event $A$ in a specified finite experiment. Equal likelihood is never automatic.\\
$\mathbb E X$ & Expected value: the probability-weighted average of $X$.\\
$\mathbf1_A$ & Indicator: $1$ when event $A$ occurs, $0$ otherwise.\\
$S_n$ & Permutations of $\{1,\ldots,n\}$ when used for symmetry; elsewhere $S_n$ may be a partial sum.\\
$\oplus$ & Bitwise exclusive-or in the Nim chapter, not ordinary addition.\\
\bottomrule
\end{longtable}
\renewcommand{\arraystretch}{1}

\section*{Conventions that prevent hidden errors}
Lengths and ordinary areas are nonnegative. A triangle is nondegenerate
unless a limiting case is expressly discussed. Every division requires
a nonzero denominator. Every square root requires its stated real domain.
The phrase \emph{without loss of generality} requires a symmetry of the
whole problem. A drawing records hypotheses; it cannot prove an
unstated relation.


\begingroup
\setlength{\parskip}{0pt}
\tableofcontents
\endgroup

% ---- front/reading ----
\chapter*{How to Learn from This Book}
\addcontentsline{toc}{chapter}{How to Learn from This Book}
\markboth{How to Learn from This Book}{}
This book develops algebra, combinatorics, Euclidean geometry, and
number theory from first proofs to advanced olympiad problem solving.
Its starting point is comfort with whole numbers, fractions, and
elementary arithmetic. If a school-algebra step is unfamiliar, work the
opening algebra chapters before more advanced chapters in any subject.

First we ask what a method is trying to achieve. Then we define its
objects, prove the precise result, and watch the method work in examples.
A successful solution contains both a choice and an argument: why this
move is promising, and why it is valid. Pay attention to both.

\section*{A route through the material}
Read Learning to Prove first. Then study the early chapters of each
subject rather than completing an entire subject before meeting the
others. Counting and divisibility reinforce induction; similarity and
area ratios reinforce algebra; polynomials and graphs make structural
thinking visible. Later chapters connect these viewpoints.

Within each subject, the chapter order supplies the prerequisites.
An extension strengthens the main route; its name alone need not be
memorized. An unfamiliar theorem is often less useful than a familiar
principle applied carefully. The final part removes the chapter title
as a clue and asks you to choose your own approach.

\section*{What the problem labels mean}
\begin{description}[style=nextline]
\item[Foundation]
Check a definition, reproduce a key argument, or make one main choice.
These are essential even when their answers look short.
\item[Bridge]
Combine ideas or adapt an example after a feature of the problem changes.
These develop transfer beyond imitation.
\item[Challenge]
Make a substantial discovery, sustain a longer argument, or classify
all possibilities. Difficulty depends on experience; these labels are
guides, not official contest ratings.
\end{description}

Hints and complete solutions are at the end. A hint should leave useful
work to do. After reading a solution, close it and reconstruct the
reasoning. Then change one hypothesis and ask which step survives.
Recognizing a written solution is a different skill from producing
one independently.

\section*{How to read a proof actively}
At each line, ask what justifies the step, which hypotheses it uses,
and what would make it fail. Before a proof, test small examples and
equality cases. Before a worked solution, predict a useful invariant,
factorization, auxiliary point, or modulus. If a prediction fails,
identify what it failed to control.

When an operation improves a configuration, specify what improves and
why the process stops. After finding candidates, substitute them into
the original conditions. Check whether a proof treats obtuse triangles,
zero values, negative variables, and boundary cases.

\section*{The advanced goal}
The advanced chapters and challenge problems support preparation for
the strongest olympiad problems. A medal is an outcome of independent
performance on unfamiliar problems. Judge progress by the proofs you
can discover, explain, and adapt, rather than the pages you have read.

The book focuses on mathematics and problem-solving decisions.
It contains no schedule or contest time-management system.


\mainmatter
\part{Learning to Prove}

% ---- chapters/foundations ----
\chapter{What a Proof Does}
\label{ch:proof-language}
\intuition A calculation can suggest a pattern. A proof explains why no
allowed example can escape it. If the claim concerns every integer,
testing a million integers still leaves infinitely many untested.
The useful question is therefore not how many cases we have checked,
but what feature all cases share.

\section{Objects, assumptions, and conclusions}
A variable stands for an object in a specified set. In
$n=2k+1$ with $k\in\ZZ$, the variable $n$ is an odd integer.
The same formula with $k\in\RR$ describes every real number. The
formula alone does not carry the domain.

A definition fixes the meaning of a word. An integer $n$ is
\emph{even} when $n=2k$ for some $k\in\ZZ$, and \emph{odd} when
$n=2k+1$ for some $k\in\ZZ$. We will prove in number theory that
every integer has exactly one of these forms. A theorem is a claim
that follows from definitions, assumptions, and already established
results.

\begin{example}[Turn a pattern into an explanation]
Prove that the sum of two odd integers is even.
\end{example}
\begin{solution}
Write the integers as $2a+1$ and $2b+1$, where $a,b\in\ZZ$.
Their sum is $2(a+b+1)$. The expression in parentheses is an integer,
which is exactly what the definition of evenness requires.
\end{solution}
\strategy We replaced a property by its definition. This is often the
best first move when a problem asks about parity, divisibility,
rationality, parallelism, or injectivity.

The \emph{hypotheses} are what we may assume; the \emph{conclusion}
is what must be established. In the statement
\[
x>1\quad\Longrightarrow\quad x^2>x,
\]
the inequality $x>1$ makes $x(x-1)>0$. The conclusion need not hold
outside that domain: at $x=1/2$ it is false.

\section{Implications, converses, and quantifiers}
An implication $P\Rightarrow Q$ does not assert its converse
$Q\Rightarrow P$. Every multiple of $4$ is even, but $6$ shows that
not every even number is a multiple of $4$. An equivalence
$P\Leftrightarrow Q$ requires both directions.

\begin{example}[A proof with two directions]
For $n\in\ZZ$, prove that $n$ is even if and only if $n^2$ is even.
\end{example}
\begin{solution}
If $n=2k$, then $n^2=2(2k^2)$ is even. Conversely, suppose $n^2$
is even. If $n$ were odd, $n=2k+1$ would give
$n^2=2(2k^2+2k)+1$, which is odd. This contradiction proves that
$n$ is even.
\end{solution}

The contrapositive of $P\Rightarrow Q$ is
$\neg Q\Rightarrow\neg P$, where $\neg$ means negation. It is
logically equivalent to the original implication. A proof by
contradiction assumes the desired conclusion false and derives an
impossibility; a contrapositive proof directly proves the alternative
implication. Keeping these structures visible improves clarity.

The order of quantifiers matters:
\[
\forall x\in\RR\ \exists y\in\RR:\ y>x
\]
is true, since $y=x+1$ works. The assertion
$\exists y\in\RR\ \forall x\in\RR:\ y>x$ is false, since a proposed
$y$ fails for $x=y+1$. In the first statement $y$ may depend on $x$;
in the second it must work before $x$ is known.

To negate a universal assertion, produce one exception. To negate an
existence assertion, show that every candidate fails. Thus
\[
\neg(\forall x\,\exists y\,P(x,y))
\quad\Longleftrightarrow\quad
\exists x\,\forall y\,\neg P(x,y).
\]
\pitfall A counterexample must satisfy every hypothesis. A right
triangle cannot refute a statement about acute triangles.

\section{Existence, uniqueness, and classification}
These are three different tasks. Existence constructs or forces an
object. Uniqueness rules out two different objects. Classification
describes every object and checks that every item on the list works.

\begin{example}[Candidates are not the answer]
Solve $\sqrt{x+2}=x$ over the real numbers.
\end{example}
\begin{solution}
The right side must be nonnegative, so $x\ge0$. Squaring gives
$x+2=x^2$, hence $(x-2)(x+1)=0$. The candidates are $2,-1$.
The domain excludes $-1$; substitution verifies $x=2$.
\end{solution}
\idea Squaring loses sign information. We recorded the lost information
before using the easier polynomial equation. In a classification,
write the reverse check as carefully as the forward deductions.

\begin{example}[Uniqueness without finding the value]
Suppose $f:\RR\to\RR$ is strictly increasing. Prove that the equation
$f(x)=0$ has at most one solution.
\end{example}
\begin{solution}
If $a<b$ were two solutions, strict increase would give
$0=f(a)<f(b)=0$, impossible.
\end{solution}
This proves no existence statement: the argument never constructs
a zero. For a simple illustration on a restricted domain, the
strictly increasing function $f:(0,1)\to\RR$, $f(x)=x$, has none.
The logical phrase \emph{at most one} includes zero.

\section{Legal transformations}
An equality may be added to another equality, multiplied by the same
quantity, or substituted into a valid expression. Reversing a step
requires checking what it discarded. Division by $x-y$ discards the
case $x=y$. Taking a square root chooses a sign. Multiplying an
inequality by a negative number reverses its direction.

\begin{example}[A convincing false proof]
Starting with $a=b\ne0$, one may derive
$a^2=ab$, then $a^2-b^2=ab-b^2$, then
$(a-b)(a+b)=b(a-b)$. Why does cancellation not prove $2=1$?
\end{example}
\begin{solution}
The cancelled factor is $a-b=0$. The last displayed equation is
$0=0$; it contains no information about $a+b$ versus $b$.
\end{solution}

An estimate also has a direction. If $A\le B$ and $B\ge C$, there is
no comparison between $A$ and $C$. Draw the number line, or invent
numbers: $1\le10$ and $10\ge7$ do not imply $1\ge7$.

\section{Practice: make the logical task explicit}
\begin{exercise}[{Foundation}]\label{f:ex01}
Prove that the product of two odd integers is odd.
\end{exercise}
\begin{exercise}[{Foundation}]\label{f:ex02}
Negate: for every integer $n$ there exists an integer $m$ such that
$m^2=n$. Decide whether the original statement is true.
\end{exercise}
\begin{exercise}[{Foundation}]\label{f:ex03}
Solve $x^2=|x|$ over $\RR$, with a complete reverse check.
\end{exercise}
\begin{exercise}[{Bridge}]\label{f:ex04}
Prove that if $x$ is rational and $y$ is irrational, then $x+y$ is
irrational. Is $xy$ always irrational?
\end{exercise}
\begin{exercise}[{Bridge}]\label{f:ex05}
Find all real $x$ satisfying $\sqrt{x+6}=x$.
\end{exercise}
\begin{exercise}[{Bridge}]\label{f:ex06}
If $a,b>0$ and $a+b=1$, prove $ab\le1/4$ and determine equality.
Explain why an example with $a+b\ne1$ is irrelevant.
\end{exercise}
\begin{exercise}[{Challenge}]\label{f:ex07}
Prove that $\sqrt2$ is irrational. Identify exactly where the
assumption that a fraction is in lowest terms enters.
\end{exercise}
\begin{exercise}[{Challenge}]\label{f:ex08}
Let $S$ be a nonempty finite set of real numbers. Suppose its arithmetic
mean equals its largest element. Prove that all its elements are equal.
Then give an example showing that the mean need not equal an element.
\end{exercise}

\chapter{Induction, Descent, and Extremal Choices}
\label{ch:induction}
\intuition Many infinite claims contain a mechanism that moves from
one case to the next. Induction proves that the mechanism never breaks.
Descent turns the mechanism around: a hypothetical counterexample
would force a smaller one forever, which positive integers cannot do.
An extremal choice selects an object whose improvement is impossible.

\section{Ordinary and strong induction}
\begin{theorem}[Induction principle]
Let $P(n)$ be a statement for integers $n\ge n_0$. If $P(n_0)$ is true
and $P(n)\Rightarrow P(n+1)$ for every $n\ge n_0$, then all $P(n)$
are true.
\end{theorem}
\begin{proof}
Otherwise the set of failing indices has a least element $k$.
The base case gives $k>n_0$. Minimality makes $P(k-1)$ true,
and the transition makes $P(k)$ true, a contradiction.
\end{proof}
We used the \emph{well-ordering principle}: every nonempty set of
nonnegative integers has a least element. This is a basic property of
the integers. It is false for positive reals; $\{1/n:n\in\NN\}$
has no least element.

\begin{example}[Find the transition rather than repeat the formula]
Prove $1+3+\cdots+(2n-1)=n^2$ for $n\ge1$.
\end{example}
\begin{solution}
For $n=1$ both sides equal $1$. If the sum through $2n-1$ is $n^2$,
the next sum is $n^2+(2n+1)=(n+1)^2$. That proves the transition.
\end{solution}
\begin{center}
\begin{tikzpicture}[scale=.48]
\foreach \x in {0,1,2,3}\foreach \y in {0,1,2,3}
{\draw[fill=soft] (\x,\y) rectangle +(1,1);}
\foreach \x in {0,1,2,3,4}{\draw[fill=accent!25] (\x,4) rectangle +(1,1);}
\foreach \y in {0,1,2,3}{\draw[fill=accent!25] (4,\y) rectangle +(1,1);}
\node at (2,-.7){$n^2$};
\node[anchor=west] at (5.6,2.5){Add $2n+1$ squares};
\node[anchor=west] at (5.6,1.5){to make $(n+1)^2$.};
\end{tikzpicture}
\end{center}
The drawing explains the identity; the induction makes the explanation
valid for every $n$.

In \emph{strong induction}, the transition for $n$ may assume all
smaller cases, not just the preceding one. This is natural when an
object splits into pieces of unpredictable sizes.

\begin{example}[Prime factorizations exist]
Every integer $n\ge2$ is a product of primes.
\end{example}
\begin{solution}
A prime is an integer greater than $1$ with no positive divisors other
than $1$ and itself. Assume every integer from $2$ to $n-1$ has a prime
factorization. If $n$ is prime, it is already a product with one factor.
If not, $n=ab$ with $2\le a,b<n$, so factor $a,b$ using the induction
hypothesis and concatenate their factors.
\end{solution}
This proves existence; uniqueness needs a different argument later.

\section{Strengthen the statement when the transition needs it}
Sometimes the desired statement is too weak to support the next step.
Prove additional structure that survives the induction.

\begin{example}[A missing square]
Remove one square from a $2^n$ by $2^n$ board. Prove that the remaining
board can be tiled by L-shaped trominoes, each covering three unit squares.
\end{example}
\begin{solution}
For $n=1$, the three remaining squares form one tromino.
For a larger board, divide it into four equal quadrants. Exactly one
quadrant contains the missing square. Put one tromino across the center,
covering one central square in each of the other three quadrants.
Each quadrant now has exactly one unavailable square. Apply induction
to the four quadrants.
\end{solution}
\strategy The stronger claim allows the missing square anywhere.
Proving only the corner-missing case would not justify the four
recursive calls.

\section{Descent: choose a quantity that cannot decrease forever}
\begin{example}[A minimal solution cannot exist]
The equation $x^2=2y^2$ has no positive integer solutions.
\end{example}
\begin{solution}
Suppose there is a solution with least positive $y$. Since $x^2$
is even, $x=2u$. Then $y^2=2u^2$, so $y=2v$. It follows that
$u^2=2v^2$, with $0<v=y/2<y$, contradicting minimality.
\end{solution}
\pitfall A strictly decreasing real quantity does not imply termination.
The sequence $1,1/2,1/4,\ldots$ decreases forever. A strictly decreasing
nonnegative integer, or a strictly decreasing choice in a finite set,
does give termination.

A descent proof must check four points: the new object exists; it
remains in the allowed domain; it satisfies the same essential
conditions; and its chosen measure is smaller. Missing any one of
these turns descent into a guess.

\section{Extremal choices and invariants}
An \emph{extremal} object maximizes or minimizes a specified quantity.
In a nonempty finite family, such an object exists. In an infinite
family, existence needs a proof.

\begin{example}[A farthest pair forces a bound]
In a finite set of at least two distinct plane points, choose $A,B$
at maximum distance. Show that every point lies in the intersection
of the closed disks of radius $AB$ centered at $A$ and $B$.
\end{example}
\begin{solution}
For any point $P$ in the set, maximality gives $AP\le AB$ and
$BP\le AB$. These are exactly the two disk-membership conditions.
\end{solution}
The conclusion is elementary, but the choice turns many comparisons
into automatic consequences. Later, a maximum matching, minimum
counterexample, or shortest segment will play the same role.

An \emph{invariant} stays unchanged under an allowed operation. A
\emph{monovariant} moves in only one direction. Either can make a
long process understandable without tracing all its states.

\begin{example}[Reachability can fail before a search begins]
A move on an integer pair $(a,b)$ replaces it by $(a+b,b)$ or
$(a,a+b)$. Starting from $(2,4)$, can one reach $(3,5)$?
\end{example}
\begin{solution}
No. The gcd is unchanged because a common divisor of $a,b$ is a
common divisor of $a+b,b$, and conversely a divisor of $a+b,b$
also divides $a$. Initially the gcd is $2$, whereas $\gcd(3,5)=1$.
\end{solution}
We will establish the gcd facts systematically in number theory.
Here the conceptual point is that an invariant provides an obstruction;
having the right invariant value does not by itself prove reachability.

\section{Practice: justify the mechanism}
\begin{exercise}[{Foundation}]\label{f:ex09}
Prove by induction that $1+\cdots+n=n(n+1)/2$ for $n\ge1$.
\end{exercise}
\begin{exercise}[{Foundation}]\label{f:ex10}
Prove $2^n\ge n+1$ for every integer $n\ge0$.
\end{exercise}
\begin{exercise}[{Bridge}]\label{f:ex11}
Prove that every integer $n\ge2$ has a prime divisor.
Give a proof using a least divisor and another using induction.
\end{exercise}
\begin{exercise}[{Bridge}]\label{f:ex12}
Prove that every integer $n\ge8$ is $3a+5b$ for some nonnegative
integers $a,b$.
\end{exercise}
\begin{exercise}[{Bridge}]\label{f:ex13}
Starting from $(1,1)$, repeatedly apply $(a,b)\mapsto(a+b,b)$
or $(a,b)\mapsto(a,a+b)$. Prove that every resulting pair is coprime.
Must the same statement hold starting from $(2,3)$?
\end{exercise}
\begin{exercise}[{Challenge}]\label{f:ex14}
For every coprime pair of positive integers $(a,b)$, prove that it
can be obtained from $(1,1)$ by the two moves in Problem~\ref{f:ex13}.
\end{exercise}
\begin{exercise}[{Challenge}]\label{f:ex15}
In a complete graph every pair of distinct vertices is joined by an
edge. A path lists distinct vertices with edges between successive
ones. Every edge of a finite complete graph is colored red or blue.
Prove that its vertices can be partitioned into a red path and a blue
path, allowing either path to be empty or a single vertex.
\end{exercise}
\begin{exercise}[{Challenge}]\label{f:ex16}
Let $a_1,\ldots,a_n$ be positive integers. Replacing an adjacent pair
$(x,y)$ by $(x+y,y)$ or $(x,x+y)$ never changes the gcd of all entries.
Prove this assertion. Explain why it is only a necessary test for
reaching a specified list.
\end{exercise}

\chapter{Finding the Key Move}
\label{ch:discovery}
\intuition A finished proof hides the search that produced it.
To learn problem solving, reconstruct the reason for its decisive
choice. What was hard to control? What new representation made that
feature simpler? This chapter develops questions that remain useful
in every subject.

\section{Work backward, then write forward}
Suppose the goal is $A\ge B$. If we recognize $A-B$ as a sum of
squares, the inequality becomes automatic. Working backward suggests
the squares; the written proof expands them and verifies the identity.
Backward reasoning is valid discovery even when it is not yet a proof.

\begin{example}[The target tells us what to build]
For real $a,b,c$, prove
$a^2+b^2+c^2\ge ab+bc+ca$.
\end{example}
\begin{solution}
The difference is
\[
\frac12\bigl((a-b)^2+(b-c)^2+(c-a)^2\bigr)\ge0.
\]
Equality holds exactly when $a=b=c$.
\end{solution}
\strategy Equality suggests the differences $a-b,b-c,c-a$.
Testing equality cases is a way to discover a proof, not merely
something to add after the proof.

\section{Change representations}
A sum can be counted, a geometric ratio can be expressed as an area,
an integer can be replaced by its prime exponents, and a repeated
operation can be viewed as a graph. The purpose is to make the
essential constraint easy to preserve.

\begin{example}[Partial sums expose a hidden difference]
Given integers $a_1,\ldots,a_n$, prove that some nonempty consecutive
block has sum divisible by $n$.
\end{example}
\begin{solution}
Set $S_0=0$ and $S_k=a_1+\cdots+a_k$ for $1\le k\le n$.
The $n+1$ partial sums occupy only $n$ residue classes modulo $n$,
so $S_i\equiv S_j\pmod n$ for some $i<j$.
Then $a_{i+1}+\cdots+a_j=S_j-S_i$ is divisible by $n$.
\end{solution}
\idea The problem asks about many intervals. Partial sums replace
every interval sum by a difference of two objects, so a collision
principle becomes available.

\section{Symmetry, normalization, and boundary cases}
If a condition is symmetric in $a,b,c$, we may relabel them to arrange
$a\ge b\ge c$. A cyclic condition may not permit every relabeling.
Always check the hypotheses and the target together.

If multiplying every variable by the same positive factor leaves
the claim unchanged, a normalization such as $a+b+c=1$ can remove
one degree of freedom. The factor must exist, and the transformed
domain must remain valid. A positive-variable problem permits
division by $a+b+c$; a general real-variable problem may not.

The boundary can expose a false conjecture or indicate where a minimum
lies. If a proposed inequality is strict but equality occurs at
$a=b$, the target is wrong. If a theorem excludes a collinear
configuration, do not silently take a limit through it without
checking what objects cease to exist.

\section{Small cases, conjectures, and counterexamples}
Small cases serve three purposes: they test a claim, suggest its
structure, and reveal exceptions. They do not prove an unrestricted
claim. When a conjecture fails, use the failure to formulate the
correct statement.

\begin{example}[A plausible polynomial pattern ends]
The expression $n^2+n+41$ is prime for many small nonnegative integers.
Prove that it is not always prime.
\end{example}
\begin{solution}
At $n=40$, its value is $40^2+40+41=41^2$, composite.
At $n=41$, it is $41(41+1)+41=41\cdot43$.
\end{solution}
\recognition For a polynomial with integer coefficients, substituting
a multiple of a known value often creates divisibility. Number theory
will turn this observation into a theorem about prime-valued polynomials.

\section{Write a proof a reader can audit}
A clear proof identifies its variables, states the decisive idea,
and supplies the intermediate implications. Distinguish what has been
proved from what is still needed. Avoid replacing an argument by a
theorem name when the theorem's hypotheses are the hard part.

For an algorithm, prove termination and correctness separately.
For a minimum or maximum, give a bound and an example attaining it.
For an existence result from an average, explain why at least one
outcome meets the average. For a construction, check every required
property rather than only the property that inspired it.

\checkpoint After any solution, ask four questions. Where did the
information enter? What was the key choice? Which step would fail if
a hypothesis changed? Can the method solve a nearby problem?

\section{Practice: choose a representation}
\begin{exercise}[{Foundation}]\label{f:ex17}
For real $a,b$, prove $a^2+b^2\ge2ab$ and state equality.
\end{exercise}
\begin{exercise}[{Foundation}]\label{f:ex18}
Find a counterexample to each statement: the sum of two primes is prime;
the product of two irrational reals is irrational; every increasing
sequence is bounded.
\end{exercise}
\begin{exercise}[{Bridge}]\label{f:ex19}
Prove that among any six integers, two have difference divisible by $5$.
Explain why the same conclusion need not hold for five integers.
\end{exercise}
\begin{exercise}[{Bridge}]\label{f:ex20}
Find the minimum of $x+1/x$ for positive real $x$, and prove that it
is attained.
\end{exercise}
\begin{exercise}[{Bridge}]\label{f:ex21}
For integers $n\ge1$, prove
$\sum_{k=1}^n1/(k(k+1))=n/(n+1)$.
\end{exercise}
\begin{exercise}[{Challenge}]\label{f:ex22}
Given $n+1$ distinct integers chosen from $\{1,\ldots,2n\}$, prove
that one of them divides another.
\end{exercise}
\begin{exercise}[{Challenge}]\label{f:ex23}
For positive real $a,b,c$ with $abc=1$, prove
$(1+a)(1+b)(1+c)\ge8$ and determine equality.
\end{exercise}
\begin{exercise}[{Challenge}]\label{f:ex24}
Start with any finite list of real numbers. Whenever two adjacent entries
are in decreasing order, exchange them. Prove that the process ends
in a nondecreasing list, regardless of which eligible pair is chosen.
\end{exercise}


\part{Algebra}

% ---- chapters/algebra ----
\chapter{Expressions, Meaning, and Legal Moves}
\label{a:expressions}

Algebra is the art of changing the form of a statement without changing
what it means. The same expression can reveal different information in
different forms. The expression \(x^2-6x+9\) looks like a sum of three
terms; the form \((x-3)^2\) immediately reveals that it is nonnegative
and vanishes only at \(x=3\). In an olympiad solution, discovering the
useful form is often the main work.

\section{Equality is a promise}

An \emph{expression} such as \(x^2+1\) names a number after \(x\) is
specified. An \emph{equation} such as \(x^2+1=2x\) asserts that two
expressions have the same value. An \emph{identity} is an equation true
for every value in its stated domain. Before manipulating a fraction,
square root, or logarithm, specify where it is defined.

For real numbers, \(\sqrt{x}\) denotes the nonnegative square root and
requires \(x\ge0\). The expression \(1/x\) requires \(x\ne0\).
For a base \(b>0\), \(b\ne1\), the real logarithm \(\log_b x\) is the
number \(t\) satisfying \(b^t=x\), and requires \(x>0\).
The absolute value \(|x|\) is \(x\) for \(x\ge0\), and \(-x\) for \(x<0\).
Thus \(\sqrt{x^2}=|x|\), not always \(x\).

\intuition Think of solving an equation as walking through doors.
Adding the same quantity to both sides is a door that opens in both
directions. Squaring opens only one way unless signs are controlled.
Dividing by an expression may close the door on the cases where it is zero.

\begin{example}[A solution lost by division]
Solve \(x(x-2)=3(x-2)\) over the real numbers.
\end{example}
\begin{solution}
Bring everything to one side:
\[
 (x-2)(x-3)=0.
\]
A product of real numbers is zero only if at least one factor is zero.
The solutions are \(2\) and \(3\). Dividing first by \(x-2\) would have
found \(3\) while discarding the valid case \(x=2\).
\end{solution}

\begin{example}[A solution introduced by squaring]
Solve \(\sqrt{x+2}=x\).
\end{example}
\begin{solution}
The left side is nonnegative, so any solution must have \(x\ge0\).
Squaring gives \(x+2=x^2\), or \((x-2)(x+1)=0\).
Only \(x=2\) satisfies the original sign condition, and it checks.
\end{solution}

\pitfall The phrase \emph{if and only if} must be earned at every step.
From \(u=v\) one may conclude \(u^2=v^2\); the reverse conclusion is
only \(u=v\) or \(u=-v\).

\section{Factor before expanding}

The basic identities
\[
 u^2-v^2=(u-v)(u+v),\qquad
 u^3-v^3=(u-v)(u^2+uv+v^2)
\]
come from multiplication. More generally, for an integer \(n\ge1\),
\[
 u^n-v^n=(u-v)\sum_{j=0}^{n-1}u^{n-1-j}v^j.
\]
Multiplying the right side causes all intermediate terms to cancel.
This is the same mechanism as telescoping a sum.

\begin{example}[Find the missing square]
Factor \(x^4+4y^4\).
\end{example}
\begin{solution}
The two terms resemble \((x^2+2y^2)^2\), except for its middle term.
Add and subtract that term:
\[
\begin{aligned}
x^4+4y^4
  &=(x^2+2y^2)^2-(2xy)^2\\
  &=(x^2-2xy+2y^2)(x^2+2xy+2y^2).
\end{aligned}
\]
The useful idea is completing a square; the final expansion is only a check.
\end{solution}

\recognition A difference of even powers suggests a difference of squares.
An alternating expression that becomes zero when \(x=y\) suggests a
factor \(x-y\). A symmetric expression suggests using sums and products
rather than expanding without a plan.

\section{Sums, products, and cancellation}

The notation \(\sum_{k=r}^{s}a_k\) means
\(a_r+a_{r+1}+\cdots+a_s\); \(\prod_{k=r}^{s}a_k\) means their product.
An empty sum is \(0\), and an empty product is \(1\).
For \(n\ge1\),
\[
 \sum_{k=1}^n k=\frac{n(n+1)}2.
\]
One proof pairs the sum with its reverse: every paired column is \(n+1\),
and there are \(n\) columns.

\begin{example}[The form that makes cancellation visible]
For a positive integer \(n\), evaluate
\(\sum_{k=1}^n1/[k(k+1)]\).
\end{example}
\begin{solution}
Because \(1/[k(k+1)]=1/k-1/(k+1)\), all interior terms cancel:
\[
 \sum_{k=1}^n\left(\frac1k-\frac1{k+1}\right)
 =1-\frac1{n+1}=\frac n{n+1}.
\]
The partial-fraction identity creates adjacent terms on purpose.
\end{solution}

\checkpoint When an expression is difficult, ask what information you
want: zeros suggest factoring, a lower bound suggests squares, and a long
sum suggests a difference of neighboring terms.

\section{Practice}
\begin{exercise}[{Foundation}]\label{a:ex01}
Factor \(x^4+4y^4\), and use the factorization to explain why
\(x^4+4\) is composite for every integer \(x>1\).
\end{exercise}
\begin{exercise}[{Foundation}]\label{a:ex02}
Solve \(1/(x-1)+1/(x+1)=1\) over the real numbers.
\end{exercise}
\begin{exercise}[{Foundation}]\label{a:ex03}
Solve \(\sqrt{x+6}=x\) over the real numbers.
\end{exercise}
\begin{exercise}[{Bridge}]\label{a:ex04}
If \(x+y+z=0\), prove \(x^3+y^3+z^3=3xyz\).
\end{exercise}
\begin{exercise}[{Bridge}]\label{a:ex05}
Evaluate \(\prod_{k=1}^{n}(1+1/k)\), and explain why there is no
corresponding rule \(\prod(1+t_k)=1+\prod t_k\).
\end{exercise}
\begin{exercise}[{Bridge}]\label{a:ex06}
For real \(x,y,z\), show that
\(x^2+y^2+z^2=xy+yz+zx\) holds exactly when \(x=y=z\).
\end{exercise}
\begin{exercise}[{Challenge}]\label{a:ex07}
Solve \(\sqrt{x+1}+\sqrt{4-x}=3\) over the real numbers.
\end{exercise}

\chapter{Equations, Discriminants, and Symmetry}
\label{a:equations}

\section{Why the quadratic discriminant appears}

\intuition A quadratic equation asks where a parabola meets the horizontal
axis. Completing the square shifts the picture until the answer is
controlled by the sign of one number. The familiar formula is the record
of this simple geometric operation.

\begin{theorem}[Quadratic formula]
For real \(a,b,c\), with \(a\ne0\), the equation \(ax^2+bx+c=0\) has real
solutions exactly when \(\Delta=b^2-4ac\ge0\). The solutions are
\[
 x=\frac{-b\pm\sqrt{\Delta}}{2a}.
\]
\end{theorem}
\begin{proof}
Multiplying by \(4a\) and completing the square gives
\((2ax+b)^2=b^2-4ac\). A real square is nonnegative; conversely,
when the right side is nonnegative, taking both signs of its square
root gives all solutions.
\end{proof}

\begin{example}[An integer obstruction]
Can integers \(x,y\) satisfy \(x+y=7\) and \(xy=11\)?
\end{example}
\begin{solution}
Both would be roots of \(t^2-7t+11=0\), whose discriminant is \(5\).
An integer root \(t\) would make \((2t-7)^2=5\), impossible.
The discriminant detects the obstruction before we search pairs.
\end{solution}

An integer-coefficient quadratic with an integer root must have a square
discriminant. The converse only guarantees rational roots: for example,
\(2t^2+t-1=0\) has roots \(1/2\) and \(-1\).

\section{Replace the variables by what the expression remembers}

An expression is \emph{symmetric} in \(x,y\) if swapping \(x,y\) does not
change it. Put \(s=x+y\) and \(p=xy\). Then
\[
 x^2+y^2=s^2-2p,\qquad x^3+y^3=s^3-3sp.
\]
The variables \(x,y\) can be recovered as the roots of
\(t^2-st+p=0\). They are real exactly when \(s^2\ge4p\).
If they are real, both are positive exactly when \(s>0\) and \(p>0\).

\begin{example}[Eliminate structure, not variables at random]
Find all real pairs satisfying \(x+y=5\), \(x^3+y^3=35\).
\end{example}
\begin{solution}
The cube identity gives \(125-15xy=35\), so \(xy=6\).
Thus \(x,y\) are the roots of \(t^2-5t+6=(t-2)(t-3)\).
The pairs are \((2,3)\) and \((3,2)\).
\end{solution}

For three variables define
\[
 p=x+y+z,\qquad q=xy+yz+zx,\qquad r=xyz.
\]
These are the elementary symmetric quantities. Direct multiplication gives
\[
 (t-x)(t-y)(t-z)=t^3-pt^2+qt-r.
\]
In particular, specifying \(p,q,r\) determines the unordered triple,
including repeated values. Not every real triple \(p,q,r\) corresponds to
three real roots; feasibility must still be checked.

\begin{example}[The equality case is the whole answer]
Find all real triples with \(x+y+z=3\), \(xy+yz+zx=3\).
\end{example}
\begin{solution}
The sum of squares is \(9-6=3\). Consequently
\[
 (x-y)^2+(y-z)^2+(z-x)^2
 =2(x^2+y^2+z^2-xy-yz-zx)=0.
\]
Each square vanishes, so \(x=y=z=1\).
\end{solution}

\section{Elimination and bounds work together}

Elimination turns a system into fewer unknowns. Bounds keep the remaining
search finite or rule out roots that do not meet the original conditions.
If a system is symmetric, one may arrange the variables in increasing
order when looking for unordered solutions. A merely cyclic system need
not permit arbitrary reordering.

\begin{example}[A finite search obtained from an inequality]
Find all positive integer triples with \(x+y+z=xyz\).
\end{example}
\begin{solution}
By symmetry assume \(x\le y\le z\). Since \(x+y+z\le3z\),
we have \(xy\le3\). If \(x\ge2\), this is impossible.
For \(x=1\), the equation becomes
\((y-1)(z-1)=2\). The positive factor pairs give \(y=2,z=3\)
under the chosen order. All permutations of \((1,2,3)\) work.
\end{solution}

\pitfall Solving the eliminated equation produces candidates, not
automatically solutions. Check signs, denominators, ordering assumptions,
and the equations that were used to eliminate variables.

\section{Practice}
\begin{exercise}[{Foundation}]\label{a:ex08}
For which real \(t\) does \(x^2-2tx+t+2=0\) have two positive real
roots, allowing a repeated root?
\end{exercise}
\begin{exercise}[{Foundation}]\label{a:ex09}
Find the real pairs with \(x+y=5\), \(xy=6\).
\end{exercise}
\begin{exercise}[\problemlevel{Bridge}]\label{a:ex10}
Find all real triples with sum \(6\), pairwise-product sum \(11\),
and product \(6\).
\end{exercise}
\begin{exercise}[\problemlevel{Bridge}]\label{a:ex11}
Show that \(x+y+z=xy+yz+zx=0\) for real variables forces
\(x=y=z=0\).
\end{exercise}
\begin{exercise}[\problemlevel{Bridge}]\label{a:ex12}
If \(x,y>0\) and \(x+y+xy=8\), determine the largest possible \(xy\).
\end{exercise}
\begin{exercise}[\problemlevel{Challenge}]\label{a:ex13}
Classify all positive integer solutions of \(x+y+z=xyz\), explaining
why the search is complete.
\end{exercise}
\begin{exercise}[\problemlevel{Challenge}]\label{a:ex14}
Let \(a,b,c\) be nonzero real numbers with \(a+b+c=0\).
Prove that all the following denominators are nonzero and that
\[
 \frac1{a^2+b^2-c^2}+
 \frac1{b^2+c^2-a^2}+
 \frac1{c^2+a^2-b^2}=0.
\]
\end{exercise}

\chapter{Sequences, Floors, and Hidden Linear Structure}
\label{a:sequences}

\section{Read a recurrence as an instruction}

A sequence \((a_n)_{n\ge1}\) assigns a number \(a_n\) to each positive
integer \(n\). A recurrence tells us how to obtain later terms from earlier
ones. Its initial values are part of the data. A formula that satisfies
a recurrence but not its initial values is not its solution.

An arithmetic progression has \(a_{n+1}=a_n+d\) and hence
\(a_n=a_1+(n-1)d\). Pairing first and last terms gives
\[
 \sum_{k=1}^n a_k=\frac n2(a_1+a_n).
\]
A geometric progression has \(a_{n+1}=qa_n\). For \(q\ne1\),
subtracting \(qS\) from \(S=a_1+\cdots+a_1q^{n-1}\) yields
\[
 S=a_1\frac{1-q^n}{1-q}.
\]
For \(q=1\), the sum is \(na_1\). A harmonic progression is a sequence
of nonzero numbers whose reciprocals form an arithmetic progression.

\begin{example}[Shift the fixed point]
Suppose \(a_1=2\) and \(a_{n+1}=3a_n+4\). Find \(a_n\).
\end{example}
\begin{solution}
The constant value \(-2\) would remain fixed under \(x\mapsto3x+4\).
Measure distance from this value: \(b_n=a_n+2\).
Then \(b_{n+1}=3b_n\), \(b_1=4\), so
\(a_n=4\cdot3^{n-1}-2\).
\end{solution}

\intuition A nonlinear formula may be linear in the wrong coordinate.
When the recurrence repeatedly introduces \(1+a_n\), try reciprocals;
when it is a fractional linear expression, try a ratio or a fixed-point
substitution.

\begin{example}[Reciprocals remove the nonlinearity]
Let \(a_1=t>0\) and \(a_{n+1}=a_n/(1+a_n)\).
\end{example}
\begin{solution}
Every term is positive, so reciprocals are legal. We obtain
\[
 \frac1{a_{n+1}}=\frac1{a_n}+1,\qquad
 a_n=\frac{t}{1+(n-1)t}.
\]
The formula shows at once that \(a_n\) decreases to \(0\).
\end{solution}

\section{Second-order recurrences}

For \(a_{n+2}=pa_{n+1}+qa_n\), trial solutions \(a_n=r^n\) lead to
the \emph{characteristic equation} \(r^2-pr-q=0\).
When \(q\ne0\), its roots are nonzero. Distinct roots \(r,s\) yield
solutions \(Ar^n+Bs^n\); a repeated root \(r\) yields
\((A+Bn)r^n\). Substitution checks these forms, and the two initial
values determine \(A,B\). There cannot be another solution with the same
initial values because the recurrence uniquely determines every later
term. Complex roots may be used in conjugate pairs to describe real
sequences. If \(q=0\), solve the resulting first-order recurrence
directly rather than dividing by a zero root.

\begin{example}[A repeated root]
If \(a_0=0,a_1=1\) and \(a_{n+2}=2a_{n+1}-a_n\), find the sequence.
\end{example}
\begin{solution}
The difference \(a_{n+1}-a_n\) is constant, so \(a_n=n\).
The characteristic polynomial \((r-1)^2\) records the same fact.
\end{solution}

For a convergent real sequence \(a_n\) with recurrence
\(a_{n+1}=F(a_n)\), continuity of \(F\) at the limit implies
that its limit \(L\) satisfies \(L=F(L)\). This only gives candidates:
convergence needs proof. A useful elementary result is that a monotone,
bounded real sequence converges. Its proof uses the least-upper-bound
property of the real numbers: an increasing sequence approaches the
supremum of its terms.

\begin{example}[Prove convergence before solving the limit equation]
Let \(a_1>0\) and \(a_{n+1}=(a_n+2/a_n)/2\). Show that
\(a_n\) converges and find its limit.
\end{example}
\begin{solution}
AM--GM gives \(a_{n+1}\ge\sqrt2\), so every term from \(a_2\) onward
is at least \(\sqrt2\). For those terms,
\[
 a_n-a_{n+1}=\frac{a_n^2-2}{2a_n}\ge0.
\]
The tail is decreasing and bounded below, so it has a limit \(L\ge\sqrt2\).
The recurrence is continuous there, giving \(L=(L+2/L)/2\), hence
\(L=\sqrt2\).
Merely solving this fixed-point equation would not prove convergence;
the bounds and monotonicity are the missing logical steps.
\end{solution}

\section{Telescoping and finite differences}

Define the forward difference \(\Delta a_n=a_{n+1}-a_n\).
Summing differences gives \(a_{n+1}-a_1\). For products, seek a ratio
of consecutive terms instead.

\begin{example}[A telescoping product]
For \(n\ge1\), evaluate
\(\prod_{k=2}^{n}(1-1/k^2)\).
\end{example}
\begin{solution}
For \(n=1\) the product is empty and equals \(1\).
For \(n\ge2\), write each factor as \((k-1)(k+1)/k^2\):
\[
 \prod_{k=2}^{n}\frac{k-1}{k}\frac{k+1}{k}
 =\frac1n\frac{n+1}{2}=\frac{n+1}{2n}.
\]
The same final formula also gives \(1\) when \(n=1\).
\end{solution}

\section{Floors separate integer and fractional information}

The floor \(\lfloor x\rfloor\) is the greatest integer not exceeding
\(x\); the ceiling \(\lceil x\rceil\) is the least integer not below it.
The fractional part is \(\{x\}=x-\lfloor x\rfloor\), so \(0\le\{x\}<1\),
including when \(x\) is negative. For example,
\(\lfloor-1.2\rfloor=-2\), \(\{-1.2\}=0.8\).

Write \(x=m+u,y=n+v\) with \(m,n\) integers and \(u,v\in[0,1)\).
Then \(\lfloor x+y\rfloor\) is either \(m+n\) or \(m+n+1\),
according as \(u+v<1\) or \(u+v\ge1\). This is a carry in addition.

\begin{example}[A floor identity]
For real \(x\), prove
\(\lfloor x\rfloor+\lfloor x+1/2\rfloor=\lfloor2x\rfloor\).
\end{example}
\begin{solution}
Write \(x=m+u\), \(0\le u<1\). Both sides contain \(2m\).
Their remaining parts are \(0\) for \(u<1/2\), and \(1\) otherwise.
\end{solution}

\pitfall The floor is not linear. Replace an unproved rule by the
inequality \(\lfloor x\rfloor\le x<\lfloor x\rfloor+1\), or work with
the fractional part explicitly.

\extension A striking application of floor counting is the following
partition theorem. If irrational \(r,s>1\) satisfy \(1/r+1/s=1\), then
the two lists
\[
 \lfloor r\rfloor,\lfloor2r\rfloor,\ldots,
 \qquad \lfloor s\rfloor,\lfloor2s\rfloor,\ldots
\]
contain each positive integer exactly once between them.
For a nonnegative integer \(m\), the number of terms in the first list
that are at most \(m\) is \(\lfloor(m+1)/r\rfloor\):
the condition is \(nr<m+1\), and the endpoint is not an integer.
The analogous count for \(s\) is \(\lfloor(m+1)/s\rfloor\).
Their unfloored arguments are nonintegers summing to \(m+1\), so the
two counts sum to \(m\).
Subtract this cumulative count at \(m-1\) from the count at \(m\):
exactly one occurrence of \(m\) appears. Each list is strictly increasing
because its spacing exceeds \(1\). Thus there are neither repetitions
between the lists nor missing positive integers.
This is Beatty's theorem. The proof illustrates a useful habit:
count how many objects lie below each threshold rather than trying
to pair their positions directly.

\section{Practice}
\begin{exercise}[\problemlevel{Foundation}]\label{a:ex15}
If \(a_1=2\), \(a_{n+1}=a_n+3\), find \(a_n\) and
\(\sum_{k=1}^n a_k\).
\end{exercise}
\begin{exercise}[\problemlevel{Foundation}]\label{a:ex16}
Solve \(a_1=1,\ a_{n+1}=2a_n+1\).
\end{exercise}
\begin{exercise}[\problemlevel{Bridge}]\label{a:ex17}
Evaluate \(\sum_{k=1}^n1/[k(k+1)(k+2)]\).
\end{exercise}
\begin{exercise}[\problemlevel{Bridge}]\label{a:ex18}
Let \(a_1=2\) and \(a_{n+1}=1+1/a_n\).
Express \(a_n\) using the Fibonacci numbers \(F_1=F_2=1\),
\(F_{n+2}=F_{n+1}+F_n\).
\end{exercise}
\begin{exercise}[\problemlevel{Challenge}]\label{a:ex19}
For every real \(t\) and positive integer \(n\), prove
\[
 \sum_{k=0}^{n-1}\floor{\frac{t+k}{n}}=\lfloor t\rfloor.
\]
\end{exercise}
\begin{exercise}[\problemlevel{Bridge}]\label{a:ex20}
For \(a_1=t>0,\ a_{n+1}=a_n/(1+a_n)\), find \(a_n\) and
\(\sum_{k=1}^{n}a_ka_{k+1}\).
\end{exercise}
\begin{exercise}[\problemlevel{Challenge}]\label{a:ex21}
Let \(a_1,a_2>0\) and \(a_{n+2}=(1+a_{n+1})/a_n\).
Prove \(a_{n+5}=a_n\) for all \(n\ge1\).
\end{exercise}

\chapter{Polynomials as Rigid Objects}
\label{a:polynomials}

\section{Roots and division}

A real polynomial is an expression \(P(x)=a_dx^d+\cdots+a_0\).
If \(a_d\ne0\), its degree is \(d\), its leading coefficient is \(a_d\),
and it is \emph{monic} when \(a_d=1\).
The zero polynomial has no finite degree. Polynomial identities are
coefficient-by-coefficient equalities.

\intuition A polynomial can bend only a limited number of times.
Its degree measures a precise form of rigidity: enough roots force
the entire polynomial to vanish.

\begin{theorem}[Factor theorem and root bound]
For a polynomial \(P\), \(P(a)=0\) if and only if \(x-a\) divides \(P(x)\).
A nonzero polynomial of degree \(d\) has at most \(d\) distinct roots.
\end{theorem}
\begin{proof}
The identity \(x^k-a^k=(x-a)(x^{k-1}+ax^{k-2}+\cdots+a^{k-1})\)
shows that \(P(x)-P(a)\) is divisible by \(x-a\).
This proves the factor theorem. Repeatedly remove factors for distinct
roots; a degree \(d\) polynomial cannot have more than \(d\) such factors.
\end{proof}

\begin{example}[Values determine a polynomial]
A polynomial \(P\) has degree at most \(2\), and
\(P(0)=1,P(1)=2,P(2)=5\). Find \(P\).
\end{example}
\begin{solution}
The pattern suggests \(Q(x)=x^2+1\). The polynomial \(P-Q\) has degree
at most \(2\) and three distinct roots, so it is zero. The guess became
a proof because of the root bound.
\end{solution}

\section{Vieta and power sums}

If a monic polynomial factors as \(\prod_{j=1}^d(x-r_j)\), expanding
shows that its coefficient of \(x^{d-k}\) is \((-1)^k\) times the sum
of all products of \(k\) different roots, counting multiplicities.
These are Vieta's relations.
Every complex polynomial of positive degree has a complex root and
therefore factors completely over \(\CC\). This is the fundamental
theorem of algebra; its proof belongs to a later course in analysis.
The elementary arguments here that specify known roots do not require it.

\begin{example}[Do not calculate roots you do not need]
Let \(\alpha,\beta\) be the roots of \(x^2-3x+1=0\).
Find \(\alpha^4+\beta^4\).
\end{example}
\begin{solution}
Vieta gives \(\alpha+\beta=3,\alpha\beta=1\).
Thus \(\alpha^2+\beta^2=7\), and
\(\alpha^4+\beta^4=7^2-2=47\).
Alternatively, \(\alpha^k+\beta^k\) satisfies
\(s_{k+2}=3s_{k+1}-s_k\), with \(s_0=2,s_1=3\).
\end{solution}

More generally, each root of
\(x^d+c_1x^{d-1}+\cdots+c_d\) satisfies its equation.
Multiplying by powers of the root and summing gives recurrences for
power sums. For the first few sums one obtains, with \(s_k=\sum r_j^k\),
\[
 s_1=-c_1,\quad s_2+c_1s_1+2c_2=0,\quad
 s_3+c_1s_2+c_2s_1+3c_3=0
\]
when the displayed coefficients exist. These are the first Newton
identities; they avoid solving the individual roots.

\section{Interpolation explains both existence and uniqueness}

Given distinct real numbers \(t_0,\ldots,t_d\) and values \(y_0,\ldots,y_d\),
define
\[
 L_j(x)=\prod_{\substack{0\le k\le d\\k\ne j}}
             \frac{x-t_k}{t_j-t_k}.
\]
The denominators are nonzero. Since \(L_j(t_i)\) is \(1\) for \(i=j\)
and \(0\) otherwise, the polynomial
\[
 P(x)=\sum_{j=0}^d y_jL_j(x)
\]
has the desired values and degree at most \(d\).
Any two such polynomials differ by a degree-at-most-\(d\) polynomial
with \(d+1\) roots, proving uniqueness.

\begin{example}[Read a coefficient instead of the whole formula]
For distinct \(t_0,\ldots,t_d\), prove
\[
 \sum_{j=0}^d\frac{t_j^m}
 {\prod_{k\ne j}(t_j-t_k)}
 =
 \begin{cases}0,&0\le m<d,\\1,&m=d.\end{cases}
\]
\end{example}
\begin{solution}
Interpolate \(x^m\) at the given nodes and compare the coefficient of
\(x^d\). Each basis polynomial has leading coefficient
\(1/\prod_{k\ne j}(t_j-t_k)\). The left side is therefore the leading
coefficient of \(x^m\), which is exactly the stated value.
\end{solution}

\section{Finite differences lower degree}

For a polynomial, define \(\Delta P(x)=P(x+1)-P(x)\).
If \(P\) has degree \(d\ge1\) and leading coefficient \(a\), the binomial
theorem shows that \(\Delta P\) has degree \(d-1\) and leading coefficient
\(da\). Repeating differences eventually leaves a constant.

\begin{example}[An equation about all real inputs]
Find all real polynomials satisfying \(P(x+1)-P(x)=2x+1\).
\end{example}
\begin{solution}
The polynomial \(x^2\) has exactly this difference.
Subtract it: \(Q=P-x^2\) satisfies \(\Delta Q=0\).
A nonconstant polynomial would have a nonzero difference of one smaller
degree. Therefore \(Q\) is constant, and \(P=x^2+C\).
\end{solution}

\recognition An identity at infinitely many distinct inputs is a
polynomial identity. An equation involving \(P(x+1)-P(x)\) suggests
degree reduction; composition suggests comparing degrees and leading
coefficients before expanding.

\section{Practice}
\begin{exercise}[\problemlevel{Foundation}]\label{a:ex22}
Find the polynomial of degree at most \(3\) with values
\(1,2,5,10\) at \(0,1,2,3\), respectively.
\end{exercise}
\begin{exercise}[\problemlevel{Foundation}]\label{a:ex23}
For roots \(\alpha,\beta\) of \(x^2-3x+1=0\), compute
\(\alpha^3+\beta^3\).
\end{exercise}
\begin{exercise}[\problemlevel{Bridge}]\label{a:ex24}
A monic quartic \(P\) vanishes at \(1,2,3\) and satisfies \(P(4)=24\).
Determine \(P\).
\end{exercise}
\begin{exercise}[\problemlevel{Bridge}]\label{a:ex25}
Determine all real polynomials with \(P(x+1)-P(x)=3x^2+3x+1\).
\end{exercise}
\begin{exercise}[\problemlevel{Challenge}]\label{a:ex26}
Find necessary and sufficient conditions on real \(a,b\) for
\(x^4+ax^2+b\ge0\) for every real \(x\).
\end{exercise}
\begin{exercise}[\problemlevel{Challenge}]\label{a:ex27}
Determine all real polynomials with \(P(x^2)=P(x)^2\).
\end{exercise}
\begin{exercise}[\problemlevel{Challenge}]\label{a:ex28}
Let a nonconstant real polynomial \(P\) have only real roots, counting
multiplicity. Prove
\[
 P'(x)^2\ge P(x)P''(x)\qquad(x\in\RR).
\]
Here the formal derivative is defined by
\((\sum a_kx^k)'=\sum_{k\ge1}ka_kx^{k-1}\).
\end{exercise}

\chapter{Integer Polynomials and Several Variables}
\label{a:integer-polynomials}

\section{Integer coefficients leave arithmetic traces}

\(\ZZ[x]\) denotes the set of polynomials whose coefficients are integers.
If \(P\in\ZZ[x]\) and \(m,n\) are integers, then
\[
 m-n\mid P(m)-P(n)\qquad(m\ne n).
\]
Indeed each difference \(m^k-n^k\) has the factor \(m-n\).
This divisibility is stronger than a numerical coincidence: it holds
for every pair of inputs. It frequently prevents interpolation data
from coming from integer coefficients.

\begin{example}[Values that cannot be interpolated over the integers]
Does a polynomial \(P\in\ZZ[x]\) exist with \(P(0)=0,P(2)=1\)?
\end{example}
\begin{solution}
If it existed, \(2\) would divide \(P(2)-P(0)=1\), a contradiction.
Over the rationals, \(P(x)=x/2\) works. The coefficient domain matters.
\end{solution}

\section{Integer-valued does not mean integer-coefficient}

Define
\[
 \binom{x}{k}=\frac{x(x-1)\cdots(x-k+1)}{k!}\quad(k\ge1),
 \qquad \binom{x}{0}=1,
\]
where \(k!=1\cdot2\cdots k\). This is a polynomial in \(x\).
For nonnegative integers \(x\), it counts \(k\)-element subsets
and is an integer. For a negative integer \(-m\), \(m\ge1\), the same
formula gives
\(\binom{-m}{k}=(-1)^k\binom{m+k-1}{k}\), also an integer.

\begin{theorem}[Newton's finite-difference expansion]
For a polynomial \(P\) of degree at most \(d\),
\[
 P(x)=\sum_{k=0}^d \Delta^kP(0)\binom{x}{k}.
\]
Consequently \(P\in\QQ[x]\) takes integer values on all integers if and
only if all the coefficients \(\Delta^kP(0)\) in this expansion are integers.
\end{theorem}
\begin{proof}
The polynomial identity
\(\Delta\binom{x}{k}=\binom{x}{k-1}\) follows either by multiplying out
or from Pascal's identity at infinitely many nonnegative integers.
The expansion can now be built by induction on degree: expand
\(\Delta P\) in this basis, integrate one difference by shifting each
basis index up by one, then choose the constant so that the value at
\(0\) is \(P(0)\). The coefficients are exactly the successive
differences at zero. Integer values give integer differences;
conversely integer coefficients in this basis give integer values.
\end{proof}

\begin{example}[A rational polynomial with integer values]
Why is \((n^3-n)/6\) an integer for every integer \(n\)?
\end{example}
\begin{solution}
Its numerator is the product of three consecutive integers
\((n-1)n(n+1)\). At least one is divisible by \(3\), and at least
one is even, so the product is divisible by \(6\).
Equivalently,
\((x^3-x)/6=\binom{x}{3}+\binom{x}{2}\).
\end{solution}

\section{Irreducibility is always relative to a coefficient domain}

A nonconstant polynomial is \emph{irreducible over \(\QQ\)} if it cannot
be written as a product of two positive-degree rational polynomials.
The polynomial \(x^2-2\) is irreducible over \(\QQ\), but factors over
\(\RR\). A quadratic or cubic over a field is reducible precisely when
it has a root in that field: any positive-degree factorization has a
linear factor. For higher degrees, absence of rational roots is not
enough; \(x^4+4\) has no rational roots but factors.

\begin{theorem}[Rational root test]
If \(a_dx^d+\cdots+a_0\in\ZZ[x]\) has root \(u/v\), where
\(\gcd(u,v)=1\) and \(v>0\), then \(u\mid a_0\) and \(v\mid a_d\).
\end{theorem}
\begin{proof}
Multiply the root equation by \(v^d\).
Reducing modulo \(u\) shows \(u\mid a_0v^d\), and coprimality gives
\(u\mid a_0\). Reducing modulo \(v\) gives the other assertion.
When \(u=0\), the conclusion \(u\mid a_0\) is interpreted as \(a_0=0\);
one may simply handle that root first.
\end{proof}

An integer polynomial is \emph{primitive} when the greatest common
divisor of its coefficients is \(1\).
The product of primitive integer polynomials is primitive: if a prime
divided every coefficient of the product, reduction modulo that prime
would give a product of two nonzero polynomials equal to zero, which
cannot happen over a field. Clearing denominators and extracting the
greatest common divisors therefore proves \emph{Gauss's lemma}: a
primitive integer polynomial reducible over \(\QQ\) factors into
positive-degree integer polynomials.

\begin{theorem}[Eisenstein's criterion]
Let \(P=a_dx^d+\cdots+a_0\) be primitive with \(d\ge1\).
If a prime \(\ell\) divides
every coefficient except \(a_d\), does not divide \(a_d\), and
\(\ell^2\nmid a_0\), then \(P\) is irreducible over \(\QQ\).
\end{theorem}
\begin{proof}
An integer factorization \(P=GH\) would reduce modulo \(\ell\) to
a factorization of \(a_dx^d\). Since leading coefficients remain
nonzero, the reductions of both positive-degree factors are monomials
of positive degree. Thus their constant coefficients are both divisible
by \(\ell\), forcing \(\ell^2\mid a_0\), a contradiction.
\end{proof}

\begin{example}[Shift the polynomial before testing]
For a prime \(\ell\), prove \(1+x+\cdots+x^{\ell-1}\) irreducible over
\(\QQ\).
\end{example}
\begin{solution}
Substitute \(x=y+1\):
\[
 1+(y+1)+\cdots+(y+1)^{\ell-1}
 =\frac{(y+1)^\ell-1}{y}
 =\sum_{j=1}^{\ell}\binom{\ell}{j}y^{j-1}.
\]
All coefficients except the leading one are divisible by \(\ell\);
the constant coefficient is \(\ell\), not divisible by \(\ell^2\).
Eisenstein applies. A factorization before the substitution would give
one after it, so the original polynomial is irreducible.
\end{solution}

Reduction modulo a prime gives another test: if a primitive polynomial
keeps its degree and becomes irreducible over the finite field
\(\FF_\ell=\ZZ/\ell\ZZ\), then it was irreducible over \(\QQ\).
An integer factorization would otherwise survive reduction with both
degrees unchanged. The converse fails; a reducible reduction may lose
arithmetic information.

\begin{example}[Reduction modulo a prime needs a factor check]
Prove \(x^4+x+1\) irreducible over \(\QQ\).
\end{example}
\begin{solution}
Reduce modulo \(2\). Neither \(0\) nor \(1\) is a root.
If a quartic over \(\FF_2\) with no linear factor were reducible,
it would split into two irreducible quadratics.
The only monic irreducible quadratic over \(\FF_2\) is
\(x^2+x+1\): a monic quadratic with constant term \(0\) has root \(0\),
and among the two with constant term \(1\), \(x^2+1\) has root \(1\).
Modulo \(x^2+x+1\), we have \(x^2=x+1,x^3=1,x^4=x\);
the quartic therefore has remainder \(1\).
It has no quadratic factor either, so it is irreducible over \(\FF_2\).
Its degree is preserved and it is primitive, so the modular criterion
proves the rational irreducibility.
\end{solution}

\section{Several variables: treat the others as coefficients}

For \(P(x,y)\), the total degree is the largest sum of exponents in a
nonzero monomial; the degree in \(x\) ignores the exponent of \(y\).
If \(P(y,y)=0\) as a polynomial in \(y\), then \(x-y\) divides \(P(x,y)\):
divide in the variable \(x\), with polynomials in \(y\) as coefficients.
The remainder is exactly \(P(y,y)\).

\begin{example}[Alternation predicts the factors]
Factor \(x^2(y-z)+y^2(z-x)+z^2(x-y)\).
\end{example}
\begin{solution}
The expression vanishes whenever two variables coincide, so it has
factors \(x-y,x-z,y-z\). These are distinct linear factors and together
have total degree \(3\), the degree of the expression. Thus their
product agrees up to a constant. Comparing the coefficient of \(x^2y\)
gives
\[
 x^2(y-z)+y^2(z-x)+z^2(x-y)=(x-y)(x-z)(y-z).
\]
\end{solution}

Every symmetric polynomial in three variables can be expressed using
\(p=x+y+z,q=xy+yz+zx,r=xyz\).
One constructive proof orders monomials lexicographically. A symmetric
leading monomial can be taken as \(x^\alpha y^\beta z^\gamma\) with
\(\alpha\ge\beta\ge\gamma\); the polynomial
\(p^{\alpha-\beta}q^{\beta-\gamma}r^\gamma\) has that leading monomial.
Subtract its multiple to remove the leading term and repeat. Among the
finitely many monomials of bounded degree, this process terminates.
This explains why \(p,q,r\) are useful coordinates rather than a lucky trick.

\section{Practice}
\begin{exercise}[\problemlevel{Foundation}]\label{a:ex29}
Prove \(m-n\mid P(m)-P(n)\) for \(P\in\ZZ[x]\) and distinct integers
\(m,n\), without using polynomial division.
\end{exercise}
\begin{exercise}[\problemlevel{Foundation}]\label{a:ex30}
Explain why no polynomial in \(\ZZ[x]\) takes values \(0,1\) at
\(0,2\), respectively, although a polynomial in \(\QQ[x]\) does.
\end{exercise}
\begin{exercise}[\problemlevel{Bridge}]\label{a:ex31}
Give a degree-two integer-valued polynomial with a noninteger
coefficient, and prove that it takes integer values at negative inputs too.
\end{exercise}
\begin{exercise}[\problemlevel{Bridge}]\label{a:ex32}
Prove \(x^4+2\) irreducible over \(\QQ\). Explain why merely checking
its rational roots would not prove this.
\end{exercise}
\begin{exercise}[\problemlevel{Bridge}]\label{a:ex33}
Find all rational roots of \(2x^3-3x^2-8x+12\).
\end{exercise}
\begin{exercise}[\problemlevel{Challenge}]\label{a:ex34}
Let \(P\in\ZZ[x]\) satisfy \(P(0)=P(1)=P(2)=1\).
Prove that \(P(n)\ne2\) for every integer \(n\).
\end{exercise}
\begin{exercise}[\problemlevel{Challenge}]\label{a:ex35}
Let \(P\in\ZZ[x]\) be nonconstant. Prove that infinitely many different
primes divide its nonzero integer values.
\end{exercise}

\chapter{Complex Numbers and Trigonometric Structure}
\label{a:complex}

\section{A number can record a rotation}

A complex number is \(z=x+iy\), where \(x,y\in\RR\) and \(i^2=-1\).
Its conjugate is \(\bar z=x-iy\), and its modulus is
\(|z|=\sqrt{x^2+y^2}\). Since \(z\bar z=|z|^2\), a nonzero number has
reciprocal \(\bar z/|z|^2\). Multiplication preserves lengths
multiplicatively: \(|zw|=|z||w|\).

For an oriented angle \(\theta\), define \(\cos\theta,\sin\theta\) to
be the horizontal and vertical coordinates of the corresponding
point on the unit circle. Thus the point has complex coordinate
\(\cos\theta+i\sin\theta\). Angles here are in radians, so a full
turn is \(2\pi\). Multiplying two such expressions and using the
angle-addition identities gives
\[
 (\cos\alpha+i\sin\alpha)(\cos\beta+i\sin\beta)
 =\cos(\alpha+\beta)+i\sin(\alpha+\beta).
\]
Thus multiplication adds angles. We sometimes abbreviate
\(\cos\theta+i\sin\theta\) as \(e^{i\theta}\); no properties of the
complex exponential beyond this identity are needed here.

\begin{figure}[htbp]
\centering
\begin{tikzpicture}[scale=1.55]
\draw[->](-1.35,0)--(1.55,0)node[right]{real};
\draw[->](0,-1.2)--(0,1.35)node[above]{imaginary};
\draw[gray](0,0)circle(1);
\coordinate(A)at(35:1);
\coordinate(B)at(105:1);
\draw[main,->](0,0)--(A)node[right]{$z$};
\draw[main,->](0,0)--(B)node[above left]{$wz$};
\draw[->,accent](35:.5)arc(35:105:.5);
\node at(72:.72){$\alpha$};
\node[below left]at(0,0){$0$};
\end{tikzpicture}
\caption{Multiplication by \(w=e^{i\alpha}\) rotates the vector for
\(z\) through \(\alpha\); if \(|w|\ne1\), it also scales the length.}
\end{figure}

\begin{example}[Powers without repeated expansion]
Compute \((1+i)^8\).
\end{example}
\begin{solution}
The number \(1+i\) has modulus \(\sqrt2\) and angle \(\pi/4\).
Its eighth power has modulus \(16\) and angle \(2\pi\), hence is \(16\).
One can also square three times: \((1+i)^2=2i\),
\((1+i)^4=-4\), \((1+i)^8=16\).
\end{solution}

\section{Roots of unity turn a regular polygon into algebra}

The \(n\)th roots of unity are the solutions of \(z^n=1\):
\[
 1,\zeta,\ldots,\zeta^{n-1},\qquad
 \zeta=\cos(2\pi/n)+i\sin(2\pi/n).
\]
They are distinct, and the degree bound proves that this is the full list.
For \(n\ge2\), the geometric-sum identity gives
\(\sum_{k=0}^{n-1}\zeta^k=0\).
Geometrically the vectors from the center to the vertices of a regular
polygon balance.

\begin{example}[A product of distances]
For \(n\ge2\), prove
\(\prod_{k=1}^{n-1}|1-\zeta^k|=n\).
\end{example}
\begin{solution}
The factor theorem and leading coefficients give
\[
 \frac{x^n-1}{x-1}=\prod_{k=1}^{n-1}(x-\zeta^k).
\]
The left side is \(1+x+\cdots+x^{n-1}\).
Set \(x=1\) and then take absolute values to obtain the result.
\end{solution}

\section{Trigonometry: choose identities that match the constraint}

The angle-addition formulas imply
\[
 \sin2t=2\sin t\cos t,\qquad
 \cos2t=2\cos^2t-1,\qquad
 \sin^2t+\cos^2t=1.
\]
For \(A+B+C=\pi\), the relation between the angles is information:
replace \(C\) by \(\pi-A-B\) when it simplifies the desired expression.
Do not treat \(A,B,C\) as independent.

\begin{example}[The constraint creates a factorization]
If \(A+B+C=\pi\), prove
\(\sin2A+\sin2B+\sin2C=4\sin A\sin B\sin C\).
\end{example}
\begin{solution}
Using sum-to-product,
\[
\begin{aligned}
\sin2A+\sin2B+\sin2C
 &=2\sin(A+B)\cos(A-B)-2\sin(A+B)\cos(A+B)\\
 &=2\sin C\,[2\sin A\sin B].
\end{aligned}
\]
The last step is the cosine-difference identity.
\end{solution}

\pitfall Tangent is \(\tan t=\sin t/\cos t\), so any formula involving
it requires \(\cos t\ne0\). Cross-multiplying can produce a valid
identity at excluded angles, but that does not define the original
tangent expression there.

\section{Chebyshev polynomials encode repeated angles}

Define \(T_0(x)=1,T_1(x)=x\), and
\(T_{n+1}(x)=2xT_n(x)-T_{n-1}(x)\).
The cosine addition formula shows inductively that
\[
 T_n(\cos\theta)=\cos(n\theta).
\]
Thus a trigonometric identity produces a polynomial of degree \(n\),
with leading coefficient \(2^{n-1}\) for \(n\ge1\).
Its \(n\) distinct roots are
\(\cos((2j-1)\pi/(2n))\), \(1\le j\le n\).

\begin{example}[A polynomial bound with a trigonometric explanation]
Show \(|4x^3-3x|\le1\) for \(-1\le x\le1\).
\end{example}
\begin{solution}
Write \(x=\cos\theta\), possible for every \(x\in[-1,1]\).
The expression is \(T_3(x)=\cos3\theta\), whose absolute value is at
most \(1\). The substitution fits the interval exactly; it would not
justify a claim for all real \(x\).
\end{solution}

\section{Practice}
\begin{exercise}[\problemlevel{Foundation}]\label{a:ex36}
Compute \((1+i)^8\) in two ways.
\end{exercise}
\begin{exercise}[\problemlevel{Bridge}]\label{a:ex37}
If \(\zeta^n=1\) and \(\zeta\ne1\), show
\(\sum_{k=0}^{n-1}k\zeta^k=n/(\zeta-1)\).
\end{exercise}
\begin{exercise}[\problemlevel{Bridge}]\label{a:ex38}
Prove the triple-angle sum identity
\(\sin2A+\sin2B+\sin2C=4\sin A\sin B\sin C\)
under \(A+B+C=\pi\).
\end{exercise}
\begin{exercise}[\problemlevel{Bridge}]\label{a:ex39}
Factor \(x^4+x^2+1\) over \(\RR\), and describe its complex roots
as roots of unity.
\end{exercise}
\begin{exercise}[\problemlevel{Challenge}]\label{a:ex40}
For \(n\ge3\), prove
\[
 \sum_{k=0}^{n-1}\cos(\theta+2\pi k/n)=0,\qquad
 \sum_{k=0}^{n-1}\cos^2(\theta+2\pi k/n)=\frac n2.
\]
Why is the restriction \(n\ge3\) needed for the second formula?
\end{exercise}
\begin{exercise}[\problemlevel{Challenge}]\label{a:ex41}
For \(n\ge2\), show
\(\prod_{k=1}^{n-1}\sin(k\pi/n)=n/2^{n-1}\).
\end{exercise}
\begin{exercise}[\problemlevel{Challenge}]\label{a:ex42}
For \(n\ge2\), prove
\[
 \sum_{k=1}^{n-1}\frac1{\sin^2(k\pi/n)}=\frac{n^2-1}{3}.
\]
You may differentiate polynomial identities using the formal derivative.
\end{exercise}

\chapter{Inequalities: From Squares to Strategy}
\label{a:inequalities}

\section{An inequality is an identity with a sign}

\intuition To prove \(A\ge B\), try to understand the excess \(A-B\).
If it becomes a sum of nonnegative terms, both the inequality and its
equality cases become visible. A remembered theorem is useful only if
its direction and equality pattern fit the problem.

The starting point is \((x-y)^2\ge0\), equivalent to
\(x^2+y^2\ge2xy\). For \(x,y\ge0\), substitute their square roots:
\((x+y)/2\ge\sqrt{xy}\). Equality holds precisely when \(x=y\).

\begin{theorem}[Arithmetic--geometric mean inequality]
For nonnegative \(x_1,\ldots,x_n\),
\[
 \frac{x_1+\cdots+x_n}{n}\ge(x_1\cdots x_n)^{1/n},
\]
with equality exactly when all \(x_i\) are equal.
\end{theorem}
\begin{proof}
The case \(n=2\) was proved above. If the result holds for \(n\), apply
it to each half of a list of \(2n\) terms, then apply the two-variable
case to their averages. It follows for all powers of two.
If it holds for \(n+1\), append the average
\(A=(x_1+\cdots+x_n)/n\) to \(n\) given terms. It gives
\((x_1\cdots x_n A)^{1/(n+1)}\le A\).
For \(A>0\), cancel \(A\) to obtain the \(n\)-term result; for \(A=0\)
all terms vanish. Descending from a sufficiently large power of two
proves every \(n\). The same steps show that equality requires all
terms equal; the all-zero case is included.
\end{proof}

Weighted AM--GM says that for positive \(x_i\) and real weights
\(\lambda_i\ge0\) summing to \(1\),
\[
 \sum_i\lambda_i x_i\ge\prod_i x_i^{\lambda_i}.
\]
For rational weights, repeat each \(x_i\) the required number of times
in ordinary AM--GM. Approximate real weights by rational ones and use
continuity of positive real powers. Equality requires all variables
with positive weights to be equal.

\begin{example}[Choose terms with the desired product]
For \(x>0\), find the minimum of \(x^2+2/x\).
\end{example}
\begin{solution}
Apply AM--GM to \(x^2,1/x,1/x\). Their product is \(1\), so the
sum is at least \(3\), attained exactly at \(x=1\).
The split into three terms is chosen to cancel the powers of \(x\).
\end{solution}

\section{Cauchy: align the numerator and denominator}

\begin{theorem}[Cauchy--Schwarz]
For real lists \(u_1,\ldots,u_n\), \(v_1,\ldots,v_n\),
\[
 \left(\sum u_iv_i\right)^2\le
 \left(\sum u_i^2\right)\left(\sum v_i^2\right).
\]
Equality holds exactly when the two lists are linearly dependent:
there exist real \(\alpha,\beta\), not both zero, with
\(\alpha u_i+\beta v_i=0\) for every \(i\).
\end{theorem}
\begin{proof}
Expanding gives the identity
\[
 \left(\sum u_i^2\right)\left(\sum v_i^2\right)
 -\left(\sum u_iv_i\right)^2
 =\sum_{i<j}(u_iv_j-u_jv_i)^2.
\]
If one list is zero, equality and the stated dependence both hold.
Otherwise choose one nonzero coordinate. Vanishing of every displayed
square makes the lists proportional.
\end{proof}

Taking \(u_i=a_i/\sqrt{b_i}\), \(v_i=\sqrt{b_i}\) for \(b_i>0\) gives
the useful fraction form
\[
 \sum_i\frac{a_i^2}{b_i}\ge
 \frac{(\sum_i a_i)^2}{\sum_i b_i}.
\]
Equality holds exactly when \(a_i/b_i\) is constant.
The \(a_i\) may be negative, but every \(b_i\) must be positive.

\begin{example}[Nesbitt's inequality, with the decision explained]
For \(a,b,c>0\), prove
\(\frac a{b+c}+\frac b{c+a}+\frac c{a+b}\ge3/2\).
\end{example}
\begin{solution}
To apply the fraction form, make each numerator a square:
\[
 \sum_{\cyc}\frac a{b+c}
 =\sum_{\cyc}\frac{a^2}{a(b+c)}
 \ge\frac{(a+b+c)^2}{2(ab+bc+ca)}\ge\frac32.
\]
The last inequality follows from
\(a^2+b^2+c^2\ge ab+bc+ca\), itself a sum-of-squares identity.
Equality throughout holds at \(a=b=c\), and the last inequality shows
this is necessary. Here \(\sum_{\cyc}\) means the three cyclic choices
\((a,b,c),(b,c,a),(c,a,b)\).
\end{solution}

\section{Ordering information: rearrangement and Chebyshev}

Suppose \(a_1\le\cdots\le a_n\) and \(b_1\le\cdots\le b_n\).
The sum \(\sum a_i b_i\) is largest when the orders agree and smallest
when they are reversed. To see this, swapping an inverted pair changes
the sum by
\[
 (a_j-a_i)(b_j-b_i)\ge0\qquad(i<j).
\]
Repeatedly remove inversions. This is the \emph{rearrangement inequality}.
It uses ordering, not positivity.

Chebyshev's inequality for similarly ordered real lists is
\[
 \frac1n\sum a_ib_i\ge
 \left(\frac1n\sum a_i\right)\left(\frac1n\sum b_i\right).
\]
Indeed
\[
 n\sum a_ib_i-\left(\sum a_i\right)\left(\sum b_i\right)
 =\sum_{i<j}(a_i-a_j)(b_i-b_j)\ge0.
\]
The direction reverses for oppositely ordered lists.

\begin{example}[A comparison using the right ordering]
For nonnegative \(a,b,c\), prove
\((a+b+c)(a^2+b^2+c^2)\le3(a^3+b^3+c^3)\).
\end{example}
\begin{solution}
The lists \(a,b,c\) and \(a^2,b^2,c^2\) have the same order, so
Chebyshev gives exactly the desired inequality.
The same claim for arbitrary real variables cannot use this argument:
squaring does not preserve their order.
\end{solution}

\section{Normalize before choosing a theorem}

An expression is homogeneous of degree \(d\) if replacing every variable
by \(t\) times itself multiplies it by \(t^d\). In a homogeneous
inequality, a common positive scale often carries no information.
If \(a,b,c>0\), one may divide by \(a+b+c\) to impose sum \(1\),
or divide by \((abc)^{1/3}\) to impose product \(1\).
The normalization must preserve the comparison being proved.

\begin{example}[A product constraint controls a sum]
If \(a,b,c>0\), \(abc=1\), prove
\(a^2+b^2+c^2\ge a+b+c\).
\end{example}
\begin{solution}
Let \(s=a+b+c\). AM--GM gives \(s\ge3\), while Cauchy gives
\(a^2+b^2+c^2\ge s^2/3\ge s\).
Equality holds exactly when \(a=b=c=1\).
The bridge \(s\ge3\) is what makes the two standard inequalities fit.
\end{solution}

\strategy First predict equality. Then check scale and signs.
If denominators are sums, try Cauchy with deliberately chosen squares.
If a product is fixed, try AM--GM. If the variables are ordered, use
that information rather than discarding it. Before adding inequalities,
check that every intermediate direction points toward the target.

\section{Practice}
\begin{exercise}[\problemlevel{Foundation}]\label{a:ex43}
Prove Nesbitt's inequality and determine equality.
\end{exercise}
\begin{exercise}[\problemlevel{Foundation}]\label{a:ex44}
For real \(x+y+z=1\), find the minimum of \(x^2+y^2+z^2\).
\end{exercise}
\begin{exercise}[\problemlevel{Bridge}]\label{a:ex45}
Pair \(1,2,3\) with \(4,5,6\) in all possible ways. Find the largest
and smallest sum of paired products, using a proof that does not
list all permutations.
\end{exercise}
\begin{exercise}[\problemlevel{Bridge}]\label{a:ex46}
If \(a,b,c>0\), \(ab+bc+ca=3\), show
\(a^2+b^2+c^2\ge3\) and \(abc\le1\).
\end{exercise}
\begin{exercise}[\problemlevel{Bridge}]\label{a:ex47}
For nonnegative \(a,b,c\), prove
\(a^2b^2+b^2c^2+c^2a^2\ge abc(a+b+c)\).
\end{exercise}
\begin{exercise}[\problemlevel{Bridge}]\label{a:ex48}
For positive \(a+b+c=1\), show
\[
 \frac{a^2}{1-a}+\frac{b^2}{1-b}+\frac{c^2}{1-c}\ge\frac12.
\]
\end{exercise}
\begin{exercise}[\problemlevel{Challenge}]\label{a:ex49}
Let \(a,b,c>0\), \(abc=1\). Prove
\(a^2+b^2+c^2\ge a+b+c\), and explain the role of the constraint.
\end{exercise}

\chapter{Convexity, Smoothing, and Sharp Bounds}
\label{a:advanced-inequalities}

\section{Convexity is a statement about averaging}

A function \(\phi\) on an interval is \emph{convex} if for all
\(x,y\) in that interval and \(0\le t\le1\),
\[
 \phi(tx+(1-t)y)\le t\phi(x)+(1-t)\phi(y).
\]
The graph lies below each joining chord. A concave function satisfies
the reverse inequality. For example, \(x^2\) is convex because the
difference between the right and left sides is
\(t(1-t)(x-y)^2\).

\begin{figure}[htbp]
\centering
\begin{tikzpicture}[x=1.4cm,y=.8cm]
\draw[->](-.15,0)--(2.7,0)node[right]{$x$};
\draw[->](0,-.1)--(0,4.8);
\draw[main,domain=.1:2.2,smooth]plot(\x,{\x*\x});
\draw(.45,.2025)--(2,4);
\draw[aux](1.15,0)--(1.15,1.9175);
\fill[accent](1.15,1.3225)circle(1.6pt);
\fill[ink](1.15,1.9175)circle(1.6pt);
\node[right]at(1.25,.85){graph};
\node[right]at(1.38,2.6){chord};
\end{tikzpicture}
\caption{For a convex function, applying the function after averaging
gives no more than averaging the function values. The point on the
chord uses the same weights as the point on the horizontal axis.}
\end{figure}

\begin{theorem}[Jensen's inequality]
If \(\phi\) is convex, \(\lambda_i\ge0\), and \(\sum_{i=1}^n\lambda_i=1\),
then
\[
 \phi\!\left(\sum_i\lambda_i x_i\right)
 \le\sum_i\lambda_i\phi(x_i).
\]
\end{theorem}
\begin{proof}
Discard zero weights. For \(n>2\), first average \(x_1,\ldots,x_{n-1}\)
using weights divided by \(1-\lambda_n\), apply the induction
hypothesis, and then combine that average with \(x_n\) using the
two-point definition of convexity. The case of a sole nonzero weight
is equality. Strict convexity gives equality only when all inputs
with positive weights are equal.
\end{proof}

\begin{example}[A tangent bound proved without calculus]
For \(u>0\), show \(\sqrt u\le(u+1)/2\).
\end{example}
\begin{solution}
The difference is \((\sqrt u-1)^2/2\ge0\).
Geometrically, the line \((u+1)/2\) touches the concave square-root
graph at \(u=1\), but the algebra proves the bound globally.
\end{solution}

More generally weighted AM--GM gives, for \(r>1\) and \(x,a>0\),
\[
 x^r\ge r a^{r-1}x-(r-1)a^r.
\]
Apply it with weights \(1/r,(r-1)/r\) to \(x^r,a^r\).
This supporting-line inequality can itself prove convexity:
at \(a=tx+(1-t)y\), apply it to \(x,y\), weight, and add.
Thus Jensen applies to \(x^r\) without needing a course in calculus.

\section{Hölder and Minkowski are controlled normalization}

\begin{theorem}[Hölder]
Let \(p,q>1\) satisfy \(1/p+1/q=1\). For real lists \(u_i,v_i\),
\[
 \sum_i|u_iv_i|\le
 \left(\sum_i|u_i|^p\right)^{1/p}
 \left(\sum_i|v_i|^q\right)^{1/q}.
\]
\end{theorem}
\begin{proof}
If either factor on the right is zero, the claim is immediate.
Otherwise divide the lists by those factors. The normalized lists
\(U,V\) have \(\sum|U_i|^p=\sum|V_i|^q=1\).
Weighted AM--GM gives
\(|U_iV_i|\le |U_i|^p/p+|V_i|^q/q\).
Sum, then undo the normalization.
\end{proof}

\begin{example}[Match the exponents to a fraction]
For \(a_i\ge0,b_i>0\), show
\[
 \sum_i\frac{a_i^3}{b_i^2}\ge
 \frac{(\sum_i a_i)^3}{(\sum_i b_i)^2}.
\]
\end{example}
\begin{solution}
In Hölder, write
\(a_i=(a_i/b_i^{2/3})b_i^{2/3}\) and take \(p=3,q=3/2\).
Cube the resulting inequality. Equality holds when the nonzero
normalized lists attain equality in weighted AM--GM, equivalently
when \(a_i/b_i\) is constant, including the all-zero numerator case.
\end{solution}

For \(p>1\), Minkowski's inequality is
\[
 \left(\sum|u_i+v_i|^p\right)^{1/p}
 \le \left(\sum|u_i|^p\right)^{1/p}
     +\left(\sum|v_i|^p\right)^{1/p}.
\]
To prove it, set \(W=(\sum|u_i+v_i|^p)^{1/p}\).
If \(W=0\), there is nothing to prove. Otherwise
\[
 W^p\le
 \sum |u_i||u_i+v_i|^{p-1}
 +\sum |v_i||u_i+v_i|^{p-1}.
\]
Hölder bounds the right side by
\((\|u\|_p+\|v\|_p)W^{p-1}\), where
\(\|u\|_p=(\sum|u_i|^p)^{1/p}\). Cancel the positive power of \(W\).
For \(p=1\), the result is the ordinary absolute-value triangle inequality.

The power means
\(M_r=(\frac1n\sum x_i^r)^{1/r}\) for positive \(x_i\) increase with
positive \(r\): if \(0<r<s\), apply Jensen to \(t^{s/r}\) at
the inputs \(x_i^r\) and take the \(s\)th root.
The geometric mean is the limiting value at \(r=0\);
AM--GM proves \(M_0\le M_r\) for positive \(r\).
Negative powers can be handled by reciprocals, and the harmonic mean
is \(M_{-1}=n/\sum(1/x_i)\).

\section{Schur: preserve ordering until the signs are visible}

\begin{theorem}[Schur's inequality]
For nonnegative \(x,y,z\) and a real exponent \(t>0\),
\[
 \sum_{\cyc}x^t(x-y)(x-z)\ge0.
\]
In particular,
\[
 x^3+y^3+z^3+3xyz
 \ge x^2y+x^2z+y^2x+y^2z+z^2x+z^2y.
\]
\end{theorem}
\begin{proof}
By symmetry order \(x\ge y\ge z\). The terms belonging to \(x,y\)
combine to
\[
 (x-y)\bigl[x^t(x-z)-y^t(y-z)\bigr]\ge0,
\]
because both factors in the first product inside brackets are at
least their nonnegative counterparts. The remaining term
\(z^t(z-x)(z-y)\) is also nonnegative.
For \(t=1\), expansion gives the second display.
\end{proof}

For \(t=1\), equality occurs when all three variables are equal, or one
is zero and the other two are equal. This can be read from the ordered
proof. In the symmetric variables \(p=x+y+z,q=xy+yz+zx,r=xyz\),
the same inequality is \(p^3+9r\ge4pq\).

\section{Smoothing needs a direction, not a slogan}

Smoothing replaces variables by a simpler tuple while controlling the
expression. For example, keeping \(x+y=s\) fixed and writing \(d=x-y\)
gives
\[
 x^4+y^4=\frac{s^4+6s^2d^2+d^4}{8}.
\]
Replacing \(x,y\) by their average makes \(d=0\), so it decreases this
sum. It also preserves their sum and nonnegativity. It generally
changes their product, so it is invalid under an additional fixed-product
constraint unless that change is accounted for.

\begin{example}[A justified two-equal-variable reduction]
Let \(a,b,c\ge0\), \(a+b+c=3\). Prove
\(a^2+b^2+c^2+abc\ge4\).
\end{example}
\begin{solution}
The next lemma applies because the expression minus \(4\) is
\(5-2q+r\), linear in \(r=abc\) when \(p,q\) are fixed.
It suffices to check a zero variable or two equal variables.
If \(c=0\), the sum of squares is at least \(9/2>4\).
If \(a=b=x,c=3-2x\), where \(0\le x\le3/2\), the difference is
\[
 2x^2+(3-2x)^2+x^2(3-2x)-4=(x-1)^2(5-2x)\ge0.
\]
Equality is possible only at \(a=b=c=1\).
\end{solution}

\begin{lemma}[A precise \(p,q,r\) reduction]\label{a:pqr}
Suppose a continuous symmetric expression in nonnegative \(a,b,c\)
has, after fixing \(p=a+b+c>0\), the form
\(A(p,q)+B(p,q)r\), where \(q=ab+bc+ca,r=abc\).
To prove it nonnegative for all such triples, it suffices to check
triples with a zero variable and triples with two equal variables.
\end{lemma}
\begin{proof}
For any fixed feasible \(p,q\), the set of nonnegative triples with
these values is closed and bounded, hence compact. The expression
attains a minimum there. If \(B(p,q)=0\), it is constant on the set,
and it is enough to find one endpoint configuration as below.
If all three variables at an extremum are positive and distinct,
consider their polynomial
\[
 X^3-pX^2+qX-r.
\]
Choose small disjoint positive intervals around its three roots.
The polynomial has opposite signs at the ends of each interval.
Sufficiently small positive or negative changes of \(r\) preserve
these sign changes, so the intermediate value theorem gives three
positive roots, one in each interval. Their sum and pairwise-product
sum remain \(p,q\), while their product has changed in either direction.
Thus such a point cannot minimize a nonconstant linear function of \(r\).
Minimizing or maximizing \(r\) itself also shows that each nonempty
fiber has a point where a variable is zero or two coincide. This
handles the constant case and proves the reduction.
\end{proof}

\pitfall The instruction to set two variables equal is not a general
theorem about symmetric inequalities. The lemma requires a verified
dependence linear in \(r\) on each fixed-\(p,q\) fiber. For an arbitrary
expression, prove a suitable smoothing lemma or use another method.

\section{Majorization: compare how spread out two lists are}

For decreasing real lists \(x_1\ge\cdots\ge x_n\) and
\(y_1\ge\cdots\ge y_n\), say that \(x\) \emph{majorizes} \(y\) when
\[
 \sum_{i=1}^k x_i\ge\sum_{i=1}^k y_i\quad(1\le k<n),
 \qquad \sum_{i=1}^n x_i=\sum_{i=1}^n y_i.
\]
The first list has more of its mass concentrated in its largest entries.
For a convex function \(\phi\) on an open interval containing all entries,
Karamata's inequality states
\(\sum\phi(x_i)\ge\sum\phi(y_i)\).

Here is a proof that explains the mechanism. A convex function at an
interior point \(t\) has a supporting slope \(s\) with
\(\phi(u)\ge\phi(t)+s(u-t)\) for every \(u\).
Indeed slopes of secants ending at \(t\) are no greater than slopes
of secants beginning at \(t\); choose \(s\) between their supremum
and infimum. These bounds are finite at an interior point.
Choose slopes \(s_i\) at \(y_i\); they are decreasing with \(i\),
choosing the same slope for repeated \(y_i\). Put
\(D_k=\sum_{i=1}^k(x_i-y_i)\), so \(D_k\ge0\), \(D_n=0\).
Then the supporting inequalities and a telescoping rearrangement give
\[
 \sum_i[\phi(x_i)-\phi(y_i)]
 \ge\sum_i s_i(x_i-y_i)
 =\sum_{k=1}^{n-1}(s_k-s_{k+1})D_k\ge0.
\]
This proves the claim without assuming that \(\phi\) is differentiable.

In polynomial inequalities, many comparisons sometimes attributed to
Muirhead can instead be proved directly by weighted AM--GM. For example,
\(a^2b\le(2a^3+b^3)/3\); summing over all six permutations gives
\[
 2(a^3+b^3+c^3)\ge
 a^2b+a^2c+b^2a+b^2c+c^2a+c^2b.
\]
The general Muirhead theorem makes this comparison systematic.
For a real exponent list \(\alpha=(\alpha_1,\ldots,\alpha_n)\) and
positive variables \(x_1,\ldots,x_n\), define its symmetrized monomial mean
\[
 M_\alpha(x)=\frac1{n!}\sum_{\pi}
                  x_{\pi(1)}^{\alpha_1}\cdots x_{\pi(n)}^{\alpha_n},
\]
where the sum runs over all permutations \(\pi\) of \(1,\ldots,n\).
If decreasing \(\alpha\) majorizes decreasing \(\beta\), then
\(M_\alpha(x)\ge M_\beta(x)\).

Here are the two ingredients of a proof. First, a majorized vector
\(\beta\) is a weighted average of permutations of \(\alpha\).
We justify this useful fact by induction on the length.
If their largest coordinates agree, remove them and apply induction.
Otherwise choose \(k<n\) with
\(\alpha_k\ge\beta_1\ge\alpha_{k+1}\).
Such a \(k\) exists unless all entries already agree.
Mix the first and \((k+1)\)st coordinates of \(\alpha\) to produce
the pair
\[
 \beta_1,\quad v=\alpha_1+\alpha_{k+1}-\beta_1.
\]
This new vector is a weighted average of \(\alpha\) and the vector
obtained by swapping those two coordinates: the first mixing weight is
\((\beta_1-\alpha_{k+1})/(\alpha_1-\alpha_{k+1})\in[0,1]\).

Its remaining coordinates majorize \(\beta_2,\ldots,\beta_n\).
To check this, for \(j<k\), the \(j\) coordinates
\(\alpha_2,\ldots,\alpha_{j+1}\) are all at least \(\beta_1\),
so their sum is at least \(\beta_2+\cdots+\beta_{j+1}\).
For \(j\ge k\), choose the remaining coordinates corresponding to the
first \(j+1\) original coordinates; their sum is
\(\alpha_1+\cdots+\alpha_{j+1}-\beta_1\), again sufficient by majorization.
The \(j\) largest remaining coordinates have at least those sums,
and total sums agree. Inductively average permutations of the remaining
coordinates to obtain the tail of \(\beta\), keeping \(\beta_1\) fixed.
Every such permuted vector is itself an average of permutations of
\(\alpha\), which completes the induction.

Second, suppose explicitly
\(\beta=\sum_j\lambda_j\alpha^{(j)}\), where the
\(\alpha^{(j)}\) are permutations of \(\alpha\), \(\lambda_j\ge0\),
and \(\sum\lambda_j=1\).
For each fixed permutation of the variables, weighted AM--GM gives
\[
 \prod_i x_i^{\beta_i}
 =\prod_j\left(\prod_i x_i^{\alpha_i^{(j)}}\right)^{\lambda_j}
 \le\sum_j\lambda_j\prod_i x_i^{\alpha_i^{(j)}}.
\]
Average this inequality over all variable permutations.
Each term on the right averages to \(M_\alpha\), proving Muirhead.
For nonnegative exponents, zero variables follow by continuity;
positive variables avoid any \(0^0\) notation.

For practical use, include every permutation, compare exponents in
decreasing order, verify equal total degree, and check every partial sum.
For a single small comparison, the explicit AM--GM decomposition
may still be the most readable proof.

\section{Practice}
\begin{exercise}[\problemlevel{Bridge}]\label{a:ex50}
Use weighted AM--GM to prove
\(\log(\lambda x+(1-\lambda)y)\ge
\lambda\log x+(1-\lambda)\log y\) for \(x,y>0\), \(0\le\lambda\le1\).
Here \(\log\) is the natural logarithm; only its increasing property
and \(\log(u^t)=t\log u\) are needed.
\end{exercise}
\begin{exercise}[\problemlevel{Bridge}]\label{a:ex51}
For \(a_i\ge0,b_i>0\), prove
\(\sum a_i^3/b_i^2\ge(\sum a_i)^3/(\sum b_i)^2\), with equality conditions.
\end{exercise}
\begin{exercise}[\problemlevel{Bridge}]\label{a:ex52}
Prove the triangle inequality
\(\sqrt{(x+u)^2+(y+v)^2}\le\sqrt{x^2+y^2}+\sqrt{u^2+v^2}\)
using Cauchy, and describe equality.
\end{exercise}
\begin{exercise}[\problemlevel{Bridge}]\label{a:ex53}
For \(r>1\), \(a,x>0\), prove
\(x^r\ge ra^{r-1}x-(r-1)a^r\), and identify equality.
\end{exercise}
\begin{exercise}[\problemlevel{Challenge}]\label{a:ex54}
Use Lemma~\ref{a:pqr} to prove
\((a+b+c)^3+9abc\ge4(a+b+c)(ab+bc+ca)\) for nonnegative \(a,b,c\).
\end{exercise}
\begin{exercise}[\problemlevel{Challenge}]\label{a:ex55}
For positive \(a,b,c\) with \(abc=1\), show
\[
 \frac1{a^3(b+c)}+\frac1{b^3(c+a)}+\frac1{c^3(a+b)}
 \ge\frac32.
\]
\end{exercise}
\begin{exercise}[\problemlevel{Challenge}]\label{a:ex56}
Prove \(\sum_{\cyc}x^t(x-y)(x-z)\ge0\) for nonnegative \(x,y,z\)
and real \(t>0\), and find all equality cases.
\end{exercise}

\chapter{Functional Equations: Learn What the Equation Allows}
\label{a:functional-equations}

\section{A function is more than its formula}

A function \(f:D\to E\) assigns to each input in its domain \(D\)
exactly one output in its codomain \(E\). Its range is the set
\(\{f(x):x\in D\}\). It is \emph{injective} if equal outputs imply equal
inputs, and \emph{surjective onto \(E\)} if its range is \(E\).
A bijection has both properties.

In a functional equation, the unknown is the entire function.
Testing \(f(x)=x\) may discover a candidate; it does not prove uniqueness.
To classify solutions, derive necessary properties, use those properties
to obtain every possible function, and substitute the final candidates
into the original equation.

\intuition A functional equation is a machine that turns legal inputs
into information. Inputs such as \(0,1,-x\), and \(f(y)\) are useful
because they cancel terms or repeat expressions. The domain determines
which buttons the machine actually has.

\begin{example}[A range statement obtained from an identity]
Suppose \(f:\RR\to\RR\) satisfies \(f(f(x))=x\). Show that \(f\) is a bijection.
\end{example}
\begin{solution}
If \(f(u)=f(v)\), apply \(f\) and obtain \(u=v\), proving injectivity.
For any real \(y\), the input \(f(y)\) has output \(y\), proving
surjectivity. Such a function is called an involution.
\end{solution}

\pitfall Never cancel a function symbol as if it were multiplication.
From \(f(U)=f(V)\), concluding \(U=V\) requires injectivity.
Likewise replacing \(f(y)\) by an arbitrary real variable requires
surjectivity onto the reals.

\section{Cauchy's equation depends on the domain}

\begin{theorem}[Additive functions]\label{a:cauchy}
Suppose \(A(x+y)=A(x)+A(y)\).
On \(\ZZ\), every solution is \(A(n)=cn\).
On \(\QQ\), every solution is \(A(q)=cq\).
On \(\RR\), every solution satisfying any one of the following conditions
is \(A(x)=cx\): continuity at one point; monotonicity; or boundedness
on some nondegenerate interval.
\end{theorem}
\begin{proof}
Set \(x=y=0\) to get \(A(0)=0\), and set \(y=-x\) to get oddness.
Repeated addition gives \(A(nx)=nA(x)\) for every integer \(n\).
For positive integer \(n\), \(nA(1/n)=A(1)\); hence rational inputs
give \(A(q)=qA(1)\).

Continuity at any point implies continuity at \(0\), since
\(A(t+h)-A(t)=A(h)\). Approximating a real \(x\) by rationals then
gives \(A(x)=xA(1)\).
If \(A\) is increasing, squeeze \(x\) between rational approximations;
monotonicity squeezes \(A(x)\) between their known values.
If \(A\) is decreasing, apply this argument to \(-A\).

Finally suppose \(|A|\le M\) on \((u,v)\), \(u<v\).
For sufficiently small \(h\), choose \(t,t+h\in(u,v)\); then
\(|A(h)|\le2M\) for \(|h|<v-u\).
For any positive integer \(N\) and \(|h|<(v-u)/N\),
\(|A(h)|=|A(Nh)|/N\le2M/N\). Choosing \(N\) arbitrarily large proves
continuity at zero. If \(M=0\), this conclusion is immediate.
\end{proof}

Here boundedness means a two-sided absolute bound on the interval.
For a real domain, additivity alone must not be replaced by linearity.
Nonlinear additive functions can be constructed by prescribing values
on a basis of \(\RR\) as a vector space over \(\QQ\); existence of such
a basis is an additional set-theoretic result. We need no construction
to see the logical issue: a proof that determines only rational inputs
has not determined all real inputs without another argument.

\begin{example}[Remove the known quadratic part]
Find all \(f:\ZZ\to\ZZ\) satisfying
\(f(m+n)=f(m)+f(n)+2mn\).
\end{example}
\begin{solution}
The term \(2mn\) is the cross-term in \((m+n)^2\).
Put \(A(n)=f(n)-n^2\). The equation becomes additive, so
\(A(n)=cn\) with \(c=A(1)\in\ZZ\).
Thus \(f(n)=n^2+cn\), and direct substitution verifies every integer \(c\).
\end{solution}

\section{Use comparisons of two routes}

In many problems one can evaluate the same expression in two ways.
For example, compute \(f(x+y+z)\) by first combining \(x,y\), and then
by first combining \(y,z\). This is the functional-equation analogue
of comparing two counts of the same set.

\begin{example}[Associativity forces a formula]
Determine \(f:\RR\to\RR\) satisfying
\(f(x)f(y)=f(x+y)+xy\) for all real \(x,y\).
\end{example}
\begin{solution}
Setting \(y=0\) gives \(f(x)(f(0)-1)=0\).
The identically zero function fails at \(x=y=1\), so \(f(0)=1\).
Compute \(f(x+y+z)\) in two groupings:
\[
\begin{aligned}
 f(x+y+z)&=[f(x)f(y)-xy]f(z)-(x+y)z,\\
 f(x+y+z)&=f(x)[f(y)f(z)-yz]-x(y+z).
\end{aligned}
\]
Subtracting yields \(xy(f(z)-1)=yz(f(x)-1)\).
With \(y=z=1\), this becomes
\(f(x)=1+cx\), where \(c=f(1)-1\).
Substitution in the original equation gives \(c^2=1\).
Both \(f(x)=1+x\) and \(f(x)=1-x\) work.
\end{solution}

\section{Multiplicative and exponential equations}

If \(f:\RR\to\RR\) satisfies \(f(x+y)=f(x)f(y)\), then either \(f=0\)
or \(f(0)=1\). In the latter case
\(f(x)=f(x/2)^2\ge0\) and \(f(x)f(-x)=1\), so \(f(x)>0\).
Then \(A(x)=\log f(x)\) is additive.
Under continuity or another appropriate regularity condition,
\(f(x)=e^{cx}\). Without regularity, the correct description is
\(f(x)=e^{A(x)}\) for an arbitrary additive \(A\).

Similarly, for \(g:(0,\infty)\to(0,\infty)\) satisfying
\(g(xy)=g(x)g(y)\), define \(A(t)=\log g(e^t)\).
It is additive, so a continuous or monotone \(g\) is a power \(g(x)=x^c\).
Positivity of \(g\) alone does not prove regularity of \(A\).

\begin{example}[The missing hypothesis changes the answer]
Find the continuous real functions satisfying
\(f(x+y)=f(x)+f(y)+1\).
\end{example}
\begin{solution}
Put \(A(x)=f(x)+1\). Then \(A\) is additive and continuous,
so \(A(x)=cx\). Hence \(f(x)=cx-1\).
Without continuity, replace \(cx\) by an arbitrary additive function.
\end{solution}

\recognition Look for the defect from a familiar equation: a term
\(2xy\) suggests subtracting a square; a constant suggests shifting
the output; products suggest logarithms when positivity and domains
permit them. Establish injectivity or range only when the equation
provides a reason.

\section{Practice}
\begin{exercise}[\problemlevel{Foundation}]\label{a:ex57}
Let \(f:\RR\to\RR\) be additive, monotone, and satisfy \(f(1)=2\).
Determine \(f\).
\end{exercise}
\begin{exercise}[\problemlevel{Bridge}]\label{a:ex58}
Classify \(f:\ZZ\to\ZZ\) with
\(f(m+n)=f(m)+f(n)+2mn\).
\end{exercise}
\begin{exercise}[\problemlevel{Bridge}]\label{a:ex59}
Classify all \(f:\RR\to\RR\) satisfying
\(f(x+y)=f(x)+f(y)+1\), first without regularity assumptions and then
under continuity.
\end{exercise}
\begin{exercise}[\problemlevel{Challenge}]\label{a:ex60}
Classify \(f:\RR\to\RR\) satisfying
\[
 f(x+y)+f(x-y)=2f(x)+2y^2
\]
without assuming continuity.
\end{exercise}
\begin{exercise}[\problemlevel{Challenge}]\label{a:ex61}
A monotone function \(f:(0,\infty)\to(0,\infty)\) satisfies
\(f(xy)=f(x)f(y)\), \(f(2)=4\). Prove \(f(x)=x^2\).
\end{exercise}
\begin{exercise}[\problemlevel{Challenge}]\label{a:ex62}
Suppose \(f:\RR\to\RR\) satisfies \(f(x+f(y))=f(x)+y\).
Describe all solutions using additive involutions. If \(f\) is
continuous, determine it explicitly.
\end{exercise}
\begin{exercise}[\problemlevel{Bridge}]\label{a:ex63}
Find all continuous \(f:\RR\to\RR\) with \(f(x+y)=f(x)f(y)\).
\end{exercise}

\chapter{Functional Equations: Iteration and Rigidity}
\label{a:advanced-fe}

\section{Iterates commute with the function}

The second iterate \(f^2\) means \(f\circ f\), not the pointwise square
\(f(x)^2\). Always distinguish \(f(f(x))\) from \([f(x)]^2\).
Every function commutes with its own iterates:
\(f(f^2(x))=f^2(f(x))\).
If the equation gives a simple formula for \(f^2\), this observation
can create a useful translation identity.

\begin{example}[An impossible integer shift]
Show that no \(f:\ZZ\to\ZZ\) satisfies \(f(f(n))=n+1\) for every integer \(n\).
\end{example}
\begin{solution}
The two ways to write the third iterate give
\[
 f(n+1)=f(f(f(n)))=f(n)+1.
\]
Thus \(f(n)=n+c\) for an integer \(c\).
Its second iterate is \(n+2c\), which cannot equal \(n+1\).
Over the reals, \(f(x)=x+1/2\) would work: the codomain caused
the final obstruction.
\end{solution}

On a finite set, an involution partitions the elements into fixed points
and pairs that are exchanged. A function whose \(k\)th iterate is the
identity is a permutation whose cycle lengths divide \(k\).
This follows by repeatedly applying the function to each element until
it returns to its start; no tail can precede a cycle because the
function is invertible, with inverse \(f^{k-1}\).

\section{A polynomial assumption is powerful information}

When the unknown is known to be a polynomial, degree comparison should
usually precede substitutions. For nonconstant polynomials \(P,Q\),
\(\deg(PQ)=\deg P+\deg Q\) and
\(\deg(P\circ Q)=(\deg P)(\deg Q)\).
These facts follow by multiplying the leading terms.

\begin{example}[Composition leaves almost no possibilities]
Find all \(P\in\RR[x]\) satisfying \(P(P(x))=P(x)^2\).
\end{example}
\begin{solution}
Constants satisfy \(c=c^2\), so give \(0,1\).
For a nonconstant degree \(d\), degrees give \(d^2=2d\), hence \(d=2\).
If \(P=ax^2+bx+c\), leading coefficients give \(a^3=a^2\), so \(a=1\).
Now \(P(P(x))=P(x)^2+bP(x)+c\).
Thus \(bP+c=0\), forcing \(b=c=0\). The remaining solution \(P=x^2\)
checks.
\end{solution}

\begin{example}[Periodicity plus polynomiality]
If \(P\in\RR[x]\) and \(P(x+1)=P(x)\) for all \(x\), prove \(P\) is constant.
\end{example}
\begin{solution}
A degree \(d\ge1\) polynomial with leading coefficient \(a\) has
\(P(x+1)-P(x)\) of degree \(d-1\) and leading coefficient \(da\ne0\).
It cannot vanish identically.
For arbitrary functions the conclusion is false: \(\sin(2\pi x)\)
has period \(1\).
\end{solution}

\section{Square-valued identities can create regularity}

\intuition An equation may contain no continuity assumption yet force
the function to preserve nonnegativity. Combined with additivity, that
produces monotonicity and therefore linearity. This is a conclusion
derived from the equation, not an assumption silently added to it.

For instance, if an additive \(A:\RR\to\RR\) also satisfies
\(A(x^2)=A(x)^2\), then \(A(t)\ge0\) for \(t\ge0\).
Consequently \(A(y)-A(x)=A(y-x)\ge0\) for \(y\ge x\);
the function is increasing and Theorem~\ref{a:cauchy} applies.
If \(A\ne0\), then \(A(1)=A(1)^2\) forces \(A(1)=1\), so \(A(x)=x\).
The zero function is the other possibility.

\begin{example}[Do not divide away an entire branch]
Classify additive functions \(A:\RR\to\RR\) satisfying
\(A(x^2)=A(x)^2\).
\end{example}
\begin{solution}
The preceding positivity argument gives \(A(x)=cx\).
Substitution yields \(c=c^2\), so \(c=0\) or \(c=1\).
Dividing the original equation by \(A(1)\) before separating its zero
case would have lost the zero function.
\end{solution}

\section{Build a proof architecture}

For a difficult two-variable functional equation, write a short
dependency chain before calculating:
\[
\begin{gathered}
 \text{special inputs}\ \longrightarrow\
 \text{range or injectivity}\ \longrightarrow\ \text{translation law}\\
 \longrightarrow\ \text{rigidity}\ \longrightarrow\ \text{verification}.
\end{gathered}
\]
There is no universal order. Surjectivity might first provide a zero
of the function; that zero may determine \(f(0)\); the resulting
iteration law may then make two evaluations comparable.

Problems~\ref{a:ex68} and \ref{a:ex69} are long enough that this
architecture matters. In each, the equation itself eventually rules
out irregular functions. A guess of \(f(x)=x\) or \(\pm x\) is only the
beginning.

\section{Practice}
\begin{exercise}[\problemlevel{Bridge}]\label{a:ex64}
Prove there is no integer-valued function on \(\ZZ\) with
\(f(f(n))=n+1\).
\end{exercise}
\begin{exercise}[\problemlevel{Bridge}]\label{a:ex65}
If \(f\) is an involution on a finite set of \(N\) elements, prove
that the number of its fixed points has the same parity as \(N\).
\end{exercise}
\begin{exercise}[\problemlevel{Challenge}]\label{a:ex66}
Classify \(f:\RR\to\RR\) with
\(f(x)f(y)=f(x+y)+xy\) for all real \(x,y\).
\end{exercise}
\begin{exercise}[\problemlevel{Bridge}]\label{a:ex67}
Classify real polynomials \(P\) with \(P(P(x))=P(x)^2\).
\end{exercise}
\begin{exercise}[\problemlevel{Challenge}]\label{a:ex68}
Determine all \(f:\RR\to\RR\) such that
\[
 f(x^2+f(y))=y+f(x)^2\qquad(x,y\in\RR).
\]
\end{exercise}
\begin{exercise}[\problemlevel{Challenge}]\label{a:ex69}
Determine all \(f:\RR\to\RR\) such that
\[
 f(xf(x)+f(y))=f(x)^2+y\qquad(x,y\in\RR).
\]
\end{exercise}
\begin{exercise}[\problemlevel{Bridge}]\label{a:ex70}
Determine all \(f:\ZZ\to\ZZ\) with
\(f(m+n)+f(m-n)=2f(m)+2f(n)\).
\end{exercise}

\chapter{Algebraic Synthesis and Hard Problems}
\label{a:synthesis}

\section{Choose a representation before choosing a theorem}

An unfamiliar algebra problem may belong to several chapters at once.
A polynomial equation may hide a recurrence; a symmetric system may
be a statement about roots; an inequality may become elementary only
after normalization. The following questions organize the first
substantial attempt.

\begin{enumerate}
\item What is invariant under relabeling, scaling, shifting, or taking
reciprocals? Preserve that structure.
\item Does the target ask for a value, a bound, or a classification?
A classification requires proving that no other cases exist.
\item What happens at equality or at a repeated root? These often
identify the sharp example.
\item Can two evaluations of the same object be compared?
\item Which assumptions have not yet been used? They often explain
why a tempting counterexample is excluded.
\end{enumerate}

These questions are not a replacement for exploration. They direct
experiments toward information that can survive in a proof.

\section{A local bound may need a global repair}

A supporting-line argument can be excellent near equality and false
far away. Rather than abandon it, separate the exceptional region
and use a cruder, stronger estimate there.

\begin{example}[Patch the region where the tangent bound fails]
Let \(a,b,c>0\), \(a+b+c=3\). Prove
\[
 \frac1{a^2}+\frac1{b^2}+\frac1{c^2}\ge a^2+b^2+c^2.
\]
\end{example}
\begin{solution}
At equality one expects \(a=b=c=1\). Seek a linear lower bound for
\(1/x^2-x^2\) that vanishes after summing. An exact calculation gives
\[
 \frac1{x^2}-x^2+4(x-1)
 =\frac{(x-1)^2(-x^2+2x+1)}{x^2}.
\]
For \(0<x\le7/3\), the quadratic factor is nonnegative:
its minimum on \([0,7/3]\) occurs at an endpoint, with values \(1,2/9\).
If all three variables are in that interval, summing proves the claim.

Otherwise, say \(c>7/3\). Then \(a+b<2/3\), so one of \(a,b\) is
less than \(1/3\). Its reciprocal square alone exceeds \(9\).
But \(a^2+b^2+c^2<(a+b+c)^2=9\), so the desired inequality is strict.
The equality case from the first region is \(a=b=c=1\).
\end{solution}

The number \(7/3\) was chosen to meet two needs: the local quadratic
remains nonnegative up to that point, and outside it a variable is
forced below \(1/3\). It is an auxiliary threshold, not a guessed
equality value.

\section{Compare a solution with a model of the same degree}

If a polynomial identity suggests a familiar family, construct that
family first. Then subtract a candidate of the same degree.
The difference often has a smaller degree that the identity cannot support.

\begin{example}[The lowest nonzero coefficient exposes a contradiction]
Determine all \(P\in\RR[x]\) with \(P(x^2)=P(x)^2\).
\end{example}
\begin{solution}
The zero polynomial works. Otherwise write \(P=x^kQ\), where \(k\ge0\)
and \(Q(0)\ne0\). Cancelling the common monomial gives
\(Q(x^2)=Q(x)^2\). Constants at zero give \(Q(0)=1\).
If \(Q\) were nonconstant, let \(j\ge1\) be its least positive exponent
with nonzero coefficient \(c\). In \(Q(x)^2\), the coefficient of
\(x^j\) is \(2c\). In \(Q(x^2)\), that coefficient is zero: for odd
\(j\) no such exponent occurs, and for even \(j\) the needed coefficient
at \(j/2\) is zero by minimality. Contradiction.
Thus \(Q=1\). The solutions are \(P=0\) and \(P=x^k\), \(k\ge0\).
\end{solution}

\section{Use a quadratic equation to produce a smaller solution}

If an integer equation is quadratic in one variable, one root may
produce another integer root through Vieta. The useful question is
whether the second root is positive and smaller. Without these checks,
there is no descent.

\begin{example}[The transformation, before the classification]
For positive integers satisfying \(x^2+y^2+z^2=3xyz\), show that
\((x,y,3xy-z)\) is another positive integer solution.
\end{example}
\begin{solution}
Regarding the equation as a quadratic in \(z\), the two roots have sum
\(3xy\) and product \(x^2+y^2\). Thus the other root is
\[
 z'=3xy-z=\frac{x^2+y^2}{z}>0.
\]
It is an integer by the sum formula, and it satisfies the same quadratic.
Problem~\ref{a:ex71} supplies the ordering and size estimates needed
to turn this transformation into a complete classification.
\end{solution}

\section{A final algebra problem set}

The problems below mix several methods. Their labels indicate the
intended level of the reasoning, not a guarantee about a particular
contest placement. Write the classification and the completeness
argument separately when that makes the logic clearer.

\begin{exercise}[\problemlevel{Challenge}]\label{a:ex71}
Show that every positive integer solution of
\(x^2+y^2+z^2=3xyz\) can be obtained from \((1,1,1)\) by permutations
and repeatedly replacing one coordinate by three times the product of
the other two minus that coordinate.
\end{exercise}
\begin{exercise}[\problemlevel{Challenge}]\label{a:ex72}
Let \(P\in\ZZ[x]\) be nonzero, with \(P(n)\ne0\) for every positive
integer \(n\). If \(P(n)\mid P(n+1)\) for every positive integer \(n\),
prove that \(P\) is constant.
\end{exercise}
\begin{exercise}[\problemlevel{Challenge}]\label{a:ex73}
For positive \(a,b,c\) with sum \(3\), prove both
\[
 \frac1{a^2}+\frac1{b^2}+\frac1{c^2}\ge a^2+b^2+c^2,
 \qquad a^2+b^2+c^2+abc\ge4,
\]
and find equality in each inequality.
\end{exercise}
\begin{exercise}[\problemlevel{Challenge}]\label{a:ex74}
Let \(a_1,\ldots,a_n\) be distinct integers, \(n\ge1\).
Prove that \(\prod_{j=1}^n(x-a_j)-1\) is irreducible over \(\QQ\).
Explain why replacing the minus sign by a plus sign does not give
a general theorem, by factoring \(x(x+1)(x+2)(x+3)+1\).
\end{exercise}
\begin{exercise}[\problemlevel{Challenge}]\label{a:ex75}
For real \(x,y,z\) with \(x+y+z=0\) and \(x^4+y^4+z^4=2\),
find the largest possible \(|xyz|\), including equality cases.
\end{exercise}
\begin{exercise}[\problemlevel{Challenge}]\label{a:ex76}
Let \(x,y,z\ge0\), \(x^2+y^2+z^2+xyz=4\).
Prove \(x+y+z\le3\) and classify equality. Explain how a
normalization converts the constraint into a previously proved inequality.
\end{exercise}
\begin{exercise}[\problemlevel{Challenge}]\label{a:ex77}
Classify all real polynomials \(P\) satisfying
\[
 P(x)^2=P(x^2-2)+2.
\]
Express the nonconstant solutions using Chebyshev polynomials, and prove
that no others exist.
\end{exercise}


\part{Combinatorics}

% ---- chapters/combinatorics ----
\chapter{Counting Is the Art of Choosing a Description}
\label{c:counting}

Combinatorics studies finite structures: subsets, arrangements, paths,
colorings, graphs, and the rules that transform them. Its central question
is often not which formula applies, but what exactly an object is,
and how to describe it without losing or repeating it.
This chapter begins with that question.

\section{Objects, decisions, and maps}

For a nonnegative integer \(n\), write \([n]=\{1,\ldots,n\}\), with
\([0]=\varnothing\). A \emph{set} records membership, not order or
repetition: \(\{1,2\}=\{2,1\}\). A \emph{sequence} records both positions
and values: \((1,2)\ne(2,1)\). A \emph{multiset} records multiplicities
but not order. Choosing which of these models a problem is already
mathematical work.

A map \(f:A\to B\) assigns one value \(f(a)\in B\) to every \(a\in A\).
It is injective if different inputs have different outputs, surjective
if every target occurs, and bijective if both hold. A bijection is a
perfect reversible translation between descriptions.

\intuition Suppose a password has three positions, each with two choices.
Draw a decision tree: each partial word has two children. The eight leaves
are the passwords. Multiplication is justified by the equal number of
continuations at each stage, not by a vague assertion of independence.
If some partial choices have fewer continuations, split into cases or
count the leaves another way.

\begin{theorem}[Sum and product rules]
If a finite set is partitioned into disjoint sets of sizes
\(a_1,\ldots,a_k\), its size is \(a_1+\cdots+a_k\).
If an object is specified uniquely by successive decisions, with exactly
\(a_i\) choices at step \(i\) for every permitted earlier history, there
are \(a_1\cdots a_k\) objects.
\end{theorem}
\begin{proof}
For the sum rule, every object belongs to exactly one case.
For the product rule, induction on the number of decisions counts the
leaves of the decision tree. Each leaf of the first \(k-1\) stages has
exactly \(a_k\) children.
\end{proof}

\begin{example}[Subsets become binary words]
How many subsets does an \(n\)-element set have?
\end{example}
\begin{solution}
Identify the set with \([n]\). Encode a subset \(S\) by
\((b_1,\ldots,b_n)\), where \(b_i=1\) if \(i\in S\) and \(b_i=0\)
otherwise. The encoding is reversible, and each position has two choices,
so the answer is \(2^n\), including the empty subset.
\end{solution}
\strategy The decisive move was to replace a variable-size object by a
fixed-length description. Later, indicator variables do the same job in
probability and linear algebra.

\section{Permutations and combinations}

A permutation of \(n\) distinct objects is an ordered list using each
exactly once. It has \(n(n-1)\cdots1=n!\) possibilities, where \(0!=1\).
An ordered selection of \(k\) distinct objects has
\(n!/(n-k)!\) possibilities. A \(k\)-element subset can be ordered in
\(k!\) ways, so
\[
 \binom nk=\frac{n!}{k!(n-k)!}\qquad(0\le k\le n).
\]
We set \(\binom nk=0\) outside this range when \(n\ge0\).
This zero convention is useful only when the upper argument is
nonnegative; an expression such as \(\binom{-2}{3}\) needs a separate
definition and must not be silently interpreted by the factorial formula.

\begin{example}[Repeated letters]
How many words contain three copies of \(A\), two copies of \(B\),
and one copy of \(C\)?
\end{example}
\begin{solution}
Temporarily distinguish the copies as \(A_1,A_2,A_3,B_1,B_2,C\).
There are \(6!\) labelled words. Each unlabelled word has
\(3!2!\) labelled versions, so the answer is \(6!/(3!2!)=60\).
Equivalently, choose the three \(A\)-positions and then the two
\(B\)-positions: \(\binom63\binom32=60\).
\end{solution}
\pitfall Division is justified only when every object has the same
number of labelled descriptions. Circular arrangements with repeated
colors often violate this condition. We will use Burnside's lemma there.

\begin{example}[Monotone paths]
A path from \((0,0)\) to \((r,s)\) uses unit steps right or up.
There are \(\binom{r+s}{r}\) such paths.
\end{example}
\begin{solution}
Record the steps as a word with \(r\) letters \(R\) and \(s\) letters
\(U\). Every such word gives one path, and every path gives one word.
Choose the \(R\)-positions. This counts paths regardless of how they
look in a drawing.
\end{solution}

\begin{center}
\begin{tikzpicture}[scale=.55]
 \draw[step=1,gray!30] (0,0) grid (5,3);
 \draw[main,-{Stealth}] (0,0)--(1,0)--(1,1)--(3,1)--(3,2)--(5,2)--(5,3);
 \node[below left] at (0,0) {\((0,0)\)};
 \node[above right] at (5,3) {\((5,3)\)};
\end{tikzpicture}
\end{center}
The picture is one of \(\binom85\) descriptions; its turns are not extra
choices after the step word has been chosen.

\section{Distributing identical objects}

\intuition In a row of seven identical stars, two separators can create
three boxes: \({**}|{}|{*****}\) means \((2,0,5)\).
Empty boxes correspond to consecutive separators or an end separator.
Thus allowing emptiness changes the encoding, not merely the wording.

\begin{theorem}[Stars and bars]
For integers \(n\ge0\) and \(k\ge1\), the number of solutions of
\(x_1+\cdots+x_k=n\) in nonnegative integers is
\(\binom{n+k-1}{k-1}\).
For positive integers \(x_i\), the number is \(\binom{n-1}{k-1}\)
when \(n\ge k\), and zero otherwise.
\end{theorem}
\begin{proof}
Use \(n\) stars and \(k-1\) bars. The \(k-1\) bar positions in a word of
length \(n+k-1\) determine all boxes. For positive variables, substitute
\(y_i=x_i-1\), reducing the total to \(n-k\).
\end{proof}

\recognition Before using a distribution formula, answer four questions:
are the objects distinct, are the boxes distinct, may a box be empty,
and does order inside a box matter? A labelled ball assigned to a box
is a map; identical balls distributed among labelled boxes give integer
solutions; unlabelled boxes give a different problem.

\section{A first problem-solving routine}
Translate the problem into a finite set. Describe an object uniquely.
Test the description on a small case. Then prove both directions:
every legal object is described, and every legal description gives one
object. Small-case agreement is a diagnostic, not the proof.

\begin{exercise}[\problemlevel{Foundation}]\label{c:ex01}
How many maps \([5]\to[3]\) are there? How many are injective? How many
injective maps \([3]\to[5]\) are there? Explain each count.
\end{exercise}
\begin{exercise}[\problemlevel{Foundation}]\label{c:ex02}
Count the words with four \(A\)'s, three \(B\)'s, and two \(C\)'s.
Then count those whose first and last letters are \(C\).
\end{exercise}
\begin{exercise}[\problemlevel{Foundation}]\label{c:ex03}
Count nonnegative integer solutions of \(x+y+z=12\), first without
extra restrictions and then with \(x\ge2,\ y\ge3,\ z\ge1\).
\end{exercise}
\begin{exercise}[\problemlevel{Bridge}]\label{c:ex04}
How many \(k\)-element subsets of \([n]\) contain no two consecutive
integers? Treat impossible parameter values explicitly.
\end{exercise}
\begin{exercise}[\problemlevel{Bridge}]\label{c:ex05}
Count length-\(n\) binary words with exactly \(k\) maximal nonempty
runs of ones. Include the case \(k=0\).
\end{exercise}
\begin{exercise}[\problemlevel{Bridge}]\label{c:ex06}
There are \(n\) distinct people, \(n\ge2\), around an oriented circle,
with rotations identified and reflections distinguished. Count all
arrangements and those in which two specified people are adjacent.
\end{exercise}
\begin{exercise}[\problemlevel{Challenge}]\label{c:ex07}
Let \(n\ge1\). Count the pairs of subsets \((A,B)\) of \([n]\)
for which \(A\subseteq B\). Then deduce
\(\sum_{k=0}^n\binom nk2^k=3^n\) by a second description.
\end{exercise}
\begin{exercise}[\problemlevel{Challenge}]\label{c:ex08}
For \(n\ge1\), count the ordered pairs \((A,B)\) of subsets of \([n]\)
such that \(A\cup B=[n]\) and \(A\cap B\ne\varnothing\).
Give both a coordinatewise argument and an argument by subtraction.
\end{exercise}

\chapter{Count the Same Structure Twice}
\label{c:double}

\section{Incidences and weighted counting}

An \emph{incidence} is a pair \((x,Y)\) expressing a specified relation,
such as a vertex belonging to an edge. Count all such pairs by
fixing the first coordinate, then by fixing the second. The two answers
must agree. This simple device often reveals a quantity that the original
statement hides.

\begin{example}[A committee with a leader]
For \(1\le k\le n\), prove
\[
 k\binom nk=n\binom{n-1}{k-1}.
\]
\end{example}
\begin{solution}
Count a \(k\)-person committee together with one leader. Choose the
committee, then its leader, for the left side. Choose the leader,
then the remaining \(k-1\) members, for the right side.
\end{solution}

\begin{example}[Splitting a selection]
For nonnegative \(r,s,k\), prove Vandermonde's identity
\[
 \binom{r+s}{k}=\sum_{j=0}^k\binom rj\binom{s}{k-j}.
\]
\end{example}
\begin{solution}
Partition a set into disjoint parts of sizes \(r\) and \(s\).
Any \(k\)-subset takes exactly \(j\) elements from the first part for
one \(j\). These cases are disjoint and exhaustive. Impossible choices
contribute zero under our convention.
\end{solution}

\strategy When a sum contains \(\binom nk\), imagine choosing a subset.
When it contains \(k\binom nk\), imagine marking an element of the subset.
When it contains \(\binom nk^2\), imagine two independent choices or a
choice split between two \(n\)-element parts.

\begin{theorem}[An incidence bound]\label{c:incidence}
Let \(X\) be a finite set and \(\mathcal F\) a finite family of its
subsets. If every member of \(\mathcal F\) has at least \(a>0\)
elements and each element of \(X\) belongs to at most \(b\) members, then
\(a|\mathcal F|\le b|X|\).
\end{theorem}
\begin{proof}
Count pairs \((x,F)\) with \(x\in F\in\mathcal F\).
The count is at least \(a|\mathcal F|\) and at most \(b|X|\).
\end{proof}
Notice that equality of two exact counts can still produce inequalities
when the hypotheses give only upper and lower estimates.

\section{Inclusion--exclusion}

\intuition To exclude objects with a defect, subtract the objects with
each individual defect. An object with two defects is then subtracted
twice, so add it back once. With three defects, continue the correction.
The alternating signs repair the multiplicity of each object.

\begin{theorem}[Inclusion--exclusion]\label{c:ie}
For subsets \(A_1,\ldots,A_m\) of a finite universe \(U\), the number
of elements in none of the \(A_i\) is
\[
 \sum_{J\subseteq[m]}(-1)^{|J|}
 \left|\bigcap_{j\in J}A_j\right|,
\]
where the intersection for \(J=\varnothing\) means \(U\).
\end{theorem}
\begin{proof}
An element lying in exactly \(t\) bad sets contributes
\(\sum_{j=0}^t(-1)^j\binom tj=(1-1)^t\).
Its contribution is one for \(t=0\) and zero otherwise.
The binomial expansion follows by choosing either \(1\) or \(-1\)
from each of \(t\) factors.
\end{proof}

\begin{example}[Derangements]
A derangement of \([n]\) is a permutation \(\pi\) with
\(\pi(i)\ne i\) for all \(i\). Its number is
\[
 D_n=\sum_{j=0}^n(-1)^j\binom nj(n-j)!
     =n!\sum_{j=0}^n\frac{(-1)^j}{j!}.
\]
\end{example}
\begin{solution}
Declare \(A_i\) to be the permutations fixing \(i\).
Fixing every point of a specified \(j\)-set leaves \((n-j)!\)
permutations. Apply inclusion--exclusion.
\end{solution}
\pitfall Inclusion--exclusion counts intersections, not intersections
with all other defects absent. Requiring exactly those fixed points
inside an intersection would subtract the same restriction twice.

\section{Surjections and set partitions}

Let \(S(n,k)\) denote the number of partitions of \([n]\) into \(k\)
nonempty, unlabelled blocks. Set \(S(0,0)=1\);
set \(S(n,0)=0\) for \(n>0\), and \(S(n,k)=0\) for \(k>n\).
Then, for \(n,k\ge1\),
\[
 S(n,k)=S(n-1,k-1)+kS(n-1,k).
\]
Indeed, \(n\) either forms a singleton block or joins one of the \(k\)
existing blocks of a partition of \([n-1]\).
Surjections \([n]\to[k]\) are partitions with their blocks labelled, so
their number is \(k!S(n,k)\), also equal to
\[
 \sum_{j=0}^k(-1)^j\binom kj(k-j)^n.
\]
Here a bad event omits a specified target. The formula is valid
for \(n\ge1\), avoiding any need to discuss \(0^0\).

\section{Involutions: cancellation without arithmetic}

An involution on a set is a map \(f\) satisfying \(f(f(x))=x\).
Its orbits have size one or two. If each object has a sign and \(f\)
reverses the sign of every nonfixed object, paired terms cancel.
Only fixed objects remain.

\begin{example}[Even and odd subsets]
For \(n\ge1\), exactly half of the subsets of \([n]\) have even size.
\end{example}
\begin{solution}
Toggle membership of the element \(1\). This reversible operation pairs
every even subset with an odd subset and has no fixed point.
There are \(2^n\) subsets, so each class has \(2^{n-1}\) members.
\end{solution}
\recognition Alternating sums suggest a sign-reversing pairing.
Parity conclusions suggest pairing everything except a recognizable
fixed family. The hard part is often choosing a canonical element to
toggle so that applying the rule twice returns to the start.

\section{Exercises}
\begin{exercise}[\problemlevel{Foundation}]\label{c:ex09}
Count \(\sum_{k=0}^n k\binom nk\) and
\(\sum_{k=0}^n k(k-1)\binom nk\) by marking subset elements.
\end{exercise}
\begin{exercise}[\problemlevel{Foundation}]\label{c:ex10}
Count the onto maps \([5]\to[3]\). Compute the same answer using
the recurrence for \(S(n,k)\).
\end{exercise}
\begin{exercise}[\problemlevel{Bridge}]\label{c:ex11}
Count integer solutions of \(x_1+x_2+x_3+x_4=10\) with
\(0\le x_i\le4\).
\end{exercise}
\begin{exercise}[\problemlevel{Bridge}]\label{c:ex12}
Prove \(\sum_{k=0}^n\binom nk^2=\binom{2n}{n}\) without expanding
factorials. Then count \(\sum_{k=0}^n k\binom nk^2\).
\end{exercise}
\begin{exercise}[\problemlevel{Bridge}]\label{c:ex13}
For \(n\ge1\), prove that \(D_n\) is even exactly when \(n\) is odd.
Use a recurrence or a parity argument.
\end{exercise}
\begin{exercise}[\problemlevel{Bridge}]\label{c:ex14}
Let \(\mathcal F\) be a family of distinct \(3\)-element subsets of
\([n]\), and suppose every pair of elements belongs to at most one
member. Prove \(3|\mathcal F|\le\binom n2\).
\end{exercise}
\begin{exercise}[\problemlevel{Challenge}]\label{c:ex15}
There are \(n\) married couples. Count the \(k\)-person committees
containing no complete couple. Give a direct formula and an
inclusion--exclusion formula, and prove that they agree by counting.
\end{exercise}
\begin{exercise}[\problemlevel{Challenge}]\label{c:ex16}
For \(n\ge1\), evaluate
\(\sum_{k=0}^n(-1)^k\binom nk^2\).
Use a sign-reversing involution on \(n\)-subsets of two labelled
\(n\)-element rows, rather than assuming a coefficient identity.
\end{exercise}

\chapter{Recurrences, Paths, and Generating Functions}
\label{c:recurrences}

\section{Remove the last decision}

A recurrence describes a large object by smaller objects of the same
kind. Its usefulness depends on a clean decomposition: the cases must be
disjoint, every object must appear, and reconstruction must be unique.
Initial conditions are part of the statement. A recurrence without
enough initial values generally does not determine a sequence.

\begin{example}[Dominoes on a strip]
Let \(t_n\) count tilings of a \(2\)-by-\(n\) board by \(1\)-by-\(2\)
dominoes. Then \(t_0=t_1=1\) and
\(t_n=t_{n-1}+t_{n-2}\) for \(n\ge2\).
\end{example}
\begin{solution}
Inspect the leftmost square in the top row. A vertical domino removes a
whole first column. A horizontal domino forces another horizontal domino
below it, removing two columns. The remaining board can be tiled
arbitrarily, and the two cases do not overlap.
\end{solution}
\begin{center}
\begin{tikzpicture}[scale=.55]
\draw (0,0) rectangle (1,2);
\draw[gray!40] (1,0) grid (5,2);
\draw[main] (0,0) rectangle (1,2);
\node[below] at (2.5,-.3) {one vertical domino};
\begin{scope}[xshift=7cm]
\draw[gray!40] (0,0) grid (5,2);
\draw[main] (0,0) rectangle (2,1);
\draw[main] (0,1) rectangle (2,2);
\node[below] at (2.5,-.3) {two horizontal dominoes};
\end{scope}
\end{tikzpicture}
\end{center}

\begin{example}[Words that avoid adjacent ones]
Let \(b_n\) count binary words of length \(n\) without adjacent ones.
Then \(b_0=1,b_1=2\), and \(b_n=b_{n-1}+b_{n-2}\).
\end{example}
\begin{solution}
A word ending in zero comes from any legal word of length \(n-1\).
A word ending in one, when \(n\ge2\), ends in \(01\), which follows
any legal word of length \(n-2\).
Although this is the same recurrence as before, the initial values differ:
\(b_n=t_{n+1}\). The recurrence alone would not reveal that shift.
\end{solution}

\section{Catalan structures and reflection}

A \emph{Dyck path} of semilength \(n\) is a word with \(n\) up-steps
\((1,1)\) and \(n\) down-steps \((1,-1)\), whose height never becomes
negative. Let \(C_n\) count these paths and put \(C_0=1\).
Think of an up-step as an opening parenthesis and a down-step as a closing
parenthesis. Nonnegative height means a closing parenthesis never appears
before its matching opening parenthesis.

\begin{theorem}[Catalan recurrence and formula]\label{c:catalan}
For \(n\ge1\),
\[
 C_n=\sum_{i=0}^{n-1}C_iC_{n-1-i}
     =\frac1{n+1}\binom{2n}{n}.
\]
\end{theorem}
\begin{proof}
For the recurrence, split a nonempty Dyck path at its first return to
height zero. It has the unique form \(U\,P\,D\,Q\), where \(P,Q\)
are Dyck paths of semilengths \(i,n-1-i\).

For the formula, there are \(\binom{2n}{n}\) unrestricted step words.
A bad one first reaches height \(-1\). Interchange up- and down-steps
up to and including that first visit. The image has \(n+1\) up-steps
and \(n-1\) down-steps. Conversely, a word ending at height \(2\)
has a first visit to height \(1\); reflecting that prefix recovers
a unique bad word. There are \(\binom{2n}{n-1}\) images, so
\[
 C_n=\binom{2n}{n}-\binom{2n}{n-1}
     =\frac1{n+1}\binom{2n}{n}.
 \]
\end{proof}
\strategy Reflection turns a global restriction into a local first
violation. Choose the \emph{first} violation so the transformation has
a recoverable stopping point. Reflecting an arbitrary negative excursion
usually destroys uniqueness.

\begin{example}[Triangulating a polygon]
The number of triangulations of a convex \((n+2)\)-gon is \(C_n\).
\end{example}
\begin{solution}
Fix a boundary edge. In any triangulation it belongs to one triangle.
The third vertex splits the remaining polygon into two smaller convex
polygons. If their parameters are \(i\) and \(n-1-i\), their
triangulations can be chosen independently. The empty side contributes
one. Thus the counts have the Catalan recurrence and initial value,
which determine them uniquely by induction.
\end{solution}

\section{A series as a bookkeeping device}

The \emph{ordinary generating function} of a sequence \(a_0,a_1,\ldots\)
is the formal expression \(A(x)=\sum_{n\ge0}a_nx^n\).
Write \([x^n]A(x)=a_n\) for its \(n\)th coefficient.
Here \(x\) need not be a real number. Equality means equality of every
coefficient, and multiplication is defined by
\[
 [x^n](A(x)B(x))=\sum_{i=0}^na_i b_{n-i}.
\]
This sum is finite. We do not need analytic convergence to justify it.
An infinite product is legitimate when the coefficient of each fixed
power depends on only finitely many factors.

\intuition Coefficients count objects, exponents record size, addition
means a choice of class, and multiplication means assembling two
independent components. A recurrence with a convolution often becomes
an algebraic equation because multiplication already performs that
convolution.

\begin{example}[The strip generating function]
For the strip counts above,
\[
 T(x)=\sum_{n\ge0}t_nx^n=\frac1{1-x-x^2}.
\]
\end{example}
\begin{solution}
Multiply the recurrence by \(x^n\) and sum for \(n\ge2\):
\[
 T(x)-1-x=x(T(x)-1)+x^2T(x).
\]
Rearranging gives \((1-x-x^2)T(x)=1\).
The denominator has constant term one, so it has a unique formal inverse:
its coefficients can be solved one at a time. No division by a possibly
zero numerical value of \(x\) occurs.
\end{solution}

\section{Partitions and products}

A partition of an integer is a multiset of positive parts whose sum
is that integer. With allowed part sizes \(1,\ldots,m\), the generating
function is
\[
 \prod_{j=1}^m(1+x^j+x^{2j}+\cdots)
 =\prod_{j=1}^m\frac1{1-x^j}.
\]
Choosing \(x^{rj}\) from factor \(j\) means using \(r\) copies of part
\(j\). If parts are distinct, replace the factor by \(1+x^j\).
If parts are ordered, this product is the wrong model: ordering changes
the object, so one needs a word or a recurrence instead.

\recognition Use generating functions when size is additive and local
choices are independent. For a modest numerical count, a recurrence may
be clearer. A clever bijection often explains more than a series
calculation; the two viewpoints should reinforce each other.

\section{Exercises}
\begin{exercise}[\problemlevel{Foundation}]\label{c:ex17}
Derive and use a recurrence for compositions of \(n\) into parts \(1\)
and \(2\), where the order of parts matters.
\end{exercise}
\begin{exercise}[\problemlevel{Foundation}]\label{c:ex18}
Let \(u_n\) count binary words of length \(n\) containing no \(111\).
Find initial values and a recurrence valid for all \(n\ge3\).
\end{exercise}
\begin{exercise}[\problemlevel{Bridge}]\label{c:ex19}
Prove that the number of balanced parenthesis words with \(n\) pairs
is \(C_n\), explaining the matching rule rather than merely naming Dyck paths.
\end{exercise}
\begin{exercise}[\problemlevel{Bridge}]\label{c:ex20}
Show that \(C(x)=\sum_{n\ge0}C_nx^n\) satisfies \(C(x)=1+xC(x)^2\).
Use this equation to compute \(C_1,\ldots,C_5\).
\end{exercise}
\begin{exercise}[\problemlevel{Bridge}]\label{c:ex21}
Give a generating function for the number of solutions of
\(a+2b+3c=n\) in nonnegative integers. Determine its coefficient at
\(x^8\) directly.
\end{exercise}
\begin{exercise}[\problemlevel{Bridge}]\label{c:ex22}
For \(r\ge s\ge0\), count words with \(r\) up-steps and \(s\)
down-steps whose height is never negative.
\end{exercise}
\begin{exercise}[\problemlevel{Challenge}]\label{c:ex23}
For \(n\ge3\), count the binary colorings of \(n\) labelled positions
on a circle that have exactly \(k\) ones, with no adjacent ones.
Rotations are distinguished. Include \(k=0\) and impossible cases.
\end{exercise}
\begin{exercise}[\problemlevel{Challenge}]\label{c:ex24}
Prove that partitions of any integer into distinct parts are equinumerous
with partitions into odd parts. Give a generating-function proof and
a reversible map using binary expansion.
\end{exercise}

\chapter{Pigeonholes, Prefixes, and Hidden Order}
\label{c:pigeonhole}

\section{Choose useful boxes}

The pigeonhole principle says that placing more than \(m\) objects in
\(m\) boxes forces a box to contain at least two objects.
More generally, \(N\) objects in \(m\ge1\) boxes force a box of size at
least \(\lceil N/m\rceil\). Its proof is only a sum: if each box has
at most \(q-1\) objects, the total is at most \(m(q-1)\).
The mathematical invention lies in the boxes.

\begin{example}[Equal remainders]
Among \(n+1\) integers, two have difference divisible by \(n\).
\end{example}
\begin{solution}
Put each integer into its residue class modulo \(n\).
Two share a class, so their difference is a multiple of \(n\).
This is true for negative integers as well: residue classes concern
congruence, not the sign of a written remainder.
\end{solution}

\begin{example}[A divisible consecutive sum]
From any sequence \(a_1,\ldots,a_n\) of integers, a nonempty
consecutive block has sum divisible by \(n\).
\end{example}
\begin{solution}
Introduce prefix sums \(s_0=0\) and
\(s_j=a_1+\cdots+a_j\) for \(1\le j\le n\).
The \(n+1\) prefixes occupy \(n\) residue classes.
If \(s_i\equiv s_j\pmod n\), \(i<j\), then
\(a_{i+1}+\cdots+a_j=s_j-s_i\) is the required sum.
\end{solution}
\strategy A consecutive block is the difference of two prefixes.
Whenever a statement asks for a consecutive segment, test whether
prefix sums, products, or parity vectors turn it into a pairwise
coincidence.

\section{Factor out the part that can vary}

\begin{example}[A divisibility pair]
Any \(n+1\) distinct members of \([2n]\) contain two of which one
divides the other.
\end{example}
\begin{solution}
Write each number uniquely as \(2^rq\), where \(q\) is odd.
There are only \(n\) possible odd parts \(q\in\{1,3,\ldots,2n-1\}\).
Two chosen numbers have the same \(q\), and their powers of two are
different; the smaller power divides the larger.
\end{solution}
\intuition Boxing by the actual number is useless. Boxing by odd part
forgets precisely the information that is allowed to change in a
divisibility chain. This is a general design principle: retain the
obstruction, discard the harmless variation.

\section{Label each object by the best structure ending there}

\begin{theorem}[Increasing or decreasing subsequences]\label{c:es}
Among \((r-1)(s-1)+1\) distinct real numbers in a sequence, there is an
increasing subsequence of length \(r\) or a decreasing subsequence of
length \(s\), where \(r,s\ge2\).
\end{theorem}
\begin{proof}
At position \(i\), let \(p_i\) be the length of a longest increasing
subsequence ending there, and \(q_i\) the analogous decreasing length.
If the claimed subsequences do not exist, the pair \((p_i,q_i)\)
belongs to \([r-1]\times[s-1]\), a set of \((r-1)(s-1)\) boxes.
The labels are distinct: if \(i<j\), then either \(a_i<a_j\), giving
\(p_j>p_i\), or \(a_i>a_j\), giving \(q_j>q_i\).
Pigeonhole gives a contradiction.
\end{proof}
\pitfall A subsequence keeps the original order but need not occupy
consecutive positions. Sorting the sequence would destroy the problem.
Distinctness is also used: equal terms extend neither a strictly
increasing nor a strictly decreasing subsequence.

\begin{example}[The bound cannot be lowered]
There is a sequence of \((r-1)(s-1)\) distinct numbers with neither
of the two subsequences in Theorem~\ref{c:es}.
\end{example}
\begin{solution}
Use \(r-1\) blocks of \(s-1\) numbers. Make each block decrease,
and make every number of a later block larger than every number of an
earlier block. An increasing subsequence takes at most one number from
each block, while a decreasing one remains within a single block.
\end{solution}

\section{Geometric and simultaneous pigeonholes}

Partitioning a square into smaller squares forces two points to be
close. Coloring integer lattice points by coordinate parities forces
an integral midpoint. In both cases boxes are vectors rather than
single numbers. If \(d\) independent coordinates each have \(m\)
possible residues, the label space has \(m^d\) possibilities.

For an irrational \(\alpha\), box the fractional parts
\(\{0\alpha\},\ldots,\{N\alpha\}\), where
\(\{x\}=x-\lfloor x\rfloor\), into \(N\) half-open intervals of
length \(1/N\). Two lie in one interval. Subtracting gives integers
\(p,q\), \(1\le q\le N\), with \(|q\alpha-p|<1/N\).
The half-open convention assigns boundary points uniquely.

\section{Exercises}
\begin{exercise}[\problemlevel{Foundation}]\label{c:ex25}
Prove that among six integers two differ by a multiple of five.
Show that five integers do not suffice.
\end{exercise}
\begin{exercise}[\problemlevel{Foundation}]\label{c:ex26}
Among five points with integer coordinates in the plane, prove that
two have a midpoint with integer coordinates.
\end{exercise}
\begin{exercise}[\problemlevel{Bridge}]\label{c:ex27}
Among \(n\) integers, show that some nonempty subset has sum
divisible by \(n\). Explain why the stronger consecutive version holds.
\end{exercise}
\begin{exercise}[\problemlevel{Bridge}]\label{c:ex28}
From \(n+1\) distinct members of \([2n]\), show that two are coprime.
Give a different boxing from the divisibility example.
\end{exercise}
\begin{exercise}[\problemlevel{Bridge}]\label{c:ex29}
Give 16 distinct numbers with no increasing or decreasing subsequence
of length five, and prove that 17 distinct numbers always force one.
\end{exercise}
\begin{exercise}[\problemlevel{Bridge}]\label{c:ex30}
Real numbers \(a_1,\ldots,a_n\) sum to zero. Prove that some cyclic
rotation of the sequence has every partial sum nonnegative.
\end{exercise}
\begin{exercise}[\problemlevel{Challenge}]\label{c:ex31}
Given real numbers \(\alpha_1,\ldots,\alpha_d\) and an integer
\(M\ge1\), prove there are integers \(q,p_1,\ldots,p_d\) such that
\(1\le q\le M^d\) and \(|q\alpha_i-p_i|<1/M\) for every \(i\).
\end{exercise}
\begin{exercise}[\problemlevel{Challenge}]\label{c:ex32}
Prove that an irrational real number \(\alpha\) has infinitely many
different rational approximations \(p/q\), \(q>0\), satisfying
\(|\alpha-p/q|<1/q^2\). Explain why merely applying pigeonhole once
does not establish infinitude.
\end{exercise}

\chapter{Invariants, Colorings, and Termination}
\label{c:invariants}

\section{Ask what a move cannot change}

An \emph{invariant} is a quantity preserved by every permitted move.
If the initial and target configurations have different values, the
target is unreachable. Examples include parity, a residue, a greatest
common divisor, a signed sum, and a permutation's sign.
An invariant is a necessary condition for reachability; it need not
classify all reachable states.

\begin{example}[Flipping two coins]
Initially an odd number of coins show heads. Each move flips exactly
two coins. Can all coins become tails?
\end{example}
\begin{solution}
The number of heads changes by \(2,0\), or \(-2\), so its parity is
invariant. All tails has zero heads, an even number; it is impossible.
Counting the number of moves would not help without a preserved feature.
\end{solution}

\begin{example}[A nonlinear-looking invariant]
Repeatedly replace two real numbers \(a,b\) on a board by
\(a+b+ab\). When one number remains, determine it from the initial
numbers \(x_1,\ldots,x_n\).
\end{example}
\begin{solution}
The identity \(1+a+b+ab=(1+a)(1+b)\) suggests tracking the product
of one plus every board entry. It is invariant. Therefore the last
number is \(\prod_{i=1}^n(1+x_i)-1\), regardless of the choices.
The argument uses no division, so entries equal to \(-1\) are allowed.
\end{solution}
\recognition For a replacement \(a,b\mapsto F(a,b)\), search for a
transform \(\phi\) with
\(\phi(F(a,b))=\phi(a)+\phi(b)\) or
\(\phi(F(a,b))=\phi(a)\phi(b)\). Simple algebra can reveal a hidden
associative operation.

\section{Colorings are weights}

Color a chessboard black and white by the parity of the sum of its
row and column indices. A domino covers one square of each color.
Thus the difference between the two color counts is unchanged when
dominoes are removed.

\begin{example}[The mutilated board]
An \(8\)-by-\(8\) board with two diagonally opposite corner squares
removed cannot be tiled by dominoes.
\end{example}
\begin{solution}
The removed corners have the same color. The remaining board has
30 squares of that color and 32 of the other. Any domino tiling would
use equal counts. Equal area alone, \(62=31\cdot2\), misses the
obstruction.
\end{solution}

\begin{center}
\begin{tikzpicture}[scale=.34]
 \foreach \i in {0,...,7}{
  \foreach \j in {0,...,7}{
   \pgfmathtruncatemacro{\p}{mod(\i+\j,2)}
   \ifnum\p=0 \fill[accent!28] (\i,\j) rectangle ++(1,1);\fi
  }}
 \fill[white] (0,0) rectangle (1,1);
 \fill[white] (7,7) rectangle (8,8);
 \draw[gray] (0,0) grid (8,8);
 \draw[red,thick] (0,0)--(1,1) (0,1)--(1,0);
 \draw[red,thick] (7,7)--(8,8) (7,8)--(8,7);
\end{tikzpicture}
\end{center}
The colors are not decoration: assigning weights \(+1,-1\) makes each
domino have total weight zero, while the whole region has nonzero weight.

More general weights can be integers, residue classes, or complex
numbers. To disprove a tiling, find a weighting for which every allowed
tile has a fixed total, but the region has an incompatible total.
To prove a tiling exists, provide an actual arrangement or a constructive
induction. A coloring obstruction alone never supplies existence.

\section{Permutation parity}

For a permutation \(a_1,\ldots,a_n\) of \([n]\), an \emph{inversion}
is a pair \(i<j\) with \(a_i>a_j\). Its number is \(I\).
Swapping adjacent entries changes \(I\) by exactly one: only their
mutual order changes. Swapping entries at positions \(i<j\) can be
performed by \(2(j-i)-1\) adjacent swaps, so every transposition
reverses the parity of \(I\).

Define the sign as \(\operatorname{sgn}(\pi)=(-1)^I\).
Consequently, a sequence of transpositions that returns to the identity
has even length. This proves parity is well defined without choosing
one preferred factorization into transpositions.
If a puzzle move is a three-cycle, it is two transpositions and hence
preserves sign. If a move also changes a blank's position, the blank
must be included in the encoding; an unproved analogy with ordinary
permutations is unsafe.

\section{Monovariants need a stopping reason}

A \emph{monovariant} changes in one direction under every move.
Strict decrease alone does not imply termination: \(1,1/2,1/3,\ldots\)
decreases forever. Two safe forms are:
\begin{itemize}
\item a nonnegative integer decreases strictly at every move;
\item a quantity changes strictly on a finite set of states.
\end{itemize}
More generally, any positive lower bound on the size of each decrease
together with a lower bound on the quantity suffices.

The inversion count gives an algorithmic example. If a permutation is
not sorted, it has an adjacent inverted pair; swap it. The inversion
count falls by one and cannot become negative. At termination there
are no adjacent inversions, so the list is increasing.
The stopping condition must be connected to the goal, not merely
declared a success.

\section{Exercises}
\begin{exercise}[\problemlevel{Foundation}]\label{c:ex33}
Which pairs of removed squares from a chessboard are ruled out by
the ordinary black--white coloring? Explain why removing opposite
colors does not by itself prove tileability.
\end{exercise}
\begin{exercise}[\problemlevel{Foundation}]\label{c:ex34}
On a row of \(n\) coins, a move flips two adjacent coins.
Classify the configurations reachable from all tails.
\end{exercise}
\begin{exercise}[\problemlevel{Bridge}]\label{c:ex35}
Starting with \(1,2,\ldots,n\), replace \(a,b\) by \(a+b+ab\)
until one number remains. Determine that number.
\end{exercise}
\begin{exercise}[\problemlevel{Bridge}]\label{c:ex36}
Starting from the integer triple \((3,5,7)\), a move replaces one
entry by that entry plus or minus another. Show that \((2,2,2)\)
is unreachable.
\end{exercise}
\begin{exercise}[\problemlevel{Bridge}]\label{c:ex37}
Prove that a knight on a chessboard cannot return to its initial square
after an odd number of moves. A knight changes coordinates by
\((\pm1,\pm2)\) or \((\pm2,\pm1)\).
\end{exercise}
\begin{exercise}[\problemlevel{Bridge}]\label{c:ex38}
Show that the minimum number of adjacent swaps needed to sort a
permutation equals its inversion count.
\end{exercise}
\begin{exercise}[\problemlevel{Challenge}]\label{c:ex39}
An \(8\)-by-\(8\) board has the corner with coordinates \((0,0)\)
removed; squares have coordinates \((i,j)\), \(0\le i,j<8\).
Show that the remaining board cannot be tiled by straight
\(1\)-by-\(3\) tiles in either orientation.
\end{exercise}
\begin{exercise}[\problemlevel{Challenge}]\label{c:ex40}
For an integer \(k\ge2\), prove that an \(m\)-by-\(n\) rectangle
can be tiled by straight \(1\)-by-\(k\) tiles in either orientation
if and only if \(k\mid m\) or \(k\mid n\).
You may use a primitive complex \(k\)th root of unity.
\end{exercise}

\chapter{Extremal Choices and Constructive Algorithms}
\label{c:extremal}

\section{Choose the best object before modifying it}

To prove existence, one can construct an object step by step, or choose
an optimal object from a finite collection and show it already has the
desired property. The second approach often makes a proposed
improvement impossible and thereby forces structure.

\begin{example}[Make every row and column nonnegative]
In a finite real matrix, one may reverse the signs of all entries of
any row or column. Show that some resulting matrix has every row sum
and column sum nonnegative.
\end{example}
\begin{solution}
There are finitely many results: choose which rows and columns to flip.
Choose one with maximum total sum. If a row had negative sum \(s\),
flipping it would increase the total by \(-2s>0\), impossible.
The same applies to columns.
\end{solution}
\intuition The goal concerns every local part. Optimize a global
quantity that a bad local part would improve. This converts many
simultaneous requirements into one maximality statement.
An arbitrary sequence of improving flips also terminates, since the
state set is finite and its total sum strictly increases.

\begin{example}[A noncrossing red--blue matching]
There are \(n\) red and \(n\) blue distinct points in the plane, with
no three collinear. Match each red point to a blue point so that the
straight connecting segments do not cross.
\end{example}
\begin{solution}
Among the \(n!\) color-respecting matchings choose one of least total
Euclidean length. If segments \(R_1B_1\) and \(R_2B_2\) cross at an
interior point \(X\), replace them by \(R_1B_2,R_2B_1\).
The strict triangle inequalities through \(X\) give
\[
 |R_1B_2|+|R_2B_1|
 <
 |R_1X|+|XB_2|+|R_2X|+|XB_1|,
 \]
which is the old total. This contradicts minimality.
No three collinear ensures the triangle inequalities are strict.
\end{solution}
\pitfall One must preserve the class of legal objects during an
exchange. Reconnecting the two red points to each other would shorten
segments in some pictures but would destroy the matching's color rule.

\section{Inductive insertion}

A tournament is an orientation of a complete graph: for every two
different vertices, exactly one arrow joins them. A directed path
visits different vertices, following arrows.

\begin{theorem}[Every tournament has a directed spanning path]
Every finite tournament has a directed path through all its vertices.
\end{theorem}
\begin{proof}
Induct on the number of vertices. Remove \(v\), and take a directed
spanning path \(x_1\to\cdots\to x_{n-1}\) of the remaining tournament.
If \(v\to x_1\), insert \(v\) first. Otherwise, if some \(x_i\)
is beaten by \(v\), choose the first such \(i\).
Then \(x_{i-1}\to v\to x_i\), so insert \(v\) there.
If none is beaten by \(v\), all \(x_i\to v\), so append \(v\).
\end{proof}
\strategy The inductive object need not retain every desired feature
after arbitrary insertion. Search for an insertion point justified by
a first transition from one type of edge to another.

\begin{example}[Greedy scheduling]
From finitely many nonempty closed intervals on a line, choose as many
as possible that are pairwise disjoint. The following greedy algorithm
is optimal: select an interval with smallest right endpoint, discard
all intervals meeting it, and repeat.
\end{example}
\begin{solution}
Let the greedy first interval be \(I\), ending at \(b\). In an optimal
disjoint family, the first interval \(J\) ends at \(d\ge b\).
Replace \(J\) by \(I\). Every later interval starts after \(d\), and
therefore after \(b\), so no new intersection is created. Thus some
optimal family starts with \(I\). Remove intervals meeting \(I\)
and apply the same argument inductively to the remaining instance.
\end{solution}
For closed intervals, sharing an endpoint counts as intersection,
which explains the strict word \emph{after}. For half-open intervals
the compatibility convention changes accordingly.

\section{Greedy, exchange, and local improvement}

A greedy rule makes the apparently best current choice. Such rules
require proof; choosing the largest available interval above can fail.
An exchange argument proves that some optimum can be modified to
contain the greedy choice. A local-improvement argument shows that
an object with no improving local move has the needed global property.
It need not be globally optimal for an unrelated objective.

When constructing an extremal example, distinguish the upper bound
from its achievement. An inequality proves that no object is better;
a construction proves that the bound is attainable. Both are needed
for a maximum or minimum problem.

\section{Exercises}
\begin{exercise}[\problemlevel{Foundation}]\label{c:ex41}
Prove that every finite graph has a partition of its vertices into two
parts such that each vertex has at least as many neighbors in the
other part as in its own part.
\end{exercise}
\begin{exercise}[\problemlevel{Bridge}]\label{c:ex42}
Prove that in every finite tournament some vertex can reach every
other vertex by a directed path of length at most two.
\end{exercise}
\begin{exercise}[\problemlevel{Challenge}]\label{c:ex43}
Let \(S\) be a finite set of noncollinear points in the real plane.
Prove that some line passes through exactly two points of \(S\).
\end{exercise}
\begin{exercise}[\problemlevel{Bridge}]\label{c:ex44}
A connected graph has positive edge weights. Prove that an edge of
strictly smaller weight than every other edge crossing some partition
of the vertices into two nonempty parts belongs to every minimum-weight
spanning tree. Define a spanning tree as a connected acyclic subgraph
using all vertices.
\end{exercise}
\begin{exercise}[\problemlevel{Bridge}]\label{c:ex45}
Give three intervals for which repeatedly choosing the longest available
interval produces fewer pairwise disjoint intervals than possible.
Explain why the earliest-finish exchange proof does not apply.
\end{exercise}
\begin{exercise}[\problemlevel{Challenge}]\label{c:ex46}
For \(n\ge3\) points in the plane, no three collinear, prove that they
can be joined in some cyclic order to form a polygon with no crossing
sides. The polygon need not be convex.
\end{exercise}
\begin{exercise}[\problemlevel{Bridge}]\label{c:ex47}
There are \(N\ge1\) indistinguishable tokens in finitely many boxes.
Whenever at least two boxes are nonempty, move one token from a box
with smaller or equal occupancy to a different box with larger or equal
occupancy. Prove that every such process terminates with all tokens
in one box.
\end{exercise}
\begin{exercise}[\problemlevel{Challenge}]\label{c:ex48}
A tournament is strongly connected if a directed path exists from every
vertex to every other. Prove that a strongly connected tournament with
at least three vertices has a directed cycle through all vertices.
\end{exercise}

\chapter{Graphs: Turn Relationships into Structure}
\label{c:graphs}

\section{The language of a graph}

Unless stated otherwise, a graph here is finite, undirected, and simple:
an edge joins two different vertices, with at most one edge per pair.
Write \(V(G)\) for its vertex set, \(E(G)\) for its edge set,
\(n=|V(G)|\), and \(e=|E(G)|\).
Vertices joined by an edge are neighbors. The degree \(d(v)\) counts a
vertex's neighbors, and \(\Delta\) is the maximum degree.
A graph is connected if every pair of vertices can be joined by a path.
A walk may repeat edges and vertices; a trail repeats no edge; a path
repeats no vertex. A cycle returns to its start without other repeated
vertices and has at least three edges.

\intuition A graph forgets irrelevant measurements while preserving
which objects are related. People knowing one another, cells sharing a
border, and sets having nonempty intersection can all produce graphs.
The model is useful only if the graph statement translates back into the
original problem.

\begin{theorem}[The handshaking identity]
\(\sum_{v\in V(G)}d(v)=2e\). In particular, the number of odd-degree
vertices is even.
\end{theorem}
\begin{proof}
Count incidences of a vertex with an edge containing it. Every edge
has exactly two ends. Reducing the equality modulo two gives the parity
assertion.
\end{proof}

\begin{example}[An average-degree argument]
A graph with \(e>kn/2\) has a vertex of degree at least \(k+1\).
\end{example}
\begin{solution}
Its average degree is \(2e/n>k\). If every degree were at most \(k\),
their average would also be at most \(k\).
The conclusion does not say every vertex has large degree.
\end{solution}

\section{Trees and Euler trails}

A tree is a connected graph containing no cycle. A forest is a graph
containing no cycle, with any number of connected components.

\begin{theorem}[Equivalent tree properties]\label{c:trees}
A tree on \(n\ge1\) vertices has \(n-1\) edges and a unique path
between every pair of vertices. Conversely, a connected graph with
\(n-1\) edges is a tree. A forest with \(c\) components has \(n-c\) edges.
\end{theorem}
\begin{proof}
Two different paths between the same vertices would contain a cycle,
so a tree path is unique. A longest path in a tree with \(n\ge2\)
has endpoints of degree one: an extra neighbor would extend the path
or create a cycle. Remove an endpoint and its edge. The remaining graph
is a tree, so induction gives \(n-1\) edges.
In any connected graph, repeatedly delete an edge of a cycle.
Connectivity survives, and eventually a spanning tree remains.
If the original graph already had \(n-1\) edges, none could have
been deleted. Apply the tree formula separately to each component for
the forest formula.
\end{proof}

\begin{theorem}[Euler's trail criterion]\label{c:eulertrail}
A connected graph with at least one edge has a closed trail using every
edge exactly once if and only if all degrees are even.
It has a trail using every edge with different endpoints if and only if
exactly two degrees are odd.
\end{theorem}
\begin{proof}
An internal visit to a vertex uses one entering and one leaving edge.
Thus a closed trail forces even degrees; an open one leaves exactly its
two endpoints odd.

For sufficiency in the even case, start at any vertex and follow unused
edges until no continuation exists. At every vertex other than the start,
used incident edges would have odd count if the trail stopped there,
leaving an unused edge because the full degree is even. Thus the trail
closes at its start. If some edges remain, connectivity gives a vertex
of the existing trail incident with an unused edge. Starting there,
construct another closed trail and splice it into the first.
Repeating uses all edges, since each splice consumes at least one new edge.
For two odd vertices \(u,v\), add a new vertex \(w\) and edges \(uw,wv\).
All degrees are now even. Remove these two consecutive edges from an
Euler circuit to obtain the desired open trail.
\end{proof}
\pitfall An Euler trail concerns edges; a Hamilton path concerns vertices.
Degree parity completely solves the former existence problem but does not
give a comparable criterion for the latter.

\section{Bipartiteness and coloring}

A proper coloring assigns colors to vertices so that adjacent vertices
have different colors. A graph is bipartite if its vertices can be split
into two parts with every edge between the parts.

\begin{theorem}\label{c:bipartite}
A graph is bipartite if and only if it has no odd cycle.
Every graph of maximum degree \(\Delta\) can be properly colored with
\(\Delta+1\) colors.
\end{theorem}
\begin{proof}
In a bipartite graph, a cycle alternates parts and has even length.
Conversely, in each connected component choose a root and color a vertex
by the parity of its distance from the root. If an edge joined two
equal-parity distances, their root paths plus the edge would give an odd
closed walk. An odd closed walk contains an odd cycle: if it repeats an
internal vertex, split it into two shorter closed walks, one of odd length,
and repeat. This contradicts the hypothesis.

For the coloring bound, order the vertices arbitrarily and color them
one at a time. At most \(\Delta\) colors are forbidden by already
colored neighbors, so one of \(\Delta+1\) colors remains.
\end{proof}

\begin{example}[Degeneracy improves greedy coloring]
If every nonempty subgraph has a vertex of degree at most \(d\) within
that subgraph, the graph can be colored with \(d+1\) colors.
\end{example}
\begin{solution}
Remove a vertex of degree at most \(d\), and repeat on what remains.
Then insert vertices in reverse removal order. Each has at most \(d\)
already colored neighbors. This is not the same as greedily coloring
in an arbitrary order: the removal procedure builds the useful order.
\end{solution}

\section{Planarity and faces}

A planar graph admits a drawing in the plane with edges meeting only at
common endpoints. Fix such a drawing. The connected regions left by its
edges are faces, including the unbounded region.
For a connected plane graph,
\[
 n-e+f=2.
\]
To see this, delete edges on cycles until a tree remains. Each deletion
merges two faces, reducing both \(e\) and \(f\) by one; a simple closed
polygon separates an inside from an outside. A tree has \(e=n-1\)
and one face, giving the formula. A dangling tree edge does not separate
faces, which also follows by successively removing leaves.

If the graph is simple and \(n\ge3\), every face boundary walk has length
at least three. Counting edge-sides gives \(3f\le2e\), even when a bridge
is traversed twice along one boundary. Consequently
\[
 e\le3n-6.
\]
If it is also bipartite, boundary lengths are at least four, so
\(e\le2n-4\). For a disconnected graph add noncrossing edges between
components first; the resulting connected simple graph has at least as
many edges, so the same upper bounds follow.

\begin{example}[Two impossible drawings]
The complete graph \(K_5\), with all ten edges between five vertices,
is not planar. Neither is the complete bipartite graph \(K_{3,3}\),
with all nine edges between two sets of three vertices.
\end{example}
\begin{solution}
For \(K_5\), the planar bound would give \(10\le3\cdot5-6=9\).
For bipartite \(K_{3,3}\), it would give \(9\le2\cdot6-4=8\).
Both are contradictions. These bounds can disprove planarity; satisfying
one does not by itself prove that a drawing exists.
\end{solution}

\section{Exercises}
\begin{exercise}[\problemlevel{Foundation}]\label{c:ex49}
Show that any simple graph with at least two vertices has two vertices
of the same degree.
\end{exercise}
\begin{exercise}[\problemlevel{Foundation}]\label{c:ex50}
Prove that a tree on at least two vertices has at least two leaves.
Can it have exactly two?
\end{exercise}
\begin{exercise}[\problemlevel{Bridge}]\label{c:ex51}
A connected graph has \(2k>0\) odd-degree vertices. Show that its
edges can be partitioned into \(k\) trails, allowing a trail to revisit
vertices but not edges.
\end{exercise}
\begin{exercise}[\problemlevel{Bridge}]\label{c:ex52}
Prove that every graph has a bipartite subgraph containing at least half
of its edges.
\end{exercise}
\begin{exercise}[\problemlevel{Bridge}]\label{c:ex53}
If every nonempty subgraph of \(G\) has average degree less than \(2d\), where
\(d\) is a positive integer, prove that \(G\) is \(2d\)-colorable.
Explain why \(d+1\) colors do not follow from that assumption alone.
\end{exercise}
\begin{exercise}[\problemlevel{Bridge}]\label{c:ex54}
Show that any finite simple planar graph contains a vertex of degree at
most five. Deduce a six-coloring theorem by an explicit induction.
\end{exercise}
\begin{exercise}[\problemlevel{Challenge}]\label{c:ex55}
A simple graph on \(n\ge3\) vertices has minimum degree at least \(n/2\).
Prove that it contains a Hamilton cycle, a cycle through every vertex.
\end{exercise}
\begin{exercise}[\problemlevel{Challenge}]\label{c:ex56}
In a tournament with outdegrees \(d_1,\ldots,d_n\), show that the
number of directed triangles is
\(\binom n3-\sum_i\binom{d_i}{2}\).
For odd \(n\), determine its largest possible value and construct a
tournament attaining it.
\end{exercise}

\chapter{Matching and Order}
\label{c:matching}

\section{A local condition for a global assignment}

In a bipartite graph with parts \(L,R\), a matching is a set of edges
with no shared vertex. It saturates \(L\) if every vertex of \(L\)
belongs to a matching edge. For \(S\subseteq L\), let \(N(S)\subseteq R\)
be all neighbors of vertices in \(S\).

\intuition Distinct representatives require enough options not just for
each individual but for every group. Three applicants can each have two
choices and still all compete for the same two jobs. Hall's condition
records exactly that collective obstruction.

\begin{theorem}[Hall's marriage theorem]\label{c:hall}
A finite bipartite graph has a matching saturating \(L\) if and only if
\(|N(S)|\ge|S|\) for every \(S\subseteq L\).
\end{theorem}
\begin{proof}
Necessity follows by looking at the different matched neighbors of \(S\).
For sufficiency, induct on \(|L|\), with the empty case immediate.
Suppose a nonempty proper subset \(S\) satisfies \(|N(S)|=|S|\).
The graph between \(S,N(S)\) satisfies Hall and has a matching by induction.
For \(T\subseteq L\setminus S\), Hall applied to \(S\cup T\) gives
\[
 |N(T)\setminus N(S)|\ge |T|.
 \]
Thus the graph between \(L\setminus S,R\setminus N(S)\) also satisfies
Hall and can be matched by induction. The two matchings are disjoint.

If no such proper \(S\) exists, every nonempty proper \(S\) has
\(|N(S)|\ge|S|+1\). Choose any \(x\in L\) and a neighbor \(y\);
one exists by Hall. Delete \(x,y\). Every nonempty subset of the
remaining left part was proper in \(L\), so losing \(y\) still leaves
at least as many neighbors as vertices. Induct and add edge \(xy\).
For \(|L|=1\), simply choose its neighbor.
\end{proof}

\begin{example}[A regular bipartite graph]
A finite bipartite graph with equal positive degree \(d\) at every
vertex has a perfect matching, one saturating both parts.
\end{example}
\begin{solution}
Count edges to obtain \(d|L|=d|R|\), so \(|L|=|R|\).
For \(S\subseteq L\), the \(d|S|\) incident edges all end in \(N(S)\),
whose vertices receive at most \(d|N(S)|\) such edges.
Thus Hall holds; a matching saturating \(L\) also saturates \(R\).
\end{solution}
\recognition Assignment, representatives, disjoint choices, and rook
placements often indicate bipartite matching. Define the two parts and
the allowed edges before testing Hall. The hard part is commonly proving
the neighborhood inequality from the problem's numerical hypotheses.

\section{Augmenting paths and a dual certificate}

Relative to a matching \(M\), an augmenting path starts and ends at
unmatched vertices and alternates edges outside and inside \(M\).
Switching membership along it gives a matching with one more edge.
Conversely, if a larger matching \(M'\) exists, the graph of edges
belonging to exactly one of \(M,M'\) is a union of alternating paths
and even cycles. Some path has one more \(M'\)-edge, so it augments \(M\).
Thus absence of an augmenting path certifies maximal \emph{size}.
A merely \emph{maximal} matching, to which no single edge can be added,
can still have an augmenting path and need not have maximum size.

A vertex cover is a set of vertices meeting every edge.
Every cover has at least as many vertices as any matching has edges.

\begin{theorem}[Konig's theorem]\label{c:konig}
In a finite bipartite graph the maximum matching size equals the
minimum vertex-cover size.
\end{theorem}
\begin{proof}
Take a maximum matching \(M\). From all unmatched left vertices,
follow edges outside \(M\) from left to right and edges of \(M\)
from right to left. Let \(Z_L,Z_R\) be the reachable vertices.
No unmatched right vertex is reachable, or an augmenting path would exist.
Then
\[
 (L\setminus Z_L)\cup Z_R
 \]
is a cover. Indeed, an edge from \(Z_L\) to \(R\setminus Z_R\)
would contradict reachability: a nonmatching edge could be followed;
for a matching edge its reachable left endpoint was reached through
that same matched right endpoint, unless it was an unmatched starting
vertex, which has no matching edge.

Each matching edge has both endpoints reachable or neither reachable.
Thus exactly one endpoint belongs to this cover. No unmatched vertex
belongs to it. The cover has \(|M|\) vertices, attaining the general
lower bound.
\end{proof}

\section{Chains and antichains}

A partially ordered set, or poset, is a set \(P\) with a reflexive,
antisymmetric, transitive relation \(\le\).
A chain is a subset in which every pair is comparable.
An antichain has no two different comparable elements.
The width is the largest antichain size.

\begin{theorem}[Dilworth's theorem]\label{c:dilworth}
In a finite poset, the smallest number of chains partitioning the
elements equals its width.
\end{theorem}
\begin{proof}
Make left and right copies of the \(n\) elements, joining \(x_L\)
to \(y_R\) when \(x<y\). A matching of size \(m\) makes directed links
\(x\to y\) on \(P\). Every vertex has at most one incoming and one
outgoing link, and no directed cycle exists because links strictly
increase the order. The links therefore form \(n-m\) disjoint chains.
Conversely, a partition into \(k\) chains gives a matching of \(n-k\)
edges by linking consecutive elements in each chain.
Take a maximum matching of size \(m\). Konig gives a cover
\(X_L\cup Y_R\) of size \(m\).
The elements outside \(X\cup Y\) form an antichain, because a comparable
pair would give an uncovered edge. There are at least \(n-m\) such
elements. Any antichain meets each of the \(n-m\) chains in at most
one element. Hence width and minimum chain count both equal \(n-m\).
\end{proof}

\begin{example}[Subsequences as chains]
For a sequence of different numbers, order positions by
\(i\prec j\) if \(i<j\) and \(a_i<a_j\).
Chains are increasing subsequences; antichains, read from left to right,
are decreasing subsequences. Dilworth therefore says that the minimum
number of increasing subsequences partitioning the sequence equals the
longest decreasing-subsequence length.
\end{example}
\begin{solution}
The relation is transitive and has no cycles, so adjoining equality makes
a poset. The asserted identifications follow directly from the two
inequalities. Apply the theorem. This is a structural strengthening of
the earlier pigeonhole argument, not a statement about consecutive runs.
\end{solution}

\section{Sperner's theorem and weighted chains}

Order the subsets of \([n]\) by inclusion. A maximal chain can be made
by taking a permutation and adding its entries one at a time.
An antichain \(\mathcal A\) meets such a chain at most once.
A particular \(k\)-set belongs to exactly \(k!(n-k)!\) of the \(n!\)
permutation chains. Counting pairs consisting of a permutation chain
and a member of \(\mathcal A\) on it gives the LYM inequality
\[
 \sum_{A\in\mathcal A}\frac1{\binom n{|A|}}\le1.
\]
Since the largest binomial coefficient is
\(\binom n{\lfloor n/2\rfloor}\), this implies
\[
 |\mathcal A|\le\binom n{\lfloor n/2\rfloor}.
\]
The bound is attained by all subsets of size \(\lfloor n/2\rfloor\).
The largest coefficient is central because
\(\binom n{k+1}/\binom nk=(n-k)/(k+1)\).

\begin{example}[An unexpectedly strong weighted bound]
If an inclusion-antichain consists of \(a_k\) sets of size \(k\), the
bound is \(\sum_k a_k/\binom nk\le1\), not merely a bound on
\(\sum_k a_k\).
\end{example}
\begin{solution}
Each set uses a proportion \(1/\binom nk\) of all permutation chains,
and no chain can serve two antichain members. A small or large set
uses more chains than a middle-size set. This explains why concentrating
at the middle level is efficient, rather than simply memorizing the
central binomial coefficient.
\end{solution}

\section{Stable assignments}

Matching with preferences is a different question from maximum cardinality.
Suppose \(n\) applicants and \(n\) colleges each strictly rank all members
of the opposite side, and each college has one place.
A matching is stable if no unmatched applicant--college pair prefer
one another to their assigned partners.

The deferred-acceptance algorithm lets a free applicant propose to their
highest-ranked college not yet tried. Each college retains its favorite
proposal so far and rejects the others, possibly releasing a former
tentative partner. There are at most \(n^2\) proposals.
At termination everyone is matched: an exhausted unmatched applicant
would imply every college is holding someone else, accounting for all
\(n\) applicants already. If an applicant prefers a college to the
final partner, that college rejected the applicant earlier and only
improved its held partner afterward. It cannot prefer that applicant
to its final partner. Thus the final matching is stable.
Strict complete preferences are important assumptions in this proof.

\section{Exercises}
\begin{exercise}[\problemlevel{Foundation}]\label{c:ex57}
Three applicants have acceptable jobs \(\{1,2\},\{1,2\},\{2,3\}\).
Find distinct jobs for them. Change just the third list to make the
assignment impossible, and identify the violating Hall set.
\end{exercise}
\begin{exercise}[\problemlevel{Bridge}]\label{c:ex58}
Prove that the edges of any \(d\)-regular bipartite graph, \(d\ge1\),
can be partitioned into \(d\) perfect matchings.
\end{exercise}
\begin{exercise}[\problemlevel{Bridge}]\label{c:ex59}
A bipartite graph has \(n\) vertices in each part and every vertex
has degree at least \(n/2\). Prove that it has a perfect matching.
\end{exercise}
\begin{exercise}[\problemlevel{Bridge}]\label{c:ex60}
Certain squares of a chessboard are allowed. Show that the maximum number
of nonattacking rooks on allowed squares equals the minimum number of
rows and columns whose union contains every allowed square.
\end{exercise}
\begin{exercise}[\problemlevel{Bridge}]\label{c:ex61}
Order \([2n]\) by divisibility. Determine its width, giving both a
largest antichain and a partition into the same number of chains.
\end{exercise}
\begin{exercise}[\problemlevel{Challenge}]\label{c:ex62}
A family of subsets of \([n]\) contains no chain of \(r+1\) distinct
sets, where \(1\le r\le n+1\). Prove that its size is at most the sum
of the \(r\) largest numbers among \(\binom n0,\ldots,\binom nn\),
and show that the bound is attainable.
\end{exercise}
\begin{exercise}[\problemlevel{Bridge}]\label{c:ex63}
Prove that the least number of antichains partitioning a finite poset
equals its longest-chain length. Deduce that a sequence with no
increasing subsequence of length \(r+1\) can be partitioned into \(r\)
decreasing subsequences, assuming its entries are distinct.
\end{exercise}
\begin{exercise}[\problemlevel{Challenge}]\label{c:ex64}
Prove that in applicant-proposing deferred acceptance, every applicant
receives their favorite partner among all stable matchings.
Do not assume that a rejected pair can never be stable without proving it.
\end{exercise}

\chapter{Extremal Bounds and the Probabilistic Method}
\label{c:probability}

\section{Forbid a small pattern, bound the whole graph}

An extremal problem asks how large a structure can be while avoiding
a specified configuration. An answer needs an upper bound and a
construction. Counting the forbidden pattern in two ways is often more
effective than trying to inspect every possible graph.

\begin{theorem}[Mantel's bound]\label{c:mantel}
A triangle-free simple graph on \(n\) vertices has at most
\(\lfloor n^2/4\rfloor\) edges. The complete bipartite graph with
parts of sizes \(\lfloor n/2\rfloor,\lceil n/2\rceil\) attains the bound.
\end{theorem}
\begin{proof}
If \(uv\) is an edge, the neighborhoods of \(u,v\) are disjoint,
since a shared neighbor would make a triangle. Hence \(d(u)+d(v)\le n\).
Summing over edges gives
\[
 \sum_v d(v)^2=\sum_{uv\in E}(d(u)+d(v))\le ne.
 \]
Cauchy's inequality gives
\(\sum_v d(v)^2\ge(\sum_vd(v))^2/n=4e^2/n\).
If \(e>0\), combine and divide to get \(e\le n^2/4\);
the edgeless case is immediate. Since \(e\) is an integer, take the
floor. The stated bipartite graph has exactly
\(\lfloor n/2\rfloor\lceil n/2\rceil=\lfloor n^2/4\rfloor\)
edges and no triangle.
\end{proof}

\begin{theorem}[A Turan bound]\label{c:turan}
For \(r\ge1\), a graph with no complete subgraph on \(r+1\) vertices
has at most \(\frac{r-1}{2r}n^2\) edges.
If \(r\mid n\), the complete \(r\)-partite graph with equal parts
attains this real-valued bound.
\end{theorem}
\begin{proof}
The case \(r=1\) is immediate. Fix \(r\ge2\), put
\(a=(r-1)/(2r)\), and induct on \(n\), including \(n=0\).
Choose a largest clique \(C\) of size \(k\le r\).
Every vertex outside \(C\) has at most \(k-1\) neighbors in \(C\),
or it would enlarge the clique. Thus
\(\sum_{v\in C}d(v)\le(k-1)n\).
The number of edges incident with \(C\) is that degree sum minus
\(\binom k2\). By induction on the remaining graph,
\[
 \begin{split}
 e&\le a(n-k)^2+(k-1)n-\frac{k(k-1)}2\\
  &=an^2-(1-k/r)(n-k/2)\le an^2.
 \end{split}
 \]
In an equal \(r\)-partite graph, only within-part pairs are missing,
so its edge count is \(r(r-1)(n/r)^2/2=an^2\).
\end{proof}
For \(r\nmid n\), equal integer part sizes are impossible.
An exchange argument supplies the exact integer bound as well.

\begin{theorem}[The exact Turan theorem]\label{c:turanexact}
Among \(K_{r+1}\)-free graphs on \(n\) vertices, a complete \(r\)-partite
graph whose part sizes differ by at most one has the most edges.
Empty parts are permitted. If \(n=qr+s\), \(0\le s<r\), its edge count is
\[
 \frac12\bigl(n^2-s(q+1)^2-(r-s)q^2\bigr).
 \]
\end{theorem}
\begin{proof}
Choose a graph with maximum edge count and, among such graphs, maximum
\(\sum_vd(v)^2\). If nonadjacent vertices \(u,v\) have different
degrees, replace the smaller-degree vertex's neighborhood by a copy of
the other's. This increases the edge count and creates no forbidden
clique: any new clique using the copy could replace it by the original
vertex. Thus nonadjacent vertices have equal degree.

Suppose such \(u,v\) have different neighborhoods. Put
\(A=N(u)\setminus N(v)\), \(B=N(v)\setminus N(u)\);
these have equal positive size \(a\).
Copying \(v\)'s neighborhood to \(u\) preserves the edge count and
changes the degree-square sum by
\[
 2\sum_{x\in B}d(x)-2\sum_{x\in A}d(x)+2a.
 \]
Copying in the opposite direction changes it by the opposite difference
plus \(2a\). The two changes sum to \(4a>0\), so one improves the
chosen optimum, a contradiction.
Therefore nonadjacent vertices have identical neighborhoods.
The relation of equality or nonadjacency is consequently transitive:
if \(u,v\) and \(v,w\) are nonadjacent, adjacency of \(u,w\)
would contradict \(N(u)=N(v)\).
Its classes are independent sets, and all edges between different
classes are present. There are at most \(r\) nonempty classes.

Allow empty classes to make \(r\) in total, with sizes \(n_i\).
The edge count is \((n^2-\sum_i n_i^2)/2\).
If \(n_i\ge n_j+2\), moving one vertex from part \(i\) to part \(j\)
decreases the square sum by \(2(n_i-n_j-1)>0\).
Thus an optimum has \(s\) parts of size \(q+1\) and \(r-s\) of
size \(q\), giving the formula. The balanced construction is itself
\(K_{r+1}\)-free, so it attains the upper bound.
\end{proof}
\strategy The first Turan proof is a short numerical induction.
The second chooses a secondary optimum so that a local replacement
forces a highly organized graph. These are complementary habits:
one estimates a quantity directly; the other simplifies an extremal
object until its quantity is easy to compute.

\begin{example}[A four-cycle bound by common neighbors]
If a graph has no cycle of length four, then
\[
 e\le\frac n4\bigl(1+\sqrt{4n-3}\bigr).
\]
\end{example}
\begin{solution}
Two different vertices have at most one common neighbor, or those
vertices and two common neighbors form a four-cycle.
Count a vertex together with an unordered pair of its neighbors:
\[
 \sum_v\binom{d(v)}2\le\binom n2.
 \]
Since \(\sum_vd(v)^2\ge4e^2/n\), this gives
\(2e^2/n-e\le n(n-1)/2\).
Solving this quadratic for nonnegative \(e\) gives the bound.
\end{solution}
\recognition Forbidden triangles suggest comparing degrees across an
edge. Forbidden four-cycles suggest counting common neighbors of a pair.
The forbidden shape tells you which small incidence to count.

\section{Ramsey arguments}

Let \(R(s,t)\) be the smallest \(N\) such that every red--blue coloring
of the edges of \(K_N\) contains a red \(K_s\) or a blue \(K_t\).
The following recurrence also proves existence:
\[
 R(s,t)\le R(s-1,t)+R(s,t-1),\qquad R(1,t)=R(s,1)=1.
 \]
Indeed, at a chosen vertex, either at least \(R(s-1,t)\) red neighbors
or at least \(R(s,t-1)\) blue neighbors exist. Apply the appropriate
smaller assertion within that neighborhood and add the chosen vertex
when necessary. Pascal's identity then gives
\[
 R(s,t)\le\binom{s+t-2}{s-1}.
 \]

\begin{example}[Why six people are enough]
In any red--blue edge coloring of \(K_6\), a monochromatic triangle exists,
and the number six is sharp.
\end{example}
\begin{solution}
At any vertex at least three edges have the same color, say red.
If one edge between their other endpoints is red, it completes a red
triangle; otherwise those three endpoints make a blue triangle.
For sharpness, color a five-cycle red and its complementary five-cycle
blue. Neither contains a triangle.
\end{solution}
\strategy A Ramsey proof forces a homogeneous substructure by repeatedly
restricting to a large same-color neighborhood. It often gives a useful
bound rather than the exact threshold. Exactness needs a separate
coloring below the proposed threshold.

\section{Expectation does not require independence}

A finite probability space assigns nonnegative probabilities to outcomes,
summing to one. The expectation of a real random variable \(X\) is the
weighted sum of its values:
\(\mathbb E X=\sum_\omega P(\omega)X(\omega)\).
An indicator \(I_A\) is one on event \(A\) and zero elsewhere; then
\(\mathbb E I_A=P(A)\).
Expanding the finite sums proves
\[
 \mathbb E(X_1+\cdots+X_m)=\sum_i\mathbb E X_i.
 \]
No independence assumption appears.
Some outcome has value at least \(\mathbb E X\), because otherwise
their weighted average would be smaller.
Also \(I_{\cup A_i}\le\sum_iI_{A_i}\), giving the union bound
\(P(\cup A_i)\le\sum_iP(A_i)\).

\begin{example}[A large cut without constructing it first]
Every graph has a cut containing at least half its edges.
\end{example}
\begin{solution}
Independently place each vertex in one of two parts with probability
one half. Each edge crosses with probability one half, so the expected
crossing count is \(e/2\). Some partition has at least that many edges.
Indicators of edges sharing a vertex need not be treated as independent:
linearity is enough.
\end{solution}

\section{Alteration: choose first, repair afterward}

A random object may violate the desired condition but still be easy
to repair. Track a score that subtracts an upper bound on the repair
cost from the original size.

\begin{example}[An independent set from random deletion]
A graph with \(n\) vertices and \(e\) edges has an independent set
of size at least \(np-ep^2\) for every real \(0\le p\le1\).
\end{example}
\begin{solution}
Select each vertex independently with probability \(p\).
The expected number \(X\) selected is \(np\); the expected number \(Y\)
of edges with both ends selected is \(ep^2\).
Some selection has \(X-Y\ge np-ep^2\).
For every surviving internal edge, delete one endpoint. At most the
original \(Y\) vertices are deleted, because eliminating an edge may
eliminate several others too. The remaining set is independent and has
size at least \(X-Y\).
\end{solution}
\pitfall It is false that \emph{every} random selection has at least
its expected score. Expectation guarantees one favorable outcome.
It is also unnecessary to delete one different vertex for every edge:
the repair cost is at most, not exactly, the edge count.

\section{Conditional expectation makes the choice explicit}

Suppose a random construction is generated by finitely many binary choices.
After some choices are fixed, consider the average score over all
remaining choices. Before the next equally likely bit, this conditional
average is the mean of the two averages obtained by fixing the bit.
Choose a bit whose average is at least the current average.
At the end, no randomness remains and the actual score is at least
the original expectation. Unequal probabilities give a weighted mean
and the same conclusion.

For the random cut, compute each edge's conditional crossing probability:
it is zero or one if both ends are assigned, and one half otherwise.
Assign each new vertex to a side that does not lower the sum.
This deterministic algorithm produces a cut of size at least \(e/2\).
The method needs a conditional expectation that can actually be evaluated;
existence of an average alone does not make an efficient algorithm.

\section{A local lemma for sparse dependencies}

The union bound can fail when there are many bad events, even if each
event interacts with only a few others. The local lemma uses this
dependency structure.

\begin{theorem}[Symmetric local lemma]\label{c:lll}
Let \(A_1,\ldots,A_m\) be events. Suppose each \(A_i\) is independent
of every joint event determined by all but at most \(d\) of the other
events, and \(P(A_i)\le p\).
If a real \(x\) with \(0<x<1\) satisfies
\[
 p\le x(1-x)^d,
 \]
then \(P(\bigcap_i\overline{A_i})>0\).
\end{theorem}
\begin{proof}
We prove by induction on \(|S|\) that, for \(i\notin S\),
\[
 P\left(A_i\,\middle|\,\bigcap_{j\in S}\overline{A_j}\right)\le x.
 \]
All conditioning events have positive probability: multiplying the
earlier conditional avoidance probabilities gives factors at least
\(1-x>0\). The empty case follows from \(p\le x\).
Split \(S\) into \(S_1\), the events allowed to depend on \(A_i\),
and the remaining \(S_2\). Then \(|S_1|\le d\).
If \(S_1\) is empty, independence gives the result.
Otherwise
\[
 P(A_i\mid\text{avoid }S)
 \le
 \frac{P(A_i\mid\text{avoid }S_2)}
 {P(\text{avoid }S_1\mid\text{avoid }S_2)}
 \le\frac p{(1-x)^{|S_1|}}\le x.
 \]
For the denominator, expose the events of \(S_1\) one by one.
Each conditional failure probability is at most \(x\) by induction,
since its conditioning set has fewer than \(|S|\) elements.
The numerator equals \(P(A_i)\) by the stated joint independence.
Finally multiply the \(m\) conditional probabilities of avoiding each
event, obtaining a product at least \((1-x)^m>0\).
\end{proof}

In independent-variable models, connect two bad events when they use
a common underlying variable. An event is then independent of all
joint conditions on nonneighbors, which verifies the theorem's stronger
joint independence requirement. Pairwise independence alone would not
justify that step.

For \(d\ge1\), choosing \(x=1/(d+1)\) gives the concrete condition
\[
 p\le\frac1{d+1}\left(\frac d{d+1}\right)^d.
 \]
The theorem is an advanced existence tool. First try an invariant,
an extremal choice, expectation, or a direct union bound; these often
give a simpler proof.

\section{Exercises}
\begin{exercise}[\problemlevel{Foundation}]\label{c:ex65}
Determine the maximum number of edges of a triangle-free graph on
seven vertices, and give a graph attaining it.
\end{exercise}
\begin{exercise}[\problemlevel{Bridge}]\label{c:ex66}
Determine the maximum edge count for a \(K_4\)-free graph on twelve
vertices. Explain why your bound is attained.
\end{exercise}
\begin{exercise}[\problemlevel{Challenge}]\label{c:ex67}
Prove \(9\le R(3,4)\le10\).
For the lower bound, color red the edges of an eight-cycle together
with its four opposite-vertex edges, and blue all remaining edges.
\end{exercise}
\begin{exercise}[\problemlevel{Challenge}]\label{c:ex68}
Prove that every finite graph has an independent set of size at least
\(\sum_v1/(d(v)+1)\). Deduce the lower bound \(n/(\bar d+1)\),
where \(\bar d=2e/n\) is the average degree and \(n>0\).
\end{exercise}
\begin{exercise}[\problemlevel{Challenge}]\label{c:ex69}
Let \(M\) be an \(n\)-by-\(n\) zero--one matrix with no four ones
at the corners of a rectangle in two different rows and columns.
If \(m\) is its total number of ones, prove
\(m\le\frac n2(1+\sqrt{4n-3})\).
\end{exercise}
\begin{exercise}[\problemlevel{Challenge}]\label{c:ex70}
For every integer \(k\ge3\), prove that there is a red--blue coloring
of the edges of \(K_{\lfloor2^{k/2}\rfloor}\) containing no
monochromatic \(K_k\).
\end{exercise}
\begin{exercise}[\problemlevel{Challenge}]\label{c:ex71}
A finite family of ten-element sets has the property that each member
intersects at most one hundred other members. Prove that the elements
can be colored red and blue so that no member is monochromatic.
\end{exercise}
\begin{exercise}[\problemlevel{Challenge}]\label{c:ex72}
A Boolean formula has \(m\) clauses, each the disjunction of three
literals on three different variables. A literal is a variable or its
negation. Prove that some truth assignment satisfies at least \(7m/8\)
clauses, and give a conditional-expectation algorithm that finds one.
\end{exercise}

\chapter{Symmetry and Winning Strategies}
\label{c:games}

\section{Count objects fixed by a symmetry}

When rotations or reflections identify arrangements, dividing by the
number of symmetries is tempting. It works only if every arrangement
has that many different images. An all-red necklace has just one
image under rotation; a necklace with a single blue bead usually has
many. We need a formula that automatically compensates for stabilizers.

Let a finite set \(G\) of reversible symmetries act on a finite set \(X\).
Assume \(G\) contains the identity and is closed under composition and
inverses. The \emph{orbit} of \(x\) consists of all its images \(g(x)\).
Write \(\operatorname{Fix}(g)=\{x:g(x)=x\}\).

\begin{theorem}[Burnside's lemma]\label{c:burnside}
The number of orbits is
\[
 \frac1{|G|}\sum_{g\in G}|\operatorname{Fix}(g)|.
\]
\end{theorem}
\begin{proof}
Count pairs \((g,x)\) with \(g(x)=x\).
For a fixed \(x\) in an orbit of size \(t\), exactly \(|G|/t\)
symmetries fix \(x\). To justify this, the symmetries sending \(x\)
to any particular orbit point are a coset of the fixing symmetries:
composition with one chosen symmetry gives a bijection between the fibers.
All \(t\) fibers have equal size and partition \(G\).
The \(t\) elements of that orbit therefore contribute
\(t(|G|/t)=|G|\) pairs. Summing over orbits proves the formula.
\end{proof}

\begin{example}[Six-bead necklaces]
Count binary colorings of six equally spaced beads, identifying
rotations but not reflections.
\end{example}
\begin{solution}
A rotation by \(j\) positions has \(\gcd(6,j)\) cycles of positions.
A fixed coloring is constant on each cycle, giving
\(2^{\gcd(6,j)}\) choices; the identity has six cycles.
Thus the count is
\[
 \frac{2^6+2^1+2^2+2^3+2^2+2^1}{6}=14.
 \]
The fixed-cycle method avoids trying to classify all necklaces by period.
\end{solution}

\begin{example}[Binary colorings of a square]
Color the four vertices of a square with two colors, identifying every
rotation and reflection. There are six classes.
\end{example}
\begin{solution}
The identity fixes \(2^4=16\) colorings. The two quarter-turns each
have one position cycle and fix two; the half-turn has two cycles and
fixes four. Two reflections through opposite vertices have three cycles
and fix eight each. Two reflections through opposite edge midpoints
have two cycles and fix four each. Burnside gives
\((16+2+2+4+8+8+4+4)/8=6\).
\end{solution}
\pitfall A reflection is not a rotation. State which motions identify
objects before counting. Also count cycles of \emph{positions}, then
assign a color to each cycle; colors are not permuted unless the problem
explicitly allows exchanging their names.

\section{Work backward from losing positions}

Consider a finite directed acyclic graph of game positions. An arrow is
a legal move, the players alternate, both have the same legal moves
from a given position, and a player with no move loses. This is
\emph{normal play} for a finite impartial game.

A position is called \(P\) if the player about to move loses under
best play, and \(N\) if that player wins.
Terminal positions are \(P\). Working backward gives:
\begin{itemize}
\item a position is \(N\) if at least one move leads to \(P\);
\item a position is \(P\) if every move leads to \(N\).
\end{itemize}
These statements must both be checked when proposing a pattern of losing
positions. A move from every proposed winning position into the proposed
losing set is not enough if moves also exist between two members of that set.

\begin{example}[Taking one or two tokens]
From a heap of \(n\) tokens, a player removes one or two. The losing
positions are exactly multiples of three.
\end{example}
\begin{solution}
Zero is terminal. A move from a multiple of three gives residue one
or two. From residue one remove one, and from residue two remove two,
returning to a multiple of three. Both conditions hold; induction on
\(n\) proves the classification.
\end{solution}

\section{Nim and binary cancellation}

In Nim, a move removes any positive number of tokens from exactly one
heap. Write \(\oplus\) for bitwise exclusive-or: add binary digits
modulo two without carrying. For example,
\(5\oplus3=101_2\oplus011_2=110_2=6\).

\begin{theorem}[Nim criterion]\label{c:nim}
A Nim position with heap sizes \(a_1,\ldots,a_m\) is losing exactly
when \(a_1\oplus\cdots\oplus a_m=0\).
\end{theorem}
\begin{proof}
Let the xor be \(S\). From \(S=0\), replacing \(a_i\) by a smaller
\(b\ne a_i\) changes it to \(a_i\oplus b\ne0\).
If \(S\ne0\), consider its highest nonzero bit. Some heap \(a_i\)
has a one in that bit. Put \(b=a_i\oplus S\).
Higher bits do not change, and that bit changes from one to zero;
hence \(b<a_i\). Replacing \(a_i\) by \(b\) changes the total xor to
\(S\oplus a_i\oplus b=0\).
The all-zero terminal position has xor zero, so the two \(P/N\)
conditions prove the theorem.
\end{proof}

\begin{example}[Finding the move, not just the label]
For heaps \(3,4,5\), determine a winning move.
\end{example}
\begin{solution}
The xor is \(3\oplus4\oplus5=2\). The highest nonzero bit of two
occurs in heap three. Replace it by \(3\oplus2=1\), leaving heaps
\(1,4,5\), whose xor is zero. Removing an arbitrary number from the
largest heap need not work; the bit argument specifies the right change.
\end{solution}

\section{The Sprague--Grundy value}

For a finite impartial normal-play game, define
\[
 g(v)=\operatorname{mex}\{g(w):v\to w\},
 \]
where mex is the smallest nonnegative integer missing from the displayed
finite set. In particular, \(\operatorname{mex}\varnothing=0\).
Define values backward along the acyclic position graph.
The defining facts are that no move preserves \(g\), and for every
integer \(t<g(v)\) some move has value \(t\).
Thus \(g(v)=0\) exactly for losing positions.

The disjoint sum of games allows one move in exactly one component.
Its losing positions are characterized by xor of the component values:
\[
 g(v_1)\oplus\cdots\oplus g(v_m)=0.
 \]
The proof is the same highest-bit argument as Nim. A move from xor zero
changes one value to a different value. From nonzero xor \(S\), choose
a component with a one at its highest nonzero bit and target the smaller
value \(g(v_i)\oplus S\); mex guarantees a move with that value.
This proves the winning criterion. For the stronger value statement,
induct on the sum of the maximum remaining play lengths of the components;
every move lowers that integer. The value of the sum is the xor \(S\).
By induction, no follower has value \(S\). For any integer \(t<S\), the highest
bit where \(S,t\) differ is one in \(S\); choose a component with one
in that bit, and change its value from \(g\) to \(g\oplus S\oplus t<g\).
The mex property supplies such a move, and the resulting xor is \(t\).
Thus every smaller value occurs and \(S\) itself does not, which is
exactly the defining mex rule.

Do not apply this theory to games with player-dependent moves, infinite
play, or a last-move-loses rule without a new proof. An algebraically
similar notation does not remove its hypotheses.

\section{Pairing and strategy stealing}

A pairing strategy responds to an opponent's move by a legal move on
its paired object. The proof must explain why the reply is still available
and why paired moves cannot interfere.

Strategy stealing proves that some winning first move exists even when
it does not identify that move. One assumes the second player has a
winning strategy, makes an extra harmless first move, and shows how to
simulate that strategy. The word \emph{harmless} needs a precise argument:
an extra move can be harmful in many games.
Problem~\ref{c:ex78} gives a setting where the inclusion of removed regions
makes the argument rigorous.

\section{Exercises}
\begin{exercise}[\problemlevel{Foundation}]\label{c:ex73}
For a prime \(p\), count binary necklaces with \(p\) labelled positions
before rotations are identified. Deduce that \(p\mid(2^p-2)\).
\end{exercise}
\begin{exercise}[\problemlevel{Bridge}]\label{c:ex74}
Count colorings of the vertices of a square with three named colors,
identifying rotations and reflections.
\end{exercise}
\begin{exercise}[\problemlevel{Bridge}]\label{c:ex75}
Count binary six-bead necklaces when reflections as well as rotations
are identified.
\end{exercise}
\begin{exercise}[\problemlevel{Bridge}]\label{c:ex76}
In a one-heap subtraction game, a move removes \(1,3\), or \(4\)
tokens. Compute the Grundy values and prove they repeat with period
seven, starting at heap size zero.
\end{exercise}
\begin{exercise}[\problemlevel{Challenge}]\label{c:ex77}
Classify the losing positions of misere Nim, where the player taking
the last token loses. Assume at least one token is present.
Explain why applying the ordinary xor criterion unchanged fails.
\end{exercise}
\begin{exercise}[\problemlevel{Challenge}]\label{c:ex78}
A rectangular chocolate board has at least two squares. The bottom-left
square is poisoned. A legal move selects a nonpoisoned remaining square
and removes it and every remaining square weakly above and weakly to its
right. A player with no legal move loses. Prove that the first player
has a winning strategy, without claiming to identify the winning move.
\end{exercise}
\begin{exercise}[\problemlevel{Bridge}]\label{c:ex79}
There are several heaps; a move removes one or two tokens from one heap.
Prove that the position is losing exactly when the xor of the residues
of the heap sizes modulo three is zero.
\end{exercise}
\begin{exercise}[\problemlevel{Challenge}]\label{c:ex80}
In Wythoff's game, a move removes any positive number from one of two
heaps, or the same positive number from both. Prove that the losing
positions, with heap sizes sorted, are
\[
 \bigl(\lfloor k\phi\rfloor,\lfloor k\phi^2\rfloor\bigr)
 \quad(k=0,1,\ldots),\qquad \phi=\frac{1+\sqrt5}{2}.
 \]
As part of the proof, show that the positive terms of these two
sequences partition the positive integers.
\end{exercise}

\chapter{Algebraic Combinatorics: Encode More Structure}
\label{c:algebraic}

The methods in this chapter replace finite configurations by vectors,
polynomials, or signed sums. They are powerful because algebra can
detect relations that ordinary counting forgets. Learn the earlier
counting and extremal methods first; an algebraic encoding should make
the decisive obstruction visible, not merely make the notation heavier.

\section{Parity vectors and Oddtown}

Let \(\FF_2=\{0,1\}\), with addition and multiplication modulo two.
Encode a subset \(A\subseteq[n]\) by its incidence vector
\(v_A=(\epsilon_1,\ldots,\epsilon_n)\), where \(\epsilon_i=1\)
exactly when \(i\in A\). Then the dot product satisfies
\[
 v_A\cdot v_B=\sum_i\epsilon_i\delta_i=|A\cap B|\pmod2.
\]
A list of vectors is linearly independent if the only zero sum of
them with coefficients in \(\FF_2\) uses all coefficients zero.
There are at most \(n\) independent vectors in \(\FF_2^n\).
Here is an elementary reason: if some vector has a last coordinate one,
add it to any other vector with last coordinate one. This preserves
independence and leaves all the other vectors in \(\FF_2^{n-1}\).
Induction bounds their number by \(n-1\). If all last coordinates
are zero, apply the induction directly.

\begin{theorem}[Oddtown]\label{c:oddtown}
If \(A_1,\ldots,A_m\subseteq[n]\) all have odd size and
\(|A_i\cap A_j|\) is even for \(i\ne j\), then \(m\le n\).
\end{theorem}
\begin{proof}
Their incidence vectors have \(v_i\cdot v_i=1\) and
\(v_i\cdot v_j=0\) for \(i\ne j\).
If \(\sum_i c_i v_i=0\), dotting with \(v_j\) gives \(c_j=0\).
Thus the vectors are independent, and there are at most \(n\).
\end{proof}

\begin{example}[Why oddness is essential]
The bound \(m\le n\) fails dramatically if all sizes are even.
\end{example}
\begin{solution}
Partition \(2r\) elements into \(r\) pairs and take all unions of pairs.
There are \(2^r\) sets. Every set and every pairwise intersection has
even size. The vectors are no longer independent: their self-dot-products
are zero, so the coefficient-isolating argument breaks exactly where
oddness was used.
\end{solution}

\section{Eventown and a dimension argument}

For a collection of vectors, its span is the set of all their linear
combinations. Choose a maximal independent subcollection, a basis;
every vector in the span is a combination of the basis, because any
vector outside its span could be added to make a larger independent
collection. A basis of size \(d\) gives exactly \(2^d\) different sums:
two equal sums would yield a nontrivial zero combination.

\begin{theorem}[Eventown]\label{c:eventown}
If a family \(\mathcal F\) of distinct subsets of \([n]\) has all
set sizes and all pairwise intersection sizes even, then
\(|\mathcal F|\le2^{\lfloor n/2\rfloor}\).
\end{theorem}
\begin{proof}
Let \(V\) be the span of its incidence vectors, with basis size \(d\).
Every two vectors in \(V\) have dot product zero, by distributing the
dot product over the original generators.
Thus \(V\) lies in the set \(V^\perp\) of vectors orthogonal to its
\(d\) basis vectors.
These are \(d\) independent linear equations in \(n\) binary variables.
Elimination uses one different pivot variable per independent equation,
leaving \(n-d\) variables free, so \(V^\perp\) has \(2^{n-d}\) members.
Since \(V\subseteq V^\perp\), \(2^d\le2^{n-d}\), hence \(d\le n/2\).
The family is a subset of \(V\), giving the bound.
\end{proof}
\begin{example}[A sharp construction]
Pair \(2\lfloor n/2\rfloor\) of the elements, leaving one unused if
\(n\) is odd. All unions of these pairs form a family of size
\(2^{\lfloor n/2\rfloor}\), attaining Eventown's bound.
\end{example}
\begin{solution}
Every union has even size. Intersecting two unions leaves another
union of full pairs, also of even size. There is one independent
include-or-exclude decision per pair.
\end{solution}
\strategy Parity constraints suggest \(\FF_2\). Intersection sizes
become dot products; unions need not become vector sums, but symmetric
difference does. State what each algebraic operation means before using it.

\section{A polynomial cannot vanish for free}

Let \(K\) be a field, such as \(\QQ,\RR\), or the residues
\(\FF_p\) modulo a prime \(p\). Nonzero elements of a field can be
divided by. This is why prime moduli will matter below.

\begin{lemma}[A coefficient test on a grid]\label{c:coefficient}
Let \(P(x,y)\) have total degree at most \(a+b\).
Let \(A,B\subset K\) have \(a+1,b+1\) distinct elements.
Then the coefficient of \(x^ay^b\) equals
\[
 \sum_{\substack{u\in A\\v\in B}}
 \frac{P(u,v)}
 {\displaystyle\prod_{\substack{u'\in A\\u'\ne u}}(u-u')
  \displaystyle\prod_{\substack{v'\in B\\v'\ne v}}(v-v')}.
 \]
In particular, if this coefficient is nonzero, \(P\) cannot vanish
on every point of \(A\times B\).
\end{lemma}
\begin{proof}
For a nonnegative integer \(r\le a\), Lagrange interpolation of \(x^r\)
on \(A\) shows that
\[
 \sum_{u\in A}\frac{u^r}{\prod_{u'\ne u}(u-u')}
 =
 \begin{cases}0,&r<a,\\1,&r=a.\end{cases}
 \]
Indeed, the interpolation polynomial agrees at \(a+1\) points and
has degree at most \(a\); their difference has too many roots unless
zero. Read its leading coefficient. The same holds on \(B\).
By linearity, test the proposed formula on monomials \(x^ry^s\).
If \((r,s)=(a,b)\), the two sums multiply to one.
If \(r+s\le a+b\) but \((r,s)\ne(a,b)\), then \(r<a\) or
\(s<b\). The corresponding sum is zero. This proves the identity.
\end{proof}
\intuition Design a polynomial that vanishes at every allegedly allowed
configuration. A coefficient that is forced to be nonzero then certifies
that the alleged collection is too small.
The grid sizes determine which coefficient to inspect.

\begin{theorem}[Cauchy--Davenport]\label{c:cd}
For nonempty \(A,B\subseteq\FF_p\), \(p\) prime, their sumset
\(A+B=\{a+b:a\in A,b\in B\}\) satisfies
\[
 |A+B|\ge\min\{p,\ |A|+|B|-1\}.
 \]
\end{theorem}
\begin{proof}
Put \(s=|A|,t=|B|\).
If \(s+t>p\), then for every \(z\in\FF_p\), the sets \(A\) and
\(z-B\) intersect, so \(z\in A+B\).
Otherwise \(s+t\le p\). Suppose \(|A+B|\le s+t-2=d\).
Choose a \(d\)-element set \(C\) containing \(A+B\), and set
\[
 P(x,y)=\prod_{c\in C}(x+y-c).
 \]
It vanishes on \(A\times B\), but its coefficient of
\(x^{s-1}y^{t-1}\) is \(\binom d{s-1}\).
Since \(d<p\), none of the factorials defining this coefficient is
divisible by \(p\), so the coefficient is nonzero in \(\FF_p\).
Lemma~\ref{c:coefficient} contradicts the vanishing.
If \(s=t=1\), the desired lower bound one is immediate, so the
contradictory assumption with \(d=0\) need not be made.
\end{proof}

\section{Nonintersecting paths and determinants}

For a \(k\)-by-\(k\) matrix \(M\), define its determinant by the finite sum
\[
 \det M=\sum_{\sigma}\operatorname{sgn}(\sigma)
             \prod_{i=1}^k M_{i,\sigma(i)},
 \]
over permutations of \([k]\). We need only this definition and the fact
that exchanging two permutation values reverses its sign.

Take a finite directed acyclic graph, distinct sources \(A_1,\ldots,A_k\),
and distinct sinks \(B_1,\ldots,B_k\), all these vertices different.
Let \(M_{ij}\) count directed paths from \(A_i\) to \(B_j\).
Expand the determinant: each signed term counts families of paths
from \(A_i\) to \(B_{\sigma(i)}\).
Every family with an intersection cancels with another family.

To define the pairing precisely, order the vertices so all arrows go
forward. Such an order exists: repeatedly remove a vertex of indegree
zero; if none existed, following predecessors would force a directed cycle.
Choose the first vertex visited by two or more paths, then the two
smallest path indices meeting there. Swap their tails after that vertex.
The endpoint permutation changes by a transposition, reversing the sign.
The chosen first intersection and the chosen pair remain the same, so
this is an involution. Acyclicity ensures the spliced paths do not
repeat a vertex. Therefore
\[
 \det M=\sum_{\substack{\text{vertex-disjoint path families}\\
                        A_i\to B_{\sigma(i)}}}
                  \operatorname{sgn}(\sigma).
 \]
This is the path form of the Lindstrom--Gessel--Viennot lemma.
If the source and sink order forces every disjoint family to have
\(\sigma\) the identity, the determinant is an ordinary nonnegative count.

\begin{example}[Two walkers that cannot meet]
Count pairs of right-and-up lattice paths, one from \((0,0)\) to
\((2,1)\) and one from \((0,1)\) to \((2,2)\), with no shared vertex.
\end{example}
\begin{solution}
The path-count matrix is
\(\begin{pmatrix}3&6\\1&3\end{pmatrix}\).
A path joining the lower start to the upper end and one joining the
upper start to the lower end must meet: as their order reverses across
the strip, a vertical or horizontal grid segment would cross, and grid
crossings occur at vertices. Thus only the identity pairing can be disjoint.
The answer is \(3\cdot3-6\cdot1=3\).
\end{solution}
\begin{center}
\begin{tikzpicture}[scale=.8]
\draw[gray!30] (0,0) grid (2,2);
\draw[accent,very thick,-{Stealth}] (0,0)--(1,0)--(2,0)--(2,1);
\draw[orange!80!black,very thick,-{Stealth}] (0,1)--(0,2)--(1,2)--(2,2);
\foreach \x/\y in {0/0,0/1,2/1,2/2}
 \fill (\x,\y) circle (1.4pt);
\node[left] at (0,0) {\(A_1\)};
\node[left] at (0,1) {\(A_2\)};
\node[right] at (2,1) {\(B_1\)};
\node[right] at (2,2) {\(B_2\)};
\end{tikzpicture}
\end{center}

\section{A permutation proof of a set-pair bound}

\begin{theorem}[Bollobas's set-pair inequality]\label{c:bollobas}
Suppose finite sets \(A_i,B_i\), \(1\le i\le m\), satisfy
\(A_i\cap B_i=\varnothing\) and
\(A_i\cap B_j\ne\varnothing\) whenever \(i\ne j\).
Then
\[
 \sum_{i=1}^m\frac1{\binom{|A_i|+|B_i|}{|A_i|}}\le1.
 \]
\end{theorem}
\begin{proof}
Order the finite union of all these sets uniformly at random.
Let \(E_i\) be the event that every element of \(A_i\) precedes every
element of \(B_i\). Among the relative orders of \(A_i\cup B_i\),
exactly a fraction \(1/\binom{|A_i|+|B_i|}{|A_i|}\) have this property.
Events \(E_i,E_j\) are disjoint for \(i\ne j\):
choose \(x\in A_i\cap B_j\) and \(y\in A_j\cap B_i\).
Event \(E_i\) forces \(x\) before \(y\), whereas \(E_j\) forces \(y\)
before \(x\). The elements are different, since otherwise
\(A_i\cap B_i\) would be nonempty. Summing disjoint event probabilities
gives the result.
\end{proof}
\begin{example}[The uniform bound is sharp]
If all \(|A_i|=a\) and \(|B_i|=b\), there are at most
\(\binom{a+b}{a}\) pairs. All \(a\)-subsets of an \((a+b)\)-set,
paired with their complements, attain this.
\end{example}
\begin{solution}
The inequality becomes \(m/\binom{a+b}{a}\le1\).
For different \(a\)-sets \(A_i,A_j\), some element of \(A_i\)
is outside \(A_j\), so \(A_i\cap B_j\ne\varnothing\);
each own pair is disjoint. The construction satisfies every hypothesis.
\end{solution}

\section{Shattering and an inductive growth bound}

A family \(\mathcal F\subseteq2^{[n]}\) \emph{shatters} \(T\subseteq[n]\)
if \(\{A\cap T:A\in\mathcal F\}=2^T\). In other words, every membership
pattern on \(T\) appears somewhere. A large family must shatter a
reasonably large set.

\begin{theorem}[Sauer--Shelah bound]\label{c:sauer}
If \(\mathcal F\) shatters no set of size \(d+1\), where \(d\ge0\),
then
\[
 |\mathcal F|\le\sum_{i=0}^d\binom ni.
 \]
\end{theorem}
\begin{proof}
For \(n=0\) the claim is immediate; for \(d\ge n\) it is the trivial
bound \(2^n\). Otherwise split by the last element:
\[
 \mathcal F_0=\{A\subseteq[n-1]:A\in\mathcal F\},\qquad
 \mathcal F_1=\{A\subseteq[n-1]:A\cup\{n\}\in\mathcal F\}.
 \]
Put \(\mathcal U=\mathcal F_0\cup\mathcal F_1\) and
\(\mathcal I=\mathcal F_0\cap\mathcal F_1\).
Then \(|\mathcal F|=|\mathcal U|+|\mathcal I|\).
Family \(\mathcal U\) shatters no \((d+1)\)-set, or the original
family would shatter that same set.
Family \(\mathcal I\) shatters no \(d\)-set, or \(\mathcal F\)
would shatter that set together with \(n\).
For \(d=0\) the latter statement means \(\mathcal I\) is empty,
since every nonempty family shatters the empty set.
For \(d\ge1\), induction and Pascal's identity give
\[
 |\mathcal F|\le
 \sum_{i=0}^d\binom{n-1}{i}
 +\sum_{i=0}^{d-1}\binom{n-1}{i}
 =\sum_{i=0}^d\binom ni.
 \]
For \(d=0\), induction on \(\mathcal U\) gives the same bound one.
\end{proof}
\recognition If a problem limits the patterns a family can produce
on small sets, split into those that contain a chosen element and those
that do not. The intersection measures patterns available in both
versions; it frequently has a stronger restriction.

\section{Exercises}
\begin{exercise}[\problemlevel{Foundation}]\label{c:ex81}
Find the largest possible size of a family of odd subsets of \([n]\)
with all pairwise intersections even, and give a construction.
\end{exercise}
\begin{exercise}[\problemlevel{Bridge}]\label{c:ex82}
Find the largest family of subsets of an eleven-element set in which
every size and every intersection size is even.
\end{exercise}
\begin{exercise}[\problemlevel{Bridge}]\label{c:ex83}
Show that Cauchy--Davenport fails if residues modulo a composite number
replace \(\FF_p\). Identify the precise algebraic step that breaks.
\end{exercise}
\begin{exercise}[\problemlevel{Bridge}]\label{c:ex84}
For a nonempty \(A\subseteq\FF_p\), let \(kA\) be the set of sums
of \(k\) not necessarily distinct members of \(A\).
Prove \(|kA|\ge\min\{p,\ k(|A|-1)+1\}\) for every \(k\ge1\).
\end{exercise}
\begin{exercise}[\problemlevel{Bridge}]\label{c:ex85}
For \(n\ge1\), count vertex-disjoint right-and-up path pairs from
\((0,0)\) to \((n,1)\) and from \((0,1)\) to \((n,2)\).
Explain why the determinant is an unsigned count in this geometry.
\end{exercise}
\begin{exercise}[\problemlevel{Bridge}]\label{c:ex86}
Pairs \((A_i,B_i)\) satisfy the hypotheses of
Theorem~\ref{c:bollobas}, with \(|A_i|=|B_i|=k\).
Determine the largest possible number of pairs and construct equality.
\end{exercise}
\begin{exercise}[\problemlevel{Challenge}]\label{c:ex87}
There are \(n\) distinct points in the plane, with no three collinear.
A closed half-plane selects the points lying in it.
Show that the number of different selected subsets is at most
\(1+n+\binom n2+\binom n3\).
\end{exercise}
\begin{exercise}[\problemlevel{Challenge}]\label{c:ex88}
For \(A\subseteq\FF_p\), \(p\) prime, \(|A|=a\ge2\), prove
\[
 |\{x+y:x,y\in A,\ x\ne y\}|\ge\min\{p,\,2a-3\}.
 \]
The factor \(x-y\) can enforce the condition \(x\ne y\) in a
vanishing polynomial. Choose the grid sizes and coefficient carefully.
\end{exercise}



% ---- chapters/convex-lattice ----
\chapter{Combinatorics in the Plane}
\label{ch:convex-lattice}
\intuition A finite geometric configuration can be studied by counting
its vertices, expressing its points as averages, or drawing a planar
graph inside it. These descriptions connect convexity to combinatorics
and area to arithmetic. Only elementary areas are needed initially;
the final challenge can be revisited after the geometry metric chapter.

\section{Convexity means that averages stay inside}
A subset $K$ of the plane is \emph{convex} if, with any two points,
it contains the entire segment between them. Repeatedly taking
segments shows that every finite weighted average
$\lambda_1P_1+\cdots+\lambda_kP_k$, with $\lambda_i\ge0$ and
$\sum_i\lambda_i=1$, stays in $K$.
The weights are called convex coefficients.

The \emph{convex hull} of a finite nonempty set $S$ is the set of
all its convex combinations. It is convex: averaging two combinations
just combines their nonnegative coefficients.
Every convex set containing $S$ contains all these combinations,
so this is precisely the smallest convex set containing $S$.
Intersecting all convex sets containing $S$ gives the same definition.

\begin{example}[A support line detects an extreme point]
Let $P$ be the unique point of $S$ with largest first coordinate.
Prove it cannot be a convex combination of the other points.
\end{example}
\begin{solution}
All other first coordinates are strictly smaller than $P$'s.
Their finite weighted average, with nonnegative weights summing to
one, is also strictly smaller. Thus it cannot equal $P$'s first coordinate.
\end{solution}
\recognition If a line puts every point on one side, it bounds the
hull. Maximizing a linear expression such as $ux+vy$ finds a
supporting line, and often an extreme point, without drawing the hull.

\section{Four points force two intersecting hulls}
\begin{theorem}[Radon's lemma in the plane]\label{l:radon}
Any four plane points can be partitioned into two nonempty groups
whose convex hulls intersect.
\end{theorem}
\begin{proof}
Write the points as $P_i=(x_i,y_i)$.
The three homogeneous equations
$\sum_i t_i=0$, $\sum_i t_ix_i=0$, $\sum_i t_iy_i=0$
in four unknowns have a nonzero real solution: elimination has
at most three pivot variables, leaving at least one free variable.
Some $t_i$ are positive and some negative, since they sum to zero.
Let $T=\sum_{t_i>0}t_i=-\sum_{t_i<0}t_i>0$.
Then
\[
\sum_{t_i>0}\frac{t_i}{T}P_i
=\sum_{t_i<0}\frac{-t_i}{T}P_i.
\]
This common point is in both hulls. Assign any zero-coefficient
points to either group; both groups remain nonempty.
\end{proof}
Geometrically, either one of four points lies in the triangle of
the other three, or the four are vertices of a convex quadrilateral
and its diagonals cross. The algebra also includes collinear cases.

\begin{center}
\begin{tikzpicture}[scale=.8]
\coordinate (A) at (0,0);\coordinate (B) at (3,0);
\coordinate (C) at (2.6,2.2);\coordinate (D) at (.3,1.7);
\draw (A)--(B)--(C)--(D)--cycle;
\draw[main] (A)--(C) (B)--(D);
\foreach \X/\pos in {A/below left,B/below right,C/above right,D/above left}
{\fill (\X) circle (1.3pt) node[\pos]{$\X$};}
\end{tikzpicture}
\end{center}
The two groups in the convex-quadrilateral case are the opposite pairs.

\begin{theorem}[Planar Caratheodory]\label{l:caratheodory}
Every point in the convex hull of a finite set is a convex combination
of at most three of its points.
\end{theorem}
\begin{proof}
Take a combination with more than three positive coefficients.
As in Radon's proof, choose a nonzero affine dependence
$\sum t_iP_i=0,\sum t_i=0$ among its points.
Replace $\lambda_i$ by $\lambda_i-c t_i$, where
$c=\min_{t_i>0}\lambda_i/t_i>0$.
The new coefficients remain nonnegative, retain sum one and the same
weighted point, and at least one positive coefficient becomes zero.
Repeating finitely many times leaves at most three.
\end{proof}
Thus the hull is a finite union of closed triangles, including
degenerate ones. For noncollinear $S$ its boundary is a convex polygon.
An extreme corner must be one of the original points: if it were
a nontrivial average of other hull points, it would not be extreme.
This explains why the rubber-band picture of a finite hull has
vertices among the points, rather than new mysterious corners.

\section{Helly: three-set tests control the whole family}
\begin{theorem}[Helly in the plane]\label{l:helly}
Let $K_1,\ldots,K_m$ be a finite family of convex plane sets,
with $m\ge3$. If every three have a common point, all have a common
point. For families with fewer than three sets, require their
intersection directly; the three-set condition would otherwise be vacuous.
\end{theorem}
\begin{proof}
Induct on $m$, the case $3$ being the hypothesis.
For $m\ge4$, choose $P_i$ in the intersection of all sets except
$K_i$, for $i=1,2,3,4$.
These intersections exist by induction on the $m-1$ remaining sets.
Radon partitions these four points into groups with a common hull
point $P$.
For any $K_j$ with $j\le4$, one group avoids $P_j$; all its
points lie in $K_j$, so convexity puts $P$ in $K_j$.
For $j>4$, all four points lie in $K_j$, so again $P\in K_j$.
Thus $P$ lies in every set.
\end{proof}
\idea The missing-set points are chosen so that each set can object
to at most one of them. Radon then gives two descriptions of the
same average, and every set accepts at least one description.

\begin{example}[Three-point covering tests]
Let $S$ be a finite set of at least three points. If every three
can be covered by a disk of radius $r$, prove all of $S$ can.
\end{example}
\begin{solution}
For each $P\in S$, the possible centers of a covering disk form
the closed radius-$r$ disk centered at $P$.
These disks are convex: the triangle inequality keeps every average
of two permitted centers within distance $r$ of $P$.
The hypothesis says every three center-disks intersect.
Helly gives one center in all of them, hence a disk covering $S$.
\end{solution}
\pitfall Pairwise intersection is insufficient in the plane.
Three unit disks with centers forming an equilateral triangle of
side $1.9$ overlap pairwise but have no common point.
If they had a common point $P$, summing squared distances to their
centers would give at most $3$.
The centroid identity gives that sum as $3PG^2+1.9^2>3$,
a contradiction.

\section{Lattice segments and elementary triangles}
A lattice point has two integer coordinates. A lattice polygon has
lattice vertices. For a segment with displacement $(a,b)\ne(0,0)$,
put $g=\gcd(|a|,|b|)$.
Its lattice points are exactly the $g+1$ equally spaced points
with displacements $j(a/g,b/g)$, $0\le j\le g$.
Indeed, if a point on the segment has displacement
$t(a/g,b/g)$ with both coordinates integral, Bezout's identity
for the coprime pair $a/g,b/g$ makes $t$ an integer.

\begin{lemma}[An elementary lattice triangle]\label{l:elementary}
A nondegenerate lattice triangle with no lattice points other than
its three vertices on or inside it has area $1/2$.
\end{lemma}
\begin{proof}
Translate a vertex to the origin and call the others $(a,b),(c,d)$.
The first edge is primitive, so $\gcd(a,b)=1$.
Choose integers $s,t$ with $sa+tb=1$.
The linear change
$(x,y)\mapsto(sx+ty,-bx+ay)$ has determinant $1$, has an integer
inverse, preserves areas, and sends $(a,b)$ to $(1,0)$.
Reflect the second coordinate if necessary, so the remaining
vertex is $(u,v)$ with integer $v>0$.
The shear $(x,y)\mapsto(x-ky,y)$, with integer $k$, preserves
the lattice and area and can arrange $0\le u<v$.

If $v\ge2$, intersect the triangle with height $1$.
Its horizontal range is
$[u/v,\,1+(u-1)/v]$.
For $u=0$ it contains the lattice point $(0,1)$;
for $1\le u<v$ it contains $(1,1)$.
Either is a nonvertex point on or inside the triangle, a contradiction.
Thus $v=1$ and its area, half the determinant, is $1/2$.
\end{proof}
\strategy Area preservation needs checking when coordinates change.
The two maps above have determinant $\pm1$ and integer inverses,
so they preserve both the lattice and the absolute area.

\section{Pick's theorem by triangulation and Euler counting}
\begin{theorem}[Pick]\label{l:pick}
For a simple lattice polygon of area $A$, with $I$ lattice points
strictly inside and $B$ on its boundary,
\[
A=I+\frac B2-1.
\]
Here simple means its edges do not cross except at their shared
consecutive endpoints; the polygon has no holes.
\end{theorem}
\begin{proof}
Triangulate the polygon, then insert every lattice point on or inside
it as a vertex, splitting adjacent triangles when necessary.
There are finitely many such points, so this terminates in triangles
with no further lattice points. Each has area $1/2$ by
Lemma~\ref{l:elementary}.

For completeness, the initial triangulation needs no new vertices.
A simple polygon has a convex corner $B$, for instance at a
noncollinear extreme supporting vertex. Let its neighbors be $A,C$.
If triangle $ABC$ contains no other polygon vertices, $AC$ is an
interior diagonal. Otherwise choose a vertex $P$ in that triangle,
other than $A,B,C$, farthest from line $AC$ toward $B$.
The open segment $BP$ is an interior diagonal: an edge crossing it
would either cross $AB$ or $BC$, or have an endpoint still farther
toward $B$. Both are impossible. Cutting along the resulting
diagonal reduces the number of vertices, so induction triangulates.
Collinear consecutive vertices may first be deleted and later restored.

Let $T$ be the number of final triangles and $E_0$ the number of
interior edges. There are exactly $B$ boundary edges, since all
boundary lattice points were inserted. Counting triangle-edge
incidences gives $3T=2E_0+B$.
Euler's planar graph formula, including the exterior face, gives
$(I+B)-(E_0+B)+(T+1)=2$, so $I-E_0+T=1$.
Eliminating $E_0$ yields $T=2I+B-2$. Therefore
$A=T/2=I+B/2-1$.
\end{proof}
\idea The minus one is an Euler characteristic in disguise.
Area counts little triangles; the same triangulation's edges and
vertices convert that count into boundary and interior points.

\begin{example}[Count without listing interior points]
Find the number of interior lattice points of the triangle
with vertices $(0,0),(6,0),(0,4)$.
\end{example}
\begin{solution}
Its area is $12$. The boundary count is the sum of the three edge
gcds, $6+4+\gcd(6,4)=12$.
Summing $g+1$ on each edge would count every vertex twice,
so the sum of gcds is the correct $B$.
Pick gives $I=12-12/2+1=7$.
\end{solution}

\begin{center}
\begin{tikzpicture}[scale=.62]
\draw[gray!30] (0,0) grid (6,4);
\foreach \x in {0,...,6}\foreach \y in {0,...,4}
{\fill[gray] (\x,\y) circle (1pt);}
\fill[accent!10] (0,0)--(6,0)--(0,4)--cycle;
\draw[main] (0,0)--(6,0)--(0,4)--cycle;
\foreach \x/\y in {1/1,2/1,3/1,4/1,1/2,2/2,1/3}
{\fill[accent] (\x,\y) circle (2pt);}
\end{tikzpicture}
\end{center}
The seven emphasized points are strictly interior; points on the
sloping edge belong to $B$.

\section{Problems}
\begin{exercise}[\problemlevel{Foundation}]\label{l:ex01}
How many lattice points lie on the segment from $(2,-1)$ to $(14,7)$,
including endpoints? Prove your count.
\end{exercise}
\begin{exercise}[\problemlevel{Foundation}]\label{l:ex02}
For positive integers $m,n$, verify Pick's formula directly for the
rectangle with vertices $(0,0),(m,0),(m,n),(0,n)$.
\end{exercise}
\begin{exercise}[\problemlevel{Bridge}]\label{l:ex03}
Prove that among any five lattice points, two have a lattice midpoint.
Deduce that a convex lattice polygon with at least five vertices has
a lattice point other than its vertices on or inside it.
\end{exercise}
\begin{exercise}[\problemlevel{Bridge}]\label{l:ex04}
Use Pick to count interior lattice points of the triangle
$(0,0),(9,0),(0,6)$.
\end{exercise}
\begin{exercise}[\problemlevel{Bridge}]\label{l:ex05}
For a finite family of closed intervals on the real line, prove
that pairwise intersection implies a common point.
Explain why this does not contradict the three-disk counterexample.
\end{exercise}
\begin{exercise}[\problemlevel{Challenge}]\label{l:ex06}
Prove that a point inside a finite planar convex hull lies in
some triangle whose vertices are original points, allowing a
degenerate triangle if necessary. Identify why nonnegative
coefficients are preserved in your reduction.
\end{exercise}
\begin{exercise}[\problemlevel{Challenge}]\label{l:ex07}
A simple lattice polygon has no interior lattice points and every
edge is primitive, with no lattice point strictly between its endpoints.
Express its area in terms of its number $n$ of vertices, and prove
that its area is an integer when $n$ is even and a half-integer when $n$ is odd.
\end{exercise}
\begin{exercise}[\problemlevel{Challenge}]\label{l:ex08}
Let $S$ be a finite set of plane points whose pairwise distances are
at most $D>0$. Prove that some disk of radius $D/\sqrt3$ contains $S$.
Show that the constant cannot be reduced. You may use the sine law
after studying the geometry metric chapter.
\end{exercise}


\part{Geometry}

% ---- chapters/geometry ----
\chapter{Seeing a Figure and Proving a Fact}
\label{g:foundations}

Geometry begins with a useful tension. A picture makes a relationship
plausible, but a proof must explain why it survives when the picture
changes. Drawings in this part are aids to reasoning. Marks for equal
lengths, right angles, and parallel lines record hypotheses or proved
facts; apparent equality on the page records neither.

\section{The objects and the questions}

Points are denoted by capital letters. The line through distinct points
$A,B$ is $AB$; the same letters in a numerical equation denote the
positive distance between them. The ray $AB$ starts at $A$ and goes
through $B$. The segment $AB$ includes its endpoints. Three distinct
points are collinear if they belong to a line. A triangle $ABC$ is
nondegenerate when its vertices are not collinear.

An ordinary angle $\angle ABC$ is the angle between the rays $BA$ and
$BC$, between $0^\circ$ and $180^\circ$. A right angle is $90^\circ$.
Parallel lines do not meet in the Euclidean plane; perpendicular lines
meet at a right angle. A circle with center $O$ and radius $r>0$ is the
set of points $X$ with $OX=r$.

\intuition Many geometry conclusions have one of four forms:
two lengths are equal; two angles are equal; three points lie on a line;
or four points lie on a circle. Each form has a useful supply of
equivalent descriptions. For instance, to show $PX=PY$, one may show
that $P$ lies on the perpendicular bisector of $XY$. The proof becomes
easier after changing the language of the conclusion.

\begin{figure}[htbp]\centering
\begin{tikzpicture}[scale=.82]
\coordinate (A) at (-2,0); \coordinate (B) at (2,0);
\coordinate (M) at (0,0); \coordinate (P) at (0,2.8);
\draw (A)--(B); \draw[main] (A)--(P)--(B);
\draw[aux] (0,-.45)--(0,3.15);
\draw (0,.22)--(.22,.22)--(.22,0);
\foreach \X/\pos in {A/below,B/below,M/below,P/above}
 \fill (\X) circle(1.3pt) node[\pos]{$\X$};
\node at (-1,-.3){$a$}; \node at (1,-.3){$a$};
\end{tikzpicture}
\caption{An equality of distances can be replaced by a locus:
$PA=PB$ exactly when $P$ lies on the perpendicular bisector of $AB$.}
\label{g:fig:bisector}
\end{figure}

\begin{theorem}[Perpendicular-bisector criterion]
For distinct $A,B$, a point $P$ satisfies $PA=PB$ if and only if it
lies on the perpendicular bisector of $AB$.
\end{theorem}
\begin{proof}
Let $M$ be the midpoint of $AB$. If $P=M$, both assertions hold.
Otherwise, if $PM\perp AB$, the right triangles $PMA,PMB$ have equal
legs, so their hypotenuses are equal. Conversely, if $PA=PB$, the
triangles $PMA,PMB$ are congruent by three sides. The adjacent angles
at $M$ are equal and sum to $180^\circ$, so each is right.
\end{proof}

We use the elementary congruence criteria SSS, SAS, and ASA. Their
meaning is that three sides, two sides and the included angle, or two
angles and a corresponding side determine a triangle up to a rigid
motion. In SAS, the word \emph{included} matters. Two sides and an
unrelated angle may describe two different triangles.

\begin{example}[The third side is not a guess]
In triangle $ABC$, suppose $AB=AC$. Prove that the segment from $A$
to the midpoint $M$ of $BC$ is perpendicular to $BC$ and bisects
$\angle BAC$.
\end{example}
\begin{solution}
The triangles $ABM,ACM$ share $AM$, have $AB=AC$, and have $BM=CM$.
SSS gives congruence. Thus $\angle BAM=\angle MAC$, and the equal
adjacent angles $\angle BMA,\angle AMC$ are right.
\end{solution}
\strategy Equal distances and a midpoint suggest constructing two
triangles which share a side. Before chasing angles, write the three
pieces of data that would establish congruence.

\section{Parallel lines create equal angles}

If a transversal meets two parallel lines, corresponding and alternate
interior angles are equal. Conversely, either suitable equality proves
parallelism. Drawing the parallel through one vertex of a triangle
shows that the three interior angles sum to $180^\circ$.

\begin{example}[Why an exterior angle is larger]
Extend $BC$ beyond $C$ to $D$. Prove
$\angle ACD=\angle CAB+\angle ABC$.
\end{example}
\begin{solution}
The left side is $180^\circ-\angle ACB$. Substitute the triangle angle
sum. In particular it is larger than each of the two remote interior
angles. This observation often excludes an impossible point order.
\end{solution}

\begin{theorem}[Midpoint theorem]
If $M,N$ are the midpoints of $AB,AC$ in a nondegenerate triangle,
then $MN\parallel BC$ and $MN=BC/2$.
\end{theorem}
\begin{proof}
Extend $MN$ beyond $N$ to $D$ with $ND=MN$. Triangles $ANM,CND$
are congruent by SAS: $AN=CN$, $MN=ND$, and the included angles
are vertical. Hence $CD=AM=BM$ and $CD\parallel AM\parallel BM$.
Thus quadrilateral $MBCD$ is a parallelogram: a pair of opposite
sides is equal and parallel. Consequently $MD=BC$ and $MD\parallel BC$.
Since $MD=2MN$, the result follows.
\end{proof}

\begin{figure}[htbp]\centering
\begin{tikzpicture}[scale=.83]
\coordinate (A) at (.8,3); \coordinate (B) at (-2,0);
\coordinate (C) at (3,0); \coordinate (M) at ($(A)!.5!(B)$);
\coordinate (N) at ($(A)!.5!(C)$); \coordinate (D) at ($(M)!2!(N)$);
\draw (A)--(B)--(C)--cycle;
\draw[main] (M)--(D)--(C);
\draw[aux] (C)--(N);
\foreach \X/\pos in {A/above,B/below left,C/below right,M/left,N/above right,D/right}
 \fill (\X) circle(1.2pt) node[\pos]{$\X$};
\end{tikzpicture}
\caption{Extend a midpoint segment by its own length. The new triangle
is congruent to an old one, producing a parallelogram.}
\end{figure}

\section{How an auxiliary construction earns its place}

An auxiliary line should have a job. A parallel may create similar
triangles. A perpendicular may create a cyclic quadrilateral. A midpoint
may turn a length into half of another. An extension may convert a
difference of lengths into a sum. State the job before making the
construction; otherwise the diagram quickly becomes a collection of
unrelated lines.

\begin{example}[A median and a right angle]
In a right triangle $ABC$ with $\angle A=90^\circ$, let $M$ be the
midpoint of $BC$. Prove $MA=MB=MC$.
\end{example}
\begin{solution}
Reflect $A$ in $M$ to obtain $D$. The diagonals $AD,BC$ of
quadrilateral $ABDC$ bisect each other, so it is a parallelogram.
Its angle at $A$ is right, so it is a rectangle. The diagonals of a
rectangle are equal: for example SAS applied to triangles $ABC,BAD$
gives $BC=AD$. Halving gives $MA=MB=MC$.
\end{solution}
\recognition A midpoint of a hypotenuse suggests completing a rectangle.
The reverse statement is equally useful: if $MA=MB=MC$ and $B,M,C$
are collinear with $M$ between $B,C$, then $\angle BAC=90^\circ$.

\pitfall A perpendicular from a vertex can meet the extension of the
opposite side. The foot of an altitude is not always inside the
triangle. A construction remains legitimate outside the original
triangle; an argument using positive areas may need modification.

\section{First geometry practice}
Problems marked Foundation check definitions and one idea. Bridge
problems combine ideas. Challenge problems ask for a significant choice
of construction or representation. Difficulty is local to the chapter:
a Challenge here is preparation for the later chapters.

\begin{exercise}[\problemlevel{Foundation}]\label{g:ex01}
Prove that the two base angles of an isosceles triangle are equal,
and prove the converse.
\end{exercise}
\begin{exercise}[\problemlevel{Foundation}]\label{g:ex02}
In a parallelogram, prove that opposite sides are equal and that
the diagonals bisect each other.
\end{exercise}
\begin{exercise}[\problemlevel{Bridge}]\label{g:ex03}
Give a second proof of the midpoint theorem using a parallel through
$M$ and similar triangles, after reading the next chapter.
Explain why trying to intersect the parallel through $A$ to $BC$
with line $BM$ gives no new point.
\end{exercise}
\begin{exercise}[\problemlevel{Bridge}]\label{g:ex04}
Prove the converse of the hypotenuse-midpoint result:
if $M$ is the midpoint of $BC$ and $MA=MB$, then
$\angle BAC=90^\circ$.
\end{exercise}
\begin{exercise}[\problemlevel{Bridge}]\label{g:ex05}
The midpoints of the four sides of an arbitrary nondegenerate
quadrilateral form a parallelogram. Describe exactly when this
parallelogram is a rectangle.
\end{exercise}
\begin{exercise}[\problemlevel{Challenge}]\label{g:ex06}
Two distinct points $B,C$ lie on the same side of a line $\ell$.
Construct the point $X\in\ell$ for which $BX+XC$ is least.
Prove uniqueness, including when the perpendicular projections of
$B,C$ onto $\ell$ coincide.
\end{exercise}

\chapter{Similarity, Angles, and Circles}
\label{g:similarity}

\section{Similarity is a controlled change of scale}

Congruence preserves size. Similarity preserves shape while multiplying
every length by one positive factor. Two triangles are similar if
their corresponding angles are equal; AA suffices because the third
angle is determined by the first two. Equivalent criteria are SSS
with proportional sides and SAS with proportional adjacent sides and
equal included angle.

\intuition A ratio is often a disguised angle equality. If a problem
contains a product $AB\cdot CD$, try rewriting it as a ratio with
two lengths on each side. The proposed correspondence of triangles
must explain \emph{which} side matches which.

\begin{theorem}[Parallel-section theorem]
If $D\in AB$, $E\in AC$, and $DE\parallel BC$, then
\[
 \frac{AD}{AB}=\frac{AE}{AC}=\frac{DE}{BC}.
\]
Conversely, if $D,E$ lie on the segments $AB,AC$ and
$AD/AB=AE/AC$, then $DE\parallel BC$.
\end{theorem}
\begin{proof}
The forward direction follows from AA for $ADE,ABC$. For the converse,
draw the parallel through $D$ meeting $AC$ at $E'$. The forward direction
gives $AE'/AC=AD/AB=AE/AC$, so $E'=E$.
\end{proof}

\begin{example}[The altitude to a hypotenuse]
Let $ABC$ be right at $A$, and let $D$ be the foot from $A$ to $BC$.
Show
\[
 AB^2=BD\cdot BC,\qquad AC^2=CD\cdot BC,\qquad AD^2=BD\cdot CD.
\]
\end{example}
\begin{solution}
The three triangles $ABC,DBA,DAC$ have the same two acute angles.
Writing their side correspondences gives
$AB/BC=BD/AB$ and $AC/BC=CD/AC$.
Finally $AD/BD=CD/AD$ from the two smaller triangles.
Multiplying the first two identities is another route to the third
after using the area identity $AB\cdot AC=AD\cdot BC$.
\end{solution}

\begin{figure}[htbp]\centering
\begin{tikzpicture}[scale=.88]
\coordinate (B) at (0,0); \coordinate (C) at (5,0);
\coordinate (A) at (1.8,2.4); \coordinate (D) at (1.8,0);
\draw (A)--(B)--(C)--cycle; \draw[main] (A)--(D);
\draw (D)++(0,.2)--++(.2,0)--++(0,-.2);
\foreach \X/\pos in {A/above,B/below left,C/below right,D/below}
 \fill (\X) circle(1.2pt) node[\pos]{$\X$};
\end{tikzpicture}
\caption{The altitude to the hypotenuse creates two smaller copies of
the original right triangle. The coordinates here satisfy
$AB\perp AC$.}
\end{figure}

\section{Inscribed angles}

\begin{theorem}[Inscribed-angle theorem]
For distinct points $A,B,P$ on a circle, the angle $\angle APB$
equals half the angular measure of the arc $AB$ which does not
contain $P$. Consequently, angles subtending the same chord from
the same side of its line are equal.
\end{theorem}
\begin{proof}
First suppose one side of the angle passes through the center $O$.
The two equal radii make an isosceles triangle; its angle sum expresses
the central angle as twice the angle at $P$. In general draw the
diameter through $P$. If it lies inside $\angle APB$, add the two
special-case equalities. If it lies outside, subtract them. These cases
cover the possible positions, and identify the arc not containing $P$.
\end{proof}

Four distinct points are \emph{concyclic} when they lie on one circle.
For a convex quadrilateral $ABCD$, concyclicity is equivalent to
$\angle ABC+\angle ADC=180^\circ$. It is also equivalent to an
appropriate pair of equal angles subtending one of its diagonals.
Directed angles give a formulation which does not require a side choice.

\begin{figure}[htbp]\centering
\begin{tikzpicture}[scale=.88]
\coordinate (O) at (0,0);
\coordinate (A) at (210:2); \coordinate (B) at (330:2);
\coordinate (P) at (100:2); \coordinate (Q) at (45:2);
\draw (O) circle(2);
\draw[main] (A)--(P)--(B); \draw (A)--(Q)--(B);
\draw[aux] (A)--(O)--(B);
\foreach \X/\pos in {A/below left,B/below right,P/above,Q/above right,O/below}
 \fill (\X) circle(1.2pt) node[\pos]{$\X$};
\end{tikzpicture}
\caption{Two points on the same arc see the opposite chord under equal
angles. Moving a point to the other arc makes the ordinary angles
supplementary.}
\end{figure}

\begin{example}[Perpendiculars make a circle]
In triangle $ABC$, let $E,F$ be the feet of the perpendiculars from
$B,C$ to the lines $CA,AB$. Prove $B,C,E,F$ are concyclic and deduce
$\angle AEF=\angle ABC$ when $ABC$ is acute.
\end{example}
\begin{solution}
The right angles $\angle BEC,\angle BFC$ show that $E,F$ lie on the
circle with diameter $BC$. In the acute configuration,
$A,E,C$ are collinear with $E$ between $A,C$, so
$\angle AEF$ is supplementary to $\angle CEF$.
The cyclic quadrilateral gives $\angle CEF+\angle CBF=180^\circ$,
and $\angle CBF=\angle CBA$. Combining yields the claim.
The circle assertion remains true for obtuse triangles; the last
ordinary-angle equality needs the stated configuration.
\end{solution}

\section{Directed angles prevent hidden case splits}

For two unoriented lines $\ell,m$, let $\angle(\ell,m)$ mean the
rotation from $\ell$ to $m$ modulo $180^\circ$. Thus
$\angle(\ell,m)+\angle(m,n)=\angle(\ell,n)$ and reversing a line
does not change its class. Write
$\dang ABC=\angle(BA,BC)$. Directed-angle equalities in this book are
always modulo $180^\circ$.

\begin{theorem}[Directed cyclicity criterion]
Let $A,B,C,D$ be distinct points with $A\ne D$ and with $B,C$ off
line $AD$. Then
\[
 A,B,C,D\text{ are concyclic}
 \quad\Longleftrightarrow\quad
 \dang ABD=\dang ACD .
\]
\end{theorem}
\begin{proof}
The forward direction is the inscribed-angle theorem, taking
supplementary angles as equal modulo $180^\circ$. For an algebraic
proof of the converse, scale and rotate so that $A=(0,0)$,
$D=(1,0)$, $B=(u,v)$ with $v\ne0$, and $C=(x,y)$ with $y\ne0$.
The dot and signed cross products of the two angle arms give
\[
 \cot\dang ABD=\frac{u^2+v^2-u}{v},\qquad
 \cot\dang ACD=\frac{x^2+y^2-x}{y}.
\]
Equality is exactly
\[
 x^2+y^2-x+\frac{u-u^2-v^2}{v}\,y=0.
\]
This is the circle through $A,D,B$. Equal cotangents characterize
equal nonzero angle classes, proving the converse.
\end{proof}

\begin{example}[A circle found from its intended use]
Let $H$ be the orthocenter of an acute triangle $ABC$, and let $E,F$
be the feet from $B,C$. Prove $A,E,H,F$ are concyclic.
\end{example}
\begin{solution}
Since $EH\perp AE$ and $FH\perp AF$, the angles at $E,F$ are right.
Both points lie on the circle with diameter $AH$. The useful thought
is to identify two right angles seeing one common diameter.
\end{solution}
\pitfall Never write an angle such as $\angle AAB$.
If a point may coincide with a vertex, either handle it directly or
use a representation that still makes sense. Directed angles solve
orientation problems, not undefined-point problems.

\section{Practice}
\begin{exercise}[\problemlevel{Foundation}]\label{g:ex07}
Two similar triangles have scale factor $k>0$. Prove that their
perimeters have ratio $k$ and their areas have ratio $k^2$.
\end{exercise}
\begin{exercise}[\problemlevel{Foundation}]\label{g:ex08}
In triangle $ABC$, the internal bisector of $\angle A$ meets $BC$
at $D$. Derive $BD/DC=AB/AC$ by comparing areas.
\end{exercise}
\begin{exercise}[\problemlevel{Bridge}]\label{g:ex09}
In a circle, two chords $AB,CD$ meet at an interior point $P$.
Prove $PA\cdot PB=PC\cdot PD$ using similar triangles.
\end{exercise}
\begin{exercise}[\problemlevel{Bridge}]\label{g:ex10}
In triangle $ABC$, let $E,F$ be the feet from $B,C$. Prove
$B,C,E,F$ are concyclic without assuming the triangle acute.
If a foot coincides with a vertex, explain the correct direct
interpretation.
\end{exercise}
\begin{exercise}[\problemlevel{Bridge}]\label{g:ex11}
In an acute triangle, prove that each altitude bisects the angle
at its foot in the triangle formed by the three altitude feet.
\end{exercise}
\begin{exercise}[\problemlevel{Challenge}]\label{g:ex12}
Let $ABC$ be acute, with orthocenter $H$. Reflect $H$ in line $BC$
to $H'$. Prove that $H'$ lies on the circumcircle of $ABC$.
Then find a proof which works for an obtuse triangle as well.
\end{exercise}

\chapter{Ratios, Areas, and Concurrency}
\label{g:ratios}

\section{Choose an area that cancels}

Write $[ABC]$ for the positive area of triangle $ABC$. If $B,D,C$
are collinear with $D$ between $B,C$, triangles $ABD,ADC$ have the
same altitude from $A$, so
\[
 \frac{[ABD]}{[ADC]}=\frac{BD}{DC}.
\]
If two triangles share two side directions and an included angle,
the formula $[ABC]=\tfrac12 AB\cdot AC\sin\angle BAC$ gives
another useful ratio. These two observations convert many
concurrency problems into multiplication and cancellation.

\begin{example}[The angle-bisector theorem]
The internal bisector $AD$ of $\angle BAC$ meets $BC$ at $D$.
Prove $BD/DC=AB/AC$.
\end{example}
\begin{solution}
Compute the area ratio twice:
\[
 \frac{BD}{DC}
 =\frac{[ABD]}{[ACD]}
 =\frac{AB\cdot AD\sin\angle BAD}
        {AC\cdot AD\sin\angle DAC}
 =\frac{AB}{AC}.
\]
Conversely, the unique point on $BC$ with this positive division ratio
must be the point where the internal bisector meets $BC$.
\end{solution}
\intuition The common side $AD$ disappears. That is why areas are
preferable to calculating the length of the bisector. Look for a
quantity that can be introduced temporarily and then canceled.

\begin{figure}[htbp]\centering
\begin{tikzpicture}[scale=.85]
\coordinate (A) at (.5,3); \coordinate (B) at (-2,0);
\coordinate (C) at (3,0); \coordinate (D) at (.1,0);
\fill[accent!10] (A)--(B)--(D)--cycle;
\fill[accent!23] (A)--(D)--(C)--cycle;
\draw (A)--(B)--(C)--cycle; \draw[main] (A)--(D);
\draw[aux] (A)--(.5,0);
\foreach \X/\pos in {A/above,B/below left,C/below right,D/below}
 \fill (\X) circle(1.2pt) node[\pos]{$\X$};
\node at (-.7,.8){$S_1$}; \node at (1.25,.8){$S_2$};
\end{tikzpicture}
\caption{A common altitude turns an area ratio into a base ratio:
$S_1/S_2=BD/DC$. The picture illustrates the area argument,
not an asserted angle bisector.}
\end{figure}

\section{Ceva's theorem: three ratios fit together}

A cevian is a line joining a vertex of a triangle to a point on the
opposite side line. Lines are concurrent when they pass through one
point. We first work with points in the open sides, so all lengths
and areas in the next proof are positive.

\begin{theorem}[Ceva, interior form]\label{g:thm:ceva}
If $D,E,F$ lie on the open sides $BC,CA,AB$, respectively, then
$AD,BE,CF$ are concurrent if and only if
\[
 \frac{BD}{DC}\frac{CE}{EA}\frac{AF}{FB}=1.
\]
\end{theorem}
\begin{proof}
Suppose they meet at $P$. Put $x=[PBC]$, $y=[PCA]$, and $z=[PAB]$.
Both $[ABD]/[ADC]$ and $[PBD]/[PDC]$ equal $BD/DC$,
by the common-altitude area formula.
Subtracting these proportional areas gives
\[
\frac{[PAB]}{[PAC]}
=\frac{[ABD]-[PBD]}{[ADC]-[PDC]}=\frac{BD}{DC}.
\]
Thus $BD/DC=z/y$. Likewise $CE/EA=x/z$ and
$AF/FB=y/x$. Their product is $1$.

Conversely, let $P=AD\cap BE$; it lies inside the triangle. Let $CP$
meet $AB$ at $F'$. The proved direction gives
$AF'/F'B=AF/FB$. A positive ratio determines a unique point of $AB$,
so $F'=F$.
\end{proof}

\begin{figure}[htbp]\centering
\begin{tikzpicture}[scale=.84]
\coordinate (A) at (0,3.4); \coordinate (B) at (-2,0);
\coordinate (C) at (3,0); \coordinate (P) at (.35,1.05);
\path[name path=ad] (A)--($(A)!3!(P)$);
\path[name path=bc] (B)--(C);
\path[name intersections={of=ad and bc,by=D}];
\path[name path=be] (B)--($(B)!3!(P)$);
\path[name path=ca] (C)--(A);
\path[name intersections={of=be and ca,by=E}];
\path[name path=cf] (C)--($(C)!3!(P)$);
\path[name path=ab] (A)--(B);
\path[name intersections={of=cf and ab,by=F}];
\draw (A)--(B)--(C)--cycle;
\draw[main] (A)--(D) (B)--(E) (C)--(F);
\foreach \X/\pos in {A/above,B/below left,C/below right,D/below,E/right,F/left,P/above right}
 \fill (\X) circle(1.2pt) node[\pos]{$\X$};
\end{tikzpicture}
\caption{Ceva tests concurrence using only the three side divisions.
The interior point need not be found explicitly.}
\end{figure}

\begin{example}[Equal-area cevians]
The medians of a triangle are concurrent.
\end{example}
\begin{solution}
Each median has side ratio $1$, so Ceva applies. Their intersection
is the centroid. To obtain its position, let $D,E$ be midpoints of
$BC,CA$. Since $DE\parallel BA$ and $DE=BA/2$, triangles $PDE,PAB$
are similar; thus $PD/PA=1/2$. Every median is divided in ratio
$2:1$, measured from its vertex.
\end{solution}

\section{Signed ratios and Menelaus}

Fix a coordinate increasing from $B$ toward $C$ on their line.
For points $X,Y$ on that line write $\overline{XY}=y-x$.
The ratio $\overline{BD}/\overline{DC}$ is positive when $D$ lies
between $B,C$ and negative otherwise. Reversing the coordinate
orientation changes neither ratio. Use the corresponding cyclic
convention on $CA$ and $AB$.

\begin{theorem}[Menelaus and signed Ceva]
For finite points $D\in BC$, $E\in CA$, $F\in AB$, none a vertex,
the points $D,E,F$ are collinear if and only if
\[
 \frac{\overline{BD}}{\overline{DC}}
 \frac{\overline{CE}}{\overline{EA}}
 \frac{\overline{AF}}{\overline{FB}}=-1.
\]
The product $+1$ instead means that $AD,BE,CF$ are concurrent or
all parallel. Interior points exclude the parallel possibility.
\end{theorem}
\begin{proof}
Use affine coordinates: $A=(1,0,0)$, $B=(0,1,0)$, $C=(0,0,1)$,
where triples represent points after normalization so their sum is $1$.
Write $D=(0,d_2,d_3)$, $E=(e_1,0,e_3)$,
$F=(f_1,f_2,0)$. The side ratios are
$d_3/d_2,e_1/e_3,f_2/f_1$. Three points are collinear precisely
when the determinant of these three rows is zero, namely
$d_2e_3f_1+d_3e_1f_2=0$. This is Menelaus.
The cevian equations are
$d_3y=d_2z$, $e_1z=e_3x$, $f_2x=f_1y$.
They have a nonzero common solution exactly when the product is $1$.
If its coordinate sum is nonzero, it represents a finite common point.
If the sum is zero, it represents their common direction: in the
ordinary plane the three lines are parallel. The same determinant
criterion for affine collinearity follows by subtracting one row
from the other two, reducing it to the familiar zero-area test.
\end{proof}
\pitfall The signs above are part of the theorem. Some books use
$\overline{CD}$ instead of $\overline{DC}$ and therefore reverse the
displayed signs. Translate conventions before multiplying ratios.

\begin{example}[An external bisector without a diagram trap]
The external bisector at $A$ meets line $BC$ at a finite point $D$.
Find its signed ratio.
\end{example}
\begin{solution}
The two triangles with vertex $A$ have positive area ratio
$AB/AC$, because the relevant angle sines are equal. The division is
external, so
$\overline{BD}/\overline{DC}=-AB/AC$.
If $AB=AC$, no finite external division point exists: the external
bisector is parallel to $BC$. This is a genuine exception, not a
point that may be divided out of the proof.
\end{solution}

\section{Trigonometric Ceva}

For interior cevians the sine area formula gives
\[
 \frac{BD}{DC}
 =\frac{AB\sin\angle BAD}{AC\sin\angle DAC}.
\]
Consequently concurrency is equivalent to
\[
 \frac{\sin\angle BAD}{\sin\angle DAC}
 \frac{\sin\angle CBE}{\sin\angle EBA}
 \frac{\sin\angle ACF}{\sin\angle FCB}=1.
\]
Notice the order at $B$: the numerator is the angle adjacent to $BC$.
This order is chosen to make the factors cancel cyclically.

\recognition Use ordinary Ceva when side ratios are accessible; use
trigonometric Ceva when angle equalities or isogonal reflections are
accessible. Do not compute side lengths merely to apply a theorem
which has an angle version.

\section{Practice}
\begin{exercise}[\problemlevel{Foundation}]\label{g:ex13}
The internal angle bisectors of a triangle are concurrent.
Prove this first by Ceva, then by equal distances to lines.
\end{exercise}
\begin{exercise}[\problemlevel{Bridge}]\label{g:ex14}
In triangle $ABC$, points $D,E$ divide $BC,CA$ with
$BD/DC=2$ and $CE/EA=3$. Find $AF/FB$ when $AD,BE,CF$
are concurrent.
\end{exercise}
\begin{exercise}[\problemlevel{Bridge}]\label{g:ex15}
Let $D,E,F$ be the midpoints of $BC,CA,AB$. The line through $D,E$
is parallel to $AB$. Explain why Menelaus cannot be applied using
an alleged finite intersection of that line with $AB$.
\end{exercise}
\begin{exercise}[\problemlevel{Bridge}]\label{g:ex16}
A line meets the side lines $BC,CA,AB$ at nonvertex points $D,E,F$.
If $BD/DC=2$ internally and $CE/EA=3$ internally, determine whether
$F$ is internal or external and find its unsigned ratio $AF/FB$.
\end{exercise}
\begin{exercise}[\problemlevel{Challenge}]\label{g:ex17}
Points $D,E,F$ lie in the open sides $BC,CA,AB$ with
$BD/DC=CE/EA=AF/FB=t>0$. The cevians $AD,BE,CF$ bound a triangle,
possibly collapsed. Show that its area divided by $[ABC]$ is
\[
 \frac{(t-1)^2}{t^2+t+1}.
\]
\end{exercise}
\begin{exercise}[\problemlevel{Challenge}]\label{g:ex18}
Let $P$ lie inside triangle $ABC$, and let $D,E,F$ be where
$AP,BP,CP$ meet opposite sides. Prove
\[
 \frac{PD}{AD}+\frac{PE}{BE}+\frac{PF}{CF}=1.
\]
Explain why replacing $PD/AD$ by $AP/PD$ would not give a constant sum.
\end{exercise}

\chapter{Triangle Centers and Their Relationships}
\label{g:centers}

\section{A center is an intersection with a reason}

The \emph{circumcenter} $O$ is the intersection of the perpendicular
bisectors. It exists because the first two bisectors give a point
equidistant from all three vertices, which therefore lies on the third.
Its common distance $R$ is the circumradius. The corresponding circle
is written $(ABC)$. The circumcenter can lie outside an obtuse triangle.

The \emph{incenter} $I$ is the intersection of the internal angle
bisectors. A point on an angle bisector is equidistant from its two
side lines: perpendicular projections give congruent right triangles.
Conversely, a point inside the angle with equal distances lies on
its bisector. Hence $I$ is equidistant from the three sides and is
the center of the inscribed circle, of radius $r$.

The \emph{orthocenter} $H$ is the intersection of the altitude lines.
Here is a useful existence proof. Through each vertex draw the
parallel to the opposite side, obtaining a larger triangle whose
side midpoints are $A,B,C$. The altitude lines of $ABC$ are the
perpendicular bisectors of the larger triangle; they concur.
The \emph{centroid} $G$ is the median intersection already constructed.

\begin{figure}[htbp]\centering
\begin{tikzpicture}[scale=.88]
\coordinate (A) at (0,3); \coordinate (B) at (-2,0);
\coordinate (C) at (3,0); \coordinate (H) at (0,2);
\coordinate (O) at (.5,.5); \coordinate (G) at (.3333,1);
\draw (A)--(B)--(C)--cycle;
\draw[aux] (A)--(0,-.2);
\draw[aux] (B)--(.7692,1.8462);
\draw[main] ($(O)!-.35!(H)$)--($(O)!1.3!(H)$);
\foreach \X/\pos in {A/above,B/below left,C/below right,H/left,O/right,G/right}
 \fill (\X) circle(1.2pt) node[\pos]{$\X$};
\end{tikzpicture}
\caption{In this triangle $O=(1/2,1/2)$, $H=(0,2)$, and
$G=(1/3,1)$. Their relation is a vector identity, so it is not
restricted to acute triangles.}
\end{figure}

\section{The Euler line and the nine-point circle}

\intuition Midpoints suggest a homothety of ratio $1/2$; centers of
circles suggest equal distances. These two thoughts explain the
nine-point circle more efficiently than nine separate angle chases.
For the algebra below a vector is a directed displacement, and
$u\cdot v$ is the Euclidean dot product.

\begin{theorem}[Euler-line identity]
Take the circumcenter as origin, and let $a,b,c$ be the position
vectors of the vertices. Then
\[
 h=a+b+c,\qquad g=\frac{a+b+c}{3}.
\]
Thus $O,G,H$ are collinear and $OG:GH=1:2$ unless all coincide.
\end{theorem}
\begin{proof}
All three vertex vectors have squared length $R^2$. Therefore
\[
 (b+c)\cdot(b-c)=|b|^2-|c|^2=0.
\]
The point with vector $a+b+c$ lies on the altitude from $A$.
The analogous calculations put it on the other two altitudes.
The centroid formula follows by its $2:1$ division of a median.
\end{proof}

\begin{theorem}[Nine-point circle]
The three side midpoints, three altitude feet, and three midpoints
of $AH,BH,CH$ lie on one circle. Its center is
$N=(O+H)/2$ and its radius is $R/2$.
\end{theorem}
\begin{proof}
Continue with $O=0$, $h=a+b+c$, $n=h/2$.
For the midpoint $M$ of $BC$, $m-n=-a/2$, so $NM=R/2$.
For the midpoint $U$ of $AH$, $u-n=a/2$, giving the same distance.
This proves six incidences. Reflect $H$ across $BC$ to $H'$.
The identity $H'\in(ABC)$ is proved below. The altitude foot $D$
is the midpoint of $HH'$. Hence the homothety with center $H$ and
ratio $1/2$ takes $H'$ to $D$ and the circumcircle to the circle
with center $N$ and radius $R/2$. This proves the remaining incidences.
\end{proof}

\begin{example}[Orthocenter reflection, including obtuse triangles]
Show that the reflection of $H$ in each side line lies on the
circumcircle.
\end{example}
\begin{solution}
Choose $B=(0,0)$, $C=(1,0)$, $A=(u,v)$ with $v\ne0$.
The altitude from $A$ is vertical, while the altitude from $B$ is
perpendicular to $(u-1,v)$. Their intersection is
$H=(u,u(1-u)/v)$. The circumcircle has equation
\[
 x^2+y^2-x+\frac{u-u^2-v^2}{v}y=0.
\]
Reflection in $BC$ changes the sign of the second coordinate.
Substitution gives zero, proving the result. This computation
also includes right triangles, exterior feet, and $u=0$ or $u=1$.
\end{solution}
\strategy Use a short coordinate proof when a synthetic angle chase
requires a distracting collection of special cases. Coordinates are
a proof method, not a concession.

\section{Contact lengths and excenters}

Put $a=BC$, $b=CA$, $c=AB$, and $s=(a+b+c)/2$.
Equal tangent segments from a point show that the incircle contacts
the sides at distances $s-a,s-b,s-c$ from the corresponding vertices.
For example, if it touches $AB,AC$ at $F,E$, write $AF=AE=x$,
$BF=BD=y$, $CE=CD=z$. Then $c=x+y$, $a=y+z$, $b=z+x$,
so $x=s-a$, $y=s-b$, $z=s-c$.

\begin{figure}[htbp]\centering
\begin{tikzpicture}[scale=.85]
\coordinate (A) at (0,3); \coordinate (B) at (-2,0);
\coordinate (C) at (3,0);
\pgfmathsetmacro{\bb}{sqrt(18)}
\pgfmathsetmacro{\cc}{sqrt(13)}
\pgfmathsetmacro{\ss}{(5+\bb+\cc)/2}
\pgfmathsetmacro{\ix}{(-2*\bb+3*\cc)/(5+\bb+\cc)}
\pgfmathsetmacro{\iy}{15/(5+\bb+\cc)}
\coordinate (I) at (\ix,\iy);
\coordinate (D) at (\ix,0);
\coordinate (E) at ($(A)!(I)!(C)$);
\coordinate (F) at ($(A)!(I)!(B)$);
\draw (A)--(B)--(C)--cycle;
\draw[main] (I) circle(\iy);
\draw[aux] (I)--(D) (I)--(E) (I)--(F);
\foreach \X/\pos in {A/above,B/below left,C/below right,I/above,D/below,E/right,F/left}
 \fill (\X) circle(1.2pt) node[\pos]{$\X$};
\end{tikzpicture}
\caption{Equal tangent lengths are linear variables. They often make
incircle problems simpler than side lengths do.}
\end{figure}

\begin{example}[Gergonne concurrence]
Join each vertex to the incircle's contact point on the opposite side.
The three lines are concurrent.
\end{example}
\begin{solution}
Their Ceva product is
\[
 \frac{s-b}{s-c}\frac{s-c}{s-a}\frac{s-a}{s-b}=1.
\]
The denominators are positive by the triangle inequalities.
The intersection is called the Gergonne point.
\end{solution}

An \emph{excenter} is the intersection of one internal and two
external angle bisectors. There are three, one opposite each vertex.
The $A$-excircle touches $BC$ and the extensions of $AB,AC$.
Its center $I_A$ has equal distances from these lines, and its
radius is $r_A=[ABC]/(s-a)$. The area equation follows by subtracting
the two small exterior triangles from the large one, or by signed
areas:
$[ABC]=\tfrac12r_A(b+c-a)=r_A(s-a)$.
Likewise $[ABC]=rs$ for the incircle, by summing three triangles
with apex $I$.

\begin{example}[A formula with an accessible meaning]
Prove $rR=abc/(4s)$.
\end{example}
\begin{solution}
The extended sine law, proved in the next chapter, gives
$[ABC]=abc/(4R)$, while the preceding area decomposition gives
$[ABC]=rs$. Equating and rearranging yields the formula.
The point is to recognize two interpretations of the same area.
\end{solution}

\pitfall A center's definition is more useful than its name.
When the problem mentions $O$, look for equal radii or perpendicular
bisectors. For $I$, look for equal distances or half angles. For $H$,
look for perpendicular directions. For $G$, look for affine averages.

\section{Practice}
\begin{exercise}[\problemlevel{Foundation}]\label{g:ex19}
Prove that the centroid divides each median in ratio $2:1$ and that
the three medians partition the triangle into six equal-area triangles.
\end{exercise}
\begin{exercise}[\problemlevel{Bridge}]\label{g:ex20}
Show that $[ABC]=rs$ and derive the contact lengths $s-a,s-b,s-c$
without assuming a diagram is to scale.
\end{exercise}
\begin{exercise}[\problemlevel{Bridge}]\label{g:ex21}
The $A$-excircle touches $BC$ at $D_A$.
Show $BD_A=s-c$, $CD_A=s-b$ and prove that
$AD_A,BD_B,CD_C$ are concurrent, where $D_B,D_C$ are the
analogous opposite-side contacts.
\end{exercise}
\begin{exercise}[\problemlevel{Bridge}]\label{g:ex22}
If $ABC$ is not equilateral, show that the circumcenter of the
medial triangle is the nine-point center of $ABC$.
Why is the assertion still true in the equilateral case?
\end{exercise}
\begin{exercise}[\problemlevel{Challenge}]\label{g:ex23}
Prove that the reflection of the orthocenter in the midpoint of
each side lies on the circumcircle. Distinguish this reflection
from reflection in the side line.
\end{exercise}
\begin{exercise}[\problemlevel{Challenge}]\label{g:ex24}
For an acute triangle with altitude feet $D,E,F$, prove that its
orthocenter is the incenter of $DEF$. Identify what changes when
the original triangle is obtuse.
\end{exercise}

\chapter{Lengths and Areas without Unnecessary Calculation}
\label{g:metric}

\section{Projection is the source of metric identities}

Keep $a=BC$, $b=CA$, $c=AB$, and write $A,B,C$ also for the
three vertex angles when no confusion is possible. Dropping a
perpendicular gives Pythagoras in a right triangle. One proof cuts
a square of side $x+y$ into four right triangles with legs $x,y$
and a central square of side $z$: comparing areas gives
$(x+y)^2=4xy/2+z^2$, hence $z^2=x^2+y^2$.

The same calculation with a signed projection gives
\[
 a^2=b^2+c^2-2bc\cos A.
\]
Indeed choose $A=(0,0)$, $B=(c,0)$, and
$C=(b\cos A,b\sin A)$. Squaring the distance $BC$ proves the
identity for acute, right, and obtuse $A$ at once.

\intuition The cosine law is a way to eliminate an angle when two
sides and the included angle occur. It is also a way to \emph{recognize}
an angle from side lengths. A sum such as $AB^2+AC^2-BC^2$
usually represents twice a dot product.

\begin{example}[A median without solving for its direction]
If $M$ is the midpoint of $BC$, prove
\[
 AM^2=\frac{2b^2+2c^2-a^2}{4}.
\]
\end{example}
\begin{solution}
Put $B=(-a/2,0)$, $C=(a/2,0)$, and $A=(u,v)$.
Then $c^2=(u+a/2)^2+v^2$ and
$b^2=(u-a/2)^2+v^2$. Adding cancels the mixed terms and gives
$b^2+c^2=2AM^2+a^2/2$. This is Apollonius's median identity.
\end{solution}

\begin{theorem}[Stewart's theorem]
If $D$ lies on segment $BC$, with $BD=m$, $DC=n$, and $AD=d$,
then
\[
 b^2m+c^2n=a(d^2+mn),\qquad a=m+n.
\]
\end{theorem}
\begin{proof}
Set $D=(0,0)$, $B=(-m,0)$, $C=(n,0)$, and $A=(u,v)$.
Then $c^2=d^2+2mu+m^2$ and
$b^2=d^2-2nu+n^2$. Multiply the latter by $m$, the former
by $n$, and add. The projection $u$ cancels.
\end{proof}

\begin{figure}[htbp]\centering
\begin{tikzpicture}[scale=.88]
\coordinate (A) at (.5,3); \coordinate (B) at (-2,0);
\coordinate (C) at (3,0); \coordinate (D) at (0,0);
\draw (A)--(B)--(C)--cycle; \draw[main] (A)--(D);
\draw[aux] (A)--(.5,0);
\foreach \X/\pos in {A/above,B/below left,C/below right,D/below}
 \fill (\X) circle(1.2pt) node[\pos]{$\X$};
\node at (-1,-.3){$m$}; \node at (1.6,-.3){$n$};
\node at (.05,1.5){$d$};
\end{tikzpicture}
\caption{Stewart eliminates the signed horizontal projection of $A$.
The cevian need not be an altitude, median, or angle bisector.}
\end{figure}

\section{The sine law and four expressions for area}

\begin{theorem}[Extended sine law]
In a triangle with circumradius $R$,
\[
 \frac{a}{\sin A}=\frac{b}{\sin B}=\frac{c}{\sin C}=2R.
\]
\end{theorem}
\begin{proof}
For chord $BC$, choose the endpoint $X$ of a diameter from $B$.
The right triangle $BCX$ gives $BC=2R\sin\angle BXC$.
By the inscribed-angle theorem, that angle is equal or supplementary
to $A$, so its sine is $\sin A$. The other identities are cyclic.
\end{proof}

Consequently, with $\Delta=[ABC]$,
\[
 \Delta=\frac12bc\sin A
        =\frac{abc}{4R}
        =rs
        =\sqrt{s(s-a)(s-b)(s-c)}.
\]
The last is Heron's formula. To prove it rather than memorize it,
use the cosine law:
\[
 16\Delta^2
 =4b^2c^2-(b^2+c^2-a^2)^2
 =((b+c)^2-a^2)(a^2-(b-c)^2).
\]
Factoring the two differences of squares gives
$16s(s-a)(s-b)(s-c)$.

\begin{example}[A formula chosen for the known data]
Find the inradius and circumradius of a triangle with side lengths
$13,14,15$.
\end{example}
\begin{solution}
Here $s=21$ and Heron gives
$\Delta=\sqrt{21\cdot8\cdot7\cdot6}=84$.
Thus $r=\Delta/s=4$ and
$R=abc/(4\Delta)=65/8$.
No altitude direction or center coordinate needs to be found.
\end{solution}
\recognition Side lengths favor Heron; two sides and an included
angle favor $\tfrac12bc\sin A$; an incircle favors $rs$; a
circumcircle favors $abc/(4R)$. Switching among these expressions
is often the whole argument.

\section{Ptolemy: a cyclic quadrilateral becomes a length identity}

\begin{theorem}[Ptolemy]
For a convex cyclic quadrilateral $ABCD$ in that order,
\[
 AC\cdot BD=AB\cdot CD+AD\cdot BC.
\]
\end{theorem}
\begin{proof}
Choose $E$ on $AC$ with $\angle ABE=\angle DBC$.
Such a ray lies inside $\angle ABC$ and meets the diagonal.
Since $\angle BAE=\angle BDC$, triangles $ABE,DBC$ are similar;
hence $AE\cdot BD=AB\cdot CD$.
Also $\angle EBC=\angle ABD$ and
$\angle BCE=\angle BDA$, so triangles $EBC,ABD$ are similar;
hence $EC\cdot BD=AD\cdot BC$. Add and use $AE+EC=AC$.
\end{proof}

\begin{figure}[htbp]\centering
\begin{tikzpicture}[scale=.9]
\coordinate (A) at (145:2); \coordinate (B) at (225:2);
\coordinate (C) at (320:2); \coordinate (D) at (55:2);
\draw (0,0) circle(2); \draw (A)--(B)--(C)--(D)--cycle;
\draw[main] (A)--(C) (B)--(D);
\foreach \X/\pos in {A/above left,B/below left,C/below right,D/above right}
 \fill (\X) circle(1.2pt) node[\pos]{$\X$};
\end{tikzpicture}
\caption{For four cyclic points, a product of diagonals becomes a sum
of products of opposite sides. The vertex order controls the plus sign.}
\end{figure}

\begin{example}[A moving point and a fixed sum]
Let $ABC$ be equilateral and $P$ lie on the arc $BC$ of its
circumcircle not containing $A$. Show $PA=PB+PC$.
\end{example}
\begin{solution}
The cyclic order is $A,B,P,C$. Ptolemy gives
$AP\cdot BC=AB\cdot PC+AC\cdot PB$.
Cancel the common positive side length.
This suggests a general habit: an addition of two lengths may be
Ptolemy after multiplication by a suitable common chord.
\end{solution}

\extension A general quadrilateral satisfies Ptolemy's inequality
$AC\cdot BD\le AB\cdot CD+AD\cdot BC$.
An elegant proof by inversion will appear in
Chapter~\ref{g:inversion}. Equality for a noncollinear quadrilateral
requires the indicated convex cyclic order.

\section{Practice}
\begin{exercise}[\problemlevel{Foundation}]\label{g:ex25}
Use the cosine law to show that $A$ is acute, right, or obtuse
according as $a^2$ is less than, equal to, or greater than $b^2+c^2$.
\end{exercise}
\begin{exercise}[\problemlevel{Bridge}]\label{g:ex26}
If the internal bisector from $A$ has length $\ell$, prove
\[
 \ell^2=bc\left(1-\frac{a^2}{(b+c)^2}\right).
\]
\end{exercise}
\begin{exercise}[\problemlevel{Bridge}]\label{g:ex27}
Prove that the sum of the squares of the three medians is
$\tfrac34(a^2+b^2+c^2)$.
\end{exercise}
\begin{exercise}[\problemlevel{Bridge}]\label{g:ex28}
For a convex cyclic quadrilateral with side lengths $a,b,c,d$ in
order and semiperimeter $s$, derive
$K^2=(s-a)(s-b)(s-c)(s-d)$, where $K$ is its area.
\end{exercise}
\begin{exercise}[\problemlevel{Challenge}]\label{g:ex29}
An interior point $P$ of an equilateral triangle $ABC$ satisfies
$PA=3$, $PB=4$, $PC=5$. Find the side length of $ABC$.
\end{exercise}
\begin{exercise}[\problemlevel{Challenge}]\label{g:ex30}
Prove that among convex quadrilaterals with four prescribed side
lengths which admit a nondegenerate quadrilateral, the maximum area
is attained by a cyclic one. Derive the maximum from those lengths.
\end{exercise}

\chapter{Power, Radical Axes, and Circle Networks}
\label{g:power}

\section{One number for every secant}

For a circle $\omega$ with center $O$ and radius $R$, define
$\operatorname{Pow}_{\omega}(P)=PO^2-R^2$.
It is positive outside the circle, negative inside, and zero on it.
Its sign is useful information.

\intuition Several lines through $P$ cut the same circle, producing
apparently unrelated segment products. They agree because all measure
one fixed quantity, the power of $P$. Introducing the circle's center
explains the invariant; using secants often avoids computing the center.

\begin{theorem}[Secant and tangent formulas]
If a line through $P$ meets $\omega$ at $X,Y$, then, using one
orientation along that line,
\[
 \overline{PX}\,\overline{PY}=\operatorname{Pow}_{\omega}(P).
\]
If $P$ is external and $PT$ is tangent, then
$PT^2=\operatorname{Pow}_{\omega}(P)$.
\end{theorem}
\begin{proof}
Take $P$ as coordinate origin on the secant and let $M$ be the
projection of $O$ onto it, with signed coordinate $m$. The intersection
coordinates satisfy
$(t-m)^2+OM^2=R^2$. Their product is
$m^2+OM^2-R^2=PO^2-R^2$ by the quadratic root-product formula.
For a tangent, $OT\perp PT$, so Pythagoras gives the same result.
\end{proof}

\begin{figure}[htbp]\centering
\begin{tikzpicture}[scale=.85]
\coordinate (O) at (0,0); \coordinate (P) at (4,0);
\coordinate (X) at (2,0); \coordinate (Y) at (-2,0);
\coordinate (T) at (1,1.73205);
\draw (O) circle(2); \draw (P)--(Y);
\draw[main] (P)--(T); \draw[aux] (O)--(T);
\foreach \X/\pos in {O/below,P/right,X/below,Y/left,T/above}
 \fill (\X) circle(1.2pt) node[\pos]{$\X$};
\end{tikzpicture}
\caption{An external point has
$PX\cdot PY=PT^2=PO^2-R^2$. For an internal point the signed
secant product is negative.}
\end{figure}

\begin{example}[The converse needs a configuration]
Distinct points $A,B,C,D$ lie on two distinct lines through $P$,
with $P\notin\{A,B,C,D\}$. Suppose
$\overline{PA}\,\overline{PB}
=\overline{PC}\,\overline{PD}$. Show they are concyclic.
\end{example}
\begin{solution}
Let the circle through $A,B,C$ meet line $PC$ a second time at $D'$,
counting a tangency twice. Its power formula gives
$\overline{PC}\,\overline{PD'}
=\overline{PC}\,\overline{PD}$. Since $P\ne C$, cancellation yields
$D'=D$. If $D'=C$, this would force $D=C$, contrary to the distinctness
assumption. Thus $D$ is genuinely the fourth point on the circle.
\end{solution}

\section{Equal powers form a line}

\begin{theorem}[Radical axis]
For two nonconcentric circles, the set of points of equal power is
a line perpendicular to the line of centers. If the circles meet
at two points, it is their common-chord line; if tangent, it is
their common tangent.
\end{theorem}
\begin{proof}
In Cartesian coordinates the circle powers have forms
$x^2+y^2+u_ix+v_iy+w_i$. Subtraction cancels the quadratic terms.
Because the centers differ, at least one linear coefficient remains.
The resulting line has normal vector parallel to the difference
of the center vectors. Intersection points have both powers zero;
the tangent case follows from the same equation.
\end{proof}

Distinct concentric circles have no radical axis in the finite plane:
their powers differ by a nonzero constant. Coincident circles have
equal power everywhere, so they do not determine a unique radical axis.

\begin{theorem}[Radical center]
For three circles with noncollinear centers, their three radical
axes pass through one point.
\end{theorem}
\begin{proof}
The radical axes of pairs $(1,2)$ and $(2,3)$ are not parallel,
because their normals are the nonparallel center differences.
Let them meet at $P$. Then
$\operatorname{Pow}_1(P)=\operatorname{Pow}_2(P)
=\operatorname{Pow}_3(P)$, so $P$ lies on the third axis.
\end{proof}

\begin{figure}[htbp]\centering
\begin{tikzpicture}[scale=.78]
\coordinate (O1) at (-1,0); \coordinate (O2) at (1,0);
\coordinate (O3) at (0,2); \coordinate (P) at (0,.75);
\draw (O1) circle(1.55); \draw (O2) circle(1.55);
\draw (O3) circle(1.55);
\draw[main] (0,-1.8)--(0,3.65);
\draw[main] (-2.8,-.65)--(2.8,2.15);
\draw[main] (-2.8,2.15)--(2.8,-.65);
\foreach \X/\lab/\pos in {O1/O_1/left,O2/O_2/right,O3/O_3/above,P/P/right}
 \fill (\X) circle(1.2pt) node[\pos]{$\lab$};
\end{tikzpicture}
\caption{For equal-radius circles the radical axes are perpendicular
bisectors of center segments. Equal power is transitive, which forces
the three axes to concur.}
\end{figure}

\begin{example}[Concurrency without locating the intersection]
Two circles meet at $A,B$. A third circle meets them at $C,D$
and $E,F$, respectively. If lines $CD$ and $EF$ meet at $P$,
prove $P\in AB$.
\end{example}
\begin{solution}
Line $CD$ is the equal-power locus for circles 1 and 3, and $EF$
is that for 2 and 3. At $P$ all powers agree. Therefore $P$
belongs to the radical axis of circles 1 and 2, namely $AB$.
The existence of the stated intersection excludes the parallel case.
\end{solution}

\section{Miquel's theorem and a hidden tangent case}

\begin{theorem}[Miquel point of a triangle]
Let $D,E,F$ lie on the open sides $BC,CA,AB$ of a triangle.
The circles $(AEF)$, $(BFD)$, $(CDE)$ have a common point.
\end{theorem}
\begin{proof}
The first two circles already meet at $F$ and are distinct, since
a circle cannot contain the three distinct collinear points $A,F,B$.
If they meet again at $M\ne F$, then $M=D$ or $M=E$ proves the
claim immediately. Otherwise the directed-angle calculation gives
\[
 \dang DME=\dang DMF+\dang FME
 =\dang DBF+\dang FAE=\dang DCE.
\]
Thus $M$ belongs to the third circle.

It remains possible that the first two circles are tangent at $F$.
Let $\ell$ be their common tangent. The tangent--chord theorem gives
\[
 \dang EFD
 =\angle(FE,\ell)+\angle(\ell,FD)
 =\angle(AE,AF)+\angle(FB,BD)=\dang ECD.
\]
So the third circle contains $F$ itself. This completes a case which
an argument starting with a second intersection would miss.
\end{proof}

The tangent--chord theorem used here says that the angle between a
tangent at $X$ and chord $XY$ equals the inscribed angle subtending
$XY$, with the corresponding directed convention. To prove it, the
radius at $X$ is perpendicular to the tangent, and the isosceles
triangle formed by the two radii expresses the chord angle as
$90^\circ$ minus a base angle. The inscribed-angle theorem identifies
the remaining half-central angle.

\begin{example}[Euler's incenter formula]
For circumcenter $O$, incenter $I$, circumradius $R$, and inradius $r$,
prove
\[
 OI^2=R(R-2r).
\]
\end{example}
\begin{solution}
Let the internal bisector from $A$ meet the circumcircle again at $L$,
the midpoint of arc $BC$ not containing $A$. Then $LB=LC$.
In triangle $BIL$, an angle chase gives
$\angle LBI=\angle BIL=90^\circ-C/2$, hence $LI=LB$.
Also $LB=2R\sin(A/2)$ by the extended sine law, while
$AI=r/\sin(A/2)$ by the right triangle from $I$ to side $AB$.
The points occur in order $A,I,L$, so the power of $I$ gives
$OI^2-R^2=-IA\cdot IL=-2Rr$.
\end{solution}
\checkpoint The formula immediately implies $R\ge2r$.
Equality gives $I=O$, which forces the three side chords to be
equidistant from $O$, hence equal. Thus equality characterizes the
equilateral triangle.

\section{Practice}
\begin{exercise}[\problemlevel{Foundation}]\label{g:ex31}
Prove that tangent segments from the same external point to a circle
have equal lengths.
\end{exercise}
\begin{exercise}[\problemlevel{Bridge}]\label{g:ex32}
A point $P$ is inside a circle, and two perpendicular chords through
$P$ have endpoints $A,B$ and $C,D$. Prove that
$PA\cdot PB=PC\cdot PD$ and
$PA^2+PB^2+PC^2+PD^2=4R^2$, where $R$ is the radius.
\end{exercise}
\begin{exercise}[\problemlevel{Bridge}]\label{g:ex33}
Two nonconcentric circles have disjoint circumferences.
Construct their radical axis using an auxiliary circle which meets
both in two points.
\end{exercise}
\begin{exercise}[\problemlevel{Challenge}]\label{g:ex34}
Two circles intersect orthogonally: their tangents are perpendicular
at an intersection. Prove that this is equivalent to
$O_1O_2^2=r_1^2+r_2^2$. Show the condition also implies that the
circles meet in two distinct points.
\end{exercise}
\begin{exercise}[\problemlevel{Challenge}]\label{g:ex35}
For three circles with pairwise intersections, each pair determines
a common-chord line. Prove that these three lines are concurrent
or parallel, unless some coincide. Describe the center configuration
in the parallel case.
\end{exercise}
\begin{exercise}[\problemlevel{Challenge}]\label{g:ex36}
Let $L$ be the midpoint of arc $BC$ not containing $A$ in triangle
$ABC$. Prove $LB=LC=LI$, where $I$ is the incenter, and prove
$LB=LI_A$, where $I_A$ is the $A$-excenter.
\end{exercise}

\chapter{Transformations and Spiral Similarity}
\label{g:transformations}

\section{Move the whole configuration}

A translation moves every point by the same vector. A rotation about
$S$ through angle $\theta$ preserves distances and turns every
direction by $\theta$. Reflection in a line preserves distances and
reverses orientation. A homothety with center $S$ and nonzero real
ratio $k$ sends a point $X$ to $X'$ with vector equation
$X'-S=k(X-S)$. For $k<0$, the image lies on the opposite ray.
Lengths scale by $|k|$, and parallelism is preserved.

\intuition A transformation can turn two separate pieces of a figure
into one piece. If a length $PX$ is hard to compare with $QY$, try
finding a rigid motion that takes $X$ to $Y$. The image of $P$ may
then reveal the desired comparison.

\begin{example}[Reflection and the shortest broken path]
Points $B,C$ lie strictly on one side of a line $\ell$. Let $B'$ be
the reflection of $B$ in $\ell$. Find the minimum of $BX+XC$
for $X\in\ell$.
\end{example}
\begin{solution}
For every such $X$, $BX=B'X$, so
$BX+XC=B'X+XC\ge B'C$ by the triangle inequality.
The segment $B'C$ crosses $\ell$ exactly once, at $X_0$.
Equality holds precisely there. The construction is valid even if
the projections of $B,C$ coincide: the crossing then lies on their
common perpendicular.
\end{solution}

\begin{figure}[htbp]\centering
\begin{tikzpicture}[scale=.8]
\coordinate (B) at (-2,2); \coordinate (Bp) at (-2,-2);
\coordinate (C) at (3,3); \coordinate (X) at (0,0);
\draw (-3,0)--(4,0) node[right]{$\ell$};
\draw[main] (B)--(X)--(C); \draw[aux] (Bp)--(C);
\draw[aux] (B)--(Bp);
\foreach \X/\lab/\pos in {B/B/above left,Bp/B'/below left,C/C/above,X/X/below right}
 \fill (\X) circle(1.2pt) node[\pos]{$\lab$};
\end{tikzpicture}
\caption{Reflect one endpoint. The shortest broken path is obtained
when its reflected version becomes a straight segment.}
\end{figure}

\begin{example}[A midpoint configuration as a homothety]
Let $M,N,P$ be the midpoints of $BC,CA,AB$. Show that the centroid
is the center of a homothety taking $ABC$ to $MNP$.
\end{example}
\begin{solution}
The centroid divides $AM,BN,CP$ in ratio $2:1$, with itself between
each vertex and opposite midpoint. Thus the homothety of center
$G$ and ratio $-1/2$ takes $A,B,C$ to $M,N,P$.
It also takes the circumcircle to the medial triangle's circumcircle,
recovering its radius $R/2$.
\end{solution}

\section{Two circles and their homothety centers}

Circles with distinct centers $O_1,O_2$ and radii $r_1,r_2$ can
be related by homotheties centered on $O_1O_2$. The internal center
divides that segment in ratio $r_1:r_2$ and uses negative scale.
If $r_1\ne r_2$, an external center divides it externally in the
same unsigned ratio and uses positive scale.

These statements follow directly from
$O_2-S=k(O_1-S)$ and $r_2=|k|r_1$.
When $r_1=r_2$, the positive-scale map is a translation, so the
external homothety center is not a finite point. The internal
center is still the midpoint and the map is a half-turn.

\begin{figure}[htbp]\centering
\begin{tikzpicture}[scale=.78]
\coordinate (O1) at (0,0); \coordinate (O2) at (4,0);
\coordinate (S) at (-4,0); \coordinate (T) at (1.3333,0);
\draw (O1) circle(1); \draw (O2) circle(2);
\draw[aux] (-4.6,0)--(6.5,0);
\draw[main] (S)--(6,2.5);
\foreach \X/\lab/\pos in {O1/O_1/below,O2/O_2/below,S/S/below,T/T/below}
 \fill (\X) circle(1.2pt) node[\pos]{$\lab$};
\end{tikzpicture}
\caption{For radii $1,2$ and center distance $4$, the external and
internal homothety centers are $S=(-4,0)$ and $T=(4/3,0)$.
The marked oblique ray is an example of corresponding point directions,
not a claimed tangent.}
\end{figure}

\section{A spiral similarity combines turning and scaling}

A direct similarity is a rotation followed by a positive homothety
with the same center. It is called a \emph{spiral similarity}.
In complex coordinates it has the particularly simple form
\[
 f(z)=uz+v,\qquad u\ne0.
\]
The modulus $|u|$ is the scale factor and the argument of $u$ is
the rotation angle. If $u\ne1$, its fixed point is
$s=v/(1-u)$; then $f(z)-s=u(z-s)$.
If $u=1$ and $v\ne0$, it is a translation and has no finite center.
If $u=1,v=0$, it is the identity and no unique center is determined.

\begin{theorem}[Two-pair spiral center]
Given $A\ne B$ and $C\ne D$, there is a unique direct similarity
taking $A$ to $C$ and $B$ to $D$. Unless it is a nontrivial
translation, it has a finite center, except that the identity
has every point fixed.
\end{theorem}
\begin{proof}
In complex notation solve
$ua+v=c$, $ub+v=d$:
$u=(d-c)/(b-a)$ and $v=c-ua$.
The fixed-point alternatives follow from the preceding equations.
\end{proof}

\begin{figure}[htbp]\centering
\begin{tikzpicture}[scale=.85]
\coordinate (S) at (0,0); \coordinate (A) at (2,0);
\coordinate (B) at (1.2,1.2);
\coordinate (C) at (60:3); \coordinate (D) at (-.6588,2.4588);
\draw[main] (S)--(A)--(B)--cycle;
\draw (S)--(C)--(D)--cycle;
\draw[aux] (A)--(C) (B)--(D);
\draw[->] (.65,0) arc(0:60:.65);
\node at (.85,.4){$\theta$};
\foreach \X/\pos in {S/below left,A/below,B/left,C/right,D/above}
 \fill (\X) circle(1.2pt) node[\pos]{$\X$};
\end{tikzpicture}
\caption{Rotation through $60^\circ$ and scale factor $3/2$ take
$A,B$ to $C,D$. The same center connects both corresponding pairs.}
\end{figure}

\begin{example}[Why a rotation creates a useful sum]
Let equilateral triangle $ABC$ have positive orientation. Rotate a
point $P$ through $60^\circ$ about $A$ to $Q$, with $B$ mapping to $C$.
Then $PQ=PA$ and $CQ=BP$.
\end{example}
\begin{solution}
The triangle $APQ$ is equilateral, so $PQ=PA$. Rotation preserves
distances and takes segment $BP$ to $CQ$, so $CQ=BP$.
The triangle inequality in $PQC$ now gives
$PC\le PA+PB$. Equality has a geometric meaning: $P,Q,C$ must be
collinear in the proper order. This construction is the source of
several equilateral-triangle identities and minimum-sum arguments.
\end{solution}

\recognition An equilateral triangle suggests $60^\circ$ rotation;
a square suggests $90^\circ$; equal lengths in a reflected angle
suggest reflection. A requested sum of distances suggests straightening
a broken line after one such move.

\section{Practice}
\begin{exercise}[\problemlevel{Foundation}]\label{g:ex37}
Prove that reflection in a line is an isometry by decomposing a
displacement into components parallel and perpendicular to the line.
\end{exercise}
\begin{exercise}[\problemlevel{Bridge}]\label{g:ex38}
Two homotheties have centers $S,T$ and ratios $k,\ell$.
Find the affine formula for their composition and determine exactly
when it is a homothety, a translation, or the identity.
\end{exercise}
\begin{exercise}[\problemlevel{Bridge}]\label{g:ex39}
Find the direct similarity taking $0$ to $1+i$ and $1$ to $2+2i$
in the complex plane. Find its center and scale factor.
\end{exercise}
\begin{exercise}[\problemlevel{Challenge}]\label{g:ex40}
Two externally tangent circles have radii $r,s$, touching at $T$.
Prove that the negative homothety centered at $T$ takes one circle
to the other, and find its ratio. Explain what happens to an
arbitrary line through $T$.
\end{exercise}
\begin{exercise}[\problemlevel{Challenge}]\label{g:ex41}
For an equilateral triangle $ABC$ and an arbitrary point $P$, prove
$PA+PB\ge PC$. Determine the full equality locus.
\end{exercise}
\begin{exercise}[\problemlevel{Challenge}]\label{g:ex42}
Construct equilateral triangles externally on the three sides of an
arbitrary triangle. Prove that their centers form an equilateral
triangle. Make the orientation choices explicit.
\end{exercise}

\chapter{Isogonal Lines and Symmedians}
\label{g:symmedians}

\section{Reflect a cevian, not the whole triangle}

Two lines through $A$ are \emph{isogonal} in angle $BAC$ when
reflection in its internal angle bisector interchanges them.
For two interior cevians $AD,AE$, this means
$\angle BAD=\angle EAC$ and $\angle DAC=\angle BAE$.

\intuition An angle bisector treats the two sides symmetrically,
but a side ratio does not. Sines translate between the two:
reflecting a cevian inverts its sine ratio while the factor $AB/AC$
remains. This explains the square in the symmedian ratio.

\begin{theorem}[Isogonal ratio identity]
Let interior cevians $AD,AE$ be isogonal. Then
\[
 \frac{BD}{DC}\frac{BE}{EC}=\frac{AB^2}{AC^2}.
\]
\end{theorem}
\begin{proof}
The sine area formula gives
$BD/DC=(AB/AC)\sin\angle BAD/\sin\angle DAC$.
For $AE$ the sine ratio is its reciprocal. Multiply.
\end{proof}

The isogonal reflection of a median is a \emph{symmedian}.
If it meets $BC$ at $K$, then
\[
 \frac{BK}{KC}=\frac{AB^2}{AC^2}.
\]
The three symmedians concur by Ceva; their intersection is called
the symmedian point or Lemoine point.

\begin{figure}[htbp]\centering
\begin{tikzpicture}[scale=.84]
\coordinate (A) at (0,3); \coordinate (B) at (-2,0);
\coordinate (C) at (3,0); \coordinate (M) at (.5,0);
\coordinate (K) at (.09677,0);
\draw (A)--(B)--(C)--cycle;
\draw[aux] (A)--(M);
\draw[main] (A)--(K);
\foreach \X/\pos in {A/above,B/below left,C/below right,M/below right,K/below left}
 \fill (\X) circle(1.2pt) node[\pos]{$\X$};
\end{tikzpicture}
\caption{The median $AM$ and symmedian $AK$ are isogonal.
Here $AB^2=13$, $AC^2=18$, so $BK/KC=13/18$.
They coincide only when $AB=AC$.}
\end{figure}

\begin{example}[Why the symmedians concur]
Prove the three symmedians have a common interior point.
\end{example}
\begin{solution}
Their division ratios are $c^2/b^2$, $a^2/c^2$, and $b^2/a^2$.
Their product is $1$, so interior Ceva applies.
The reasoning also shows why the squares are compatible: every side
occurs once in a numerator and once in a denominator.
\end{solution}

\section{The tangent construction}

\begin{theorem}[Tangents locate a symmedian]
Suppose the tangents at $B,C$ to $(ABC)$ meet at a finite point $T$.
Then line $AT$ is the $A$-symmedian.
\end{theorem}
\begin{proof}
The tangent segments satisfy $TB=TC$. The tangent--chord theorem
and the sine law give
\[
 \frac{[ABT]}{[ACT]}
 =\frac{AB\sin\angle ABT}{AC\sin\angle ACT}
 =\frac{c\sin C}{b\sin B}=\frac{c^2}{b^2}.
\]
The circle lies inside the angle bounded by rays $TB,TC$.
Since $A\ne B,C$, ray $TA$ lies strictly between these rays, so
line $AT$ meets the open segment $BC$ at $K$.
The area ratio is also $BK/KC$, by comparing triangles $ABK,ACK$
and then scaling their common altitude from $K$ to $T$ along $AK$.
Thus $BK/KC=c^2/b^2$.
The isogonal ratio identity characterizes the symmedian.
\end{proof}

\begin{figure}[htbp]\centering
\begin{tikzpicture}[scale=.8]
\coordinate (O) at (0,0); \coordinate (B) at (210:2);
\coordinate (C) at (330:2); \coordinate (A) at (105:2);
\coordinate (T) at (0,-4);
\draw (O) circle(2); \draw (A)--(B)--(C)--cycle;
\draw[main] (B)--(T)--(C) (A)--(T);
\foreach \X/\pos in {A/above,B/left,C/right,T/below}
 \fill (\X) circle(1.2pt) node[\pos]{$\X$};
\end{tikzpicture}
\caption{The two tangents supply equal lengths, while the tangent--chord
theorem supplies angles. Together they produce the squared side ratio.}
\end{figure}

\pitfall Tangents at endpoints of a diameter are parallel. In that
case their intersection is not a finite point, and the theorem should
be read as a direction statement or replaced by the ratio definition.
Do not divide by the alleged distance to an infinite point.

\section{Isogonal conjugates}

For an interior point $P$, reflect lines $AP,BP,CP$ in the three
internal angle bisectors. Their reflected cevians concur at an
interior point $Q$, called the \emph{isogonal conjugate} of $P$.
Indeed each new side ratio is its old reciprocal multiplied by
the appropriate squared side ratio. The product remains $1$.
Reflecting again returns $P$, so the operation is an involution.

\begin{example}[Two familiar centers are isogonal]
In an acute triangle, the circumcenter and orthocenter are
isogonal conjugates.
\end{example}
\begin{solution}
At $A$, the angle between $AB$ and $AH$ is $90^\circ-B$,
while the angle between $AO$ and $AC$ is also $90^\circ-B$,
using the isosceles triangle $AOC$ and central-angle relations.
The analogous equality holds on the other side, so $AO,AH$ are
isogonal lines. Apply the same reasoning at $B,C$.
If one center lies on a side or equals a vertex, line-based
interpretations need direct checking; the clean interior formulation
here avoids those degeneracies.
\end{solution}

\begin{example}[A chord ratio from a tangent intersection]
Let $D\ne A$ be the second intersection of $AT$ with $(ABC)$
in the tangent construction. Show $DB/DC=AB/AC$.
\end{example}
\begin{solution}
For a chord through $A$, the inscribed-angle theorem gives
$\sin\angle BAD/\sin\angle DAC=BD/DC$.
The symmedian's sine ratio is $AB/AC$, by its ratio identity.
Thus the result follows. The equality is in unsigned lengths;
the directed-angle version gives the same positive sine ratio.
\end{solution}

\recognition Equal ratios of two chords often hide an isogonal pair.
Tangents at two vertices often hide a symmedian. If a median is
present together with angle reflection, expect a square of a side ratio.

\section{Practice}
\begin{exercise}[\problemlevel{Foundation}]\label{g:ex43}
Derive the symmedian side-ratio formula directly from the sine
area formula and the definition of isogonality.
\end{exercise}
\begin{exercise}[\problemlevel{Bridge}]\label{g:ex44}
Show that the isogonal conjugate of the centroid is the symmedian
point.
\end{exercise}
\begin{exercise}[\problemlevel{Bridge}]\label{g:ex45}
Prove that the internal bisector, median, altitude, and symmedian
from one vertex coincide if and only if the two adjacent sides
are equal.
\end{exercise}
\begin{exercise}[\problemlevel{Challenge}]\label{g:ex46}
For an interior point with positive barycentric coordinates
$(x:y:z)$, prove that its isogonal conjugate has coordinates
$(a^2/x:b^2/y:c^2/z)$. Define the coordinates before using them,
or return after Chapter~\ref{g:coordinates}.
\end{exercise}
\begin{exercise}[\problemlevel{Challenge}]\label{g:ex47}
Let $AD,BE,CF$ be concurrent interior cevians. Replace each by its
isogonal reflection. Prove concurrence using trigonometric Ceva,
without using barycentric coordinates.
\end{exercise}
\begin{exercise}[\problemlevel{Challenge}]\label{g:ex48}
In a triangle with $BC$ not a diameter of its circumcircle, tangents
at $B,C$ meet at $T$. Let $K=AT\cap BC$, and let $D\ne A$ be
the second intersection of $AT$ with the circumcircle. Prove
$BK/KC=(DB/DC)^2$ and determine the common value in terms
of $AB,AC$.
\end{exercise}

\chapter{Coordinates as a Geometric Language}
\label{g:coordinates}

\section{Choose coordinates to erase data}

Coordinates are most effective when they remove conditions before
calculation begins. A side placed on the horizontal axis makes its
altitudes vertical. A midpoint placed at the origin makes its endpoints
opposite vectors. A circle made into the unit circle turns conjugation
into reciprocation. The wrong coordinates retain every variable and
hide every geometric relation.

For vectors $u=(u_1,u_2)$, $v=(v_1,v_2)$ define
\[
 u\cdot v=u_1v_1+u_2v_2,\qquad
 u\times v=u_1v_2-u_2v_1.
\]
The dot product detects perpendicularity and gives
$(X-Y)\cdot(X-Y)=XY^2$. The signed cross product detects parallelism
and gives twice oriented area. Thus points $A,B,C$ are collinear
if and only if $(B-A)\times(C-A)=0$.

\begin{example}[The foot formula is one linear equation]
Find the projection $F$ of $P$ onto line $AB$, with $A\ne B$.
\end{example}
\begin{solution}
Write $F=A+t(B-A)$. The perpendicular condition is
$(P-F)\cdot(B-A)=0$, giving
\[
 t=\frac{(P-A)\cdot(B-A)}{|B-A|^2}.
\]
The denominator is positive. A value $t<0$ or $t>1$ simply says the
foot lies outside the segment; no new formula is required.
\end{solution}

\begin{figure}[htbp]\centering
\begin{tikzpicture}[scale=.8]
\coordinate (A) at (-2,0); \coordinate (B) at (2,1);
\coordinate (P) at (1,3); \coordinate (F) at ($(A)!(P)!(B)$);
\draw[->] (-2.7,0)--(3.5,0) node[right]{$x$};
\draw[->] (0,-.6)--(0,3.6) node[above]{$y$};
\draw[main] ($(A)!-.1!(B)$)--($(A)!1.2!(B)$);
\draw[aux] (P)--(F);
\foreach \X/\pos in {A/below left,B/below right,P/above,F/right}
 \fill (\X) circle(1.2pt) node[\pos]{$\X$};
\end{tikzpicture}
\caption{A projection belongs to a line and satisfies one dot-product
equation. This formulation automatically allows exterior feet.}
\end{figure}

\section{Circles and the Simson line}

A circle has an equation
$x^2+y^2+ux+vy+w=0$; completing squares identifies its center
$(-u/2,-v/2)$ and radius. Three noncollinear points determine
$u,v,w$ by three linear equations.

\begin{theorem}[Simson line]\label{g:thm:simson}
For a triangle $ABC$ and a point $P$, the perpendicular projections
of $P$ onto the three side lines are collinear if and only if
$P\in(ABC)$. Coincident feet are allowed; collinear means contained
in a common line.
\end{theorem}
\begin{proof}
Set $B=(0,0)$, $C=(1,0)$, $A=(u,v)$ with $v\ne0$, and
$P=(s,t)$. Put
$L=u^2+v^2$, $M=(u-1)^2+v^2$,
$U=su+tv$, $V=(s-1)(u-1)+tv$.
The projections onto $BC,AB,AC$ are
\[
 X=(s,0),\qquad
 Z=(uU/L,vU/L),\qquad
 Y=(1+(u-1)V/M,vV/M).
\]
Expanding the signed area gives
\[
 (Y-X)\times(Z-X)
 =-\frac{v^3}{LM}
   \left(s^2+t^2-s+\frac{u-u^2-v^2}{v}t\right).
\]
All denominators are nonzero. The expression in parentheses is
the circumcircle equation evaluated at $P$. It vanishes exactly
when the feet are collinear, including all projection degeneracies.
\end{proof}

\begin{figure}[htbp]\centering
\begin{tikzpicture}[scale=.84]
\coordinate (A) at (100:2); \coordinate (B) at (210:2);
\coordinate (C) at (330:2); \coordinate (P) at (35:2);
\coordinate (X) at ($(B)!(P)!(C)$);
\coordinate (Y) at ($(C)!(P)!(A)$);
\coordinate (Z) at ($(A)!(P)!(B)$);
\draw (0,0) circle(2); \draw (A)--(B)--(C)--cycle;
\draw[aux] (P)--(X) (P)--(Y) (P)--(Z);
\draw[main] ($(X)!-.15!(Z)$)--($(X)!1.15!(Z)$);
\foreach \X/\pos in {A/above,B/left,C/below right,P/right,X/below,Y/above left,Z/above left}
 \fill (\X) circle(1.2pt) node[\pos]{$\X$};
\end{tikzpicture}
\caption{The three pedal feet of a point on the circumcircle lie on
its Simson line. The feet may lie on side extensions.}
\end{figure}

\strategy The determinant calculation is a complete proof, but it
should be read structurally: every foot is linear in $s,t$, so the
collinearity equation is quadratic. A circle equation is exactly the
kind of expression one should hope to obtain.

\section{Complex coordinates: turning becomes multiplication}

Identify $(x,y)$ with $z=x+iy$, where $i^2=-1$.
The conjugate is $\bar z=x-iy$, and
$z\bar z=|z|^2=x^2+y^2$.
Multiplication by a unit complex number rotates the plane.
For nonzero vectors $u,v$, $u/v$ is real precisely when they are
parallel; it is purely imaginary precisely when they are perpendicular.

On the unit circle $|z|=1$, one has $\bar z=1/z$.
For distinct unit points $a,b$, their line has equation
\[
 z+ab\bar z=a+b.
\]
Indeed it holds at $a,b$ and is a real linear equation in $x,y$.
The perpendicular projection of a point $p$ to that line is
\[
 f=\frac{a+b+p-ab\bar p}{2}.
\]
Its reflected partner is $a+b-ab\bar p$; the midpoint with $p$
is the foot. Substituting into the line equation verifies incidence,
and the quotient of its displacement from $p$ by $b-a$ is purely
imaginary, verifying perpendicularity.

\begin{example}[The orthocenter becomes a sum]
For unit complex vertices $a,b,c$, prove $h=a+b+c$.
\end{example}
\begin{solution}
The direction $h-a=b+c$ is perpendicular to $b-c$, since
\[
 \frac{b+c}{b-c}
 +\overline{\left(\frac{b+c}{b-c}\right)}=0
\]
after replacing $\bar b,\bar c$ by $1/b,1/c$.
Thus $h$ lies on the altitude from $a$, and similarly on the other two.
This is the complex version of the earlier vector identity.
\end{solution}

\section{Barycentric coordinates: preserve the triangle}

Relative to noncollinear $A,B,C$, a finite point $P$ has
\emph{barycentric coordinates} $(x:y:z)$ if
\[
 P=\frac{xA+yB+zC}{x+y+z},\qquad x+y+z\ne0.
\]
Multiplying all three numbers by the same nonzero scalar changes
nothing. Normalize to $x+y+z=1$ when convenient.
The side $BC$ has $x=0$; on that side $BD/DC=z/y$ with the
signed convention. Interior points have all coordinates positive.
Their normalized coordinates are the area ratios
$[PBC]:[PCA]:[PAB]$.

\begin{figure}[htbp]\centering
\begin{tikzpicture}[scale=.8]
\coordinate (A) at (0,3); \coordinate (B) at (-2,0);
\coordinate (C) at (3,0); \coordinate (P) at (.4,1);
\fill[accent!10] (P)--(B)--(C)--cycle;
\fill[accent!22] (P)--(C)--(A)--cycle;
\fill[accent!34] (P)--(A)--(B)--cycle;
\draw (A)--(B)--(C)--cycle;
\draw (A)--(P)--(B) (P)--(C);
\foreach \X/\pos in {A/above,B/left,C/right,P/above right}
 \fill (\X) circle(1.2pt) node[\pos]{$\X$};
\node at (.5,.35){$x$}; \node at (1.1,1.2){$y$};
\node at (-.8,1.1){$z$};
\end{tikzpicture}
\caption{Normalized barycentric coordinates are opposite area ratios.
They make side ratios and concurrence linear.}
\end{figure}

\begin{example}[The incenter without intersecting two bisectors]
Show that the incenter has barycentric coordinates $(a:b:c)$.
\end{example}
\begin{solution}
Its distances from all three sides equal $r$, so the areas
$[IBC],[ICA],[IAB]$ are respectively $ar/2,br/2,cr/2$.
Their ratio is $a:b:c$.
The centroid is $(1:1:1)$, and the symmedian point is
$(a^2:b^2:c^2)$ by its side-division ratios.
\end{solution}

A line has a homogeneous equation $ux+vy+wz=0$. Two distinct
points determine such a line; substitute their coordinate triples
and solve the two linear equations. This is often cleaner than
memorizing a determinant formula. Barycentrics excel at concurrency,
collinearity, and ratios. Perpendicularity and distances need a metric
formula involving the side lengths, so Cartesian or complex coordinates
may then be simpler.

\pitfall Homogeneous coordinates with $x+y+z=0$ do not represent a
finite affine point. If a calculation produces this sum, determine
whether the geometric statement is about parallel lines before
trying to normalize.

\section{Practice}
\begin{exercise}[\problemlevel{Foundation}]\label{g:ex49}
With the circumcenter as origin, prove $h=a+b+c$ by dot products
and use it to obtain the nine-point center.
\end{exercise}
\begin{exercise}[\problemlevel{Bridge}]\label{g:ex50}
Derive the unit-circle line equation and projection formula above.
Verify both incidence and perpendicularity.
\end{exercise}
\begin{exercise}[\problemlevel{Bridge}]\label{g:ex51}
Let $P=(x:y:z)$ be interior. Prove that the line $AP$ meets $BC$
in the ratio $z:y$ and that $PD/AD=x/(x+y+z)$.
\end{exercise}
\begin{exercise}[\problemlevel{Challenge}]\label{g:ex52}
Prove the orthologic-triangles theorem: if the lines through
$A_1,B_1,C_1$ perpendicular to $B_2C_2,C_2A_2,A_2B_2$
are concurrent, then the lines through $A_2,B_2,C_2$
perpendicular to $B_1C_1,C_1A_1,A_1B_1$ are concurrent.
Both triangles are nondegenerate.
\end{exercise}
\begin{exercise}[\problemlevel{Challenge}]\label{g:ex53}
If two points $P,Q$ are antipodal on $(ABC)$, prove that their
Simson lines are perpendicular. Include cases when some feet coincide.
\end{exercise}
\begin{exercise}[\problemlevel{Challenge}]\label{g:ex54}
Let $P$ lie on $(ABC)$. Reflect $P$ in the three side lines.
Prove that its three reflected points are collinear, and prove the
converse for an arbitrary point $P$.
\end{exercise}

\chapter{Inversion and Tangency}
\label{g:inversion}

\section{Reciprocate distance, preserve a ray}

Fix a center $O$ and a radius $k>0$. Inversion sends every point
$X\ne O$ to the point $X'$ on ray $OX$ satisfying
$OX\cdot OX'=k^2$. The center itself has no finite image.
Inverting twice returns the original point.
With $O$ at the origin the vector formula is
\[
 x'=\frac{k^2x}{|x|^2}.
\]
This is not a similarity: scale depends on position.

\intuition A circle through the inversion center becomes a line.
Thus a dense collection of circles through one point can become
a sparse collection of lines. The best inversion center is usually
dictated by incidences, not by a desire to calculate reciprocals.

\begin{figure}[htbp]\centering
\begin{tikzpicture}[scale=1.35]
\draw[main] (0,0) circle(1);
\draw[->] (-.3,0)--(2.7,0);
\coordinate (X) at (2.3,0); \coordinate (Xp) at (.4348,0);
\coordinate (O) at (0,0);
\foreach \X/\lab/\pos in {O/O/below left,X/X/below,Xp/X'/below}
 \fill (\X) circle(1pt) node[\pos]{$\lab$};
\node at (-.2,1.2){$k=1$};
\end{tikzpicture}
\caption{Inside and outside exchange places. Points on the inversion
circle stay fixed, and $OX\cdot OX'=1$ in this unit example.}
\end{figure}

\begin{theorem}[Images of lines and circles]
Under inversion, a line not through $O$ becomes a circle through $O$,
with $O$ deleted. A circle through $O$ becomes a line not through $O$.
A circle not through $O$ becomes another circle. A line through $O$
maps to itself, with $O$ deleted.
\end{theorem}
\begin{proof}
Scale to $k=1$. A generalized circle has equation
\[
 A(x^2+y^2)+Bx+Cy+D=0,
\]
with $(A,B,C)$ not all zero. Substitute
$x=x'/(x'^2+y'^2)$, $y=y'/(x'^2+y'^2)$, valid away from $O$.
Multiplication by $x'^2+y'^2$ gives
\[
 D(x'^2+y'^2)+Bx'+Cy'+A=0.
\]
The constant and quadratic coefficients have exchanged roles.
The stated cases follow. An ordinary circle not through $O$ has
$A,D\ne0$, so its image is again an ordinary circle.
\end{proof}

\begin{figure}[htbp]\centering
\begin{tikzpicture}[scale=1.15]
\coordinate (O) at (0,0); \coordinate (A) at (2,0);
\coordinate (B) at (1,1); \coordinate (Ap) at (.5,0);
\coordinate (Bp) at (.5,.5);
\draw (1,0) circle(1);
\draw[main] (.5,-1.3)--(.5,1.5);
\draw[aux] (O)--(B);
\foreach \X/\lab/\pos in {O/O/below left,A/A/right,B/B/above right,Ap/A'/below right,Bp/B'/above left}
 \fill (\X) circle(1.1pt) node[\pos]{$\lab$};
\end{tikzpicture}
\caption{The circle $(x-1)^2+y^2=1$ passes through the unit inversion
center and becomes the line $x'=1/2$.}
\end{figure}

\section{The distance formula and angle behavior}

\begin{theorem}[Inversion distance formula]
For distinct noncentral points $X,Y$,
\[
 X'Y'=\frac{k^2\,XY}{OX\cdot OY}.
\]
\end{theorem}
\begin{proof}
Square the vector difference:
\[
 \left|\frac{k^2x}{|x|^2}-\frac{k^2y}{|y|^2}\right|^2
 =\frac{k^4(|x|^2+|y|^2-2x\cdot y)}{|x|^2|y|^2}.
\]
Take positive square roots.
\end{proof}

Inversion preserves the magnitude of an angle between smooth curves
at a noncentral intersection, while reversing orientation. One
calculus verification differentiates its vector formula:
\[
 Df_x(v)=\frac{k^2}{|x|^2}
 \left(v-2\frac{x\cdot v}{|x|^2}x\right).
\]
The bracket is reflection in the line perpendicular to $x$, so
dot-product angles are preserved and orientation reverses.
For readers avoiding calculus, the same conclusion follows from
the distance formula applied to small triangles and then taking
the limiting tangent directions. Tangencies away from $O$ remain
tangencies. At $O$ there is no finite image point, so angle claims
there must be replaced by statements about the images' directions.

\begin{example}[Ptolemy's inequality becomes a triangle inequality]
Prove for four distinct points
$AC\cdot BD\le AB\cdot CD+AD\cdot BC$.
\end{example}
\begin{solution}
Invert at $A$. The triangle inequality says
$B'D'\le B'C'+C'D'$. Use the distance formula and multiply by
the positive product $AB\cdot AC\cdot AD/k^2$ to obtain the claim.
If the four points are not collinear, equality means $B',C',D'$
are collinear with $C'$ between $B',D'$. Inverting the line back
gives a circle through $A$ in the appropriate cyclic order.
\end{solution}

\begin{example}[A circle fixed by inversion]
Suppose a circle $\omega$ of radius $\rho$ has center $C$ with
$OC^2=k^2+\rho^2$. Show that inversion of radius $k$ about $O$
maps $\omega$ to itself.
\end{example}
\begin{solution}
Its equation is $|x|^2-2c\cdot x+k^2=0$.
Substituting $x=k^2x'/|x'|^2$ returns the same equation after scaling.
The condition is also exactly orthogonality to the inversion circle,
by Problem~\ref{g:ex34}.
\end{solution}

\section{A tangency configuration developed from its metric core}

The $A$-mixtilinear incircle is the nonzero circle inside angle $A$
tangent to lines $AB,AC$ and internally tangent to $(ABC)$.
It need not be contained in triangle $ABC$.
Let its center be $J$, its radius $\rho$, and its contacts with
$AB,AC$ be $F,E$. The ordinary incircle has center $I$ and radius $r$.

\begin{theorem}[Contact-chord midpoint lemma]
The midpoint of $EF$ is $I$.
\end{theorem}
\begin{proof}
Put $A=(0,0)$ and the internal angle bisector on the positive
$x$-axis. Write $\theta=A/2$, $J=(t,0)$, $I=(d,0)$,
and let the circumcenter be $(u,v)$ with radius $R$.
Then $\rho=t\sin\theta$, $r=d\sin\theta$, and $u^2+v^2=R^2$.
Internal tangency gives
\[
 (t-u)^2+v^2=(R-t\sin\theta)^2,
 \quad\text{hence}\quad
 t\cos^2\theta=2(u-R\sin\theta).
\]
Euler's formula $OI^2=R^2-2Rr$ similarly gives
$d=2(u-R\sin\theta)$, after canceling $d>0$.
The projection of $(t,0)$ onto either side ray has $x$-coordinate
$t\cos^2\theta=d$, and the two $y$-coordinates are opposites.
Their midpoint is therefore $(d,0)=I$.
\end{proof}

\begin{figure}[htbp]\centering
\begin{tikzpicture}[scale=1.65]
\coordinate (A) at (0,0); \coordinate (B) at (1.5,-.866025);
\coordinate (C) at (1.5,.866025); \coordinate (I) at (1,0);
\coordinate (J) at (1.333333,0);
\coordinate (E) at (1,.57735); \coordinate (F) at (1,-.57735);
\draw (1,0) circle(1);
\draw (A)--(B)--(C)--cycle;
\draw[main] (J) circle(.666667);
\draw[main] (E)--(F);
\draw[aux] (A)--(2,0);
\draw[aux] (J)--(E) (J)--(F);
\foreach \X/\pos in {A/left,B/below,C/above,I/below left,J/above right,E/above left,F/below left}
 \fill (\X) circle(.8pt) node[\pos]{$\X$};
\end{tikzpicture}
\caption{An equilateral example makes the midpoint lemma visible.
The mixtilinear circle extends beyond side $BC$, while remaining inside
the circumcircle and angle $A$. The general proof does not assume symmetry.}
\end{figure}

\begin{example}[Choosing an inversion radius to swap two points]
If $AB\cdot AC=q>0$ and $B,C$ lie on the same ray from $A$,
what inversion exchanges them?
\end{example}
\begin{solution}
Use center $A$ and radius $\sqrt q$. Then the defining equation
$AX\cdot AX'=q$ sends $B$ to $C$ and conversely.
This is a practical choice, not a numerical trick: swapping known
points can turn a desired circle into an already known circle.
\end{solution}

\recognition Invert at a common point of circles to create lines;
at a common tangency point to create parallel lines; or at a vertex
when products such as $AB\cdot AC$ are fixed. Write down the images
of the important objects before adding further calculations.

\section{Practice}
\begin{exercise}[\problemlevel{Foundation}]\label{g:ex55}
Under unit inversion about the origin, find the images of
$x=2$ and $(x-2)^2+y^2=1$, stating all excluded points.
\end{exercise}
\begin{exercise}[\problemlevel{Bridge}]\label{g:ex56}
Prove that a circle orthogonal to the inversion circle is fixed as
a set. Conversely, show that a fixed circle whose center differs
from the inversion center is orthogonal to it. Identify the additional
fixed circle when the centers coincide.
\end{exercise}
\begin{exercise}[\problemlevel{Challenge}]\label{g:ex57}
Give a straightedge-and-compass construction of the $A$-mixtilinear
incircle using the ordinary incenter and the contact-chord midpoint
lemma. Prove existence as well as correctness.
\end{exercise}
\begin{exercise}[\problemlevel{Challenge}]\label{g:ex58}
For distinct $A,B,C,D$ with no three collinear, prove that equality
in Ptolemy's inequality is equivalent to their lying on a circle
in the order $A,B,C,D$ or its reverse.
\end{exercise}
\begin{exercise}[\problemlevel{Challenge}]\label{g:ex59}
Two circles are internally tangent at $T$.
Show that inversion at $T$ sends them to distinct parallel lines,
and compute the distance between the lines in terms of their
original radii $R>r$ and the inversion radius $k$.
\end{exercise}
\begin{exercise}[\problemlevel{Challenge}]\label{g:ex60}
Four circles with radii $0<r_i<R$ are internally tangent to an
outer circle of radius $R$ at distinct points $P_1,P_2,P_3,P_4$
in cyclic order. Let $t_{ij}$ be the distance between the contact
points on a common external tangent of circles $i,j$.
Prove the precisely scoped Casey identity
\[
 t_{12}t_{34}+t_{23}t_{14}=t_{13}t_{24}.
\]
Do not assume or assert an unrestricted converse.
\end{exercise}

\chapter{Projective Tools and Conics}
\label{g:projective}

\section{Incidence survives a change of viewpoint}

Parallel railway tracks appear to meet in a perspective picture.
Projective geometry formalizes that useful change of viewpoint:
parallel lines meet at an ideal point, and every pair of distinct
lines meets. Lengths and angles usually do not survive, but
incidence and cross ratio do.

In homogeneous plane coordinates, a point is a nonzero triple
$(x:y:z)$, with scalar multiples identified. Finite Cartesian points
have $z\ne0$ and coordinates $(x/z,y/z)$. Points with $z=0$ are
ideal points. A line has an equation $ux+vy+wz=0$.
An invertible $3$ by $3$ linear transformation of the triples induces
a \emph{projective transformation}; it preserves incidences because
substitution transforms linear equations into linear equations.

\begin{figure}[htbp]\centering
\begin{tikzpicture}[scale=.82]
\coordinate (S) at (0,3.5);
\coordinate (A) at (-2,0); \coordinate (B) at (-.8,0);
\coordinate (C) at (.6,0); \coordinate (D) at (2,0);
\draw (-2.5,0)--(2.5,0); \draw (-2,1.2)--(2,1.2);
\foreach \X in {A,B,C,D}{
 \draw[aux] (S)--(\X);
 \coordinate (\X p) at ($(S)!.6571429!(\X)$);
 \fill (\X) circle(1.2pt) node[below]{$\X$};
 \fill (\X p) circle(1.2pt) node[above left]{$\X'$};
}
\fill (S) circle(1.2pt) node[above]{$S$};
\end{tikzpicture}
\caption{Central projection generally changes lengths and ordinary
ratios, but preserves the cross ratio of four collinear points.}
\end{figure}

\section{Cross ratio and harmonic division}

For four distinct collinear finite points, choose a signed coordinate
and define
\[
 (A,B;C,D)
 =\frac{c-a}{c-b}\frac{d-b}{d-a}.
\]
Changing the coordinate by a nonzero affine scale leaves this number
unchanged. More generally, a projectivity of a line has the form
$t\mapsto(\alpha t+\beta)/(\gamma t+\delta)$, with
$\alpha\delta-\beta\gamma\ne0$, and preserves this expression:
subtracting two transformed values produces a common numerator
factor and two denominators, all of which cancel in the cross ratio.
Values at infinity are obtained by cancellation or limits.

A range is \emph{harmonic} when its cross ratio is $-1$.
If $A=-1$, $B=1$, then $(A,B;C,D)=-1$ is equivalent to $cd=1$.
Thus internal and external division are paired by reciprocation.

\begin{example}[Three pairs determine a projectivity]
Why are three distinct point images enough to determine a projectivity
of a line?
\end{example}
\begin{solution}
First send the three source points to $0,1,\infty$, and similarly
normalize the target points. A fractional linear map fixing
$0,\infty$ has form $t\mapsto kt$; fixing $1$ forces $k=1$.
Undo the two normalizations. This argument concerns a projectivity
of one line. It does not say that an arbitrary degree-$d$ moving-point
construction is determined by $d+1$ sample positions.
\end{solution}

\section{Poles and polars}

For the circle $x^2+y^2=R^2$, define the \emph{polar} of a noncentral
point $P=(p,q)$ by
\[
 px+qy=R^2.
\]
If $P$ is external, its two tangent contact points satisfy this
equation because $OP\cdot OT=OT^2=R^2$ at contact $T$.
Thus the polar is the chord of contact. For an interior noncentral
point, the same equation still defines a real line, even though
there are no real tangents from $P$.

\begin{theorem}[La Hire reciprocity]
$P$ lies on the polar of $Q$ if and only if $Q$ lies on the polar
of $P$.
\end{theorem}
\begin{proof}
Both statements are the symmetric equation $p\cdot q=R^2$.
\end{proof}

\begin{figure}[htbp]\centering
\begin{tikzpicture}[scale=1.05]
\coordinate (O) at (0,0); \coordinate (P) at (3,0);
\coordinate (T) at (.75,1.299038); \coordinate (U) at (.75,-1.299038);
\coordinate (Q) at (.75,0); \coordinate (A) at (-1.5,0);
\coordinate (B) at (1.5,0);
\draw (O) circle(1.5); \draw[main] (P)--(T) (P)--(U);
\draw[main] (.75,-1.8)--(.75,1.8);
\draw[aux] (-1.8,0)--(3.3,0);
\foreach \X/\pos in {O/below left,P/right,T/above,U/below,Q/above right,A/left,B/below right}
 \fill (\X) circle(1.1pt) node[\pos]{$\X$};
\end{tikzpicture}
\caption{The polar of an external point is its chord of contact.
On diameter $AB$, $OP\cdot OQ=R^2$, so $(A,B;Q,P)=-1$.}
\end{figure}

\begin{example}[Concurrency becomes collinearity]
Let three noncentral points $P,Q,S$ have concurrent polars at a
finite point $T\ne O$. Prove $P,Q,S$ are collinear.
\end{example}
\begin{solution}
By reciprocity each of $P,Q,S$ lies on the polar of $T$, so they
are collinear. The converse follows the same way when the pole of
the line is finite. A line through the center has an ideal pole,
and its point-polars are parallel; this is the correct finite-plane
exception.
\end{solution}

\section{Conics and a usable proof of Pascal}

An ellipse is the locus with sum of distances from two foci equal to
a fixed number $2a$ larger than the focal separation. Taking the
foci as $(\pm c,0)$ and squaring twice gives
\[
 \frac{x^2}{a^2}+\frac{y^2}{b^2}=1,\qquad b^2=a^2-c^2>0.
\]
The sign condition in the distance definition must be retained when
reversing the squaring. A hyperbola uses a fixed nonzero absolute
difference of focal distances. A parabola consists of points
equidistant from a focus and a directrix line.

\begin{figure}[htbp]\centering
\begin{tikzpicture}[scale=.78]
\coordinate (F1) at (-2.236068,0); \coordinate (F2) at (2.236068,0);
\coordinate (P) at (1.5,1.732051);
\draw[main] (0,0) ellipse(3 and 2);
\draw (F1)--(P)--(F2);
\draw[aux] (-3.4,0)--(3.4,0);
\foreach \X/\lab/\pos in {F1/F_1/below,F2/F_2/below,P/P/above right}
 \fill (\X) circle(1.2pt) node[\pos]{$\lab$};
\node at (0,-2.4){$PF_1+PF_2=2a$};
\end{tikzpicture}
\caption{The focal definition supplies geometric meaning; the quadratic
equation supplies an efficient calculation tool.}
\end{figure}

Projectively, every nondegenerate real conic with real points can
be represented by
$P(t)=(t^2:t:1)$, together with $P(\infty)=(1:0:0)$.
To see the normalization, choose a point of the conic as the ideal
point $(1:0:0)$ and its tangent as $z=0$. A quadratic equation then
has no $x^2$ or $xy$ terms; invertible coordinate changes complete
the remaining square and scale it to $xz=y^2$.

\begin{theorem}[Pascal]
Let $A,B,C,D,E,F$ be six distinct points on a nondegenerate conic.
If the intersections are well-defined, then
\[
 X=AB\cap DE,\quad Y=BC\cap EF,\quad Z=CD\cap FA
\]
are collinear, with ideal intersections allowed.
\end{theorem}
\begin{proof}
A fractional linear change of parameter can send the parameters of
$A,B,C$ to $\infty,0,1$. This is induced by a plane projectivity:
substituting $t=(\alpha u+\beta)/(\gamma u+\delta)$ makes the
three homogeneous coordinates quadratic polynomials in $u$,
an invertible linear change of $(u^2:u:1)$.
Thus normalize $A=(1:0:0)$, $B=(0:0:1)$, $C=(1:1:1)$.
Write the remaining parameters as distinct finite values $d,e,f$.

The line joining finite parameters $s,t$ has coefficients
$(1,-s-t,st)$. Directly solving the line pairs now gives
\[
\begin{split}
X&=(-de:0:1),\\
Y&=(ef:ef:e+f-1),\\
Z&=((1+d)f-d:f:1).
\end{split}
\]
Their collinearity determinant expands along the first row to
\[
\begin{vmatrix}
-de&0&1\\
ef&ef&e+f-1\\
(1+d)f-d&f&1
\end{vmatrix}
=def(f-1)+def(1-f)=0.
\]
The computation uses homogeneous coordinates, so zero last
coordinates already include ideal intersections. Undoing the
projectivity proves the original assertion.
\end{proof}

\begin{figure}[htbp]\centering
\begin{tikzpicture}[scale=.68]
\coordinate (A) at (26:2); \coordinate (B) at (56:2);
\coordinate (C) at (186:2); \coordinate (D) at (289:2);
\coordinate (E) at (324:2); \coordinate (F) at (354:2);
\coordinate (X) at (2.812993,-.291347);
\coordinate (Y) at (3.181660,2.897816);
\coordinate (Z) at (2.492463,-3.064092);
\draw (0,0) circle(2);
\draw (A)--(B)--(C)--(D)--(E)--(F)--cycle;
\draw[aux] (A)--(X) (D)--(X) (B)--(Y) (E)--(Y) (C)--(Z) (F)--(Z);
\draw[main] (Z)--(Y);
\foreach \X/\pos in {A/right,B/above,C/left,D/below,E/below right,F/above right,X/right,Y/right,Z/right}
 \fill (\X) circle(1.5pt) node[\pos]{$\X$};
\end{tikzpicture}
\caption{Pascal turns six conic incidences into a line.
The hexagon need not be regular, and intersections may be exterior.}
\end{figure}

\begin{theorem}[Desargues]
Suppose triangles $ABC,A'B'C'$ are perspective from a point $S$:
$AA',BB',CC'$ concur at $S$. When the corresponding side
intersections are defined, $AB\cap A'B'$, $BC\cap B'C'$,
$CA\cap C'A'$ are collinear.
\end{theorem}
\begin{proof}
If $S$ is a vertex, say $S=A$, then $B'\in AB$ and $C'\in AC$.
The intersections of the first and third corresponding side pairs
are $B',C'$, while the remaining intersection lies on $B'C'$.
This proves the assertion whenever the stated lines and intersections
are defined. The case where $S$ is a primed vertex follows by swapping
the triangles. Now suppose $S$ is none of their vertices.
Use nonzero homogeneous vector representatives and scale them so
$A'=A+\alpha S$, $B'=B+\beta S$, $C'=C+\gamma S$.
This is possible when the paired vertices are distinct from $S$.
Their corresponding side intersections have representatives
$X=\beta A-\alpha B$,
$Y=\gamma B-\beta C$,
$Z=\alpha C-\gamma A$.
The relation $\gamma X+\alpha Y+\beta Z=0$ proves collinearity.
Cases with a zero coefficient are included when the stated
intersections remain well-defined.
\end{proof}

\begin{example}[What a projective theorem can and cannot do]
Six points on a circle produce the three intersections in Pascal.
Suppose one needs to prove those intersections collinear.
\end{example}
\begin{solution}
The task is incidence-only, so Pascal is a natural candidate.
List the desired intersection pairs, then order the six circle
points so that those pairs are opposite sides of the hexagon.
If the desired conclusion instead contains a midpoint or a right
angle, projective invariance alone will not establish it; those
metric properties are not preserved.
\end{solution}

\section{Practice}
\begin{exercise}[\problemlevel{Foundation}]\label{g:ex61}
Prove directly that a fractional linear transformation with nonzero
determinant preserves cross ratio, including a point sent to infinity.
\end{exercise}
\begin{exercise}[\problemlevel{Bridge}]\label{g:ex62}
For the circle $x^2+y^2=R^2$, find the polar of $(p,q)$ and the
pole of line $ux+vy=w$, with $w\ne0$. Explain the cases $w=0$
and $(p,q)=(0,0)$.
\end{exercise}
\begin{exercise}[\problemlevel{Challenge}]\label{g:ex63}
In a triangle $ABC$, let $D\in BC$ be nonvertex and let $E\in AD$.
Put $F=BE\cap AC$, $G=CE\cap AB$, and $X=FG\cap BC$,
assuming these are distinct finite points. Prove
$(B,C;D,X)=-1$.
\end{exercise}
\begin{exercise}[\problemlevel{Challenge}]\label{g:ex64}
Use Pascal to prove that the three intersections of opposite sides
of any six distinct points on a circle are collinear. Then explain
why a parallel pair of opposite sides makes that line a direction
statement rather than invalidating the theorem.
\end{exercise}
\begin{exercise}[\problemlevel{Challenge}]\label{g:ex65}
Six distinct tangents to a circle form a hexagon with distinct consecutive
vertices. Prove that its three main diagonals concur, allowing a
common ideal point. This is Brianchon's theorem in the circle case.
\end{exercise}
\begin{exercise}[\problemlevel{Challenge}]\label{g:ex66}
For the ellipse $x^2/a^2+y^2/b^2=1$, prove that the tangent at
$(x_0,y_0)$ is $xx_0/a^2+yy_0/b^2=1$.
Show that its two axis intercepts are reciprocally related to the
contact coordinates. Handle zero contact coordinates separately.
\end{exercise}

\chapter{Geometric Optimization and Synthesis}
\label{g:synthesis}

\section{An inequality needs an equality picture}

Before proving a geometric inequality, ask when equality should
hold. If the expected case is equilateral, try a sum of squares
measuring unequal sides. If it is collinearity, try a triangle
inequality after a transformation. If it involves perpendicularity,
try a projection or a squared dot product.

\begin{example}[Weitzenbock's inequality]
For every triangle,
\[
 a^2+b^2+c^2\ge4\sqrt3\,\Delta,
\]
with equality exactly for an equilateral triangle.
\end{example}
\begin{solution}
Set $B=(-x,0)$, $C=(x,0)$ with $x>0$, and $A=(u,v)$
with $v>0$. Then
$a^2+b^2+c^2=6x^2+2u^2+2v^2$ and $\Delta=xv$.
Their difference is
\[
 2u^2+2(v-\sqrt3x)^2\ge0.
\]
Equality means $u=0$ and $v=\sqrt3x$, the equilateral configuration.
The choice of a side midpoint as origin made the equality case a
simple square.
\end{solution}

\begin{figure}[htbp]\centering
\begin{tikzpicture}[scale=.78]
\coordinate (B) at (-2,0); \coordinate (C) at (2,0);
\coordinate (A) at (.7,2.7); \coordinate (E) at (0,3.4641);
\draw (A)--(B)--(C)--cycle;
\draw[aux] (B)--(E)--(C) (0,0)--(E);
\draw[<->] (0,2.7)--(.7,2.7);
\node at (.35,2.95){$u$};
\foreach \X/\pos in {A/right,B/below left,C/below right,E/above}
 \fill (\X) circle(1.2pt) node[\pos]{$\X$};
\end{tikzpicture}
\caption{The proof measures horizontal displacement from the equilateral
vertex and vertical displacement from its required height.}
\end{figure}

\begin{example}[Viviani's identity]
In an equilateral triangle of altitude $h$, the sum of the
perpendicular distances from an interior point $P$ to its three
sides is $h$.
\end{example}
\begin{solution}
If the common side length is $a$ and the distances are $x,y,z$,
sum the areas of $PBC,PCA,PAB$:
$\Delta=\tfrac12a(x+y+z)=\tfrac12ah$.
Cancel $a/2$. The identity explains why a naive attempt to minimize
this sum cannot select a unique interior point: it is constant.
\end{solution}

\section{The Fermat point and the meaning of stationarity}

Suppose $F$ lies inside $ABC$ and
$\angle AFB=\angle BFC=\angle CFA=120^\circ$.
Then $FA+FB+FC$ is the smallest possible sum of distances to the
three vertices. Here is a global proof, not just a local observation.

Put $u_A=(F-A)/FA$ and similarly $u_B,u_C$.
Three unit vectors separated by $120^\circ$ have sum zero.
For every $P$ the dot-product inequality gives
$|P-A|\ge(P-A)\cdot u_A$; adding yields
\[
 PA+PB+PC\ge FA+FB+FC
 +(P-F)\cdot(u_A+u_B+u_C)=FA+FB+FC.
\]
Equality forces each corresponding direction to agree, which
determines $P=F$ for a nondegenerate triangle.

Such an interior point exists exactly when all three vertex angles
are less than $120^\circ$. One existence argument minimizes the
continuous distance sum on the closed triangle, a compact set.
A minimizer cannot be in an open side: a small inward displacement
decreases the distance to the opposite vertex to first order while
the two along-side distances have zero first-order change.
At a vertex of angle $A<120^\circ$, moving into its internal bisector
changes the sum to first order by $1-2\cos(A/2)<0$.
Thus the minimum is interior. For a nonzero vector $w$,
\[
 \left.\frac{d}{dt}|w+tv|\right|_{t=0}
 =\frac{w\cdot v}{|w|},
\]
which follows by multiplying the difference of square roots by its
conjugate expression. At an interior minimizer all directional
changes vanish, so $u_A+u_B+u_C=0$, equivalent to the three
$120^\circ$ angles. This optional use of elementary differentiation
is confined to the existence argument.

If, for example, $\angle A\ge120^\circ$, let $e_B,e_C$ be unit
vectors from $A$ toward $B,C$ and put $p=P-A$. Then
$|e_B+e_C|=2\cos(A/2)\le1$. The inequalities
$PB\ge AB-p\cdot e_B$ and
$PC\ge AC-p\cdot e_C$ imply
$PA+PB+PC\ge AB+AC$. Thus the minimum is at $A$.

\begin{example}[The minimum can be expressed from the sides]
When all angles are below $120^\circ$, show that the minimum $L$
satisfies
\[
 L^2=\frac{a^2+b^2+c^2+4\sqrt3\,\Delta}{2}.
\]
\end{example}
\begin{solution}
Write $x=FA$, $y=FB$, $z=FC$. The cosine law at $F$ gives
$a^2=y^2+z^2+yz$ and its cyclic versions.
Also $\Delta=\sqrt3(xy+yz+zx)/4$, by summing the three areas.
Substitute and simplify to obtain $(x+y+z)^2=L^2$.
\end{solution}

\section{Locus problems require both directions}

Fix distinct points $A,B$ and $k>0$. The locus $PA/PB=k$ is the
perpendicular bisector if $k=1$, and otherwise is an
\emph{Apollonius circle}. With $A=(0,0)$ and $B=(1,0)$, its equation is
\[
 x^2+y^2=k^2((x-1)^2+y^2).
\]
For $k\ne1$, completing squares gives center
$(k^2/(k^2-1),0)$ and radius $k/|k^2-1|$.
Because both distances are nonnegative and $k>0$, squaring has not
introduced an extraneous branch; checking $P=B$ excludes it automatically.

\begin{figure}[htbp]\centering
\begin{tikzpicture}[scale=1.4]
\coordinate (A) at (0,0); \coordinate (B) at (1,0);
\coordinate (O) at (1.333333,0);
\draw[aux] (-.3,0)--(2.3,0);
\draw[main] (O) circle(.666667);
\coordinate (P) at (1.333333,.666667);
\draw (A)--(P)--(B);
\foreach \X/\pos in {A/below,B/below left,O/below right,P/above}
 \fill (\X) circle(.9pt) node[\pos]{$\X$};
\node at (1.6,-.95){$PA=2PB$};
\end{tikzpicture}
\caption{A ratio locus becomes a circle. A correct locus proof shows
both that every admissible point lies on the circle and that every
point on the circle satisfies the original ratio.}
\end{figure}

\section{A method-selection conversation}

Suppose the goal is a line through three complicated points.
Before naming a theorem, inspect how the points were made.
Intersections on side lines suggest Menelaus. Equal-power intersections
suggest radical axes. Perpendicular projections from a common point
suggest Simson. Intersections of opposite chords suggest Pascal.
Midpoints suggest vectors or a homothety. These cues are hypotheses
to test, not automatic prescriptions.

For a hard synthetic problem, work backward from the conclusion:
which angle equality would prove the desired circle? Which second
circle would supply that equality? Which construction would put
the needed four points on it? Then work forward from the data to
see whether the proposed intermediate claim is accessible.
The proof should retain this logic even if its final order is forward.

\pitfall A successful advanced technique can be inefficient in the
wrong problem. Inversion may destroy a midpoint condition;
barycentrics may make a perpendicular condition expensive;
angle chasing may hide a one-line area identity. Compare what each
representation makes simple and what it makes difficult.

\section{Feuerbach's tangencies from weighted vectors}
\intuition A difficult circle-tangency theorem can become a distance
identity. Two circles are internally tangent when the center distance
is the difference of their radii, and externally tangent when it is
their sum. Weighted vectors will compute precisely those distances.

\begin{theorem}[Feuerbach]\label{g:feuerbach}
The nine-point circle of a nonequilateral triangle is internally
tangent to its incircle and externally tangent to its three excircles.
For an equilateral triangle it coincides with the incircle; its
external tangencies to the excircles still hold.
\end{theorem}
\begin{proof}
Take the circumcenter as vector origin. Write the vertices as
$\mathbf A,\mathbf B,\mathbf C$, each of length $R$.
The orthocenter has vector $\mathbf H=\mathbf A+\mathbf B+\mathbf C$:
for example $(\mathbf H-\mathbf A)\cdot(\mathbf B-\mathbf C)=0$.
Thus the nine-point center is $\mathbf N=\mathbf H/2$ and its
radius is $R/2$.

For any real weights $\alpha,\beta,\gamma$, expansion gives
\[
\begin{split}
|\alpha\mathbf A+\beta\mathbf B+\gamma\mathbf C|^2
={}&R^2(\alpha+\beta+\gamma)^2\\
&-\alpha\beta c^2-\beta\gamma a^2-\gamma\alpha b^2.
\end{split}
\tag{*}
\]
The incenter has barycentric weights $a:b:c$, because its three
subtriangle areas are $ar/2,br/2,cr/2$.
Applying (*) to $\mathbf I=(a\mathbf A+b\mathbf B+c\mathbf C)/(2s)$
gives
$OI^2=R^2-abc/(2s)=R^2-2Rr$.
We used $\Delta=rs=abc/(4R)$.
Consequently $R\ge2r$. Equality makes $O=I$; equal distances
from the circumcenter to all three chords then make the sides equal.

Put $x=s-a$, $y=s-b$, $z=s-c$, so $s=x+y+z$ and
$a=y+z,b=z+x,c=x+y$. The weights for $\mathbf I-\mathbf N$ are
$-x,-y,-z$, divided by $2s$. Hence (*), Heron's formula
$\Delta^2=sxyz$, and expansion give
\[
\begin{split}
IN^2&=\frac{R^2}{4}
-\frac{xyc^2+yza^2+zxb^2}{4s^2}\\
&=\frac{R^2}{4}
-\frac{sabc-4sxyz}{4s^2}
=\frac{R^2}{4}-Rr+r^2.
\end{split}
\]
For the middle identity, expand with $s=x+y+z$; both numerators
are $s^2(xy+yz+zx)-5sxyz$.
Thus $IN=R/2-r$, proving the internal tangency unless the two
radii are equal, the equilateral coincidence already described.

For the $A$-excenter the signed barycentric weights are $-a:b:c$.
Their sum is $2x$, so the weights of
$\mathbf I_A-\mathbf N$ are $-s,z,y$, divided by $2x$.
Using (*) again,
\[
I_AN^2=\frac{R^2}{4}
+\frac{szc^2+syb^2-yza^2}{4x^2}.
\]
Expansion gives
$szc^2+syb^2-yza^2=xabc+4sxyz$.
Since $r_A=\Delta/x$, this yields
$I_AN^2=R^2/4+Rr_A+r_A^2=(R/2+r_A)^2$.
The positive center distance equals the sum of the radii, proving
external tangency. Relabeling proves the other two cases.
\end{proof}
\strategy The central move was to translate tangency into a distance.
The same weighted norm identity proved Euler's radius inequality
and both kinds of Feuerbach tangency. The algebra was organized by
the contact-length variables $s-a,s-b,s-c$, rather than expanded
blindly in the side lengths.

\begin{example}[A right triangle is included]
For vertices $A=(0,3)$, $B=(0,0)$, $C=(4,0)$, locate the internal
Feuerbach contact.
\end{example}
\begin{solution}
Here $I=(1,1)$, $r=1$, $O=(2,3/2)$, and $R=5/2$.
The orthocenter is $B$, so $N=(1,3/4)$ and the nine-point radius
is $5/4$. The center distance is $1/4=5/4-1$.
The two circles touch at $(1,2)$, on the ray from $N$ through $I$.
\end{solution}
\begin{figure}[htbp]\centering
\begin{tikzpicture}[scale=1]
\coordinate (A) at (0,3);\coordinate (B) at (0,0);
\coordinate (C) at (4,0);\coordinate (I) at (1,1);
\coordinate (N) at (1,.75);\coordinate (T) at (1,2);
\draw (A)--(B)--(C)--cycle;
\draw[main] (I) circle (1);
\draw[dashed,accent] (N) circle (1.25);
\draw[aux] (N)--(T);
\foreach \X/\pos in {A/above,B/below left,C/below right}
{\fill (\X) circle (1.2pt) node[\pos]{$\X$};}
\fill (T) circle (1.2pt);
\draw[gray] (T)--(.45,2.3) node[left,text=black]{$T$};
\fill (I) circle (1.2pt);
\draw (I)--(1.6,1.25) node[right]{$I$};
\fill (N) circle (1.2pt);
\draw (N)--(1.6,.5) node[right]{$N$};
\end{tikzpicture}
\caption{A right triangle's incircle and nine-point circle touch at
$T$. The theorem does not assume an acute triangle.}
\end{figure}

\section{Synthesis problems}
These problems combine chapters. A theorem is not identified in the
statement. Draw at least two configurations before committing to a
general argument.

\begin{exercise}[\problemlevel{Bridge}]\label{g:ex67}
Prove $a^2+b^2+c^2\ge4\sqrt3\,\Delta$ by an algebraic route from
Heron's formula, and determine equality. Compare with the coordinate
proof in the text.
\end{exercise}
\begin{exercise}[\problemlevel{Challenge}]\label{g:ex68}
For $P\in(ABC)$, prove that the line through its reflections in the
three side lines contains the orthocenter $H$. Deduce that the
Simson line of $P$ contains the midpoint of $PH$.
\end{exercise}
\begin{exercise}[\problemlevel{Challenge}]\label{g:ex69}
Let $M$ be the midpoint of a chord $PQ$ of a circle.
Two other chords $AB,CD$ pass through $M$. Let
$X=AD\cap PQ$, $Y=BC\cap PQ$, assuming distinct endpoints and
finite distinct intersections. Prove $MX=MY$.
This is the butterfly theorem.
\end{exercise}
\begin{exercise}[\problemlevel{Challenge}]\label{g:ex70}
Determine the minimum of $PA+PB+PC$ for an arbitrary point $P$
in the plane, including triangles with an angle at least
$120^\circ$. Justify why a minimizer outside the triangle cannot
improve the minimum found inside.
\end{exercise}
\begin{exercise}[\problemlevel{Challenge}]\label{g:ex71}
Project $A,B,C$ perpendicularly onto an arbitrary line $\ell$,
obtaining $X,Y,Z$. Prove that the lines through $X,Y,Z$
perpendicular to $BC,CA,AB$, respectively, concur.
Show that if $\ell$ is translated perpendicularly by a vector $v$,
its concurrence point is translated by the same vector.
\end{exercise}
\begin{exercise}[\problemlevel{Challenge}]\label{g:ex72}
For an acute triangle $ABC$, let $D,E,F$ be its altitude feet on
$BC,CA,AB$. Among all triangles $XYZ$ with
$X\in BC$, $Y\in CA$, $Z\in AB$, prove
\[
 XY+YZ+ZX\ge DE+EF+FD.
\]
Find the equality case. This is Fagnano's problem.
\end{exercise}


\part{Number Theory}

% ---- chapters/number-theory ----
\chapter{Divisibility, the Euclidean Algorithm, and Structure}

\label{ch:nt-divisibility}
\intuition Number theory asks for exact arithmetic structure. An
approximate quotient cannot tell us whether one integer divides another.
The remainder can. The Euclidean algorithm repeatedly replaces a large
pair by a smaller pair without changing the information we care about:
their common divisors.

\section{Division with a remainder}
\begin{definition}
For integers $a,b$, we write $a\mid b$ if $b=ak$ for an integer $k$.
A positive common divisor of $a,b$ divides both; their greatest common
divisor is denoted $\gcd(a,b)$, provided they are not both zero.
\end{definition}
If $d$ divides $a,b$, it divides every integer linear combination
$ua+vb$. This simple observation powers much of the subject.
The converse must be used carefully: dividing one combination does
not generally mean dividing each summand.

\begin{theorem}[Division algorithm]\label{n:division}
For every integer $a$ and positive integer $b$, there are unique integers
$q,r$ with $a=bq+r$ and $0\le r<b$.
\end{theorem}
\begin{proof}
Among the nonnegative integers of the form $a-bq$, choose the least,
called $r$. Such integers exist by taking $q$ sufficiently negative.
If $r\ge b$, then $r-b$ is a smaller nonnegative member, a contradiction.
For uniqueness, two representations give $b(q-q')=r'-r$.
The right side has absolute value less than $b$, so the only multiple
of $b$ it can equal is zero.
\end{proof}
For example, $-17=5(-4)+3$. The remainder is $3$, not $-2$.
Remainders are representatives of residue classes, not a convention
that positive inputs alone can satisfy.

\section{Why Euclid preserves the gcd}
\begin{theorem}[Euclidean reduction]\label{n:euclid}
If $a=bq+r$, then $\gcd(a,b)=\gcd(b,r)$ whenever these gcds are defined.
\end{theorem}
\begin{proof}
A divisor of $a,b$ divides $r=a-bq$. A divisor of $b,r$ divides
$a=bq+r$. Thus the sets of common positive divisors coincide.
\end{proof}
Successive nonzero remainders strictly decrease, so repeated division
terminates. The last nonzero remainder is the gcd. This separates two
claims: preservation tells us the output is correct, and decrease
tells us there will be an output.

\begin{example}[Compute and remember the history]
Find $\gcd(252,198)$ and express it as $252u+198v$.
\end{example}
\begin{solution}
The divisions are
\[
252=198+54,\quad198=3\cdot54+36,\quad54=36+18,\quad36=2\cdot18.
\]
Thus the gcd is $18$. Reverse the substitutions:
$18=54-36=4\cdot54-198=4\cdot252-5\cdot198$.
\end{solution}

\begin{theorem}[Bezout's identity]\label{n:bezout}
For integers $a,b$ not both zero, their gcd is an integer linear
combination $au+bv$. Every common divisor divides that combination.
\end{theorem}
\begin{proof}
The Euclidean algorithm expresses every remainder as an integer
linear combination of $a,b$, starting with $a,b$ themselves.
Back substitution therefore expresses the final nonzero remainder.
Signs are handled by replacing initial inputs by absolute values and
then restoring their signs in the coefficients.
\end{proof}

\section{Coprimality allows cancellation}
\begin{lemma}[Euclid's cancellation lemma]\label{n:coprime-cancel}
If $\gcd(a,b)=1$ and $a\mid bc$, then $a\mid c$.
\end{lemma}
\begin{proof}
Choose integers $u,v$ with $au+bv=1$. Multiplying by $c$ gives
$c=auc+bvc$, and both terms on the right are divisible by $a$.
\end{proof}
\pitfall From $a\mid bc$ alone, concluding $a\mid b$ or $a\mid c$
is false: $6\mid2\cdot3$. It becomes true for prime $a$ because a
prime is coprime to every integer it does not divide.

\begin{example}[Remove the shared factor]
For positive integers $a,b$, prove
\[
\gcd(a,b)\lcm(a,b)=ab.
\]
\end{example}
\begin{solution}
Write $a=gx,b=gy$, where $g=\gcd(a,b)$ and $\gcd(x,y)=1$.
A common multiple $M$ is divisible by $gx$; write $M=gxt$.
Since $gy\mid M$, we have $y\mid xt$, hence $y\mid t$ by cancellation.
Thus every common multiple is a multiple of $gxy$.
Conversely $gxy$ is a common multiple, so $\lcm(a,b)=gxy$.
\end{solution}
\recognition Shared factors obscure a problem. Substituting
$a=gx,b=gy$ exposes coprime parts, where products can be separated.

\section{Linear Diophantine equations}
\begin{theorem}\label{n:linear-dio}
Let $a,b$ be nonzero integers and $g=\gcd(a,b)$. The equation
$ax+by=c$ has integer solutions exactly when $g\mid c$.
If $(x_0,y_0)$ is one solution, all are
\[
x=x_0+\frac bg\,t,\qquad y=y_0-\frac ag\,t,\qquad t\in\ZZ.
\]
\end{theorem}
\begin{proof}
Necessity follows because $g$ divides the left side. For sufficiency,
multiply a Bezout identity by $c/g$. Subtracting a particular solution
from any other gives $(a/g)(x-x_0)=-(b/g)(y-y_0)$.
The coefficients are coprime, so $b/g\mid x-x_0$. Substitution gives
the displayed form; directly checking it proves the converse.
\end{proof}

\begin{example}[Integer solutions and positive solutions differ]
Find all nonnegative integer solutions to $7x+11y=100$.
\end{example}
\begin{solution}
Modulo $7$, $4y\equiv2$, so $y\equiv4\pmod7$.
Also $0\le y\le9$. Thus $y=4$ is the only possibility and $x=8$.
The unrestricted equation has infinitely many integer solutions,
$x=8+11t,y=4-7t$, but the nonnegative interval permits only $t=0$.
\end{solution}

\section{A first structural bound: the coin problem}
\intuition With coprime denominations, every sufficiently large amount
is possible. The residue modulo one denomination chooses the number
of the other coins; size then makes the remaining coefficient positive.

\begin{theorem}[Two-denomination Frobenius theorem]\label{n:frobenius}
For coprime integers $a,b>1$, the largest integer not representable as
$ax+by$ with $x,y\in\NNz$ is $ab-a-b$.
\end{theorem}
\begin{proof}
If $ab-a-b=ax+by$, then $ab=a(x+1)+b(y+1)$.
Reduction modulo $a$ gives $a\mid y+1$, so $y+1\ge a$; the right
side is then at least $ab+a$, impossible.
Now let $N>ab-a-b$. The residues $0,b,\ldots,(a-1)b$ modulo $a$
are distinct by cancellation, so choose $0\le y\le a-1$ with
$by\equiv N\pmod a$. Then $N-by=ax$ for an integer $x$.
Since
$N-by>ab-a-b-(a-1)b=-a$,
the multiple $ax$ is greater than $-a$, hence $x\ge0$.
\end{proof}
The final inference uses integrality. A real number greater than
$-a$ need not be nonnegative; a multiple of $a$ in that interval does.

\section{Problems}
\begin{exercise}[\problemlevel{Foundation}]\label{n:ex01}
Compute $\gcd(2026,144)$ and find integers $u,v$ with $2026u+144v=\gcd(2026,144)$.
\end{exercise}
\begin{exercise}[\problemlevel{Foundation}]\label{n:ex02}
Prove $\gcd(n,n+1)=1$ for every integer $n$.
\end{exercise}
\begin{exercise}[\problemlevel{Bridge}]\label{n:ex03}
For every integer $n$, prove $\gcd(n^2+n+1,n+1)=1$.
\end{exercise}
\begin{exercise}[\problemlevel{Bridge}]\label{n:ex04}
Find all integer solutions of $18x+30y=42$.
\end{exercise}
\begin{exercise}[\problemlevel{Bridge}]\label{n:ex05}
If $m,n$ are coprime positive integers and each divides $A$, prove $mn\mid A$.
\end{exercise}
\begin{exercise}[\problemlevel{Bridge}]\label{n:ex06}
Find all positive pairs $a\le b$ with $\gcd(a,b)=6$ and $\lcm(a,b)=180$.
\end{exercise}
\begin{exercise}[\problemlevel{Challenge}]\label{n:ex07}
For coprime $a,b>1$, prove that exactly one of $N$ and $ab-a-b-N$ is representable as $ax+by$ with $x,y\ge0$, for every integer $0\le N\le ab-a-b$.
\end{exercise}
\begin{exercise}[\problemlevel{Challenge}]\label{n:ex08}
Let $a,b$ be coprime positive integers. Prove that $\gcd(a+b,a^2+b^2)$ is $2$ when $a,b$ are both odd and $1$ otherwise.
\end{exercise}
\chapter{Primes, Factorization, and Arithmetic Functions}

\label{ch:nt-primes}
\intuition A prime factorization is a coordinate system for
multiplication: products add exponents, divisibility compares them,
and squares require them to be even. Before using this coordinate
system, we must prove that its coordinates are unique.

\section{Primes and unique factorization}
A prime is an integer $p>1$ whose only positive divisors are $1,p$.
An integer $n>1$ that is not prime is composite.
If a prime $p$ divides $ab$ and does not divide $a$, then
$\gcd(p,a)=1$ and Lemma~\ref{n:coprime-cancel} gives $p\mid b$.
This is Euclid's lemma for primes.

\begin{theorem}[Fundamental theorem of arithmetic]\label{n:fta}
Every positive integer $n>1$ is a product of primes, uniquely up to
the order of the factors.
\end{theorem}
\begin{proof}
Existence follows by strong induction: a composite number splits
into two smaller factors. For uniqueness, suppose
$p_1\cdots p_r=q_1\cdots q_s$ are two prime factorizations.
Euclid's lemma makes $p_1$ divide some $q_j$, hence equal it.
Cancel the common prime and repeat. Eventually no primes remain
on either side, since a nonempty product of primes exceeds $1$.
\end{proof}
We include $1$ separately as the empty product. It is not prime;
otherwise uniqueness would be destroyed by inserting arbitrary factors
of $1$. Negative integers have an additional factor $-1$.

\begin{example}[A coprime product that is a square]
If positive coprime $a,b$ have square product $ab$, prove that
$a,b$ are both squares.
\end{example}
\begin{solution}
No prime occurs in both factors. Its exponent in $ab$ is therefore
its exponent in just one factor. Every exponent of a square is even,
so every exponent in each factor is even.
\end{solution}
\pitfall Without coprimality this fails: $2\cdot8=16$.

\section{There are infinitely many primes}
\begin{theorem}[Euclid]\label{n:infinite-primes}
There are infinitely many primes.
\end{theorem}
\begin{proof}
If $p_1,\ldots,p_k$ listed all primes, the integer
$N=p_1\cdots p_k+1>1$ would have a prime divisor $q$.
No listed prime divides $N$, since it would then divide $1$.
This contradicts the list's completeness.
\end{proof}
\idea The product-plus-one construction does not claim that $N$
itself is prime. It forces at least one new prime divisor.
For example $2\cdot3\cdot5\cdot7\cdot11\cdot13+1=30031=59\cdot509$.

\begin{example}[Arbitrarily long gaps without primes]
For every positive integer $k$, construct $k$ consecutive composite
integers.
\end{example}
\begin{solution}
Put $M=(k+1)!$. The integers $M+2,M+3,\ldots,M+k+1$
are divisible by $2,3,\ldots,k+1$, respectively, and each exceeds
its corresponding divisor. Thus all $k$ are composite.
\end{solution}
Infinitely many primes and arbitrarily long composite runs coexist.
A statement about how often primes occur is different from either.

\section{Divisor counts and divisor sums}
Write $n=p_1^{e_1}\cdots p_r^{e_r}$ with distinct primes and
$e_i\ge1$. A positive divisor chooses independently one exponent
$f_i\in\{0,\ldots,e_i\}$ for each prime. Consequently
\[
\tau(n)=\prod_{i=1}^r(e_i+1),\qquad
\sigma(n)=\prod_{i=1}^r(1+p_i+\cdots+p_i^{e_i}).
\]
Here $\tau(n)$ counts positive divisors and $\sigma(n)$ sums them.
The second formula follows by expanding the product: every choice
of exponents contributes exactly one divisor.

\begin{example}[A condition on the number of divisors]
Classify positive integers with exactly three positive divisors.
\end{example}
\begin{solution}
The product $\prod(e_i+1)=3$ has just one factor larger than $1$,
equal to $3$. Thus $n=p^2$ for a prime $p$.
Conversely $p^2$ has divisors $1,p,p^2$, so every such number works.
\end{solution}

\section{Totients and multiplicative functions}
Euler's totient $\varphi(n)$ counts $1\le a\le n$ with
$\gcd(a,n)=1$. Among $1,\ldots,p^e$, exactly $p^{e-1}$
are multiples of $p$, giving
$\varphi(p^e)=p^e-p^{e-1}$.
More generally, excluding multiples of each distinct prime divisor
gives
\[
\varphi(n)=n\prod_{p\mid n}\left(1-\frac1p\right).
\]
To justify the exclusions, expand the product. For each subset of
prime divisors, $n/\prod p$ counts the multiples of all primes in
that subset; inclusion--exclusion counts each coprime integer once
and every other integer zero times.

A function $f$ on positive integers is \emph{multiplicative} if
$f(mn)=f(m)f(n)$ whenever $\gcd(m,n)=1$, and $f(1)=1$.
The prime-factor formulas show that $\tau,\sigma,\varphi$ are
multiplicative. They are not all multiplicative without coprimality:
$\tau(2)\tau(2)=4$ but $\tau(4)=3$.

\section{Mobius inversion as cancellation}
Define $\mu(1)=1$, $\mu(n)=0$ if a prime square divides $n$,
and $\mu(n)=(-1)^r$ if $n$ is a product of $r$ distinct primes.
Then
\[
\sum_{d\mid n}\mu(d)=
\begin{cases}1&n=1,\\0&n>1.\end{cases}
\]
Indeed, only squarefree divisors contribute, and for $r>0$
their contributions sum to $(1-1)^r=0$.

\begin{theorem}[Mobius inversion]\label{n:mobius}
If $F(n)=\sum_{d\mid n}f(d)$ for all positive $n$, then
$f(n)=\sum_{d\mid n}\mu(d)F(n/d)$.
\end{theorem}
\begin{proof}
Expand the latter sum and collect the coefficient of $f(k)$:
it is $\sum_{d\mid n/k}\mu(d)$, equal to $1$ for $k=n$ and
$0$ otherwise. All sums are finite, so rearranging is legitimate.
\end{proof}
\recognition Use inversion when a counting function includes objects
whose periods, gcds, or primitive components divide a specified integer.
First identify what the divisor sum counts; the formula should remove
overcounting, not conceal it.

\begin{example}[A useful totient identity]
Prove $\sum_{d\mid n}\varphi(d)=n$.
\end{example}
\begin{solution}
Partition $k\in\{1,\ldots,n\}$ by $g=\gcd(k,n)$.
Writing $k=gu,n=gd$ leaves $1\le u\le d$ and $\gcd(u,d)=1$,
so the class has $\varphi(d)$ members. The classes partition $n$
integers, proving the identity.
\end{solution}

\section{Problems}
\begin{exercise}[\problemlevel{Foundation}]\label{n:ex09}
Find $\tau(72),\sigma(72),\varphi(72)$.
\end{exercise}
\begin{exercise}[\problemlevel{Foundation}]\label{n:ex10}
Prove that a positive integer has an odd number of positive divisors exactly when it is a square.
\end{exercise}
\begin{exercise}[\problemlevel{Bridge}]\label{n:ex11}
For every $k\ge1$, construct $k$ consecutive composite integers, all greater than $k$.
\end{exercise}
\begin{exercise}[\problemlevel{Bridge}]\label{n:ex12}
Classify positive $n$ with exactly six positive divisors.
\end{exercise}
\begin{exercise}[\problemlevel{Challenge}]\label{n:ex13}
Prove that infinitely many primes are congruent to $3$ modulo $4$.
\end{exercise}
\begin{exercise}[\problemlevel{Bridge}]\label{n:ex14}
Prove that $\sigma$ is multiplicative, and give a counterexample when the inputs are not required to be coprime.
\end{exercise}
\begin{exercise}[\problemlevel{Bridge}]\label{n:ex15}
Using $\sum_{d\mid n}\varphi(d)=n$, express $\varphi(n)$ as a sum involving $\mu$ and derive the prime-product formula.
\end{exercise}
\begin{exercise}[\problemlevel{Challenge}]\label{n:ex16}
Let $n>1$. Prove $\prod_{d\mid n}d=n^{\tau(n)/2}$, interpreting the right side as a positive real power.
\end{exercise}
\chapter{Congruences, Inverses, and the Chinese Remainder Theorem}

\label{ch:nt-congruences}
\intuition Congruences keep the remainder and discard the quotient.
This loses information deliberately: an equation may be hard to solve
in integers but impossible already modulo a small number. When several
remainders are known, the Chinese remainder theorem reconstructs their
combined information.

\section{Arithmetic with residues}
For $m>0$, the congruence $a\equiv b\pmod m$ means $m\mid a-b$.
It is an equivalence relation. If $a\equiv b$ and $c\equiv d$,
then $a+c\equiv b+d$ and $ac\equiv bd$, all modulo $m$.
For multiplication, write $ac-bd=a(c-d)+d(a-b)$.
Consequently integer polynomials preserve congruence.

Cancellation needs an extra condition:
\[
ac\equiv bc\pmod m
\quad\Longleftrightarrow\quad
a\equiv b\pmod{m/\gcd(c,m)}.
\]
To prove this, divide $m,c$ by their gcd and cancel the now coprime
factor using Lemma~\ref{n:coprime-cancel}. The modulus may shrink.
For example $2\cdot1\equiv2\cdot4\pmod6$, although $1\not\equiv4\pmod6$.

\section{Inverses and linear congruences}
\begin{theorem}\label{n:inverse}
An integer $a$ has a multiplicative inverse modulo $m$ exactly when
$\gcd(a,m)=1$. For $g=\gcd(a,m)$, the congruence $ax\equiv b\pmod m$
is solvable exactly when $g\mid b$. If solvable, it has exactly $g$
solutions modulo $m$, forming one residue class modulo $m/g$.
\end{theorem}
\begin{proof}
For $a=0$, the linear congruence is solvable exactly when $m\mid b$,
and then every one of the $m$ classes is a solution. An inverse
exists in this case only for modulus $1$, whose single residue class
satisfies the defining congruence. For $a\ne0$, an inverse exists
precisely when $ax+my=1$ is solvable, by
Theorem~\ref{n:linear-dio}. Likewise $ax\equiv b\pmod m$
is $ax+my=b$ for some integer $y$. Divide by $g$ and use the inverse
of $a/g$ modulo $m/g$. One class modulo $m/g$ has $g$ representatives
among $0,\ldots,m-1$.
\end{proof}

\begin{example}[A modulus changes after cancellation]
Solve $14x\equiv21\pmod{35}$.
\end{example}
\begin{solution}
The gcd is $7$, dividing $21$. Dividing the entire divisibility
statement by $7$ gives $2x\equiv3\pmod5$, hence $x\equiv4\pmod5$.
The solutions modulo $35$ are $4,9,14,19,24,29,34$.
\end{solution}
\pitfall Dividing the coefficients while leaving modulus $35$
would silently discard valid solutions.

\section{Coprime remainders can be assembled}
\begin{theorem}[Chinese remainder theorem]\label{n:crt}
If positive integers $m_1,\ldots,m_r$ are pairwise coprime, any
specified residues $a_i$ determine exactly one class $x$ modulo
$M=\prod_i m_i$ with $x\equiv a_i\pmod{m_i}$.
\end{theorem}
\begin{proof}
Put $M_i=M/m_i$. Since $\gcd(M_i,m_i)=1$, choose $u_i$ with
$M_i u_i\equiv1\pmod{m_i}$. The integer
$x=\sum_i a_iM_i u_i$ has the required residues: all terms except
the $i$th vanish modulo $m_i$. If $x,y$ are two solutions, every
$m_i$ divides $x-y$, so their coprime product divides $x-y$.
\end{proof}
\idea The quantities $M_i u_i$ act like switches: one in a selected
modulus and zero in all the others.

\begin{example}[Build a controlled pattern]
Find an integer that is $2$ modulo $3$, $3$ modulo $5$, and $2$ modulo $7$.
\end{example}
\begin{solution}
The first and third conditions give $x=2+21t$. Modulo $5$, this
requires $2+t\equiv3$, so $t\equiv1\pmod5$.
Thus $x\equiv23\pmod{105}$. Direct checking verifies all three residues.
\end{solution}

\section{Noncoprime moduli: compatibility comes first}
\begin{theorem}[Generalized CRT]\label{n:general-crt}
The pair $x\equiv a\pmod m$, $x\equiv b\pmod n$ is solvable
exactly when $\gcd(m,n)\mid b-a$. If so, its solutions form one
class modulo $\lcm(m,n)$.
\end{theorem}
\begin{proof}
Set $x=a+mt$. The other condition becomes $mt\equiv b-a\pmod n$.
Theorem~\ref{n:inverse} supplies exactly the compatibility criterion
and one class of $t$ modulo $n/g$, where $g=\gcd(m,n)$.
Therefore $x$ has one class modulo $mn/g$.
\end{proof}
For several noncoprime conditions, merge them successively, updating
the modulus to the lcm at each step. Test compatibility before searching.

\section{Digits and rational expansions}
In base $B\ge2$, an integer has form
$N=d_0+d_1B+\cdots+d_kB^k$ with $0\le d_i<B$.
Modulo $B-1$, this is its digit sum. Modulo $B+1$, it is the
alternating sum $d_0-d_1+d_2-\cdots$. These tests come from
substitution $B\equiv1$ or $-1$, not a special property of base ten.

\begin{example}[A decimal expansion is a finite-state process]
Prove that a rational number has an eventually periodic base-$B$
expansion.
\end{example}
\begin{solution}
It suffices to divide $0\le a<m$ by $m$. Long division replaces a
remainder $r$ by the remainder of $Br$ modulo $m$, recording the
quotient digit. There are only $m$ possible remainders.
If remainder zero occurs, the expansion terminates. Otherwise a
remainder repeats, and the deterministic rule repeats all subsequent
digits. Thus a nonrepeating initial segment is followed by a cycle.
\end{solution}
Conversely a repeating block represents a geometric series and hence
a rational number; the algebra chapters develop that calculation.

\section{Square residues and modular obstructions}
Before using a modulus, list what the relevant expressions can be.
Squares modulo $4$ are $0,1$, and modulo $8$ they are $0,1,4$.
These tiny lists can rule out infinitely many integer possibilities.

\begin{example}[Choose a modulus that distinguishes squares]
Prove $x^2+y^2=4z+3$ has no integer solution.
\end{example}
\begin{solution}
Modulo $4$, each square is $0$ or $1$, so their sum can be
$0,1,2$, never $3$.
\end{solution}
\recognition Look for a small modulus in which the allowed powers
occupy few classes and the target occupies a missing class.
A congruence obstruction proves nonexistence; absence of an obstruction
does not prove existence.

\section{Problems}
\begin{exercise}[\problemlevel{Foundation}]\label{n:ex17}
Solve $18x\equiv30\pmod{42}$, listing all classes modulo $42$.
\end{exercise}
\begin{exercise}[\problemlevel{Foundation}]\label{n:ex18}
Solve $x\equiv5\pmod{12}$ and $x\equiv11\pmod{18}$.
\end{exercise}
\begin{exercise}[\problemlevel{Bridge}]\label{n:ex19}
Show that $x\equiv1\pmod4$, $x\equiv2\pmod6$ has no solution.
\end{exercise}
\begin{exercise}[\problemlevel{Bridge}]\label{n:ex20}
List all possible residues of an integer square modulo $8$ and prove the list is complete.
\end{exercise}
\begin{exercise}[\problemlevel{Bridge}]\label{n:ex21}
Find the last two decimal digits of $7^{2026}$.
\end{exercise}
\begin{exercise}[\problemlevel{Bridge}]\label{n:ex22}
Prove the base-$B$ digit-sum test modulo $B-1$ and alternating-digit-sum test modulo $B+1$.
\end{exercise}
\begin{exercise}[\problemlevel{Challenge}]\label{n:ex23}
For every $k\ge1$, prove there are $k$ consecutive positive integers each divisible by the square of some prime.
\end{exercise}
\begin{exercise}[\problemlevel{Challenge}]\label{n:ex24}
Let $a_n$ be defined by $a_0=0$, $a_1=1$, and $a_{n+2}=a_{n+1}+a_n$. Prove that for every $m\ge2$, the sequence of pairs $(a_n,a_{n+1})$ modulo $m$ is periodic from the first pair, and some positive $a_n$ is divisible by $m$.
\end{exercise}
\chapter{Powers, Orders, and Primitive Roots}

\label{ch:nt-orders}
\intuition A power sequence modulo $m$ has only finitely many possible
values. For an invertible base, multiplication can be undone, so the
sequence cycles back to $1$. The length of that cycle is the order.
An order converts exponential divisibility into divisibility of
exponents.

\section{Fermat's and Euler's theorems}
\begin{theorem}[Fermat and Euler]\label{n:fermat-euler}
If $p$ is prime and $p\nmid a$, then $a^{p-1}\equiv1\pmod p$.
More generally, if $\gcd(a,m)=1$, then
$a^{\varphi(m)}\equiv1\pmod m$.
\end{theorem}
\begin{proof}
Multiplication by $a$ permutes the invertible residue classes:
if $ax\equiv ay$, multiply by an inverse of $a$ to get $x\equiv y$.
Let their representatives be $u_1,\ldots,u_{\varphi(m)}$.
Multiplying the permuted classes gives
$a^{\varphi(m)}\prod_i u_i\equiv\prod_i u_i$.
The product is invertible, so cancel it. For prime $p$, every
nonzero class is invertible and $\varphi(p)=p-1$.
\end{proof}
\pitfall The coprimality cannot be dropped: $2^{\varphi(4)}=4$
is not $1$ modulo $4$. Euler's theorem does not say
$a^{\varphi(m)}\equiv1$ for arbitrary $a$.

\begin{example}[A large exponent becomes small]
Find the residue of $3^{1000}$ modulo $13$.
\end{example}
\begin{solution}
Fermat reduces the exponent modulo $12$, giving $3^4=81\equiv3$.
A shorter cycle is visible: $3^3=27\equiv1$, so reducing $1000$
modulo $3$ also gives $3$.
\end{solution}
Fermat supplies a period, not necessarily the least period.

\section{Order is the exact period}
\begin{definition}
For $m\ge2$ and $\gcd(a,m)=1$, $\ord_m(a)$ is the smallest
positive $d$ with $a^d\equiv1\pmod m$.
\end{definition}
\begin{theorem}\label{n:order-divides}
If $d=\ord_m(a)$, then $a^k\equiv1\pmod m$ exactly when $d\mid k$.
In particular $d\mid\varphi(m)$.
\end{theorem}
\begin{proof}
Write $k=qd+r$, $0\le r<d$. Then $a^k\equiv a^r$.
If it is $1$, minimality forces $r=0$ (unless $r=0$ already).
The converse follows by taking a power of $a^d\equiv1$.
Apply Euler's theorem for the final assertion.
\end{proof}
The same criterion holds for negative $k$ when negative powers
mean powers of the inverse.

\begin{example}[Constrain a prime divisor]
Let $q$ be a prime divisor of $2^{11}-1$. Prove $q\equiv1\pmod{22}$.
\end{example}
\begin{solution}
The prime $q$ is odd. The order of $2$ modulo $q$ divides $11$,
so it is $1$ or $11$. Order $1$ would mean $q\mid2-1=1$,
impossible. Thus $11\mid q-1$ by Fermat. Since $q-1$ is even,
$22\mid q-1$.
\end{solution}
\recognition A prime divisor of $a^n-1$ often carries an order
dividing $n$. If proper divisors of $n$ can be excluded, the
prime divisor itself lies in a constrained progression.

\section{Wilson's theorem}
\begin{theorem}[Wilson]\label{n:wilson}
An integer $p>1$ is prime if and only if
$(p-1)!\equiv-1\pmod p$.
\end{theorem}
\begin{proof}
For prime $p$, pair each nonzero residue with its inverse.
An unpaired residue satisfies $x^2\equiv1$, hence
$(x-1)(x+1)\equiv0$. Primality gives $x\equiv1$ or $-1$.
For odd $p$, their product is $-1$ and all other pairs contribute $1$.
For $p=2$, the formula holds directly.
Conversely, if composite $p$ had the congruence, any prime divisor
$q<p$ would divide $(p-1)!$ and $p$, so reducing modulo $q$
would give $0\equiv-1$, impossible.
\end{proof}

\begin{example}[A square root of minus one]
If prime $p\equiv1\pmod4$, show that $-1$ is a square modulo $p$.
\end{example}
\begin{solution}
Pair $k$ with $p-k$ in $(p-1)!$. This gives
\[
(p-1)!\equiv(-1)^{(p-1)/2}
\left(\left(\frac{p-1}{2}\right)!\right)^2\pmod p.
\]
The exponent is even, and Wilson makes the left side $-1$.
Thus the half-factorial is a square root of $-1$ modulo $p$.
\end{solution}

\section{Why a primitive root exists modulo a prime}
A \emph{primitive root} modulo $p$ has order $p-1$; its powers
list every nonzero residue once.
This is a theorem, not an automatic property of finite modular systems.
Modulo $8$, every odd square is $1$, so there is no element of order $4$.

\begin{lemma}\label{n:field-root-bound}
A nonzero degree-$d$ polynomial modulo a prime $p$ has at most $d$ roots.
\end{lemma}
\begin{proof}
If $a$ is a root, polynomial division gives a factor $(x-a)$.
For another root $b\ne a$, the nonzero factor $b-a$ is invertible,
so the remaining polynomial vanishes at $b$. Induct on the degree.
A nonzero constant has no roots.
\end{proof}

\begin{theorem}\label{n:primitive-root}
For every prime $p$, a primitive root modulo $p$ exists.
\end{theorem}
\begin{proof}
For $p=2$, residue $1$ works. Let
$p-1=\prod_i q_i^{e_i}$ be the prime factorization.
The polynomial $X^{(p-1)/q_i}-1$ has fewer than $p-1$ roots,
so choose a nonzero $a_i$ that is not a root.
Its order $d_i$ divides $p-1$ but not $(p-1)/q_i$,
so $q_i^{e_i}\mid d_i$.
The element $b_i=a_i^{d_i/q_i^{e_i}}$ has order $q_i^{e_i}$:
its $k$th power is $1$ precisely when $d_i$ divides
$k\,d_i/q_i^{e_i}$, or $q_i^{e_i}\mid k$.

If commuting elements $u,v$ have coprime orders $r,s$, their product
has order $rs$. Indeed, $(uv)^k=1$ implies $u^k=v^{-k}$,
an element whose order divides both $r$ and $s$, hence is $1$.
Thus $r,s\mid k$, so $rs\mid k$; the converse is immediate.
Applying this repeatedly shows $\prod_i b_i$ has order $p-1$.
\end{proof}
The proof constructs prime-power pieces of a full cycle and combines
them. The polynomial root bound is what prevents a missing piece.

\section{Power congruences become linear congruences}
\begin{theorem}\label{n:power-congruence}
For prime $p$ and positive integer $k$, the nonzero solutions of
$x^k\equiv a\pmod p$
number either zero or $\gcd(k,p-1)$. In particular,
$x^k\equiv1\pmod p$ has exactly $\gcd(k,p-1)$ solutions.
\end{theorem}
\begin{proof}
If $p\mid a$, there are no nonzero solutions. Otherwise choose
a primitive root $g$ and write $x=g^t$, $a=g^s$.
The equation is $kt\equiv s\pmod{p-1}$.
Theorem~\ref{n:inverse} gives the count and solvability criterion.
For $a=1$, take $s=0$, which always satisfies the criterion.
\end{proof}
\strategy Primitive roots are most useful for structural counting,
not necessarily for computing a large discrete logarithm. A small
order calculation is often faster in a specific problem.

\section{Problems}
\begin{exercise}[\problemlevel{Foundation}]\label{n:ex25}
Find $\ord_{13}(3)$ and use it to compute $3^{2026}\pmod{13}$.
\end{exercise}
\begin{exercise}[\problemlevel{Bridge}]\label{n:ex26}
Prove that if $n>1$ and $(n-1)!\equiv-1\pmod n$, then $n$ is prime, without assuming the forward direction of Wilson's theorem.
\end{exercise}
\begin{exercise}[\problemlevel{Bridge}]\label{n:ex27}
If $p$ is an odd prime and $q$ is a prime divisor of $2^p-1$, prove $q\equiv1\pmod{2p}$.
\end{exercise}
\begin{exercise}[\problemlevel{Bridge}]\label{n:ex28}
For prime $p$ and positive $k$, find the number of solutions to $x^k\equiv1\pmod p$.
\end{exercise}
\begin{exercise}[\problemlevel{Challenge}]\label{n:ex29}
Prove that there are infinitely many primes congruent to $1$ modulo $4$.
\end{exercise}
\begin{exercise}[\problemlevel{Foundation}]\label{n:ex30}
For every integer $a$ and prime $p$, prove $p\mid a^p-a$.
\end{exercise}
\begin{exercise}[\problemlevel{Challenge}]\label{n:ex31}
For prime $p$ and positive integer $r$, evaluate $\sum_{a=1}^{p-1}a^r$ modulo $p$.
\end{exercise}
\begin{exercise}[\problemlevel{Challenge}]\label{n:ex32}
For $n\ge1$, prove that every prime divisor $q$ of $2^{2^n}+1$ satisfies $q\equiv1\pmod{2^{n+1}}$.
\end{exercise}
\chapter{Valuations, Factorials, and Lifting Exponents}

\label{ch:nt-valuations}
\intuition Divisibility by $p$ is a yes-or-no question. Valuation asks
how many copies of $p$ are present. This turns products into sums,
making large factorials and exponential differences manageable.
The main danger is addition: equal valuations can cancel and reveal
additional powers of $p$.

\section{Prime exponents as coordinates}
For a prime $p$ and nonzero integer $a$, define $v_p(a)$ to be
the exponent of $p$ in $|a|$. Thus $v_2(-40)=3$.
For nonzero rationals define $v_p(a/b)=v_p(a)-v_p(b)$; this is
well defined because cancelling a common factor changes both exponents
equally. Unique factorization gives
\[
v_p(ab)=v_p(a)+v_p(b),\qquad
v_p(\gcd(a,b))=\min\{v_p(a),v_p(b)\}
\]
for positive gcd inputs. When $a+b\ne0$,
\[
v_p(a+b)\ge\min\{v_p(a),v_p(b)\},
\]
with equality if the two valuations differ. To see this, factor out
the smaller power; the remaining sum is not divisible by $p$ if
only one summand is a multiple of $p$.

\begin{example}[Detect an impossible equality]
Show that $x^2=2y^2$ has no nonzero integer solution using valuations.
\end{example}
\begin{solution}
The equality would imply $2v_2(x)=1+2v_2(y)$.
An even integer cannot equal an odd one.
\end{solution}
This is the prime-exponent version of the descent argument.

\section{Factorial exponents}
\begin{theorem}[Legendre's formula]\label{n:legendre-factorial}
For prime $p$ and $n\ge0$,
\[
v_p(n!)=\sum_{j\ge1}\floor{\frac{n}{p^j}}
=\frac{n-s_p(n)}{p-1},
\]
where $s_p(n)$ is the sum of the base-$p$ digits of $n$.
\end{theorem}
\begin{proof}
Each multiple of $p$ contributes one factor; each multiple of $p^2$
contributes one extra, and so on. This gives the sum, which is finite.
If $n=\sum_i d_i p^i$, then
$\sum_{j\ge1}\floor{n/p^j}=\sum_i d_i(1+p+\cdots+p^{i-1})
=(n-\sum_i d_i)/(p-1)$.
\end{proof}

\begin{example}[Trailing zeroes]
Find the number of trailing decimal zeroes in $100!$.
\end{example}
\begin{solution}
A zero needs a factor $10=2\cdot5$. There are more factors of $2$
than $5$, so the answer is
$v_5(100!)=\floor{100/5}+\floor{100/25}=20+4=24$.
\end{solution}
\pitfall Counting multiples of $5$ alone misses the extra copies
inside $25,50,75,100$.

\section{Odd-prime lifting}
\begin{theorem}[LTE for an odd prime]\label{n:lte-odd}
Let $p$ be odd, $a,b$ distinct positive integers with $p\mid a-b$
and $p\nmid ab$, and $n\ge1$. Then
\[
v_p(a^n-b^n)=v_p(a-b)+v_p(n).
\]
\end{theorem}
\begin{proof}
If $p\nmid n$, divide $a^n-b^n$ by $a-b$.
The quotient is a sum of $n$ terms, each congruent to
$b^{n-1}$ modulo $p$. It is nonzero modulo $p$, proving no
additional factor appears.

For the key $p$-step, write $a=b+p^t u$ with $t\ge1$ and $p\nmid u$.
In the binomial expansion of $a^p-b^p$, the first term is
$p^{t+1}u b^{p-1}$, with valuation $t+1$. The intermediate
terms have valuations at least $2t+1\ge t+2$ because
$p\mid\binom pj$ for $1\le j\le p-1$.
The last term has valuation $pt\ge t+2$ for odd $p$.
Thus the valuation is exactly $t+1$.
Finally write $n=p^s m$, $p\nmid m$. Apply the $p$-step $s$
times, then the coprime-exponent case for $m$.
\end{proof}
\strategy Factoring exposes at least one factor $a-b$.
LTE tells us exactly how the quotient contributes additional prime
powers. Its hypotheses prevent unwanted cancellation.

\begin{example}[Choose the correct base]
Find $v_3(2^{60}-1)$.
\end{example}
\begin{solution}
We cannot apply the theorem to $2-1$, which is not divisible by $3$.
Write $2^{60}-1=4^{30}-1$. Now
$v_3(4^{30}-1)=v_3(4-1)+v_3(30)=2$.
\end{solution}

\section{The prime two needs a different formula}
\begin{theorem}[LTE at two]\label{n:lte-two}
For distinct positive odd $a,b$ and $n\ge1$,
\[
v_2(a^n-b^n)=
\begin{cases}
v_2(a-b),&n\text{ odd},\\
v_2(a-b)+v_2(a+b)+v_2(n)-1,&n\text{ even}.
\end{cases}
\]
\end{theorem}
\begin{proof}
For odd $n$, the quotient $(a^n-b^n)/(a-b)$ is a sum of an odd
number of odd terms, so is odd. For even $n=2^s m$, $m$ odd,
first discard the odd exponent by the same observation.
Factor
\[
a^{2^s}-b^{2^s}
=(a-b)(a+b)\prod_{j=1}^{s-1}(a^{2^j}+b^{2^j}).
\]
For $j\ge1$, an odd square is $1$ modulo $8$, so each sum in the
product is $2$ modulo $8$ and has valuation $1$.
There are $s-1$ such factors, giving the formula.
\end{proof}
Do not insert $p=2$ into the odd-prime formula. For example,
$v_2(3^2-1)=3$, not $v_2(3-1)+v_2(2)=2$.

\section{Lucas: binomial coefficients digit by digit}
\begin{theorem}[Lucas]\label{n:lucas}
For prime $p$, write $n=\sum_i n_i p^i$ and $k=\sum_i k_i p^i$
in base $p$. Then
\[
\binom nk\equiv\prod_i\binom{n_i}{k_i}\pmod p.
\]
\end{theorem}
\begin{proof}
All intermediate coefficients of $(1+X)^p$ are divisible by $p$,
so $(1+X)^p\equiv1+X^p\pmod p$. Iterating and using the digits,
\[
(1+X)^n\equiv\prod_i(1+X^{p^i})^{n_i}\pmod p.
\]
In each factor the chosen exponent is $j_i p^i$ with
$0\le j_i\le n_i<p$. Thus its sum has those $j_i$ as its
unique base-$p$ digits. Comparing coefficients of $X^k$ gives
the product; if a digit $k_i>n_i$, that coefficient is zero.
\end{proof}

\begin{example}[Count the odd entries of a Pascal row]
Prove that the number of odd binomial coefficients in row $n$ is
$2^{s_2(n)}$.
\end{example}
\begin{solution}
Lucas modulo $2$ says $\binom nk$ is odd exactly when each binary
digit of $k$ is at most the corresponding digit of $n$.
A zero digit forces zero; each one digit permits either choice.
The choices are independent, giving $2^{s_2(n)}$.
\end{solution}

\section{Kummer: carries count prime factors}
\begin{theorem}[Kummer]\label{n:kummer}
For $0\le k\le n$, $v_p\binom nk$ equals the number of carries
when adding $k$ and $n-k$ in base $p$, counting every carry position.
\end{theorem}
\begin{proof}
Legendre's formula gives
$v_p\binom nk=(s_p(k)+s_p(n-k)-s_p(n))/(p-1)$.
Each carry replaces $p$ units at one digit by one at the next,
decreasing the total digit sum by $p-1$. This counts all carries,
including a chain through several positions, and proves the formula.
\end{proof}
\recognition A factorial quotient suggests Legendre; a binomial
coefficient modulo a prime suggests Lucas; its exact prime exponent
suggests Kummer. These are three views of the same digit structure.

\section{Problems}
\begin{exercise}[\problemlevel{Foundation}]\label{n:ex33}
Find $v_2(100!)$ and the number of trailing decimal zeroes of $100!$.
\end{exercise}
\begin{exercise}[\problemlevel{Bridge}]\label{n:ex34}
Find $v_3(2^{60}-1)$ and explain why applying LTE directly with $a=2,b=1$ is invalid.
\end{exercise}
\begin{exercise}[\problemlevel{Bridge}]\label{n:ex35}
Find $v_2(3^{100}-1)$.
\end{exercise}
\begin{exercise}[\problemlevel{Bridge}]\label{n:ex36}
For $m\ge1$, prove $v_2\binom{2^m}{2^{m-1}}=1$.
\end{exercise}
\begin{exercise}[\problemlevel{Bridge}]\label{n:ex37}
Compute $\binom{100}{50}$ modulo $7$.
\end{exercise}
\begin{exercise}[\problemlevel{Bridge}]\label{n:ex38}
How many binomial coefficients in row $2026$ are odd?
\end{exercise}
\begin{exercise}[\problemlevel{Challenge}]\label{n:ex39}
For odd prime $p$ and positive $n$, prove $v_p((1+p)^n-1)=1+v_p(n)$, and deduce that $(1+p)^{p^r}$ has exact order $p^{s-r-1}$ modulo $p^s$ when $0\le r<s-1$.
\end{exercise}
\begin{exercise}[\problemlevel{Challenge}]\label{n:ex40}
For prime $p$, $r\ge1$, and $1\le k<p^r$, prove $v_p\binom{p^r}{k}=r-v_p(k)$.
\end{exercise}
\chapter{Diophantine Equations: Factor, Bound, and Parametrize}

\label{ch:nt-diophantine}
\intuition An integer equation contains two kinds of rigidity:
exact divisibility and the gaps between consecutive integers.
Factorization exposes the first; inequalities expose the second.
Parametrization is what we seek when there are infinitely many
solutions but only a few independent choices.

\section{A systematic first inspection}
Before manipulating, specify the domain: positive, nonnegative, or
arbitrary integers. Check zero and sign cases. Then ask:
can the equation factor; can a shared gcd be removed; can a small
modulus obstruct it; can one variable be bounded after ordering;
and does the equation contain a symmetry that produces another solution?

These are questions, not a fixed algorithm. Their order should follow
the visible structure. A product of neighboring expressions invites a
gcd computation; a square one unit from another square invites a bound.

\begin{example}[Complete a product]
Find positive integer solutions of $xy=x+y+5$.
\end{example}
\begin{solution}
Add $1$ after moving $x,y$: $(x-1)(y-1)=6$.
Neither $x$ nor $y$ can be $1$, so both new factors are positive.
The factor pairs $(1,6),(2,3),(3,2),(6,1)$ give
$(x,y)=(2,7),(3,4),(4,3),(7,2)$. Substituting back verifies them.
\end{solution}
\strategy Completing a product is the multiplicative analogue of
completing a square. The constant added must be added to both sides.
If negative integers were permitted, negative factor pairs would also
need consideration.

\section{Trap an expression between consecutive squares}
\begin{example}[A bound can finish a classification]
Prove that $n^4+2n^2+2$ is not a square for any nonnegative integer $n$.
\end{example}
\begin{solution}
It is $(n^2+1)^2+1$, strictly greater than $(n^2+1)^2$.
It is smaller than $(n^2+2)^2$, since their difference is
$2n^2+2>0$. There is no integer square strictly between two
consecutive squares.
\end{solution}
The same idea works with consecutive cubes or powers, but the gaps
must be computed rather than assumed.

\begin{example}[A hidden factorization]
Determine positive integers $n$ for which $n^4+4$ is prime.
\end{example}
\begin{solution}
The identity
\[
n^4+4=(n^2-2n+2)(n^2+2n+2)
\]
is obtained by adding and subtracting $4n^2$.
For $n\ge2$, both factors exceed $1$; the first is $(n-1)^2+1$.
At $n=1$, the number is $5$, prime. Thus only $n=1$ works.
\end{solution}
\recognition A sum of squares can factor after adding and subtracting
a middle term. The factorization matters only after both factors are
proved nontrivial in the permitted domain.

\section{Primitive Pythagorean triples}
\begin{theorem}\label{n:pythagorean}
All primitive positive integer solutions of $x^2+y^2=z^2$, after
interchanging $x,y$ if necessary, are
\[
x=m^2-n^2,\qquad y=2mn,\qquad z=m^2+n^2,
\]
where $m>n\ge1$, $\gcd(m,n)=1$, and $m,n$ have opposite parity.
All positive solutions are positive integer multiples of primitive ones.
\end{theorem}
\begin{proof}
Divide a solution by $\gcd(x,y)$ to obtain a primitive one; its
hypotenuse is divisible by that gcd because its square is.
Two odd legs would have square sum $2$ modulo $4$, impossible.
Two even legs contradict primitivity, so exactly one leg, say $y$,
is even; $x,z$ are odd.

The integers $U=(z+x)/2$, $V=(z-x)/2$ are positive, and
$UV=(y/2)^2$. Their gcd divides both $U+V=z$ and $U-V=x$.
The latter two are coprime, since a prime dividing them would divide
$y$ as well. Thus $\gcd(U,V)=1$, and each is a square:
$U=m^2,V=n^2$. This yields the displayed formulas.
Their parameters are coprime; if both were odd, $x,y,z$ would all
be even, so they have opposite parity.

Conversely the formulas satisfy the equation by expansion.
A common odd prime divisor of $m^2-n^2$ and $2mn$ would divide
both $m,n$; and the odd leg excludes a common factor $2$.
Thus the triple is primitive.
\end{proof}
\idea First turn the equation into a product of coprime factors.
Only then can a square product be split into square factors.
Skipping the coprimality calculation is a common source of incorrect
parametrizations.

\begin{example}[The parameters simplify an extra condition]
Classify primitive right triangles with one leg of length $12$.
\end{example}
\begin{solution}
The even leg is $2mn=12$, so $mn=6$.
The pairs $m>n>0$ are $(6,1)$ and $(3,2)$; both are coprime
and of opposite parity. They give legs $(35,12)$ with hypotenuse $37$,
and $(5,12)$ with hypotenuse $13$.
\end{solution}

\section{Reciprocal equations become divisor problems}
\begin{example}[Two unit fractions]
Classify positive integers $x,y$ satisfying $1/x+1/y=1/N$ for
a fixed positive integer $N$.
\end{example}
\begin{solution}
The equation gives $N(x+y)=xy$, or
$(x-N)(y-N)=N^2$.
Since each positive summand is smaller than their sum, $x,y>N$.
Thus every solution is
\[
x=N+d,\qquad y=N+\frac{N^2}{d},\qquad d\mid N^2,\ d>0.
\]
Conversely these values are positive integers and satisfy the equation.
There are $\tau(N^2)$ ordered solutions.
\end{solution}
An infinite search in two variables becomes a finite divisor list.
The exact same completion may help in problems involving harmonic
means, rectangular areas, or reciprocal inequalities.

\section{Necessary conditions are not classifications}
Parity, congruences, and valuations can remove candidates. They may
leave candidates that still fail. To finish a classification, either
construct every survivor or give an argument that there are no others.
For example, many integers are $0$ or $1$ modulo $4$ without being
squares. A residue test is a filter, not a square-detection algorithm.

\section{Problems}
\begin{exercise}[\problemlevel{Foundation}]\label{n:ex41}
Find all nonnegative integer pairs $(x,y)$ with $x^2-y^2=2025$.
\end{exercise}
\begin{exercise}[\problemlevel{Foundation}]\label{n:ex42}
Find all positive integer solutions of $xy=x+y+5$.
\end{exercise}
\begin{exercise}[\problemlevel{Bridge}]\label{n:ex43}
Determine every positive integer $n$ for which $n^4+4$ is prime.
\end{exercise}
\begin{exercise}[\problemlevel{Bridge}]\label{n:ex44}
Find all primitive Pythagorean triples with hypotenuse $65$.
\end{exercise}
\begin{exercise}[\problemlevel{Bridge}]\label{n:ex45}
Show that $x^2+y^2=3z^2$ has no positive integer solution.
\end{exercise}
\begin{exercise}[\problemlevel{Bridge}]\label{n:ex46}
Find all positive integer solutions of $1/x+1/y=1/6$ with $x\le y$.
\end{exercise}
\begin{exercise}[\problemlevel{Challenge}]\label{n:ex47}
Prove that $x^2+y^2+z^2=2xyz$ has no positive integer solution.
\end{exercise}
\begin{exercise}[\problemlevel{Challenge}]\label{n:ex48}
Let $a,b>0$ satisfy $a\mid b^2+1$ and $b\mid a^2+1$. Prove that $\gcd(a,b)=1$ and $ab\mid a^2+b^2+1$.
\end{exercise}
\chapter{Descent and Vieta Jumping}

\label{ch:nt-descent}
\intuition Some integer equations create a second solution from the
first. If a carefully chosen new solution is smaller, a least
counterexample becomes impossible. The hard work is proving that
the new solution is still a positive integer; a formal quadratic
root alone is not enough.

\section{A descent checklist}
Choose a positive integer measure, often a largest variable or the
sum of variables. Assuming a counterexample, choose one minimizing
that measure. Construct the new object and verify integrality,
positivity, preservation of the equation, and strict decrease.
The argument ends at a genuine contradiction, not at an expectation
that repeated steps will probably shrink.

\begin{example}[An obstruction that repeats]
Prove again that $x^2+y^2=3z^2$ has no positive solution.
\end{example}
\begin{solution}
Choose a solution with least $z$. Squares modulo $3$ force
$3\mid x,y$; the equation then forces $3\mid z$.
Dividing by $3$ gives a solution with smaller positive $z$.
\end{solution}
The modulus does not directly contradict the equation. It forces a
smaller copy of the same problem, which descent rules out.

\section{Vieta's other root}
Suppose an equation is quadratic in $y$:
$y^2-B(x)y+C(x)=0$.
If $y$ is an integer root and $B(x)$ is an integer, the other root
$y'=B(x)-y$ is an integer. Its product with $y$ is $C(x)$.
The sum formula proves integrality; the product formula is often
better for positivity and size.

\begin{theorem}[A complete Vieta-jump classification]\label{n:markov-chain}
All positive solutions of $x^2+y^2+1=3xy$, after ordering $x\le y$,
are consecutive terms of
\[
u_0=u_1=1,\qquad u_{k+1}=3u_k-u_{k-1}\quad(k\ge1).
\]
The pairs begin $(1,1),(1,2),(2,5),(5,13),(13,34)$.
\end{theorem}
\begin{proof}
The other root in $y$ is $y'=3x-y=(x^2+1)/y$, a positive integer.
If $x\ge2$ and $y\ge x$, it is at most $x+1/x<x+1$,
so $y'\le x$. Equality would force $xy=x^2+1$,
impossible since $x\nmid1$. Thus $0<y'<x$.
Replacing the ordered pair by $(y',x)$ reduces its sum.
Eventually $x=1$, and the equation becomes
$(y-1)(y-2)=0$.

Conversely the pair $(1,1)$ works, and replacing $(a,b)$ by
$(b,3b-a)$ preserves the equation by direct substitution.
Starting from $(1,1)$, the next entry is $2$; after that, $b>a>0$
implies $3b-a>b$. Thus all forward pairs are positive.
Every downward step is the inverse of this unique upward step
on an ordered pair, so the sequence lists all solutions.
\end{proof}
\strategy A descent can prove both nonexistence and classification.
Here it reaches a small seed, and reversing the operation builds the
entire solution tree, which happens to be a single chain.

\section{A quotient that must be a square}
\begin{example}[A difficult jump needs a sign argument]
Let $a,b$ be positive integers such that $ab+1$ divides $a^2+b^2$.
Prove the quotient is a square.
\end{example}
\begin{solution}
Write $a^2+b^2=k(ab+1)$ with positive integer $k$.
Suppose a nonsquare $k$ occurs, and choose such a pair for this
fixed $k$ with least sum; arrange $a\le b$.
The other quadratic root is
$b'=ka-b$, with $bb'=a^2-k$.

If $b'=0$, then $k=a^2$, contrary to the choice.
If $b'=-c<0$, then $bc=k-a^2<k$, while
$b=ka+c\ge k+1$, so $bc\ge b>k$, a contradiction.
Thus $b'>0$.
Moreover $b'=(a^2-k)/b<a^2/b\le a$.
The new positive pair $(b',a)$ satisfies the same equation with
the same quotient $k$, but has smaller sum. This contradiction
proves the quotient square.
\end{solution}
\idea The sign argument is the center of the proof. Writing down a
new integer root is easy; making it a smaller allowed root is the
discovery. To see that square quotients really occur, take
$(a,b)=(t,t^3)$: the quotient is $t^2$ for every positive $t$.

\section{A longer descent: a right triangle cannot have square area}
\begin{theorem}[Fermat's right-triangle theorem]\label{n:fermat-area}
No right triangle with positive integer side lengths has area equal
to a positive integer square.
\end{theorem}
\begin{proof}
Assume one exists with smallest hypotenuse. It is primitive:
dividing all sides by a common factor $g>1$ gives an integer right
triangle, which has an even leg and hence integer area.
Thus if the original area is $t^2$, then $g^2\mid t^2$,
so $g\mid t$ by prime exponents and the new area is the integer
square $(t/g)^2$.
Parametrize its legs as $m^2-n^2,2mn$, where $m>n$ are coprime
and of opposite parity. Its area is
$mn(m+n)(m-n)$.

The four positive factors $m,n,m+n,m-n$ are pairwise coprime.
For the last pair, a common divisor divides $2m,2n$, and both
$m+n,m-n$ are odd, so that gcd is $1$.
The other gcds reduce to $\gcd(m,n)=1$.
Since their product is a square, write
$m=u^2,n=v^2,m+n=w^2,m-n=t^2$.
If $m$ were even and $n$ odd, their being squares would give
$m\equiv0,n\equiv1\pmod4$, hence $m-n\equiv3\pmod4$,
impossible for a square. Thus $m$ is odd, $n$ even, and $v$ is even.
Also $w,t$ are odd and coprime.

Set $A=(w+t)/2$, $B=(w-t)/2$.
They are positive, coprime, and of opposite parity, since $A+B=w$
is odd. Their product is
$AB=(w^2-t^2)/4=v^2/2=2(v/2)^2$.
The odd factor is a square and the even factor is twice a square,
so $\{A,B\}=\{r^2,2s^2\}$ for coprime positive $r,s$.
Now
\[
u^2=m=\frac{w^2+t^2}{2}=A^2+B^2=r^4+4s^4.
\]
Thus $r^2,2s^2,u$ are integer sides of a new right triangle,
whose area is $(rs)^2$. Its hypotenuse $u$ is smaller than the
original $m^2+n^2$, because $u^2=m$ and $m\ge2$.
This contradicts minimality.
\end{proof}
The repeated use of coprime square factors is deliberate. A long
descent often consists of several short parametrizations, each removing
one kind of ambiguity.

\section{Choosing descent rather than using it automatically}
Descent is promising when an equation has a symmetry producing another
solution, when a modular condition forces common divisibility, or when
a parametrization replaces a large object by smaller parameters.
It is less promising when the transformation can leave the domain.
Always test the smallest and equality cases before announcing a
strict inequality.

\section{Problems}
\begin{exercise}[\problemlevel{Foundation}]\label{n:ex49}
Verify directly that $(1,1),(1,2),(2,5),(5,13)$ solve $x^2+y^2+1=3xy$, and compute the next ordered solution in the chain.
\end{exercise}
\begin{exercise}[\problemlevel{Bridge}]\label{n:ex50}
Prove that every positive solution of $x^2+y^2+1=3xy$ has $\gcd(x,y)=1$.
\end{exercise}
\begin{exercise}[\problemlevel{Bridge}]\label{n:ex51}
Show that $x^2+y^2=3z^2$ has no nonzero integer solution, allowing negative entries and zero entries.
\end{exercise}
\begin{exercise}[\problemlevel{Challenge}]\label{n:ex52}
Let $x,y$ be positive integers and suppose $xy\mid x^2+y^2+1$. Prove that the quotient equals $3$.
\end{exercise}
\begin{exercise}[\problemlevel{Challenge}]\label{n:ex53}
Prove that no positive integers $u>v,w$ satisfy $u^4-v^4=w^2$.
\end{exercise}
\begin{exercise}[\problemlevel{Challenge}]\label{n:ex54}
Classify positive integer pairs $(a,b)$ satisfying both $a\mid b^2+1$ and $b\mid a^2+1$.
\end{exercise}
\begin{exercise}[\problemlevel{Challenge}]\label{n:ex55}
Let $a,b$ be positive integers with $ab+1\mid a^2+b^2$. Show that every nonnegative integer square occurs as the quotient except $0$, and explain why $0$ cannot occur.
\end{exercise}
\begin{exercise}[\problemlevel{Bridge}]\label{n:ex56}
Suppose a proposed descent produces a nonnegative integer solution with a smaller sum, but the original domain requires all entries positive. Explain why this does not finish the proof, and identify where the worked Vieta-jump proof deals with this issue.
\end{exercise}
\chapter{Quadratic Residues and Sums of Two Squares}

\label{ch:nt-quadratic}
\intuition A square equation first asks a local question: can its
right side be a square modulo a prime? Quadratic reciprocity lets us
exchange a large modulus for a smaller one. The sum-of-two-squares
theorem then shows how local prime restrictions can combine into a
complete global classification.

\section{The Legendre symbol and Euler's criterion}
For odd prime $p$, define
$\left(\frac ap\right)$ to be $0$ if $p\mid a$, $1$ if $a$ is
a nonzero square modulo $p$, and $-1$ otherwise.
Among the $p-1$ nonzero classes, exactly $(p-1)/2$ are squares:
$x^2\equiv y^2$ means $x\equiv\pm y$, and the nonzero roots
pair without fixed points.

\begin{theorem}[Euler's criterion]\label{n:euler-criterion}
If $p\nmid a$, then $a^{(p-1)/2}\equiv\left(\frac ap\right)\pmod p$.
Also the Legendre symbol is multiplicative in its numerator.
\end{theorem}
\begin{proof}
Choose primitive root $g$ and write $a=g^t$.
It is a square precisely when $t$ is even.
The element $g^{(p-1)/2}$ has order $2$, so is $-1$ modulo $p$.
Thus $a^{(p-1)/2}\equiv(-1)^t$, the desired symbol.
Multiplying numerators adds the exponents $t$, proving multiplicativity;
if a numerator is divisible by $p$, both sides of the multiplicative
identity are zero.
\end{proof}
Immediately $\left(\frac{-1}{p}\right)=(-1)^{(p-1)/2}$.
Thus minus one is a square exactly when $p\equiv1\pmod4$.

\section{Gauss's lemma turns a symbol into a parity count}
\begin{lemma}[Gauss]\label{n:gauss-lemma}
For $p\nmid a$, let $N$ count those least positive residues of
$a,2a,\ldots,((p-1)/2)a$ that exceed $p/2$.
Then $\left(\frac ap\right)=(-1)^N$.
\end{lemma}
\begin{proof}
Put $m=(p-1)/2$. Replace every residue $r>p/2$ by $p-r$.
The resulting integers lie in $\{1,\ldots,m\}$ and are all distinct:
a repetition would give $ai\equiv\pm aj$, hence $i\equiv\pm j$,
which in this range forces $i=j$.
Their product is therefore $m!$. Comparing products modulo $p$,
$a^m m!\equiv(-1)^N m!$. Cancel the invertible factorial and
apply Euler's criterion.
\end{proof}

\begin{corollary}[The second supplement]\label{n:two-symbol}
For odd prime $p$,
$\left(\frac2p\right)=(-1)^{(p^2-1)/8}$.
Thus $2$ is a square precisely for $p\equiv1,7\pmod8$.
\end{corollary}
\begin{proof}
In $2,4,\ldots,p-1$, the terms exceeding $p/2$ number
$(p-1)/2-\floor{p/4}$. Its parity is even for $p\equiv1,7$
modulo $8$ and odd for $p\equiv3,5$. The exponent $(p^2-1)/8$
has the same parity in these four cases.
\end{proof}

\section{Quadratic reciprocity, with its counting proof}
\begin{lemma}\label{n:eisenstein-parity}
For positive odd $a$ with $p\nmid a$,
\[
\left(\frac ap\right)=
(-1)^{\sum_{i=1}^{(p-1)/2}\floor{ai/p}}.
\]
\end{lemma}
\begin{proof}
Write $ai=pq_i+r_i$, $1\le r_i\le p-1$.
After replacing the $N$ large residues by their negatives as in
Gauss's lemma, the numbers become a permutation of $1,\ldots,m$.
Modulo $2$, this means $\sum r_i\equiv\sum i+N$, since $p$ is odd.
But summing $ai=pq_i+r_i$ modulo $2$, with $a$ odd, gives
$\sum i\equiv\sum q_i+\sum r_i$.
Hence $N\equiv\sum q_i\pmod2$, and Gauss's lemma finishes.
\end{proof}

\begin{theorem}[Quadratic reciprocity]\label{n:quadratic-reciprocity}
For distinct odd primes $p,q$,
\[
\left(\frac qp\right)\left(\frac pq\right)
=(-1)^{((p-1)/2)((q-1)/2)}.
\]
\end{theorem}
\begin{proof}
Put $m=(p-1)/2,n=(q-1)/2$. In the integer rectangle
$1\le i\le m$, $1\le j\le n$, no point lies on $qi=pj$,
because $p,q$ are distinct primes and the coordinates are smaller
than their respective moduli.
The number below the line $pj=qi$ is
$S=\sum_{i=1}^m\floor{qi/p}$.
The number above is $T=\sum_{j=1}^n\floor{pj/q}$.
The bounds of the rectangle are correct because
$qi/p<q/2$ and $pj/q<p/2$.
Thus $S+T=mn$. Apply Lemma~\ref{n:eisenstein-parity} to the two
symbols and multiply.
\end{proof}
\idea Reciprocity is a parity count in a rectangle. Its minus sign
is not an arbitrary convention: it records the odd number of lattice
points when both primes are $3$ modulo $4$.

\begin{center}
\begin{tikzpicture}[scale=.55]
\draw[->] (0,0)--(6,0) node[right]{$i$};
\draw[->] (0,0)--(0,5) node[above]{$j$};
\foreach \x in {1,2,3,4,5}\foreach \y in {1,2,3,4}
{\fill (\x,\y) circle (2pt);}
\draw[main] (0,0)--(5.8,4.67);
\node[fill=white] at (4.6,1.1){$pj<qi$};
\node[fill=white] at (1.5,3.5){$pj>qi$};
\end{tikzpicture}
\end{center}
The sketch illustrates the partition; the proof uses a rectangle
with the exact prime-dependent dimensions.

\begin{example}[Reduce before computing]
Determine whether $5$ is a square modulo $11$.
\end{example}
\begin{solution}
Since $5\equiv1\pmod4$, reciprocity introduces no minus sign:
$\left(\frac5{11}\right)=\left(\frac{11}5\right)
=\left(\frac15\right)=1$. Indeed $4^2\equiv5\pmod{11}$.
\end{solution}
\strategy Factor the numerator, remove square factors, use the
supplements for $-1,2$, then exchange numerator and denominator.
Every exchange must account for the sign before reducing.

\section{Composite moduli and the Jacobi warning}
For odd positive $n=\prod p_i^{e_i}$, define the Jacobi symbol
$\left(\frac an\right)=\prod_i\left(\frac a{p_i}\right)^{e_i}$.
If a number coprime to $n$ is a square modulo $n$, its symbol is $1$.
The converse fails: $\left(\frac2{15}\right)=
\left(\frac23\right)\left(\frac25\right)=1$,
but $2$ is not a square even modulo $3$.
The individual prime-power conditions, assembled by CRT, decide
solvability; their product discards information.

\section{Why primes congruent to three obstruct two squares}
\begin{lemma}\label{n:three-mod-four}
If prime $q\equiv3\pmod4$ divides $x^2+y^2$, then $q\mid x,y$.
\end{lemma}
\begin{proof}
If $q\nmid y$, its inverse gives $(xy^{-1})^2\equiv-1\pmod q$,
contrary to the first supplement. Thus $q\mid y$, and then
$q\mid x^2$ implies $q\mid x$.
\end{proof}
Repeated division shows that such a prime has even exponent in
any positive integer represented as a sum of two squares.

\section{Gaussian integers supply the converse}
Define the Gaussian integers $\ZZ[i]=\{a+bi:a,b\in\ZZ\}$.
Their norm is $N(a+bi)=a^2+b^2$, and direct multiplication gives
$N(zw)=N(z)N(w)$.
A \emph{unit} has a multiplicative inverse in $\ZZ[i]$.
Its norm must be $1$, and conversely norm $1$ gives an inverse
by conjugation. Thus the units are $1,-1,i,-i$.
For nonzero $\beta$, choose a Gaussian integer $q$ by rounding the
real and imaginary parts of $\alpha/\beta$ to nearest integers.
Then $r=\alpha-q\beta$ satisfies
$N(r)\le N(\beta)/2<N(\beta)$.
Thus Euclid's algorithm and Bezout's identity work in $\ZZ[i]$,
with norm replacing absolute value.

\begin{lemma}\label{n:prime-two-squares}
If prime $p\equiv1\pmod4$, then $p=a^2+b^2$ for integers $a,b$.
\end{lemma}
\begin{proof}
Choose $t$ with $t^2\equiv-1\pmod p$ and let
$\delta=\gcd(p,t+i)$ in $\ZZ[i]$.
This gcd cannot be a unit: otherwise a Bezout identity would permit
cancelling $t+i$ from $p\mid(t+i)(t-i)$, giving $p\mid t-i$,
impossible since its imaginary coefficient is $-1$.
Nor is $\delta$ an associate of $p$, because $p$ does not divide
$t+i$. Since $\delta\mid p$, its norm divides $p^2$.
It is neither $1$ nor $p^2$; in the latter case the quotient
$p/\delta$ would have norm $1$, making them associates.
Therefore $N(\delta)=p$, the required representation.
\end{proof}

\begin{theorem}[Two-square theorem]\label{n:two-squares}
A positive integer is a sum of two integer squares, allowing zero,
exactly when each prime divisor congruent to $3$ modulo $4$
has even exponent.
\end{theorem}
\begin{proof}
Necessity follows from Lemma~\ref{n:three-mod-four}.
For sufficiency, factor the integer into primes.
Use $2=1^2+1^2$, Lemma~\ref{n:prime-two-squares} for primes
$1$ modulo $4$, and $q^{2e}=(q^e)^2+0^2$ for the others.
Products preserve the property by
\[
(a^2+b^2)(c^2+d^2)=(ac-bd)^2+(ad+bc)^2.
\]
\end{proof}

\begin{example}[A local obstruction and a global construction]
Decide whether $65$ and $21$ are sums of two integer squares.
\end{example}
\begin{solution}
Since $65=5\cdot13=(1^2+2^2)(2^2+3^2)$, multiplication gives
$65=(-4)^2+7^2$. Another sign choice gives $1^2+8^2$.
The factorization $21=3\cdot7$ contains primes $3$ modulo $4$
to odd exponents, so no representation exists.
\end{solution}

\section{Problems}
\begin{exercise}[\problemlevel{Foundation}]\label{n:ex57}
List the nonzero squares modulo $11$ and verify $\left(\frac5{11}\right)=1$.
\end{exercise}
\begin{exercise}[\problemlevel{Challenge}]\label{n:ex58}
For prime $p>3$, show that $3$ is a square modulo $p$ exactly when $p\equiv1$ or $11\pmod{12}$.
\end{exercise}
\begin{exercise}[\problemlevel{Bridge}]\label{n:ex59}
If $q\equiv3\pmod4$ is prime and $q\mid x^2+y^2$, prove that $q^2\mid x^2+y^2$.
\end{exercise}
\begin{exercise}[\problemlevel{Bridge}]\label{n:ex60}
Find all nonnegative integer pairs $x\ge y$ with $x^2+y^2=2025$.
\end{exercise}
\begin{exercise}[\problemlevel{Bridge}]\label{n:ex61}
Compute the Jacobi symbol $\left(\frac2{15}\right)$ and explain why $x^2\equiv2\pmod{15}$ is nevertheless impossible.
\end{exercise}
\begin{exercise}[\problemlevel{Bridge}]\label{n:ex62}
Use Gaussian Euclid to find a gcd of $13$ and $5+i$, and use its norm to represent $13$ as two squares.
\end{exercise}
\begin{exercise}[\problemlevel{Challenge}]\label{n:ex63}
Show that the circle $x^2+y^2=3$ has no point with both coordinates rational.
\end{exercise}
\begin{exercise}[\problemlevel{Challenge}]\label{n:ex64}
Prove that an odd prime represented as a sum of two squares has a unique representation up to signs and interchanging the coordinates.
\end{exercise}
\chapter{Approximation, Continued Fractions, and Pell Equations}

\label{ch:nt-pell}
\intuition Integer points near an irrational line create good rational
approximations. For a quadratic irrational, the error can be encoded
by an integer norm $x^2-Dy^2$. Small errors then force repeated norm
values, which produce exact Pell solutions. Continued fractions
organize the approximations and make a practical search possible.

\section{Pigeonhole approximation}
\begin{theorem}[Dirichlet's elementary approximation]\label{n:dirichlet-approx}
For irrational $\alpha$ and integer $N\ge1$, there are integers
$p,q$ with $1\le q\le N$ and
$0<|q\alpha-p|<1/N$.
Consequently infinitely many fractions satisfy
$|\alpha-p/q|<1/q^2$.
\end{theorem}
\begin{proof}
Place the $N+1$ fractional parts of
$0,\alpha,\ldots,N\alpha$ in the $N$ half-open intervals
$[j/N,(j+1)/N)$. Two lie in the same interval.
Their difference is $q\alpha-p$ with $1\le q\le N$ and absolute
value less than $1/N$, nonzero by irrationality.
Since $q\le N$, dividing gives error less than $1/q^2$.
The denominators cannot remain bounded as $N$ grows: for a finite
set of denominators, the nonzero distances from $q\alpha$ to the
nearest integer have a positive minimum, eventually larger than $1/N$.
\end{proof}
\idea We did not ask for a rational approximation directly.
We asked two fractional parts to be close, then subtracted.
This is the same collision-to-difference idea used for consecutive sums.

\section{The continued-fraction algorithm}
For a positive irrational $\alpha$, put $\alpha_0=\alpha$,
$a_n=\floor{\alpha_n}$, and
$\alpha_{n+1}=1/(\alpha_n-a_n)$.
Then $a_0\ge0$ and $a_n\ge1$ for $n\ge1$.
The notation $\alpha=[a_0;a_1,a_2,\ldots]$ records the repeated
integer-part extraction. For a rational input the algorithm ends
when a remainder is zero; it is Euclid's algorithm applied to its
numerator and denominator.

The finite truncations $P_n/Q_n=[a_0;\ldots,a_n]$ are its
\emph{convergents}. Define
\[
\begin{array}{c|cc}
 &P&Q\\ \hline
-2&0&1\\
-1&1&0
\end{array}
\qquad
P_n=a_nP_{n-1}+P_{n-2},\quad
Q_n=a_nQ_{n-1}+Q_{n-2}.
\]
These recurrences follow by substituting one final fraction at a time.
They avoid expanding a long nested expression.

\begin{theorem}[Convergent identities]\label{n:convergents}
For $n\ge0$,
\[
P_nQ_{n-1}-P_{n-1}Q_n=(-1)^{n+1},
\qquad
\alpha=\frac{P_n\alpha_{n+1}+P_{n-1}}
{Q_n\alpha_{n+1}+Q_{n-1}}.
\]
In particular the convergents are reduced, and
\[
\left|\alpha-\frac{P_n}{Q_n}\right|
<\frac1{Q_nQ_{n+1}}.
\]
\end{theorem}
\begin{proof}
The determinant is $-1$ at $n=0$ and changes sign at each
recurrence step. The expression involving the remaining tail
$\alpha_{n+1}$ follows by substituting
$\alpha_n=a_n+1/\alpha_{n+1}$ inductively.
Subtract $P_n/Q_n$ and use the determinant to obtain absolute
error $1/[Q_n(Q_n\alpha_{n+1}+Q_{n-1})]$.
Since $\alpha_{n+1}>a_{n+1}$, the second factor exceeds
$a_{n+1}Q_n+Q_{n-1}=Q_{n+1}$.
The determinant also makes any common divisor of $P_n,Q_n$
divide $1$.
\end{proof}
The denominators increase, apart from a possible equality
$Q_0=Q_1=1$, and tend to infinity.

\begin{example}[Approximate a square root without decimals]
Compute early convergents of $\sqrt2$.
\end{example}
\begin{solution}
Since $1/(\sqrt2-1)=\sqrt2+1=2+(\sqrt2-1)$, the continued
fraction is $[1;2,2,\ldots]$.
The recurrences give
$1/1,3/2,7/5,17/12,41/29,\ldots$.
Their norms $P_n^2-2Q_n^2$ alternate between $-1$ and $1$.
The small norm explains the particularly good approximation.
\end{solution}

\section{Why exceptionally close fractions are convergents}
\begin{lemma}[Best error in the numerator]\label{n:best-approx}
If $0<q<Q_{n+1}$, then
$|q\alpha-p|\ge|Q_n\alpha-P_n|$ for integers $p,q$.
Equality can occur, in particular at the convergent.
\end{lemma}
\begin{proof}
The determinant $\pm1$ allows us to write
$(p,q)=r(P_n,Q_n)+s(P_{n-1},Q_{n-1})$ with integers $r,s$:
solving the two linear equations introduces no denominator other
than $\pm1$.
Set $e_j=Q_j\alpha-P_j$. The tail identity gives
$e_{n-1}=-\alpha_{n+1}e_n$.
If $|q\alpha-p|<|e_n|$, then $|r-s\alpha_{n+1}|<1$.
For $s\ge1$, this forces $r\ge s a_{n+1}$ and hence
$q=rQ_n+sQ_{n-1}\ge Q_{n+1}$.
For $s\le-1$, it forces $r\le s a_{n+1}$ and hence $q<0$.
For $s=0$, it forces $r=0$, so $q=0$.
All cases contradict the assumed denominator range.
\end{proof}

\begin{theorem}[Legendre's approximation criterion]\label{n:legendre-approx}
If a reduced fraction $p/q$, $q>0$, satisfies
$|\alpha-p/q|<1/(2q^2)$, it is a convergent of $\alpha$.
\end{theorem}
\begin{proof}
Choose $n$ with $Q_n\le q<Q_{n+1}$, choosing the later index
if two initial denominators coincide. If $p/q\ne P_n/Q_n$,
their difference has absolute value at least $1/(qQ_n)$.
By Lemma~\ref{n:best-approx}, $|e_n|\le|q\alpha-p|$.
Thus the triangle inequality gives
\[
\left|\frac pq-\frac{P_n}{Q_n}\right|
\le\left|\frac pq-\alpha\right|+\frac{|e_n|}{Q_n}
<\frac{1+q/Q_n}{2q^2}
\le\frac1{qQ_n},
\]
a contradiction.
\end{proof}
\strategy The denominator bound says that an exceptionally good
approximation cannot hide between consecutive convergents.
The proof depends on integer determinants, not numerical evidence.

\section{Pell solutions exist}
\begin{theorem}[Existence of Pell solutions]\label{n:pell-existence}
For a positive nonsquare integer $D$, the equation
$x^2-Dy^2=1$ has a solution with positive integers $x,y$.
\end{theorem}
\begin{proof}
Dirichlet provides infinitely many positive $p,q$ with
$|\sqrt D-p/q|<1/q^2$; discard finitely many if necessary
to ensure positivity. Their nonzero integer norms satisfy
\[
|p^2-Dq^2|
<p/q+\sqrt D<2\sqrt D+1.
\]
Thus one fixed nonzero norm $m$ occurs for infinitely many pairs.
Among these, infinitely many have the same two residues modulo $|m|$.
Choose two such pairs with
$u_2=p_2+q_2\sqrt D>u_1=p_1+q_1\sqrt D>0$.
Their quotient has form
\[
\frac{u_2}{u_1}
=\frac{p_2p_1-Dq_2q_1}{m}
+\frac{q_2p_1-p_2q_1}{m}\sqrt D=X+Y\sqrt D.
\]
Both coefficients are integers: the congruences make the first
numerator congruent to $p_1^2-Dq_1^2=m$, and the second to zero,
modulo $|m|$. The quotient has norm $1$.
It exceeds $1$, and its conjugate is its positive reciprocal.
Therefore $X=(u_2/u_1+u_1/u_2)/2>0$ and
$Y=(u_2/u_1-u_1/u_2)/(2\sqrt D)>0$.
\end{proof}
\idea Repeated small norms turn an approximation into an exact
identity. The extra residue condition removes denominators from the
quotient. Without it, the coefficients need not be integers.

\section{The fundamental solution generates every solution}
\begin{theorem}\label{n:pell-classification}
Let $u=x_1+y_1\sqrt D>1$ be the least value among positive integer
solutions of $x^2-Dy^2=1$. Every positive solution is
$x+y\sqrt D=u^k$ for some integer $k\ge1$.
\end{theorem}
\begin{proof}
Existence was just proved. A least value exists by choosing the
least positive $y$: $x=\sqrt{1+Dy^2}$ increases with $y$.
For another positive solution $v$, choose $k\ge0$ so that
$1\le v/u^k<u$. Such a $k$ exists because powers of $u>1$
grow without bound. Since $u^{-1}=x_1-y_1\sqrt D$,
the quotient is $X+Y\sqrt D$ with integer coefficients and norm $1$.
If it is greater than $1$, its reciprocal conjugate makes $X,Y>0$,
contradicting minimality of $u$. Hence it equals $1$, so $v=u^k$.
Conversely all positive powers have integer coefficients and norm $1$;
their coefficients are positive by multiplication.
\end{proof}

\begin{example}[Complete classification for $D=2$]
Find all positive solutions of $x^2-2y^2=1$.
\end{example}
\begin{solution}
The smallest possible $y=1$ fails, while $y=2$ gives $x=3$.
Thus the fundamental unit is $3+2\sqrt2$ and all solutions are
its positive powers. Equivalently,
$x_{k+1}=3x_k+4y_k$, $y_{k+1}=2x_k+3y_k$, starting at $(3,2)$.
These recurrences generate $(17,12),(99,70),(577,408),\ldots$.
\end{solution}

\section{Continued fractions make the search systematic}
If $x^2-Dy^2=1$, then $\gcd(x,y)=1$ and
\[
\left|\frac xy-\sqrt D\right|
=\frac1{y^2(x/y+\sqrt D)}<\frac1{2y^2},
\]
since $D\ge2$ and $x/y>\sqrt D$.
Legendre's criterion makes $x/y$ a convergent of $\sqrt D$.
Generate convergents until norm $1$ appears. Existence guarantees
this happens; the first such convergent with $y>0$ has the least
$y$, hence is fundamental.

For quadratic irrationals these continued fractions are eventually
periodic. Here is the finite-state reason for square roots.
Write each tail as $(\sqrt D+m)/d$, with integers $d>0$ and
$d\mid D-m^2$. The next state is
$m'=ad-m$, $d'=(D-m'^2)/d$.
After the first step the conjugate tail lies between $-1$ and $0$;
this property is preserved by $t\mapsto1/(t-a)$ for $a\ge1$.
Together with the actual tail being greater than $1$, it implies
$-\sqrt D<m<\sqrt D$ and $0<d<2\sqrt D$.
Thus only finitely many integer states occur, and a repeated state
repeats all subsequent digits. The recurrence preserves integrality
because $m'\equiv-m\pmod d$.
No decimal approximation is needed.

\section{Problems}
\begin{exercise}[\problemlevel{Foundation}]\label{n:ex65}
Compute the finite continued fraction $[1;2,2,2]$.
\end{exercise}
\begin{exercise}[\problemlevel{Bridge}]\label{n:ex66}
List the first four positive solutions of $x^2-2y^2=1$ in increasing $y$, and justify that none are skipped.
\end{exercise}
\begin{exercise}[\problemlevel{Challenge}]\label{n:ex67}
Classify all positive solutions of $x^2-2y^2=-1$.
\end{exercise}
\begin{exercise}[\problemlevel{Foundation}]\label{n:ex68}
Explain why the infinitely-many-approximations conclusion for irrational $\alpha$ requires unbounded denominators.
\end{exercise}
\begin{exercise}[\problemlevel{Bridge}]\label{n:ex69}
Classify positive solutions of $x^2-3y^2=1$ and list the first three.
\end{exercise}
\begin{exercise}[\problemlevel{Bridge}]\label{n:ex70}
Prove that every continued-fraction convergent $P_n/Q_n$ is in lowest terms, without using an error estimate.
\end{exercise}
\begin{exercise}[\problemlevel{Challenge}]\label{n:ex71}
Prove that no rational $p/q$, with $q>0$, satisfies $|\sqrt2-p/q|<1/(100q^2)$. Thus Dirichlet's constant cannot be replaced by an arbitrarily small constant for every irrational.
\end{exercise}
\begin{exercise}[\problemlevel{Challenge}]\label{n:ex72}
Classify right triangles with consecutive positive integer legs, and list the first three.
\end{exercise}
\chapter{Prime Powers, Integer Polynomials, and Further Constructions}

\label{ch:nt-lifting}
\intuition A solution modulo $p$ is a first approximation to a
solution modulo $p^2,p^3,\ldots$. Lifting chooses one new base-$p$
digit at a time. Integer polynomials also let us force prime divisors
and residue patterns by choosing their inputs deliberately.

\section{A polynomial has modular memory}
If $P\in\ZZ[X]$ and $a\equiv b\pmod m$, then
$P(a)\equiv P(b)\pmod m$.
Thus its values modulo $m$ repeat with period $m$, though the
least period may be smaller.
If $P$ has rational coefficients and takes integer values on integers,
choose a positive common coefficient denominator $D$.
Then $DP\in\ZZ[X]$, and expanding powers shows
$DP(n+mD)-DP(n)$ is divisible by $mD$.
Dividing by $D$ proves $P(n+mD)\equiv P(n)\pmod m$.
The period must account for denominators.

\begin{example}[Integer-valued does not mean period $m$]
For $P(n)=n(n-1)/2$, compare $P(n)$ and $P(n+2)$ modulo $2$.
\end{example}
\begin{solution}
Their difference is $2n+1$, odd, so period $2$ fails.
The denominator argument supplies period $4$:
$P(n+4)-P(n)=4n+6$ is even.
\end{solution}

\section{Hensel lifting: solve for one more digit}
For $f(X)=\sum c_jX^j\in\ZZ[X]$, its formal derivative is
$f'(X)=\sum_{j\ge1}j c_j X^{j-1}$.
This is an algebraic definition; no calculus is required.
Expanding each power by the binomial theorem gives, for $k\ge1$,
\[
f(r+tp^k)\equiv f(r)+tp^k f'(r)\pmod{p^{k+1}},
\]
because all terms with two factors $p^k$ are divisible by $p^{k+1}$.

\begin{theorem}[Simple-root Hensel lifting]\label{n:hensel}
If prime $p$ satisfies $f(r)\equiv0\pmod{p^k}$ and
$f'(r)\not\equiv0\pmod p$, exactly one residue
$r+tp^k$ with $0\le t<p$ is a root modulo $p^{k+1}$.
The nonzero-derivative condition persists through further lifts.
\end{theorem}
\begin{proof}
Divide the expansion by $p^k$. The required condition is
$f(r)/p^k+t f'(r)\equiv0\pmod p$.
The derivative is invertible, so this determines exactly one
class $t$ modulo $p$. The lifted residue equals $r$ modulo $p$,
so its derivative has the same nonzero residue.
\end{proof}

\begin{example}[Lift a square root]
Starting from $3^2\equiv2\pmod7$, find the root above it modulo $343$.
\end{example}
\begin{solution}
For $f(X)=X^2-2$, at $r=3,k=1$ the digit condition is
$1+6t\equiv0\pmod7$, so $t=1$ and the root modulo $49$ is $10$.
Next $f(10)/49=2$ and $f'(10)\equiv6\pmod7$, giving $t=2$.
The root is $10+2\cdot49=108$ modulo $343$.
Indeed $108^2-2=34\cdot343$.
\end{solution}
\pitfall A singular root, with $f'(r)\equiv0$, may have many lifts
or none. For $f(X)=X^2$, the root $0$ modulo $p$ has all $p$
lifts $tp$ modulo $p^2$. For $f(X)=X^2-p$, that same initial root
has no lift, since $(tp)^2-p$ is not divisible by $p^2$.

\section{Primitive roots at odd prime powers}
\begin{theorem}\label{n:primitive-prime-power}
For every odd prime $p$, there is an integer that is a primitive
root modulo every $p^k$, $k\ge1$.
\end{theorem}
\begin{proof}
Choose $g$ primitive modulo $p$. If $g^{p-1}\not\equiv1\pmod{p^2}$,
keep $g$. Otherwise use $h=g+p$:
\[
h^{p-1}\equiv g^{p-1}+(p-1)g^{p-2}p\not\equiv1\pmod{p^2}.
\]
This new $h$ is still primitive modulo $p$.
In either case choose $h$ with $v_p(h^{p-1}-1)=1$.
If $h^t\equiv1\pmod{p^k}$, reduction modulo $p$ forces $t=(p-1)s$.
LTE gives $v_p(h^{(p-1)s}-1)=1+v_p(s)$.
The least permitted exponent is $(p-1)p^{k-1}=\varphi(p^k)$,
the exact order.
\end{proof}

For a general modulus $m\ge2$, a primitive root means an invertible
residue of order $\varphi(m)$, whose powers therefore list every
invertible residue exactly once.
For completeness, the moduli $m\ge2$ admitting primitive roots
are exactly $2,4,p^k,2p^k$ for odd primes $p$.
We have proved the $p^k$ case; choosing an odd representative
gives the $2p^k$ case by CRT. Moduli $2,4$ are direct.
For the converse, if all invertible residues are powers of one
primitive root and $\varphi(m)$ is even, the equation $x^2=1$
has exactly two solutions: its exponent must be $0$ or
$\varphi(m)/2$ modulo $\varphi(m)$.
If $m$ contains two distinct odd prime factors, CRT supplies at
least four sign choices for $x^2=1$.
The same holds if $4\mid m$ and an odd prime divides $m$.
For $2^k$ with $k\ge3$, the four distinct classes
$1,-1,1+2^{k-1},-1+2^{k-1}$ square to $1$.
These obstructions eliminate precisely the other factorizations of $m$.

\section{Prime values versus prime divisors}
\begin{example}[A nonconstant polynomial cannot always be prime]
Prove no nonconstant $P\in\ZZ[X]$ takes positive prime values
at every positive integer.
\end{example}
\begin{solution}
Assume it does, and choose $a>0$ with $P(a)=q$ prime.
For every $t\ge0$, $P(a+qt)\equiv P(a)\equiv0\pmod q$.
These values would all have to equal $q$.
A nonconstant polynomial cannot equal one fixed value at infinitely
many inputs, by the root bound applied to $P-q$.
\end{solution}
The failure of prime values does not mean that only finitely many
primes divide the values. In fact the opposite holds.

\begin{theorem}[Schur's prime-divisor theorem]\label{n:schur-primes}
A nonconstant integer polynomial has infinitely many prime divisors
among its nonzero integer values.
\end{theorem}
\begin{proof}
Suppose the primes were $p_1,\ldots,p_k$, and put $M=\prod p_i$
(take $M=1$ if the list is empty).
Choose an integer $a$ with $A=P(a)\ne0$.
The polynomial $Q(T)=P(a+AT)/A$ has integer coefficients:
the constant term is $1$, and every nonconstant term contains a
factor $A$ before division. Thus $Q(T)=1+TR(T)$ and is nonconstant.
For $T=Mt$, it is $1$ modulo every listed prime.
For sufficiently large $|t|$, $|Q(Mt)|>1$, so it has a prime
divisor outside the list. That prime also divides
$P(a+AMt)=A Q(Mt)$, a contradiction.
\end{proof}
\idea Choose an input that makes the new value congruent to $1$
modulo all the supposedly available prime divisors. This is Euclid's
construction adapted to a polynomial.

\section{Cyclotomic factors and exact order}
The $n$th cyclotomic polynomial $\Phi_n(X)$ is the monic
polynomial whose roots are the primitive $n$th complex roots of unity,
those with exact multiplicative order $n$.
Every $n$th root has a unique exact order dividing $n$, so
\[
X^n-1=\prod_{d\mid n}\Phi_d(X).
\]
These polynomials have integer coefficients: start with
$\Phi_1=X-1$ and use induction, dividing the integer monic
polynomial $X^n-1$ by the product of the already known monic factors.
Monic long division has integer quotient and remainder; the complex
factorization forces the remainder to be zero.
Examples are $\Phi_2=X+1$, $\Phi_3=X^2+X+1$,
$\Phi_4=X^2+1$, and $\Phi_6=X^2-X+1$.

\begin{theorem}\label{n:cyclotomic-order}
If prime $q$ divides $\Phi_n(a)$ and $q\nmid an$, then
$\ord_q(a)=n$.
\end{theorem}
\begin{proof}
The order $d$ divides $n$, since $\Phi_n(a)\mid a^n-1$.
If $d<n$, the factorization of $X^d-1$ makes $a$ a root of
some $\Phi_e$ with $e\mid d$ modulo $q$.
It is also a root of $\Phi_n$.
Thus $a$ would be a repeated root of $X^n-1$ modulo $q$.
But the derivative $nX^{n-1}$ is nonzero at the nonzero root $a$,
so such repetition is impossible.
\end{proof}
\pitfall The exception $q\mid n$ is real:
$3\mid\Phi_6(2)=3$, yet $\ord_3(2)=2$, not $6$.
Advanced primitive-divisor theorems strengthen these ideas, but for
many problems a direct factorization and an order check are enough.

\section{Repunits and finite states}
Let $R_k=(B^k-1)/(B-1)=1+B+\cdots+B^{k-1}$ in base $B$.
If $\gcd(B,m)=1$, some $R_k$ with $1\le k\le m$ is divisible
by $m$. If no $R_k$ among $R_1,\ldots,R_m$ were zero modulo $m$,
two would share one of the $m-1$ nonzero residues.
Their difference is $B^jR_{i-j}$; cancelling the invertible $B^j$
gives a shorter divisible repunit, a contradiction.
This proof also works directly by including remainder $0$ in the
finite-state process.

\section{Selecting a number-theory method}
Try a modulus to obstruct, a gcd to separate, valuations to measure,
an order to constrain exponents, factorization to enumerate,
an inequality to bound, and descent to shrink.
When a method nearly works, identify the exact missing condition:
coprimality, a simple derivative, a nonzero residue, or a positive
other root. Often the main discovery is how to create that condition.

\section{Problems}
\begin{exercise}[\problemlevel{Bridge}]\label{n:ex73}
Find all roots of $x^2\equiv2\pmod{343}$.
\end{exercise}
\begin{exercise}[\problemlevel{Foundation}]\label{n:ex74}
For prime $p$, prove that $x^2\equiv p\pmod{p^2}$ has no solution, although it has a root modulo $p$.
\end{exercise}
\begin{exercise}[\problemlevel{Challenge}]\label{n:ex75}
For $k\ge3$, classify all roots of $x^2\equiv1\pmod{2^k}$.
\end{exercise}
\begin{exercise}[\problemlevel{Bridge}]\label{n:ex76}
Find the order of $3$ modulo $7^k$ for every $k\ge1$.
\end{exercise}
\begin{exercise}[\problemlevel{Bridge}]\label{n:ex77}
Prove that a nonconstant integer polynomial cannot have absolute value prime at every integer input.
\end{exercise}
\begin{exercise}[\problemlevel{Bridge}]\label{n:ex78}
If the base-ten repunit $R_n=(10^n-1)/9$ is prime, prove that $n$ is prime. Does the converse hold?
\end{exercise}
\begin{exercise}[\problemlevel{Challenge}]\label{n:ex79}
Let $P\in\ZZ[X]$ be nonconstant with $P(0)=1$. Give a short proof that infinitely many primes divide its nonzero values.
\end{exercise}
\begin{exercise}[\problemlevel{Challenge}]\label{n:ex80}
If $a>1$ and prime $q\ne3$ divides $a^2-a+1$, prove $q\equiv1\pmod6$.
\end{exercise}


\part{Connecting the Subjects}

% ---- chapters/capstones ----
\chapter{Choosing a Method without a Chapter Title}
\label{ch:method-selection}
\intuition A topical chapter gives away part of the solution: one
expects a valuation in a valuation chapter and an inversion in an
inversion chapter. Unfamiliar problems do not supply that clue.
The next step is to identify what the problem's information can control.

\section{Diagnose the obstacle}
Start with the exact domain and target. An existence problem needs
one object, while a classification needs all of them. An inequality
needs a bound and equality conditions. A construction needs every
required property checked. These different logical tasks call for
different stopping points.

Next name the obstacle in concrete terms. Are there too many cases?
Look for a symmetry or a coarser representation. Does a sum refuse
to factor? Ask whether it can be counted or completed to a product.
Does a geometric relation involve several circles? Seek a shared power,
angle, or transformation. Are variables unbounded integers? Try a
gcd, congruence, valuation, or descent measure.

Do not run through theorem names without a purpose. A useful sentence
is: \emph{I will compare these residues because a collision would make
their difference divisible by the desired modulus.} It states both
the intended move and its payoff.

\section{A representation can bridge subjects}
\begin{example}[Rational points on a circle]
Parametrize the rational points on $x^2+y^2=1$.
\end{example}
\begin{solution}
The fixed point $(-1,0)$ is rational. Draw a line of rational slope
$t$ through it: $y=t(x+1)$. Substitution into the circle equation
gives the known intersection $x=-1$ and the other one
\[
x=\frac{1-t^2}{1+t^2},\qquad y=\frac{2t}{1+t^2}.
\]
These coordinates are rational for rational $t$, and direct substitution
verifies the circle equation. Conversely, any rational point other than
$(-1,0)$ has rational slope $t=y/(x+1)$ and therefore has this form.
The exceptional point must be added separately.
\end{solution}
\idea Geometry supplies a line through a known point; algebra solves
its second intersection; arithmetic turns rational slopes into rational
coordinates. This is also the source of the Pythagorean parametrization
after clearing denominators.

\begin{center}
\begin{tikzpicture}[scale=1.25]
\draw[->,gray] (-1.4,0)--(1.5,0) node[right]{$x$};
\draw[->,gray] (0,-1.25)--(0,1.4) node[above]{$y$};
\draw[main] (0,0) circle (1);
\draw (-1,0)--(1.2,1.1);
\fill (-1,0) circle (1.2pt) node[below left]{$(-1,0)$};
\fill (.6,.8) circle (1.2pt) node[above right]{$(x,y)$};
\node[fill=white,inner sep=1pt] at (.5,1.22){$y=t(x+1)$};
\end{tikzpicture}
\end{center}
The plotted example has slope $t=1/2$ and point $(3/5,4/5)$.

\section{Extremal choices create structure}
\begin{example}[A shortest assignment does not cross]
There are $n$ red and $n$ blue points in the plane, all distinct,
with no three collinear. Prove that they can be joined in red--blue
pairs so that the joining segments do not cross.
\end{example}
\begin{solution}
Among the finitely many bijections between colors, choose one minimizing
total segment length. Suppose segments $R_1B_1$ and $R_2B_2$
cross at an interior point $X$. Replacing them by
$R_1B_2,R_2B_1$ preserves the pairing and gives, by strict triangle
inequalities,
\[
R_1B_2+R_2B_1
<(R_1X+XB_2)+(R_2X+XB_1)
=R_1B_1+R_2B_2.
\]
Strictness follows because the relevant triples are noncollinear.
This contradicts minimality, so the chosen matching does not cross.
\end{solution}
\strategy The minimum itself need not be computed. We choose it for
the property that no local improvement can exist. Combinatorics
supplies the finite family, and Euclidean distance supplies the
strict improvement.

\section{Check a proposed proof against nearby cases}
An argument that works only for a pictured acute triangle needs a
separate obtuse argument or a more robust representation. An integer
root that becomes zero cannot be used in a positive-integer descent.
A bound that is never attained is an infimum, not necessarily a minimum.
A random construction with positive expected surplus needs an
outcome selected from that average before it becomes an existence proof.

These checks are part of solving, not cosmetic corrections at the end.
They can reveal the decisive extra hypothesis or a stronger true result.

\section{From solved problems to reusable knowledge}
After finishing, extract the smallest useful principle. Record what
made it recognizable, not only the final algebra. For example:
\begin{itemize}
\item A consecutive sum is a difference of partial sums.
\item A coprime square product splits into square factors.
\item Two circles invite a radical axis; a common point can turn powers
into products of directed lengths.
\item A quadratic integer equation may have another integer root.
\item A minimum configuration cannot admit a strict local improvement.
\end{itemize}
Each principle needs a domain and a failure case. This makes it
available in a new setting rather than tied to the original wording.

\chapter{A Polynomial Bridge to Zero Sums}
\label{ch:zero-sum-bridge}
\intuition A finite-field polynomial can act as a detector: zero on
the configurations we want and nonzero elsewhere. Summing that
detector over all inputs can force the desired configurations to exist.
This chapter combines polynomials, modular arithmetic, and counting
to prove a strong zero-sum theorem from elementary ingredients.

\section{Power sums over a prime field}
Fix a prime $p$, and work in the field $\FF_p$ of residues modulo $p$.
For a nonnegative integer $e$,
\[
\sum_{t\in\FF_p}t^e=
\begin{cases}
0,&e=0,\\
0,&e>0,\ p-1\nmid e,\\
-1,&e>0,\ p-1\mid e.
\end{cases}
\quad\text{in }\FF_p.
\]
For $e=0$, each term means the constant polynomial $1$, including
at $t=0$, so the sum is $p=0$ in the field.
For positive $e$, use a primitive root and sum a geometric series,
as in the number-theory power-sum problem.

Consequently, if a multivariable monomial has even one variable
with exponent less than $p-1$, its sum over all inputs in $\FF_p^n$
is zero. The sum factors into one-variable sums.

\section{Chevalley--Warning in the form we need}
\begin{theorem}\label{m:chevalley-warning}
Let $F_1,\ldots,F_r$ be polynomials over $\FF_p$ in $n\ge1$
variables, with $\sum_j\deg F_j<n$.
The number of their common zeros in $\FF_p^n$ is divisible by $p$.
Zero polynomials may simply be omitted from the list.
\end{theorem}
\begin{proof}
At an input $x$, Fermat's theorem makes
\[
I(x)=\prod_{j=1}^r\bigl(1-F_j(x)^{p-1}\bigr)
\]
equal to $1$ when every $F_j(x)=0$, and to $0$ otherwise.
Thus $\sum_x I(x)$ is the common-zero count modulo $p$.
After expanding $I$, every monomial has total degree less than
$n(p-1)$, so not all of its $n$ exponents can be at least $p-1$.
One exponent is smaller, and its power sum vanishes.
Every monomial therefore sums to zero, and so does $I$.
\end{proof}
\idea We counted solutions modulo $p$ without finding them individually.
If the all-zero input is already one solution, divisibility can
force another. The degree condition is the reason the counting
identity cancels; it cannot be omitted.

\begin{example}[A guaranteed nonzero solution]
Let $a,b,c\in\FF_p$. Prove that
$ax^2+by^2+cz^2=0$ has a solution other than $(0,0,0)$.
\end{example}
\begin{solution}
If the polynomial is identically zero, every triple is a solution.
Otherwise its degree is at most $2<3$, so the common-zero count
is divisible by $p$. Since the zero triple is one solution, it
cannot be the only one. This works also for $p=2$.
\end{solution}
This finite-field assertion is different from a positive-integer
solution of an equality over $\ZZ$. The domain is essential.

\section{The prime zero-sum case}
\begin{theorem}\label{m:egz-prime}
Among any $2p-1$ integers, where $p$ is prime, some $p$ of them
have sum divisible by $p$.
\end{theorem}
\begin{proof}
Reduce the integers to $a_1,\ldots,a_{2p-1}\in\FF_p$.
In $2p-1$ variables use
\[
F_1=\sum_i X_i^{p-1},\qquad
F_2=\sum_i a_iX_i^{p-1}.
\]
Their degree sum is at most $2p-2<2p-1$.
For a common zero, the set $S$ of nonzero coordinates has
$|S|\equiv0\pmod p$ and $\sum_{i\in S}a_i=0$.
Since $0\le|S|\le2p-1$, its size is $0$ or $p$.
The empty support contributes one zero.
If $A$ is the number of $p$-element supports with zero sum, each
has $(p-1)^p$ assignments of nonzero coordinates. Hence the
common-zero count is
$1+(p-1)^p A\equiv1-A\pmod p$.
It is divisible by $p$, so $A\equiv1\pmod p$; in particular $A>0$.
\end{proof}
\strategy The exponent $p-1$ turns every nonzero coordinate into $1$.
One equation controls support size; the other controls its weighted sum.
Their combined degree fits just below the number of variables.

\section{Build the composite case from smaller cases}
\begin{theorem}[Erdos--Ginzburg--Ziv]\label{m:egz}
Among any $2n-1$ integers, some $n$ have sum divisible by $n$,
for every positive integer $n$.
\end{theorem}
\begin{proof}
Use strong induction on $n$. The case $1$ is immediate, and prime
$n$ was proved. For composite $n=ab$, $a,b>1$, take $2ab-1$
integers. Repeatedly use the theorem for $a$ to choose
$2b-1$ disjoint blocks of $a$ integers, each with sum divisible by $a$.
Before choosing the final block there remain
$2ab-1-(2b-2)a=2a-1$ integers, so every choice is justified.
Write the block sums as $ac_1,\ldots,ac_{2b-1}$.
By the theorem for $b$, some $b$ of the $c_j$ have sum divisible
by $b$. Their blocks contain $ab=n$ original entries and their
combined sum is divisible by $ab$.
\end{proof}
\pitfall Choosing blocks whose sums are merely equal modulo $a$
would not automatically allow division by $a$. We selected sums
equal to zero modulo $a$, then divided exact integer sums.

\begin{example}[The bound is sharp]
For $n\ge2$, show that $2n-2$ entries do not always suffice.
\end{example}
\begin{solution}
Take $n-1$ zeroes and $n-1$ ones.
Any choice of $n$ entries contains between $1$ and $n-1$ ones,
so its sum is not divisible by $n$.
\end{solution}
This example completes an extremal result: we proved a bound,
then constructed an obstruction immediately below it.

\chapter{Mixed Problems}
\label{ch:mixed-problems}
The subjects are deliberately interleaved. Before selecting a method,
state the domain and the logical target. The hints identify a direction;
the solutions show both discovery and verification.

\begin{exercise}[\problemlevel{Bridge}]\label{m:ex01}
Classify real polynomials $P$ such that $P(P(x))=P(x)^2$ for every real $x$.
\end{exercise}
\begin{exercise}[\problemlevel{Bridge}]\label{m:ex02}
Find all functions $f:\RR\to\RR$ satisfying
$f(x^2+y)=xf(x)+f(y)$ for every real $x,y$.
\end{exercise}
\begin{exercise}[\problemlevel{Challenge}]\label{m:ex03}
If an integer $n>1$ divides $2^n+1$, prove that $3\mid n$.
\end{exercise}
\begin{exercise}[\problemlevel{Bridge}]\label{m:ex04}
Every entry of a $6$ by $6$ array is colored black or white.
Prove that two rows and two columns contain four corners of the same color.
\end{exercise}
\begin{exercise}[\problemlevel{Bridge}]\label{m:ex05}
Diagonals $AC,BD$ of a convex quadrilateral intersect at $P$.
Prove that the quadrilateral is cyclic if and only if
$PA\cdot PC=PB\cdot PD$.
\end{exercise}
\begin{exercise}[\problemlevel{Bridge}]\label{m:ex06}
Find all rational points of $x^2+y^2=1$, including any point missed
by a slope parametrization.
\end{exercise}
\begin{exercise}[\problemlevel{Bridge}]\label{m:ex07}
For positive $a,b,c$ with $a+b+c=1$, prove
$a^2/(1-a)+b^2/(1-b)+c^2/(1-c)\ge1/2$ and find equality.
\end{exercise}
\begin{exercise}[\problemlevel{Challenge}]\label{m:ex08}
Classify polynomials $P\in\ZZ[X]$ such that for every positive integer
$n$, $P(n)\ne0$ and $P(n)\mid n$.
\end{exercise}
\begin{exercise}[\problemlevel{Challenge}]\label{m:ex09}
Show that every red--blue edge coloring of the complete graph on nine
vertices contains a red triangle or a blue complete graph on four vertices.
Construct a coloring on eight vertices containing neither.
\end{exercise}
\begin{exercise}[\problemlevel{Challenge}]\label{m:ex10}
For positive integers $a,b$ with $ab>1$, suppose $ab-1$ divides
$a^2+b^2$. Prove that the quotient is $5$.
\end{exercise}
\begin{exercise}[\problemlevel{Challenge}]\label{m:ex11}
Let $S$ be a finite noncollinear set of points in the Euclidean plane.
Prove that some line passes through exactly two points of $S$.
\end{exercise}
\begin{exercise}[\problemlevel{Challenge}]\label{m:ex12}
Prove that $n$ red and $n$ blue distinct points, no three collinear,
admit a red--blue pairing with no crossing segments. Explain why
choosing a minimum is justified.
\end{exercise}
\begin{exercise}[\problemlevel{Bridge}]\label{m:ex13}
Prove that no right triangle with positive rational side lengths
has area $1$.
\end{exercise}
\begin{exercise}[\problemlevel{Challenge}]\label{m:ex14}
In any triangle with circumradius $R$ and inradius $r$, prove
$R\ge2r$ and classify equality. Derive the identity
$OI^2=R(R-2r)$ rather than citing it.
\end{exercise}
\begin{exercise}[\problemlevel{Bridge}]\label{m:ex15}
Prove the zero-sum theorem's bound $2n-1$ is best possible for every
$n\ge2$, and explain why the usual partial-sums argument alone does
not choose exactly $n$ entries.
\end{exercise}
\begin{exercise}[\problemlevel{Bridge}]\label{m:ex16}
A convex $n$-gon, $n\ge4$, has no three diagonals concurrent in
its interior. Prove that its diagonals have exactly $\binom n4$
interior intersection points.
\end{exercise}


\appendix

% ---- front/reals ----
\chapter{Real-Number Background for Advanced Methods}
\label{app:reals}
\intuition Most of the book uses finite algebraic and combinatorial
arguments. A few advanced reductions use continuity and the existence
of a maximum. This appendix supplies the real-number principles they
need. It may be read when such a result first occurs.

\section{Completeness: a property of the real numbers}
An \emph{upper bound} of a set $S\subseteq\RR$ is a number $M$
with $s\le M$ for every $s\in S$. A \emph{least upper bound},
written $\sup S$, is an upper bound no larger than any other.
The completeness axiom for $\RR$ says that every nonempty set
bounded above has a least upper bound.
This distinguishes $\RR$ from $\QQ$: the positive rational numbers
with square less than $2$ have no rational least upper bound.
Completeness is a defining background property of the real-number
system, not an olympiad theorem about integers.

A sequence $x_n$ tends to $L$ if, for every $\eps>0$, there is an
index $N$ such that $|x_n-L|<\eps$ for all $n\ge N$.
The index may depend on $\eps$. The definition controls an entire
tail, not just infinitely many individual terms.

\begin{theorem}[Bounded monotone sequences]\label{real:monotone}
An increasing real sequence bounded above converges to the supremum
of its values. A decreasing sequence bounded below converges similarly.
\end{theorem}
\begin{proof}
Let $L=\sup\{x_n:n\ge1\}$. For any $\eps>0$, the number $L-\eps$
is not an upper bound, so some $x_N>L-\eps$.
Monotonicity gives $L-\eps<x_n\le L$ for every $n\ge N$,
which is exactly convergence. Negating the terms proves the other case.
\end{proof}
For example, the recursive sequence $x_{n+1}=\sqrt{2+x_n}$ with
$x_0=0$ is increasing and bounded by $2$. Once convergence is proved,
continuity permits $L=\sqrt{2+L}$, giving $L=2$.
Solving the fixed-point equation alone would not prove convergence.

\section{Continuity and the intermediate value theorem}
A function $f$ is continuous at $x$ if for every $\eps>0$ some
$\delta>0$ makes $|f(y)-f(x)|<\eps$ whenever $y$ is in its domain
and $|y-x|<\delta$. Polynomial functions are continuous, as follows
by expanding $f(y)-f(x)$ into factors containing $y-x$ and bounding
the remaining factors near $x$.
Sums and products preserve continuity; a quotient does so where
its denominator is nonzero.

\begin{theorem}[Intermediate value theorem]\label{real:ivt}
If $f$ is continuous on $[a,b]$ and $f(a)<0<f(b)$, then
$f(c)=0$ for some $c\in(a,b)$.
\end{theorem}
\begin{proof}
Bisect the interval repeatedly. If a midpoint is a zero, stop.
Otherwise retain a half with negative left endpoint and positive
right endpoint. Let these endpoints be $a_n,b_n$.
They are respectively increasing and decreasing, bounded, and
$b_n-a_n=(b-a)/2^n\to0$.
By Theorem~\ref{real:monotone}, they converge to the same $c$.
Continuity gives $f(a_n)\to f(c)$ and $f(b_n)\to f(c)$.
The first limit is nonpositive and the second nonnegative, so
$f(c)=0$. The endpoint signs exclude $c=a,b$.
\end{proof}
Replacing $f$ by $f-t$ proves every value between $f(a)$ and
$f(b)$ is attained. Without continuity a jump can skip values.

\section{Subsequences, closed sets, and finite-dimensional bounds}
A subsequence chooses terms with strictly increasing indices.
A set $K\subseteq\RR^d$ is \emph{bounded} if all coordinates of
its points are bounded in absolute value by one constant.
It is \emph{closed} if every convergent sequence of points of $K$
has its limit in $K$. In $\RR^d$, convergence means coordinatewise
convergence, equivalently convergence in Euclidean distance.

\begin{theorem}[Bolzano--Weierstrass]\label{real:bw}
Every bounded sequence in $\RR^d$ has a convergent subsequence.
Consequently every sequence in a closed bounded set has a
subsequence converging to a point of that set.
\end{theorem}
\begin{proof}
Enclose the sequence in a closed cube. Bisect every side,
partitioning the cube into finitely many closed smaller cubes.
At least one contains infinitely many terms. Retain it and repeat.
Choose a term with larger index at each stage from the retained cube.
The cubes are nested and their side lengths tend to zero.
For each coordinate, the left endpoints increase and the right
endpoints decrease to a common value by completeness.
The chosen terms converge to the vector of these values.
If the original set is closed, its definition puts that limit in it.
\end{proof}
The subsequence property is the form of \emph{compactness} used in
this book. In finite-dimensional real space it is equivalent to the
usual closed-and-bounded criterion. No general topology is required.

\section{Why a continuous expression attains its extrema}
\begin{theorem}[Extreme value theorem]\label{real:extreme}
A continuous real function on a nonempty closed bounded subset
$K\subseteq\RR^d$ attains both a maximum and a minimum.
\end{theorem}
\begin{proof}
First it is bounded. Otherwise choose $x_n\in K$ with
$|f(x_n)|>n$. By Theorem~\ref{real:bw}, a subsequence tends to
some $x\in K$. Continuity would make its function values tend to
the finite number $f(x)$, a contradiction.
Let $M=\sup f(K)$. For each $n$, choose $x_n\in K$ with
$f(x_n)>M-1/n$. A convergent subsequence tends to $x\in K$,
and continuity gives $f(x)=M$. Apply the same proof to $-f$
to obtain a minimum.
\end{proof}
\strategy To use this theorem, verify all three ingredients:
nonempty closed bounded domain, a real-valued expression continuous
on it, and no denominator vanishing on its boundary.
The positive triples with $a+b+c=1$ are not closed because a
coordinate can tend to zero. Their nonnegative closure is compact;
using it requires extending the expression to the boundary or treating
boundary limits separately.

\section{The complex root theorem}
The fundamental theorem of algebra says that every nonconstant
complex polynomial has a complex root, and hence factors into linear
factors with multiplicity. It is accepted as background where used
in the algebra part. Its proof belongs to a further course in complex
analysis or topology; the elementary real-number facts above do not
prove it. Polynomial division, root bounds, Vieta's formulas, and
finite-field arguments in this book are proved algebraically and
do not require a proof of that theorem.


\chapter{Hints}
\label{app:hints}

% ---- hints/foundations ----
\section{Learning to Prove}
\paragraph{Problem \ref{f:ex01}.} Replace each odd integer by $2k+1$.
\paragraph{Problem \ref{f:ex02}.} Reverse the quantifiers and negate the
equality; try a negative $n$.
\paragraph{Problem \ref{f:ex03}.} Set $t=|x|$ and remember $x^2=t^2$.
\paragraph{Problem \ref{f:ex04}.} If the sum were rational, subtract $x$.
For the product, check $x=0$ before dividing.
\paragraph{Problem \ref{f:ex05}.} Record $x\ge0$ before squaring.
\paragraph{Problem \ref{f:ex06}.} Expand $(a-b)^2$ and use $(a+b)^2=1$.
\paragraph{Problem \ref{f:ex07}.} A square is even only when its base is
even. Apply that fact to numerator and denominator.
\paragraph{Problem \ref{f:ex08}.} Compare each entry to the maximum and
sum the nonnegative differences.
\paragraph{Problem \ref{f:ex09}.} Add $n+1$ to the formula for $n$ and
factor.
\paragraph{Problem \ref{f:ex10}.} Doubling the previous bound gives
$2n+2\ge n+2$.
\paragraph{Problem \ref{f:ex11}.} A least divisor larger than $1$ cannot
have a smaller divisor larger than $1$.
\paragraph{Problem \ref{f:ex12}.} Establish three consecutive base
cases; then add $3$.
\paragraph{Problem \ref{f:ex13}.} Prove the equality of the two sets
of common divisors under one move.
\paragraph{Problem \ref{f:ex14}.} Reverse a move by subtracting the
smaller entry from the larger. Use the sum as the descent measure.
\paragraph{Problem \ref{f:ex15}.} Induct with two existing paths.
Try adding the new vertex at one endpoint of either path. If both
attempts fail, examine the edge between those endpoints.
\paragraph{Problem \ref{f:ex16}.} A common divisor of the whole list
is unchanged in both directions by the replacement. The moves only
increase positive entries.
\paragraph{Problem \ref{f:ex17}.} Subtract $2ab$.
\paragraph{Problem \ref{f:ex18}.} Use $2,3$; $\sqrt2$; and the integers
in their natural order.
\paragraph{Problem \ref{f:ex19}.} Classify the integers by their
remainders after division by $5$.
\paragraph{Problem \ref{f:ex20}.} The equality case suggests $x=1$.
Rewrite $x+1/x-2$ as a square divided by a positive number.
\paragraph{Problem \ref{f:ex21}.} Write $1/(k(k+1))=1/k-1/(k+1)$.
\paragraph{Problem \ref{f:ex22}.} Strip all powers of $2$ from each
chosen integer. How many odd parts are possible?
\paragraph{Problem \ref{f:ex23}.} Use $1+t\ge2\sqrt t$ for each
positive variable.
\paragraph{Problem \ref{f:ex24}.} Count pairs of positions whose
entries are in the wrong order. What does an adjacent exchange do
to that integer?



% ---- hints/algebra ----
\section{Algebra}

\subsection*{Expressions and equations}
\begin{description}[style=nextline]
\item[Problem~\ref{a:ex01}.] Complete the square
\((x^2+2y^2)^2\), then subtract \((2xy)^2\). For \(y=1\), check
that both resulting integer factors exceed \(1\) when \(x>1\).
\item[Problem~\ref{a:ex02}.] Exclude \(x=\pm1\), combine the fractions,
and solve the resulting quadratic. Check the candidates.
\item[Problem~\ref{a:ex03}.] A square root is nonnegative. Record \(x\ge0\)
before squaring.
\item[Problem~\ref{a:ex04}.] Expand
\((x+y+z)(x^2+y^2+z^2-xy-yz-zx)\).
\item[Problem~\ref{a:ex05}.] Write \(1+1/k=(k+1)/k\).
For the proposed false rule, test just two factors.
\item[Problem~\ref{a:ex06}.] Twice the difference is a sum of three squares.
\item[Problem~\ref{a:ex07}.] The domain is \([-1,4]\). After the first
squaring, isolate the product of the two square roots.
\item[Problem~\ref{a:ex08}.] Require a nonnegative discriminant, positive
sum of roots, and positive product of roots.
\item[Problem~\ref{a:ex09}.] Make \(x,y\) the roots of one quadratic.
\item[Problem~\ref{a:ex10}.] Build the cubic with these elementary
symmetric quantities, then factor it.
\item[Problem~\ref{a:ex11}.] Compute \(x^2+y^2+z^2\) from the two data.
\item[Problem~\ref{a:ex12}.] Put \(p=xy\), so \(x+y=8-p>0\).
Use \((x+y)^2\ge4xy\), retaining the positivity condition.
\item[Problem~\ref{a:ex13}.] Order the variables. The equation gives
\(xy\le3\); after \(x=1\), factor \((y-1)(z-1)\).
\item[Problem~\ref{a:ex14}.] Replace \(c^2\) by \((a+b)^2\).
Each denominator becomes a nonzero pairwise product times \(-2\).
\end{description}

\subsection*{Sequences and polynomials}
\begin{description}[style=nextline]
\item[Problem~\ref{a:ex15}.] The terms are \(2,5,8,\ldots\).
Pair the first and last terms for the sum.
\item[Problem~\ref{a:ex16}.] Add \(1\) to every term.
\item[Problem~\ref{a:ex17}.] Subtract consecutive terms of
\(1/[k(k+1)]\).
\item[Problem~\ref{a:ex18}.] Consecutive Fibonacci ratios obey
\(F_{m+2}/F_{m+1}=1+F_m/F_{m+1}\).
\item[Problem~\ref{a:ex19}.] Write \(t=m+u\), \(m\in\ZZ\), \(0\le u<1\).
The fractional \(u\) does not change any floor. Divide \(m\) by \(n\).
\item[Problem~\ref{a:ex20}.] Take reciprocals. Also
\(a_k-a_{k+1}=a_ka_{k+1}\).
\item[Problem~\ref{a:ex21}.] Starting from \(a,b\), compute five more
terms symbolically and cancel factors. Positivity keeps all divisions legal.
\item[Problem~\ref{a:ex22}.] The values suggest \(x^2+1\); prove uniqueness
by applying the root bound to the difference.
\item[Problem~\ref{a:ex23}.] Use
\((\alpha+\beta)^3-3\alpha\beta(\alpha+\beta)\).
\item[Problem~\ref{a:ex24}.] Divide out the three known linear factors.
The remaining factor is monic and linear.
\item[Problem~\ref{a:ex25}.] Subtract \(x^3\); a periodic polynomial is constant.
\item[Problem~\ref{a:ex26}.] Put \(t=x^2\ge0\). The quadratic in \(t\)
has its minimum either at \(0\) or at \(-a/2\).
\item[Problem~\ref{a:ex27}.] Remove the largest power of \(x\) that
divides \(P\). For the remaining factor, inspect its lowest positive exponent.
\item[Problem~\ref{a:ex28}.] Away from the roots,
\(P'/P=\sum_j1/(x-r_j)\). Differentiate this rational identity.
At a root the desired inequality is immediate.
\item[Problem~\ref{a:ex29}.] Apply the difference-of-powers factorization
to each monomial.
\item[Problem~\ref{a:ex30}.] Compare the values modulo \(2\).
\item[Problem~\ref{a:ex31}.] Try \(\binom{x}{2}\), and remember that
two consecutive negative integers still have an even product.
\item[Problem~\ref{a:ex32}.] Eisenstein with prime \(2\) applies.
A quartic can split into two quadratics without rational roots.
\item[Problem~\ref{a:ex33}.] Group the first two and last two terms.
\item[Problem~\ref{a:ex34}.] Factor \(P(x)-1\) by its three known roots.
If it took value \(1\), a product of three consecutive integers would divide \(1\).
\item[Problem~\ref{a:ex35}.] Choose \(a\) with \(c=P(a)\ne0\).
The polynomial \(Q(t)=P(a+ct)/c\) lies in \(\ZZ[t]\) and has \(Q(0)=1\).
Evaluate \(Q\) at multiples of the product of any proposed finite prime list.
\end{description}

\subsection*{Complex numbers and inequalities}
\begin{description}[style=nextline]
\item[Problem~\ref{a:ex36}.] Use either repeated squaring or modulus
\(\sqrt2\) and angle \(\pi/4\).
\item[Problem~\ref{a:ex37}.] Subtract the sum from \(\zeta\) times itself.
Use \(\sum_{k=0}^{n-1}\zeta^k=0\).
\item[Problem~\ref{a:ex38}.] Combine the first two sines, and replace
\(C\) by \(\pi-A-B\).
\item[Problem~\ref{a:ex39}.] Write the polynomial as
\((x^2+1)^2-x^2\), or use \((x^6-1)/(x^2-1)\).
\item[Problem~\ref{a:ex40}.] Take real parts of a geometric sum.
For the squares, use \(\cos^2t=(1+\cos2t)/2\).
\item[Problem~\ref{a:ex41}.] Evaluate
\((x^n-1)/(x-1)\) at \(x=1\), then use
\(|1-e^{2i\theta}|=2|\sin\theta|\).
\item[Problem~\ref{a:ex42}.] Put \(Q(x)=1+x+\cdots+x^{n-1}\).
Compute \(Q'(1)/Q(1)\) and
\((Q'(1)/Q(1))^2-Q''(1)/Q(1)\) from its factorization.
\item[Problem~\ref{a:ex43}.] Replace \(a/(b+c)\) by
\(a^2/[a(b+c)]\), then apply Cauchy.
\item[Problem~\ref{a:ex44}.] Apply Cauchy to \((x,y,z)\) and \((1,1,1)\).
\item[Problem~\ref{a:ex45}.] Swapping an inversion increases the sum.
Use equally ordered and oppositely ordered pairs.
\item[Problem~\ref{a:ex46}.] A sum of squared differences proves
the first assertion. Apply AM--GM to \(ab,bc,ca\) for the second.
\item[Problem~\ref{a:ex47}.] Apply \(u^2+v^2+w^2\ge uv+vw+wu\)
to \(u=ab,v=bc,w=ca\).
\item[Problem~\ref{a:ex48}.] The three positive denominators sum to \(2\).
\item[Problem~\ref{a:ex49}.] First get \(s=a+b+c\ge3\); then compare
\(\sum a^2\) with \(s^2/3\).
\item[Problem~\ref{a:ex50}.] Take logarithms of weighted AM--GM.
\item[Problem~\ref{a:ex51}.] Write \(a_i=(a_i/b_i^{2/3})b_i^{2/3}\)
and use Hölder with exponents \(3,3/2\).
\item[Problem~\ref{a:ex52}.] Square both nonnegative sides and use
\(xu+yv\le\sqrt{x^2+y^2}\sqrt{u^2+v^2}\).
\item[Problem~\ref{a:ex53}.] Apply weighted AM--GM to \(x^r,a^r\)
with weights \(1/r,(r-1)/r\).
\item[Problem~\ref{a:ex54}.] The expression is linear in \(r=abc\)
for fixed \(p,q\). With \(a=b=t,c=s\), it becomes \(s(t-s)^2\).
\item[Problem~\ref{a:ex55}.] Use Cauchy with numerator \(1/a\) and
denominator \(a(b+c)\). Under \(abc=1\), \(\sum1/a=ab+bc+ca\).
\item[Problem~\ref{a:ex56}.] Order \(x\ge y\ge z\), then combine the
\(x\)-term and \(y\)-term before considering signs.
\end{description}

\subsection*{Functional equations and synthesis}
\begin{description}[style=nextline]
\item[Problem~\ref{a:ex57}.] First determine rational inputs; monotonicity
squeezes every real input between them.
\item[Problem~\ref{a:ex58}.] Subtract the known quadratic \(n^2\).
\item[Problem~\ref{a:ex59}.] Add \(1\) to the output, not to the input.
\item[Problem~\ref{a:ex60}.] Put \(g(x)=f(x)-x^2-f(0)\).
Show \(g(2x)=2g(x)\), then replace \(x,y\) by
\((u+v)/2,(u-v)/2\).
\item[Problem~\ref{a:ex61}.] Use \(A(t)=\log f(e^t)\).
Its monotonicity follows from that of \(f\).
\item[Problem~\ref{a:ex62}.] The case \(x=0\) gives surjectivity and
injectivity. Compare the period from \(y=0\) with injectivity to find \(f(0)\).
\item[Problem~\ref{a:ex63}.] Separate \(f=0\). Otherwise prove positivity
before taking logarithms.
\item[Problem~\ref{a:ex64}.] Evaluate the third iterate in two ways.
\item[Problem~\ref{a:ex65}.] Every nonfixed element belongs to a pair.
\item[Problem~\ref{a:ex66}.] Evaluate \(f(x+y+z)\) in two groupings.
\item[Problem~\ref{a:ex67}.] Compare degrees first, then leading coefficients.
\item[Problem~\ref{a:ex68}.] Let \(c=f(0)\). Obtain
\(f(f(y))=y+c^2\) and surjectivity. Compare two nested substitutions
to get \(f(v+f(u)^2)=f(v)+u^2+c^2\). Use an input with zero output.
\item[Problem~\ref{a:ex69}.] Use a zero of \(f\) to prove \(f(0)=0\).
Then compare the translated equation for \(x\) and for \(f(x)\);
it gives \(f(x)^2=x^2\). Rule out mixing the two signs.
\item[Problem~\ref{a:ex70}.] First find \(f(0)\) and \(f(-n)\).
Set \(m=n\) to avoid confusion only if helpful; the simpler recurrence
comes from setting the second variable to \(1\).
\item[Problem~\ref{a:ex71}.] Order \(x\le y\le z\). For the quadratic
in \(z\), evaluate it at \(y\). Treat \(x=y=1\) separately.
\item[Problem~\ref{a:ex72}.] For large \(n\), the integer
\(P(n+1)/P(n)\) is positive and approaches \(1\).
\item[Problem~\ref{a:ex73}.] For the first inequality, use a tangent
bound up to \(7/3\), then a crude reciprocal estimate outside that region.
For the second, use the justified \(p,q,r\) reduction.
\item[Problem~\ref{a:ex74}.] Suppose the monic polynomial is \(GH\),
with \(\deg G\le n/2\). At every \(a_j\), \(G(a_j)=\pm1\).
Compare \(G^2-1\) with \(\prod(x-a_j)\).
\item[Problem~\ref{a:ex75}.] Under sum zero,
\(\sum x^4=(\sum x^2)^2/2\). Arrange the sign so two variables
are nonnegative; put \(s=x+y,t=xy\).
\item[Problem~\ref{a:ex76}.] If \(s=x+y+z>3\), scale the triple by \(3/s\).
Each nonzero term of the constraint strictly decreases.
\item[Problem~\ref{a:ex77}.] Define \(C_n(x)=2T_n(x/2)\).
It satisfies the equation. Subtract the candidate of the same degree:
if the difference has degree \(k<n\), compare \(n+k\) with \(2k\).
\end{description}



% ---- hints/combinatorics ----
\section{Combinatorics}
Use a hint to restart your own reasoning. In particular, try to prove
that an encoding is reversible before reading a counting formula.

\begin{description}[style=nextline]
\item[\ref{c:ex01}] Count choices for successive images. Injection
prohibits reusing an image.
\item[\ref{c:ex02}] Label repeated letters temporarily. In the second
question the two \(C\)'s have already been placed.
\item[\ref{c:ex03}] Subtract the required lower bounds before placing bars.
\item[\ref{c:ex04}] If the selected positions are \(a_1<\cdots<a_k\),
compress the gaps by replacing \(a_i\) with \(a_i-i+1\).
\item[\ref{c:ex05}] Add an imaginary zero at both ends; choose the
\(2k\) boundaries at which a run starts or ends.
\item[\ref{c:ex06}] Fix one person to remove rotation. For \(n\ge3\),
treat the specified adjacent pair as a block. Check \(n=2\) separately.
\item[\ref{c:ex07}] Each element is outside \(B\), in \(B\setminus A\),
or in \(A\). For the sum, fix \(B\).
\item[\ref{c:ex08}] First enforce the union condition, then remove the
pairs in which no element belongs to both sets.
\item[\ref{c:ex09}] Mark one element, or an ordered pair of different
elements, inside a subset.
\item[\ref{c:ex10}] Omit one or two target values. The recurrence gives
\(S(5,3)=25\).
\item[\ref{c:ex11}] A bad coordinate has at least five units. Two bad
coordinates leave total zero.
\item[\ref{c:ex12}] Select \(n\) elements from two \(n\)-element parts.
For the weighted sum, mark one selected element in the first part.
\item[\ref{c:ex13}] Derive \(D_n=(n-1)(D_{n-1}+D_{n-2})\), with
\(D_0=1,D_1=0\), or reduce the inclusion--exclusion formula modulo two.
\item[\ref{c:ex14}] Count incidences between triples and their contained pairs.
\item[\ref{c:ex15}] Choose the represented couples, then one spouse from
each. For inclusion--exclusion, prescribe complete couples.
\item[\ref{c:ex16}] Pair the two rows into columns. Swap the selected
entry in the first column containing exactly one selected entry.
\end{description}

\subsection*{Recurrences and pigeonholes}
\begin{description}[style=nextline]
\item[\ref{c:ex17}] Inspect the first part. The empty composition counts once.
\item[\ref{c:ex18}] A legal word ends in \(0\), \(01\), or \(011\).
\item[\ref{c:ex19}] Match each closing parenthesis to the most recent
unmatched opening parenthesis. Negative height is exactly failure.
\item[\ref{c:ex20}] Translate the first-return decomposition into multiplication.
\item[\ref{c:ex21}] Use one factor for each variable and list the possible \(c\).
\item[\ref{c:ex22}] Reflect the prefix through the first visit to height \(-1\).
How many down-steps does the reflected word have?
\item[\ref{c:ex23}] Count pairs consisting of a legal coloring and one
marked one. Read the positive zero-gaps cyclically after the mark.
\item[\ref{c:ex24}] In each odd part \(u\), expand its multiplicity in base two.
\item[\ref{c:ex25}] Use the five residue classes. For sharpness, take
one representative from each.
\item[\ref{c:ex26}] Record both coordinate parities.
\item[\ref{c:ex27}] Include the empty prefix \(s_0=0\).
\item[\ref{c:ex28}] Use the pairs \(\{1,2\},\{3,4\},\ldots\).
\item[\ref{c:ex29}] Use four decreasing blocks of length four, with
successive blocks having larger values.
\item[\ref{c:ex30}] Start immediately after a prefix of minimum sum.
\item[\ref{c:ex31}] Use \(M^d\) boxes for vectors of fractional parts,
and \(M^d+1\) vectors.
\item[\ref{c:ex32}] First obtain \(|q\alpha-p|<1/N\), \(q\le N\).
If denominators stayed bounded, the nonzero distances to integers
would have a positive minimum.
\end{description}
\subsection*{Invariants and extremal choices}
\begin{description}[style=nextline]
\item[\ref{c:ex33}] Compare the number of black and white squares.
Necessary conditions need not be sufficient.
\item[\ref{c:ex34}] Parity gives necessity. For sufficiency, correct the
coins from left to right without changing an already corrected coin.
\item[\ref{c:ex35}] Track \(\prod(1+x)\).
\item[\ref{c:ex36}] The gcd of all entries survives adding or subtracting one.
\item[\ref{c:ex37}] A knight changes the parity of the coordinate sum.
\item[\ref{c:ex38}] Each adjacent swap changes one inversion.
\item[\ref{c:ex39}] Color square \((i,j)\) by \(i+j\pmod3\).
Count the three classes before deleting the corner.
\item[\ref{c:ex40}] Give square \((i,j)\) weight \(\omega^{i+j}\).
Factor the total weight of the rectangle.
\item[\ref{c:ex41}] Maximize the number of edges between the parts.
\item[\ref{c:ex42}] Choose a vertex with greatest outdegree. What happens
if it cannot reach a vertex in two steps?
\item[\ref{c:ex43}] Minimize the positive distance from a point to a
line through two other points. If that line contains three points,
choose two on the same side of the perpendicular foot.
\item[\ref{c:ex44}] Add the edge to a minimum tree, then remove a different
edge of the resulting cycle that crosses the same partition.
\item[\ref{c:ex45}] Use one long interval meeting two short disjoint intervals.
\item[\ref{c:ex46}] Minimize the perimeter over all cyclic orders.
Uncrossing must reconnect the two remaining paths into one cycle.
\item[\ref{c:ex47}] The sum of squares of the occupancies strictly increases
and is at most \(N^2\).
\item[\ref{c:ex48}] Take a longest directed cycle. An outside vertex
can be inserted unless all its arrows to the cycle point the same way.
Use strong connectivity to bridge the two remaining types.
\end{description}

\subsection*{Graphs, matching, and order}
\begin{description}[style=nextline]
\item[\ref{c:ex49}] Degrees lie between zero and \(n-1\), but those
two extremes cannot both occur.
\item[\ref{c:ex50}] Examine the endpoints of a longest path.
\item[\ref{c:ex51}] Add one new vertex adjacent to all odd-degree vertices,
find an Euler circuit, and remove the new vertex.
\item[\ref{c:ex52}] Maximize a cut, or color vertices independently with
two equally likely colors.
\item[\ref{c:ex53}] A subgraph of average degree less than \(2d\)
has a vertex of degree at most \(2d-1\).
\item[\ref{c:ex54}] Apply \(e\le3n-6\), then remove a low-degree vertex.
\item[\ref{c:ex55}] For a longest path \(v_1,\ldots,v_k\), compare
the index sets where \(v_1v_{i+1}\) and \(v_iv_k\) are edges.
\item[\ref{c:ex56}] A transitive triple has one vertex beating the other two.
For the maximum, balance the outdegrees.
\item[\ref{c:ex57}] In the impossible version let all three lists be
contained in \(\{1,2\}\).
\item[\ref{c:ex58}] Remove a perfect matching and induct on \(d\).
\item[\ref{c:ex59}] Separate Hall subsets of size at most \(n/2\) from
larger subsets. In the latter case inspect a vertex outside \(N(S)\).
\item[\ref{c:ex60}] Make rows the left part and columns the right part.
\item[\ref{c:ex61}] Use odd parts for a chain partition and the upper
half of \([2n]\) for an antichain.
\item[\ref{c:ex62}] Each permutation chain meets the family at most \(r\)
times. Then optimize the resulting weighted bound subject to the
capacity of each level.
\item[\ref{c:ex63}] Give an element the length of the longest chain ending
there. Equal labels form antichains.
\item[\ref{c:ex64}] Suppose the first rejection of any pair belonging to
a stable matching occurs. Compare the preferred proposer with his partner
in that matching.
\end{description}

\subsection*{Extremal, probabilistic, and algebraic methods; games}
\begin{description}[style=nextline]
\item[\ref{c:ex65}] Apply Mantel and use a balanced complete bipartite graph.
\item[\ref{c:ex66}] Apply the proved Turan bound with \(r=3\), and
use three equal independent parts.
\item[\ref{c:ex67}] Use the Ramsey recurrence for the upper bound.
In an eight-cycle, an independent set of size four must alternate.
\item[\ref{c:ex68}] Order the vertices randomly. Keep a vertex if it
comes before all its neighbors.
\item[\ref{c:ex69}] Count pairs of ones in a column by the pair of rows
they occupy, then use Cauchy.
\item[\ref{c:ex70}] Color each edge randomly. Count monochromatic
\(k\)-subsets in expectation, retaining the factorial in \(\binom nk\).
\item[\ref{c:ex71}] Each bad event has probability \(1/512\).
Try \(x=1/101\) in the local lemma and use \((1+1/100)^{100}<3\).
\item[\ref{c:ex72}] A clause fails with probability \(1/8\).
After partial assignment, its success probability is directly computable.
\item[\ref{c:ex73}] A nonidentity rotation on \(p\) positions is one cycle.
\item[\ref{c:ex74}] Reuse the square's cycle types, replacing \(2^j\)
by \(3^j\).
\item[\ref{c:ex75}] There are three reflections fixing two opposite
beads and three fixing no bead.
\item[\ref{c:ex76}] Compute \(g_0,\ldots,g_{10}\), then induct on
the recurrence \(g_n=\operatorname{mex}\{g_{n-1},g_{n-3},g_{n-4}\}\).
\item[\ref{c:ex77}] First classify positions whose heaps all have size
zero or one. If exactly one heap is larger, adjust it to leave an
odd number of ones.
\item[\ref{c:ex78}] Remove only the top-right square first. Any later
legal move would already have removed that corner in the original board.
\item[\ref{c:ex79}] Show that the one-heap Grundy value is \(n\bmod3\).
\item[\ref{c:ex80}] The two Beatty sequences partition positive integers,
and their \(k\)th terms differ by \(k\). Match the difference of the
two heaps, or use the sequence containing the smaller heap.
\item[\ref{c:ex81}] Apply Oddtown; singletons attain equality.
\item[\ref{c:ex82}] Pair ten elements and leave the eleventh unused.
\item[\ref{c:ex83}] In residues modulo eight, try \(A=B=\{0,4\}\).
\item[\ref{c:ex84}] Apply Cauchy--Davenport to \((k-1)A+A\).
\item[\ref{c:ex85}] The path matrix is
\(\left(\begin{smallmatrix}n+1&\binom{n+2}{2}\\1&n+1\end{smallmatrix}\right)\).
\item[\ref{c:ex86}] Use all \(k\)-subsets of a \(2k\)-set with their complements.
\item[\ref{c:ex87}] No four points can be shattered. Separate the
convex-quadrilateral case from the case of one point inside a triangle.
\item[\ref{c:ex88}] In the small-sumset case use
\((x-y)\prod_{c\in C}(x+y-c)\), with \(|C|=2a-4\),
on a grid of sizes \(a,a-1\). For large \(a\), use an intersection count.
\end{description}




% ---- hints/convex-lattice ----
\subsection*{Combinatorics in the Plane}
\paragraph{Problem \ref{l:ex01}.} The displacement is $(12,8)$.
Use its gcd to find the primitive direction.
\paragraph{Problem \ref{l:ex02}.} Count the interior rows and columns;
avoid counting corner points twice on the boundary.
\paragraph{Problem \ref{l:ex03}.} There are four coordinate-parity classes.
\paragraph{Problem \ref{l:ex04}.} The edge gcds give the boundary count.
\paragraph{Problem \ref{l:ex05}.} Use the largest left endpoint and
the smallest right endpoint.
\paragraph{Problem \ref{l:ex06}.} With more than three positive weights,
perturb them along an affine dependence until one becomes zero.
\paragraph{Problem \ref{l:ex07}.} Here $B=n$ and $I=0$.
\paragraph{Problem \ref{l:ex08}.} Cover each triple first, then use
Helly on the possible centers. For an acute triangle, use its largest
angle in the sine-law formula for its circumradius.



% ---- hints/geometry ----
\section{Geometry}

These hints indicate a first productive move. If a hint unlocks a
problem, close this appendix and finish the argument before consulting
its solution.

\begin{description}[style=nextline,leftmargin=2.2em]
\item[\ref{g:ex01}] Compare $ABC$ with $ACB$. For the converse, use
ASA with their common side $BC$.
\item[\ref{g:ex02}] Draw a diagonal and use alternate angles. For the
midpoint statement, compare the two triangles on opposite sides of
the diagonal intersection.
\item[\ref{g:ex03}] The parallel through the midpoint $M$ meets $AC$
at some $N'$. Its scale factor is already known.
\item[\ref{g:ex04}] Since $MA=MB=MC$, use the two isosceles triangles
with apex $M$ and the angle sum of $ABC$.
\item[\ref{g:ex05}] Join one diagonal of the original quadrilateral.
The midpoint theorem gives both parallelism and half lengths.
\item[\ref{g:ex06}] Reflect $B$ in $\ell$, then use a straight segment.
\item[\ref{g:ex07}] A common scale multiplies both base and altitude.
\item[\ref{g:ex08}] Compare $[ABD]/[ACD]$ once using bases on $BC$
and once using the sine area formula at $A$.
\item[\ref{g:ex09}] Compare triangles $PAC,PDB$, with the cyclic
angles choosing the correspondence.
\item[\ref{g:ex10}] Both feet see $BC$ under a right angle. A repeated
foot should be treated as one point, not as an undefined angle.
\item[\ref{g:ex11}] Two of the small quadrilaterals have a diameter
joining an original vertex to $H$.
\item[\ref{g:ex12}] For the acute picture, chase $\angle BH'C$.
For every picture, put $BC$ on the horizontal axis.
\item[\ref{g:ex13}] Multiply the three angle-bisector ratios.
For the second proof, distances from a point on two bisectors agree
for all three side lines.
\item[\ref{g:ex14}] The missing ratio is the reciprocal of the product
of the known ratios.
\item[\ref{g:ex15}] Menelaus requires an actual finite side-line
intersection in its finite version.
\item[\ref{g:ex16}] The product in Menelaus is negative.
\item[\ref{g:ex17}] Affine coordinates let you replace $ABC$ by
$A=(0,0)$, $B=(1,0)$, $C=(0,1)$ without changing area ratios.
The three intersections are cyclic permutations of one barycentric triple.
\item[\ref{g:ex18}] Use normalized barycentric coordinates or compare
the area $[PBC]$ with $[ABC]$.
\item[\ref{g:ex19}] The midpoint segment is parallel to the third
side; compare triangles with the common median intersection.
\item[\ref{g:ex20}] Give the three pairs of equal tangent segments
three temporary variables.
\item[\ref{g:ex21}] On an extended side, a tangent length is a sum,
not a difference. Then use Ceva.
\item[\ref{g:ex22}] Three known points determine the medial
triangle's circumcircle; all lie on the nine-point circle.
\item[\ref{g:ex23}] With circumcenter at the origin, the midpoint
reflection of $h=a+b+c$ in the midpoint of $BC$ is $-a$.
\item[\ref{g:ex24}] Use Problem~\ref{g:ex11}. In an obtuse triangle
track whether the altitude line bisects an internal or external angle.
\item[\ref{g:ex25}] The cosine law isolates $\cos A$ with a positive
coefficient.
\item[\ref{g:ex26}] Insert $BD=ac/(b+c)$ and $DC=ab/(b+c)$
into Stewart's theorem.
\item[\ref{g:ex27}] Add the three median identities.
\item[\ref{g:ex28}] Split the quadrilateral along a diagonal.
The opposite included angles have equal sines and opposite cosines.
\item[\ref{g:ex29}] Rotate $P$ through $60^\circ$ about $A$, sending
$B$ to $C$. The resulting triangle has side lengths $3,4,5$.
The interior condition determines which orientation to use.
\item[\ref{g:ex30}] For a general quadrilateral, the same calculation
as in the cyclic case leaves a nonnegative correction involving
$\cos^2((A+C)/2)$.
\item[\ref{g:ex31}] The radius is perpendicular to each tangent;
apply Pythagoras to the two right triangles.
\item[\ref{g:ex32}] Put $P$ at the origin and the two perpendicular
chords on the coordinate axes. Use root sums and products.
\item[\ref{g:ex33}] The auxiliary common-chord lines meet at a point
on the unknown radical axis. Its perpendicular direction is already
known from the two centers.
\item[\ref{g:ex34}] At an intersection, the two radii are perpendicular.
For the converse compare $|r_1-r_2|$, center distance, and $r_1+r_2$.
\item[\ref{g:ex35}] Equal powers are transitive. If two axes are
parallel, their normals constrain the centers.
\item[\ref{g:ex36}] Use the internal and external angle bisectors
at $B$ and the fact that $A,I,L,I_A$ are collinear.
\item[\ref{g:ex37}] Reflection keeps the parallel component and
negates the perpendicular component.
\item[\ref{g:ex38}] Write the maps as $x\mapsto kx+(1-k)S$ and
$x\mapsto\ell x+(1-\ell)T$.
\item[\ref{g:ex39}] Solve $f(z)=uz+v$ at the two source points.
\item[\ref{g:ex40}] The centers lie on opposite rays from their
contact point.
\item[\ref{g:ex41}] Rotate as in the worked example. Equality requires
the image point to lie on the segment joining the other two.
\item[\ref{g:ex42}] Use complex coordinates and let
$\omega=e^{i\pi/3}$. Write the center of each exterior equilateral
triangle as the midpoint of its side plus the correctly oriented
perpendicular displacement.
\item[\ref{g:ex43}] The median's side ratio is $1$, so its sine ratio
is $AC/AB$. Reflection reciprocates it.
\item[\ref{g:ex44}] Reflect the three median lines.
\item[\ref{g:ex45}] Compare the angle-bisector ratio with the median
ratio. One equality already forces the adjacent sides equal.
\item[\ref{g:ex46}] On $BC$, a barycentric triple $(0:y:z)$ gives
the ratio $z/y$.
\item[\ref{g:ex47}] Reflection replaces each sine ratio by its reciprocal.
\item[\ref{g:ex48}] Combine the tangent symmedian ratio with the
inscribed-angle interpretation of a chord ratio.
\item[\ref{g:ex49}] Compute $(b+c)\cdot(b-c)$.
\item[\ref{g:ex50}] The proposed reflection is $a+b-ab\bar p$.
Check that its midpoint with $p$ lies on the chord.
\item[\ref{g:ex51}] Normalize $x+y+z=1$ and write $P=xA+(1-x)D$.
\item[\ref{g:ex52}] A concurrency condition is the vanishing of one
cyclic dot-product sum. Reverse the subscripts.
\item[\ref{g:ex53}] In unit-circle complex coordinates, a direction
of the Simson line is $(a-b)(p-c)/p$.
\item[\ref{g:ex54}] Each reflected point is twice its foot minus $P$.
This is one homothety, not three unrelated reflections.
\item[\ref{g:ex55}] Substitute $x=x'/(x'^2+y'^2)$ and
$y=y'/(x'^2+y'^2)$ in each equation.
\item[\ref{g:ex56}] Compare the coefficients of the original and
transformed circle equations.
\item[\ref{g:ex57}] Draw the perpendicular to $AI$ at $I$; its
intersections with the side rays are the intended contacts.
Intersect their perpendiculars to the sides to find $J$.
\item[\ref{g:ex58}] Invert at $A$. Equality means the image of $C$
lies between the images of $B,D$.
\item[\ref{g:ex59}] A circle tangent at the inversion center has a
diameter on the common normal. Its opposite endpoint locates the
image line.
\item[\ref{g:ex60}] Put $P_i=Ru_i$, $|u_i|=1$. Show
$t_{ij}=\sqrt{(R-r_i)(R-r_j)}\,P_iP_j/R$, then use Ptolemy.
\item[\ref{g:ex61}] Factor the difference between two fractional
linear images before multiplying the four differences.
\item[\ref{g:ex62}] Compare the line equation with $p x+q y=R^2$.
\item[\ref{g:ex63}] Use barycentrics with $D=(0:d:e)$ and
$E=(u:d:e)$ after an appropriate scaling.
\item[\ref{g:ex64}] Order the six points according to the desired
opposite-side intersections. An ideal intersection fixes the
direction of the Pascal line.
\item[\ref{g:ex65}] Apply Pascal to the six contact points and use
pole--polar reciprocity to dualize the conclusion.
\item[\ref{g:ex66}] Send the ellipse to the unit circle by
$(x,y)\mapsto(x/a,y/b)$.
\item[\ref{g:ex67}] Let $x=a^2,y=b^2,z=c^2$ and square both
nonnegative sides. The residual is a sum of squared differences.
\item[\ref{g:ex68}] On the unit circle, a reflection of $p$ in $BC$
is $b+c-bc/p$ and $h=a+b+c$. Prove the relevant quotient real.
\item[\ref{g:ex69}] Scale $PQ$ to endpoints $(-1,0),(1,0)$.
The circle has equation $x^2+y^2-2ky=1$.
Opposite points on a chord through the origin have product of
signed distances $-1$.
\item[\ref{g:ex70}] A nearest point of the closed triangle is closer
to every vertex than an exterior point is. Then use the Fermat
point argument or the $120^\circ$ vertex argument.
\item[\ref{g:ex71}] Orthogonal projection onto a line is an affine map
whose linear part is symmetric. Use the cyclic concurrency sum
from Problem~\ref{g:ex52}.
\item[\ref{g:ex72}] Reflect $X$ in $AB$ and $AC$. A broken path
becomes a path between the two reflected points, whose distance
is $2AX\sin A$. Then minimize $AX$ on $BC$.
\end{description}



% ---- hints/number-theory ----
\section{Number Theory}
\subsection*{Divisibility, the Euclidean Algorithm, and Structure}
\paragraph{Problem \ref{n:ex01}.} Keep the remainders, then substitute backward.
\paragraph{Problem \ref{n:ex02}.} A common divisor divides their difference.
\paragraph{Problem \ref{n:ex03}.} Subtract $n(n+1)$ from the first entry.
\paragraph{Problem \ref{n:ex04}.} Divide by the gcd, find one solution, and use Theorem~\ref{n:linear-dio}.
\paragraph{Problem \ref{n:ex05}.} Write $A=mk$ and cancel $m$ from $n\mid mk$.
\paragraph{Problem \ref{n:ex06}.} Set $a=6x,b=6y$. Determine $xy$ and require $\gcd(x,y)=1$.
\paragraph{Problem \ref{n:ex07}.} They cannot both be representable. For existence, take the representative $0\le y<a$ with $by\equiv N\pmod a$ and examine the sign of $(N-by)/a$.
\paragraph{Problem \ref{n:ex08}.} A common divisor divides $2b^2$, but is coprime to $b$.
\subsection*{Primes, Factorization, and Arithmetic Functions}
\paragraph{Problem \ref{n:ex09}.} Use $72=2^3\,3^2$.
\paragraph{Problem \ref{n:ex10}.} Use the divisor-count formula, or pair $d$ with $n/d$.
\paragraph{Problem \ref{n:ex11}.} Use one factorial divisible by every integer from $2$ to $k+1$.
\paragraph{Problem \ref{n:ex12}.} The factors $e_i+1$ have product $6$.
\paragraph{Problem \ref{n:ex13}.} Assume there are finitely many, and consider four times their product minus one. Track prime factors modulo $4$.
\paragraph{Problem \ref{n:ex14}.} A divisor of a coprime product uniquely splits into one divisor of each factor.
\paragraph{Problem \ref{n:ex15}.} Apply Mobius inversion to $F(n)=n$.
\paragraph{Problem \ref{n:ex16}.} Multiply the identity $d(n/d)=n$ over all divisors.
\subsection*{Congruences, Inverses, and the Chinese Remainder Theorem}
\paragraph{Problem \ref{n:ex17}.} Divide by $6$, including the modulus.
\paragraph{Problem \ref{n:ex18}.} Write $x=5+12t$ and reduce the second congruence.
\paragraph{Problem \ref{n:ex19}.} Compare parity, or apply the compatibility condition.
\paragraph{Problem \ref{n:ex20}.} Separate odd integers from the two types of even integers.
\paragraph{Problem \ref{n:ex21}.} Look for a short cycle modulo $100$.
\paragraph{Problem \ref{n:ex22}.} Substitute $B\equiv1$ and $B\equiv-1$ into the digit polynomial.
\paragraph{Problem \ref{n:ex23}.} Choose $k$ distinct primes and prescribe one square-divisibility condition for each position.
\paragraph{Problem \ref{n:ex24}.} The transition $(u,v)\mapsto(v,u+v)$ is invertible modulo $m$.
\subsection*{Powers, Orders, and Primitive Roots}
\paragraph{Problem \ref{n:ex25}.} Compute the first few powers.
\paragraph{Problem \ref{n:ex26}.} Reduce the congruence modulo a proper prime divisor if $n$ is composite.
\paragraph{Problem \ref{n:ex27}.} The order of $2$ modulo $q$ is a divisor of the prime $p$.
\paragraph{Problem \ref{n:ex28}.} Write $x$ as a power of a primitive root.
\paragraph{Problem \ref{n:ex29}.} Consider $(2p_1\cdots p_k)^2+1$ and the order of $2p_1\cdots p_k$ modulo a prime divisor.
\paragraph{Problem \ref{n:ex30}.} Separate the case $p\mid a$ from the invertible case.
\paragraph{Problem \ref{n:ex31}.} Use a primitive root and sum a geometric progression; handle ratio $1$ separately.
\paragraph{Problem \ref{n:ex32}.} The order divides $2^{n+1}$ but not $2^n$.
\subsection*{Valuations, Factorials, and Lifting Exponents}
\paragraph{Problem \ref{n:ex33}.} Use Legendre with $p=2$ and $p=5$.
\paragraph{Problem \ref{n:ex34}.} Replace the base $2$ by $4$.
\paragraph{Problem \ref{n:ex35}.} The exponent is even, so both $3-1$ and $3+1$ matter.
\paragraph{Problem \ref{n:ex36}.} Use Legendre's formula or count the single carry.
\paragraph{Problem \ref{n:ex37}.} The base-$7$ expansions are $100=(202)_7$ and $50=(101)_7$.
\paragraph{Problem \ref{n:ex38}.} Count the one digits in the binary expansion of $2026$.
\paragraph{Problem \ref{n:ex39}.} Apply LTE, then characterize which positive exponents yield valuation at least $s$.
\paragraph{Problem \ref{n:ex40}.} Use $\binom{p^r}{k}=(p^r/k)\binom{p^r-1}{k-1}$ and Lucas on the second factor.
\subsection*{Diophantine Equations: Factor, Bound, and Parametrize}
\paragraph{Problem \ref{n:ex41}.} Use positive odd factors $(x-y)(x+y)=2025$, ordered by size.
\paragraph{Problem \ref{n:ex42}.} Complete $(x-1)(y-1)$ and check the signs.
\paragraph{Problem \ref{n:ex43}.} Add and subtract $4n^2$.
\paragraph{Problem \ref{n:ex44}.} Solve $m^2+n^2=65$ with $m>n>0$ and the primitive parameter conditions.
\paragraph{Problem \ref{n:ex45}.} Squares modulo $3$ force both $x,y$ to be divisible by $3$.
\paragraph{Problem \ref{n:ex46}.} Use $(x-6)(y-6)=36$.
\paragraph{Problem \ref{n:ex47}.} Divide out the largest common power of $2$, then inspect the odd entries modulo $4$.
\paragraph{Problem \ref{n:ex48}.} Any common divisor divides $b$ and $b^2+1$. Then combine two coprime divisibilities.
\subsection*{Descent and Vieta Jumping}
\paragraph{Problem \ref{n:ex49}.} Use the recurrence $u_{k+1}=3u_k-u_{k-1}$.
\paragraph{Problem \ref{n:ex50}.} A common divisor divides all three terms except the constant $1$.
\paragraph{Problem \ref{n:ex51}.} Use absolute values and the same division by $3$; if one entry is zero, do not assume positivity.
\paragraph{Problem \ref{n:ex52}.} Fix the integer quotient and descend using the other root. A minimal pair must have its smaller entry $1$.
\paragraph{Problem \ref{n:ex53}.} Construct a right triangle whose area is a square and use Theorem~\ref{n:fermat-area}.
\paragraph{Problem \ref{n:ex54}.} Combine Problem~\ref{n:ex48} with Problem~\ref{n:ex52} and the chain theorem.
\paragraph{Problem \ref{n:ex55}.} For a proposed quotient $t^2$, test $(a,b)=(t,t^3)$.
\paragraph{Problem \ref{n:ex56}.} Zero is not allowed. The other root may be integral without being positive.
\subsection*{Quadratic Residues and Sums of Two Squares}
\paragraph{Problem \ref{n:ex57}.} Square $1,\ldots,5$; negatives give the same squares.
\paragraph{Problem \ref{n:ex58}.} Use reciprocity with $3$ and the sign from $p$ modulo $4$.
\paragraph{Problem \ref{n:ex59}.} First show that $q$ divides each coordinate.
\paragraph{Problem \ref{n:ex60}.} Remove two common factors of $3$, then solve a sum of squares equal to $25$.
\paragraph{Problem \ref{n:ex61}.} Examine the prime moduli separately.
\paragraph{Problem \ref{n:ex62}.} Subtract twice $5+i$ from $13$.
\paragraph{Problem \ref{n:ex63}.} Clear denominators and use either a valuation or descent.
\paragraph{Problem \ref{n:ex64}.} In Gaussian integers, a factor of norm a prime is irreducible. Use Bezout to obtain its prime-divisor property.
\subsection*{Approximation, Continued Fractions, and Pell Equations}
\paragraph{Problem \ref{n:ex65}.} Work from the innermost denominator, or use the convergent recurrences.
\paragraph{Problem \ref{n:ex66}.} Use powers of the fundamental unit $3+2\sqrt2$.
\paragraph{Problem \ref{n:ex67}.} Divide $x+y\sqrt2$ by $1+\sqrt2$ and use the norm-$1$ classification.
\paragraph{Problem \ref{n:ex68}.} For finitely many denominators, use their positive distances to integers.
\paragraph{Problem \ref{n:ex69}.} The pair $(2,1)$ is already fundamental.
\paragraph{Problem \ref{n:ex70}.} Use the determinant with the previous convergent.
\paragraph{Problem \ref{n:ex71}.} If $0<p/q\le2$, use the nonzero integer $|p^2-2q^2|$. Treat the outside range directly.
\paragraph{Problem \ref{n:ex72}.} The equation $t^2+(t+1)^2=h^2$ becomes a negative Pell equation.
\subsection*{Prime Powers, Integer Polynomials, and Further Constructions}
\paragraph{Problem \ref{n:ex73}.} Start with the two roots modulo $7$ and lift them uniquely.
\paragraph{Problem \ref{n:ex74}.} Every root modulo $p$ is zero.
\paragraph{Problem \ref{n:ex75}.} For odd $x$, the even neighbors $x-1,x+1$ have gcd $2$.
\paragraph{Problem \ref{n:ex76}.} Check the order modulo $7$ and the valuation of $3^6-1$.
\paragraph{Problem \ref{n:ex77}.} Choose an input with absolute value greater than $1$ and revisit its residue class modulo that value.
\paragraph{Problem \ref{n:ex78}.} Factor the geometric series when $n=ab$ with $a,b>1$.
\paragraph{Problem \ref{n:ex79}.} If all prime divisors were in a finite list, evaluate at multiples of their product.
\paragraph{Problem \ref{n:ex80}.} Use $(a+1)(a^2-a+1)=a^3+1$, then rule out order $2$.



% ---- hints/capstones ----
\section{Mixed Problems}
\paragraph{Problem \ref{m:ex01}.} Compare degrees, then the leading
coefficient. For a quadratic, compare $P(P(x))$ to $P(x)^2$ directly.
\paragraph{Problem \ref{m:ex02}.} First find $f(0)$. The square
increment can be any nonnegative real. Extend to negative increments,
then compare $f((x+1)^2)$ in two ways.
\paragraph{Problem \ref{m:ex03}.} Let $p$ be the least prime divisor
of $n$ and study the order of $2$ modulo $p$.
\paragraph{Problem \ref{m:ex04}.} Count pairs of same-color entries
within a row, recording the two column positions and the color.
\paragraph{Problem \ref{m:ex05}.} For the converse, intersect ray
$PD$ with the circle through $A,B,C$ and use power.
\paragraph{Problem \ref{m:ex06}.} Use lines through $(-1,0)$.
\paragraph{Problem \ref{m:ex07}.} Cauchy in quotient form uses
denominator sum $(1-a)+(1-b)+(1-c)=2$.
\paragraph{Problem \ref{m:ex08}.} A polynomial of degree at least two
eventually exceeds $n$ in absolute value. For $an+b$, the divisibility
also makes $an+b$ divide $b$.
\paragraph{Problem \ref{m:ex09}.} A vertex with four red neighbors
settles the first case. Otherwise parity gives a vertex with at most
two red neighbors. For sharpness, use an eight-cycle plus opposite chords.
\paragraph{Problem \ref{m:ex10}.} Fix the quotient $k$ and minimize
$a+b$. Evaluate the quadratic in the second variable at its smaller
entry to locate its other root.
\paragraph{Problem \ref{m:ex11}.} Minimize a positive distance from
a point of $S$ to a line determined by two others. Three points on
that line would give a smaller positive altitude.
\paragraph{Problem \ref{m:ex12}.} Minimize total matching length,
then uncross two crossing segments.
\paragraph{Problem \ref{m:ex13}.} Clear all side-length denominators.
The new integer triangle would have square area.
\paragraph{Problem \ref{m:ex14}.} Let the internal bisector $AI$
meet the circumcircle again at $L$. Show $IL=LB$ and compute $AI\cdot IL$.
\paragraph{Problem \ref{m:ex15}.} Use $n-1$ zeroes and $n-1$ ones.
A consecutive block from partial sums has uncontrolled length.
\paragraph{Problem \ref{m:ex16}.} Choose four vertices and their
cyclic order; their two diagonals cross once. Reverse the correspondence.


\chapter{Solutions}
\label{app:solutions}

% ---- solutions/foundations ----
\section{Learning to Prove}
\subsection*{Language and logical structure}
\paragraph{Problem \ref{f:ex01}.}
If the odd integers are $2a+1,2b+1$, their product is
$2(2ab+a+b)+1$. The expression in parentheses is an integer, so
the product is odd. Negative values of $a,b$ cause no exception.

\paragraph{Problem \ref{f:ex02}.}
The negation is: there exists $n\in\ZZ$ such that for every
$m\in\ZZ$, $m^2\ne n$. The original statement is false:
$n=-1$ works in its negation because every real square is nonnegative.
One counterexample is sufficient; classifying all nonsquares is
unnecessary.

\paragraph{Problem \ref{f:ex03}.}
Let $t=|x|\ge0$. Then $x^2=t^2$, so $t^2=t$ gives
$t(t-1)=0$. Thus $t=0$ or $1$, hence $x=0,1,-1$.
Each satisfies the original equation. Keeping $t\ge0$ explicit
also prevents a spurious negative value of the absolute value.

\paragraph{Problem \ref{f:ex04}.}
If $x+y$ were rational, subtracting rational $x$ would make $y$
rational, a contradiction. The product is not always irrational:
$x=0$ makes it zero. If $x\ne0$, rationality of $xy$ would give
$y=(xy)/x\in\QQ$, again a contradiction. The correct strengthened
statement includes the nonzero hypothesis.

\paragraph{Problem \ref{f:ex05}.}
The equation implies $x\ge0$. Squaring gives
$x^2-x-6=(x-3)(x+2)=0$. Only $x=3$ satisfies the sign requirement.
It does satisfy the equation since $\sqrt{3+6}=3$.

\paragraph{Problem \ref{f:ex06}.}
The identity $(a+b)^2-4ab=(a-b)^2\ge0$ gives $4ab\le1$.
Equality means $a=b$, so the sum condition gives $a=b=1/2$.
These values are allowed and attain the bound. A proposed
counterexample with another sum fails a hypothesis and therefore
cannot refute the conditional statement.

\paragraph{Problem \ref{f:ex07}.}
Suppose $\sqrt2=a/b$ with positive integers $a,b$ and
$\gcd(a,b)=1$. Then $a^2=2b^2$, so $a$ is even, say $a=2u$.
It follows that $b^2=2u^2$, so $b$ is also even. This contradicts
coprimality. Lowest terms enters only in the last contradiction:
an arbitrary fraction may perfectly well have two even entries.
We may assume lowest terms because repeatedly cancelling common
factors terminates at a positive denominator with no common factor.

\paragraph{Problem \ref{f:ex08}.}
Write the entries as $x_1,\ldots,x_k$, and their maximum as $M$.
Every difference $M-x_i$ is nonnegative, while
$\sum_i(M-x_i)=kM-\sum_i x_i=0$. A sum of nonnegative reals is
zero only if every term is zero; hence all $x_i=M$.
The mean of the set $\{0,2\}$ is $1$, which is not a member.
For a set with no repetitions, equality of all entries means the
set has just one element.

\subsection*{Induction, descent, and choices}
\paragraph{Problem \ref{f:ex09}.}
The base $n=1$ gives $1=1\cdot2/2$.
If the formula holds for $n$, adding the next entry gives
$n(n+1)/2+(n+1)=(n+1)(n+2)/2$. This is the formula for $n+1$,
so induction applies.

\paragraph{Problem \ref{f:ex10}.}
For $n=0$, $1=0+1$. If $2^n\ge n+1$, then
$2^{n+1}\ge2n+2\ge n+2$, since $n\ge0$. The second comparison
is part of the transition, not an automatic simplification.

\paragraph{Problem \ref{f:ex11}.}
Choose the least divisor $d>1$ of $n$, which exists since $n\mid n$.
If $d=uv$ with $1<u<d$, then $u\mid n$, contradicting the choice.
Thus $d$ is prime.
Alternatively, use strong induction. A prime $n$ divides itself.
A composite $n=ab$ with $1<a<n$ has, by induction, a prime divisor
of $a$, which also divides $n$. The base $n=2$ is immediate.

\paragraph{Problem \ref{f:ex12}.}
The bases are $8=3+5$, $9=3+3+3$, and $10=5+5$.
For $n\ge11$, the smaller integer $n-3\ge8$ is representable
by strong induction; adding one $3$ represents $n$.
Three bases are needed because the transition reduces the index
by $3$, not by $1$.

\paragraph{Problem \ref{f:ex13}.}
An integer divides $a,b$ exactly when it divides $a+b,b$:
one direction uses addition, the other subtraction. The analogous
fact holds for $a,a+b$. Thus each move preserves the gcd.
It is initially $1$ for both $(1,1)$ and $(2,3)$, so all resulting
pairs in either process are coprime.

\paragraph{Problem \ref{f:ex14}.}
If $a=b$ and the pair is coprime, then $a=b=1$.
Otherwise subtract the smaller positive entry from the larger.
The entries stay positive and coprime, and their sum strictly
decreases. Consequently this reverse process reaches an equal
pair, necessarily $(1,1)$. Reversing its finite sequence of steps
is a sequence of the required forward moves. This proves
reachability, which the gcd invariant alone could not prove.

\paragraph{Problem \ref{f:ex15}.}
Induct on the number of vertices; an empty or single-vertex graph
is immediate. Suppose all vertices except $v$ are partitioned into
a red path $R$ and a blue path $B$. If either is empty, put $v$
in that empty path, so assume both are nonempty and choose endpoints
$r,b$, respectively.
If $rv$ is red, append $v$ to $R$. If $bv$ is blue, append it to $B$.
Otherwise $rv$ is blue and $bv$ is red.
If $rb$ is red, remove $b$ from the end of $B$ and append $b,v$
to the end $r$ of $R$. If $rb$ is blue, remove $r$ from $R$
and append $r,v$ to the end $b$ of $B$. In each case the paths
remain monochromatic, disjoint, and cover every vertex.
Deleting the only vertex of a path is allowed by the empty-path convention.

\paragraph{Problem \ref{f:ex16}.}
A positive integer $d$ divides all original entries exactly when
it divides all new entries: only $x$ changes to $x+y$, and
$d\mid x,y$ is equivalent to $d\mid x+y,y$.
The same argument applies to the other replacement. Thus the
sets of common divisors, and therefore their largest members,
are identical. Equality of gcds is insufficient for reachability:
from $(1,2)$ neither coordinate can decrease, so $(1,1)$ is
unreachable although both gcds are $1$.

\subsection*{Representation and discovery}
\paragraph{Problem \ref{f:ex17}.}
The difference is $(a-b)^2\ge0$. Equality holds exactly when $a=b$.
No positivity hypothesis is needed for this identity.

\paragraph{Problem \ref{f:ex18}.}
The primes $3,5$ have composite sum $8$.
The irrational numbers $\sqrt2,\sqrt2$ have rational
product $2$. The increasing sequence $a_n=n$ has no real upper
bound. Each example satisfies the relevant hypotheses and
violates the claimed conclusion.

\paragraph{Problem \ref{f:ex19}.}
There are five residues $0,1,2,3,4$ modulo $5$.
Two of six integers have the same residue, so their difference
is divisible by $5$. The five integers $0,1,2,3,4$ have pairwise
different residues, and no nonzero difference is divisible by $5$.

\paragraph{Problem \ref{f:ex20}.}
For $x>0$,
$x+1/x-2=(x-1)^2/x\ge0$. Hence the minimum is at least $2$.
The allowed value $x=1$ gives $2$, proving both the bound and
attainment. Equality in the identity occurs only at $x=1$.

\paragraph{Problem \ref{f:ex21}.}
Partial fractions give $1/(k(k+1))=1/k-1/(k+1)$.
Summing causes every interior term to cancel, leaving
$1-1/(n+1)=n/(n+1)$. The formula also follows by induction,
but cancellation explains why the answer has this shape.

\paragraph{Problem \ref{f:ex22}.}
Each positive integer has a unique expression $2^r u$ with
$r\ge0$ and $u$ odd, by repeatedly dividing by $2$.
For integers at most $2n$, the odd part belongs to the $n$-element
set $\{1,3,\ldots,2n-1\}$. Two of the $n+1$ chosen integers
therefore have the same odd part: $2^r u,2^s u$, with $r\ne s$
because the integers are distinct. The smaller exponent gives
a number dividing the other.

\paragraph{Problem \ref{f:ex23}.}
For $t>0$, $(\sqrt t-1)^2\ge0$ gives $1+t\ge2\sqrt t$.
Multiply this inequality for $a,b,c$ to obtain
$(1+a)(1+b)(1+c)\ge8\sqrt{abc}=8$.
All factors are positive, and equality requires each factor
inequality to be an equality, so $a=b=c=1$.

\paragraph{Problem \ref{f:ex24}.}
An inversion is a pair of positions $i<j$ with $a_i>a_j$.
The number of inversions is a nonnegative integer at most
$n(n-1)/2$. Exchanging an adjacent inverted pair $x>y$
removes the inversion between them. For any other entry, the
total of its comparisons with $x,y$ is unchanged, since it lies
before both or after both. Thus exactly one inversion is removed.
The process must terminate. At termination no adjacent pair
decreases, which means the entire list is nondecreasing by
transitivity. Repeated entries cause no exception because
inversions use a strict inequality.



% ---- solutions/algebra ----
\section{Algebra}

\subsection*{Expressions and equations}

\paragraph{Problem~\ref{a:ex01}.}
Completing a square gives
\[
 x^4+4y^4=(x^2-2xy+2y^2)(x^2+2xy+2y^2).
\]
At \(y=1\), the factors are \((x-1)^2+1\) and \((x+1)^2+1\).
For an integer \(x>1\), both are integers greater than \(1\);
therefore their product is composite. The restriction excludes \(x=1\),
where the first factor is \(1\) and the product is the prime \(5\).

\paragraph{Problem~\ref{a:ex02}.}
The domain excludes \(1,-1\). Combining the fractions gives
\(2x/(x^2-1)=1\), equivalently \(x^2-2x-1=0\).
Its roots are \(1+\sqrt2\) and \(1-\sqrt2\), neither excluded.
Since we multiplied only by a nonzero denominator, both are solutions.

\paragraph{Problem~\ref{a:ex03}.}
Any solution has \(x\ge0\), and squaring gives
\(x^2-x-6=(x-3)(x+2)=0\). Only \(3\) is nonnegative,
and \(\sqrt{3+6}=3\). Thus \(x=3\).

\paragraph{Problem~\ref{a:ex04}.}
Direct multiplication yields
\[
 x^3+y^3+z^3-3xyz
 =(x+y+z)(x^2+y^2+z^2-xy-yz-zx).
\]
Under the given hypothesis, the right side is zero, proving the claim.
This factorization is also a useful recognition cue whenever
the sum of three variables vanishes.

\paragraph{Problem~\ref{a:ex05}.}
The product is
\[
 \frac21\cdot\frac32\cdot\frac43\cdots\frac{n+1}{n}=n+1.
\]
The cancellation is between factors in numerators and denominators;
it has nothing to do with distributing a product over addition.
For example, \((1+1)(1+1)=4\), whereas \(1+1\cdot1=2\).

\paragraph{Problem~\ref{a:ex06}.}
Twice the difference of the two sides is
\[
 (x-y)^2+(y-z)^2+(z-x)^2.
\]
It vanishes if and only if all three squares vanish, which is
equivalent to \(x=y=z\). Conversely equal variables clearly satisfy
the original equality.

\paragraph{Problem~\ref{a:ex07}.}
The domain is \(-1\le x\le4\). Squaring gives
\[
 5+2\sqrt{(x+1)(4-x)}=9,
\]
so \((x+1)(4-x)=4\). Thus \(-x^2+3x=0\) and \(x=0\) or \(3\).
Both values lie in the domain and give the original sum \(3\),
so both and only those values work.

\paragraph{Problem~\ref{a:ex08}.}
The discriminant is \(4(t^2-t-2)=4(t-2)(t+1)\).
For real roots we need \(t\le-1\) or \(t\ge2\).
Their sum is \(2t\), so positivity of both requires \(t>0\).
For \(t\ge2\), their product \(t+2\) is positive; with positive sum
and real roots, both roots are positive. The answer is \(t\ge2\),
including the repeated positive root at \(t=2\).

\paragraph{Problem~\ref{a:ex09}.}
The two numbers are roots of \(X^2-5X+6=(X-2)(X-3)\).
The pairs are \((2,3),(3,2)\), and both check.

\paragraph{Problem~\ref{a:ex10}.}
The monic polynomial with roots \(x,y,z\) is
\[
 X^3-6X^2+11X-6=(X-1)(X-2)(X-3).
\]
Its three roots are distinct, so the triple consists of \(1,2,3\)
in any order. Conversely every permutation has the stated symmetric
quantities.

\paragraph{Problem~\ref{a:ex11}.}
Since \(x^2+y^2+z^2=(x+y+z)^2-2(xy+yz+zx)=0\), all three
real squares vanish. Thus the only triple is \((0,0,0)\).

\paragraph{Problem~\ref{a:ex12}.}
Set \(p=xy>0\). Then \(x+y=8-p>0\), so \(p<8\).
Reality of \(x,y\) gives
\[
 (8-p)^2\ge4p,\qquad (p-4)(p-16)\ge0.
\]
Together with \(p<8\), this implies \(p\le4\).
At \(x=y=2\), the constraint holds and \(p=4\).
Thus the largest product is \(4\), attained only at \(x=y=2\).

\paragraph{Problem~\ref{a:ex13}.}
Order \(x\le y\le z\). Dividing \(xyz=x+y+z\le3z\) by \(z>0\)
gives \(xy\le3\). Hence \(x=1\). The equation becomes
\((y-1)(z-1)=2\). Neither factor can be zero, so under \(y\le z\)
the only factor pair is \(1,2\), giving \(y=2,z=3\).
Reordering gives exactly the permutations of \((1,2,3)\), each verified
by \(1+2+3=1\cdot2\cdot3\).

\paragraph{Problem~\ref{a:ex14}.}
Because \(c=-a-b\),
\(a^2+b^2-c^2=-2ab\ne0\). The other denominators are \(-2bc,-2ca\),
also nonzero. The sum is
\[
 -\frac12\left(\frac1{ab}+\frac1{bc}+\frac1{ca}\right)
 =-\frac{a+b+c}{2abc}=0.
\]

\subsection*{Sequences and floors}

\paragraph{Problem~\ref{a:ex15}.}
Induction on the recurrence gives \(a_n=3n-1\).
Therefore
\[
 \sum_{k=1}^n a_k
 =3\frac{n(n+1)}2-n=\frac{n(3n+1)}2.
\]
The formula has the correct initial term and each increase is \(3\).

\paragraph{Problem~\ref{a:ex16}.}
Let \(b_n=a_n+1\). Then \(b_{n+1}=2b_n\), \(b_1=2\),
so \(b_n=2^n\) and \(a_n=2^n-1\).

\paragraph{Problem~\ref{a:ex17}.}
The identity
\[
 \frac1{k(k+1)(k+2)}
 =\frac12\left(\frac1{k(k+1)}-\frac1{(k+1)(k+2)}\right)
\]
gives the sum
\[
 \frac14-\frac1{2(n+1)(n+2)}.
\]
Both the first and the last uncancelled term must be retained.

\paragraph{Problem~\ref{a:ex18}.}
The answer is \(a_n=F_{n+2}/F_{n+1}\).
It gives \(F_3/F_2=2\) initially, and
\[
 1+\frac1{a_n}
 =\frac{F_{n+2}+F_{n+1}}{F_{n+2}}
 =\frac{F_{n+3}}{F_{n+2}}.
\]
All Fibonacci denominators here are positive. Induction proves the formula.

\paragraph{Problem~\ref{a:ex19}.}
Write \(t=m+u\), where \(m=\lfloor t\rfloor\in\ZZ\) and \(0\le u<1\).
For any integer \(r\), adding \(u<1\) to \(r\) does not change which
multiple-of-\(n\) interval contains it. More explicitly, divide
\(m=nq+r\), with integers \(q,r\) and \(0\le r<n\). Then
\[
 \floor{\frac{t+k}{n}}=q+\floor{\frac{r+k+u}{n}}.
\]
The last floor is \(1\) for exactly the \(r\) values
\(k=n-r,\ldots,n-1\), and \(0\) otherwise. The sum is
\(nq+r=m=\lfloor t\rfloor\). This division also covers negative \(m\).

\paragraph{Problem~\ref{a:ex20}.}
Reciprocals give \(a_n=t/[1+(n-1)t]\).
Also the recurrence implies
\[
 a_k-a_{k+1}
 =a_k-\frac{a_k}{1+a_k}
 =\frac{a_k^2}{1+a_k}=a_ka_{k+1}.
\]
Thus the sum is \(a_1-a_{n+1}=t-t/(1+nt)\).

\paragraph{Problem~\ref{a:ex21}.}
Starting from \(a_1=a,a_2=b>0\), compute
\[
 a_3=\frac{1+b}{a},\quad
 a_4=\frac{a+b+1}{ab},\quad
 a_5=\frac{a+1}{b},\quad a_6=a,\quad a_7=b.
\]
Every denominator is positive, and cancellation proves each identity.
The recurrence uniquely determines the future from two adjacent terms.
Since the pair \((a_6,a_7)\) repeats \((a_1,a_2)\), the entire sequence
repeats after five steps. Its least period may be a divisor of \(5\),
but the requested five-step identity holds in every case.

\subsection*{Polynomial rigidity}

\paragraph{Problem~\ref{a:ex22}.}
The polynomial \(x^2+1\) has all four required values.
If \(P\) were another degree-at-most-\(3\) polynomial with those values,
\(P-(x^2+1)\) would have four distinct roots and degree at most \(3\).
It must vanish identically. Hence \(P=x^2+1\).

\paragraph{Problem~\ref{a:ex23}.}
Vieta gives \(\alpha+\beta=3,\alpha\beta=1\), so
\(\alpha^3+\beta^3=3^3-3\cdot1\cdot3=18\).

\paragraph{Problem~\ref{a:ex24}.}
Factoring the known roots gives
\(P(x)=(x-1)(x-2)(x-3)(x-t)\).
At \(x=4\), \(6(4-t)=24\), so \(t=0\).
Thus \(P=x(x-1)(x-2)(x-3)\), which satisfies every requirement.

\paragraph{Problem~\ref{a:ex25}.}
Since \((x+1)^3-x^3=3x^2+3x+1\), subtracting \(x^3\) produces
a polynomial \(Q\) with \(\Delta Q=0\).
The degree-drop property excludes a nonconstant \(Q\).
Therefore \(P=x^3+C\), for an arbitrary real constant \(C\);
substitution proves sufficiency.

\paragraph{Problem~\ref{a:ex26}.}
Put \(t=x^2\), so the condition is \(t^2+at+b\ge0\) for all \(t\ge0\).
Necessarily \(b\ge0\). If \(a\ge0\), all terms are nonnegative,
so this condition is sufficient. If \(a<0\), completing the square gives
\[
 t^2+at+b=(t+a/2)^2+b-a^2/4.
\]
Its minimum on \(t\ge0\) occurs at \(t=-a/2>0\), so we additionally
need \(b\ge a^2/4\). The combined answer is
\(b\ge0\) and \(a\ge-2\sqrt b\).
The argument establishes both necessity and sufficiency.

\paragraph{Problem~\ref{a:ex27}.}
The zero polynomial works. Otherwise factor off the exact power of \(x\):
\(P=x^kQ\), \(k\ge0,Q(0)\ne0\). The identity gives \(Q(x^2)=Q(x)^2\),
so \(Q(0)=Q(0)^2\) and \(Q(0)=1\).
If \(j\) is the smallest positive exponent with nonzero coefficient
\(c\) in \(Q\), the coefficient of \(x^j\) in \(Q(x)^2\) is \(2c\).
In \(Q(x^2)\) it is zero, because either \(j\) is odd or its source
coefficient at \(j/2\) vanishes by minimality. This contradiction shows
\(Q=1\). Thus all solutions are \(0\) and \(x^k\), \(k\ge0\).

\paragraph{Problem~\ref{a:ex28}.}
Write \(P(x)=c\prod_{j=1}^d(x-r_j)\), \(c\ne0\), with real roots listed
with multiplicities. For \(x\) not equal to any root, differentiating
the product and dividing by \(P\) gives
\[
 \frac{P'(x)}{P(x)}=\sum_{j=1}^d\frac1{x-r_j}.
\]
Differentiate this rational identity and rearrange:
\[
 P'(x)^2-P(x)P''(x)
 =P(x)^2\sum_{j=1}^d\frac1{(x-r_j)^2}>0.
\]
The ordinary product and quotient differentiation rules used here follow
algebraically from the formal derivative rule. At a root, the left side
is \(P'(x)^2\ge0\). Thus the inequality holds for every real \(x\),
including repeated roots.

\subsection*{Integer and multivariable polynomials}

\paragraph{Problem~\ref{a:ex29}.}
Write \(P(x)=\sum_{k=0}^d a_kx^k\), \(a_k\in\ZZ\).
For every \(k\ge1\),
\[
 m^k-n^k=(m-n)(m^{k-1}+m^{k-2}n+\cdots+n^{k-1}).
\]
Taking the integer linear combination of these identities proves
\(m-n\mid P(m)-P(n)\).

\paragraph{Problem~\ref{a:ex30}.}
The divisibility in Problem~\ref{a:ex29} would give
\(2\mid P(2)-P(0)=1\), which is impossible.
The rational-coefficient polynomial \(x/2\) has the two requested
values, illustrating exactly where the obstruction occurs.

\paragraph{Problem~\ref{a:ex31}.}
Take \(P(x)=x(x-1)/2\). One of two consecutive integers \(n,n-1\)
is even, whether positive, zero, or negative; therefore \(P(n)\in\ZZ\)
for every integer \(n\). Its coefficient of \(x^2\) is \(1/2\).
Thus integer values do not force integer coefficients.

\paragraph{Problem~\ref{a:ex32}.}
The polynomial is primitive. The prime \(2\) divides every
nonleading coefficient, does not divide the leading coefficient,
and \(4\) does not divide its constant coefficient \(2\).
Eisenstein proves irreducibility over \(\QQ\).
Checking rational roots alone would only rule out linear factors.
A quartic may factor into two quadratic factors, as \(x^4+4\) does.

\paragraph{Problem~\ref{a:ex33}.}
Grouping gives
\[
 2x^3-3x^2-8x+12=x^2(2x-3)-4(2x-3)
 =(2x-3)(x-2)(x+2).
\]
The roots are \(3/2,2,-2\). The factorization proves there are no others.

\paragraph{Problem~\ref{a:ex34}.}
The integer polynomial \(P(x)-1\) has roots \(0,1,2\), so
\[
 P(x)-1=x(x-1)(x-2)Q(x)
\]
with \(Q\in\ZZ[x]\). Successive division by monic linear factors
preserves integer coefficients. If \(P(n)=2\), then
\(n(n-1)(n-2)Q(n)=1\).
For \(n=0,1,2\), the left side is zero; otherwise the product of three
consecutive integers is a nonzero multiple of \(6\), so it cannot divide \(1\).
This excludes every integer \(n\).

\paragraph{Problem~\ref{a:ex35}.}
Choose an integer \(a\) with \(c=P(a)\ne0\); a nonzero polynomial
has only finitely many roots. Expanding each monomial of \(P(a+ct)\)
shows that
\[
 Q(t)=\frac{P(a+ct)}c
\]
is an integer-coefficient nonconstant polynomial and \(Q(0)=1\).
For the constant coefficient this is the definition of \(c\);
every coefficient of positive degree contains a factor \(c\)
before division.

Suppose the prime divisors of all nonzero values of \(P\) formed the
finite list \(p_1,\ldots,p_s\), and put \(N=p_1\cdots p_s\), with
empty product \(1\). For every integer \(k\), \(Q(Nk)\equiv1\pmod N\).
Because \(Q\) is nonconstant, choose \(k\) with \(|Q(Nk)|>1\).
A prime divisor of \(Q(Nk)\) is coprime to \(N\), hence absent from
the list, but divides the nonzero value
\(P(a+cNk)=cQ(Nk)\). Contradiction.

\subsection*{Complex numbers and trigonometry}

\paragraph{Problem~\ref{a:ex36}.}
Repeated squaring gives \((1+i)^2=2i\), \((1+i)^4=-4\),
and \((1+i)^8=16\).
Alternatively, \(1+i=\sqrt2\,(\cos(\pi/4)+i\sin(\pi/4))\).
Its eighth power has modulus \(16\) and angle \(2\pi\), giving the same answer.

\paragraph{Problem~\ref{a:ex37}.}
Let \(S=\sum_{k=0}^{n-1}k\zeta^k\).
Since \(\zeta^n=1\),
\[
 (\zeta-1)S=(n-1)-\sum_{j=1}^{n-1}\zeta^j=n.
\]
Here \(\sum_{j=1}^{n-1}\zeta^j=-1\), by the geometric-sum identity
and \(\zeta\ne1\). Dividing gives \(S=n/(\zeta-1)\).
No assumption that \(\zeta\) is a primitive root is needed.

\paragraph{Problem~\ref{a:ex38}.}
Because \(A+B=\pi-C\),
\[
\begin{aligned}
\sin2A+\sin2B+\sin2C
 &=2\sin(A+B)\cos(A-B)-2\sin(A+B)\cos(A+B)\\
 &=2\sin C\,[\cos(A-B)-\cos(A+B)]\\
 &=4\sin A\sin B\sin C.
\end{aligned}
\]
The formulas remain valid even when the angles are not angles of a triangle;
only their sum was used.

\paragraph{Problem~\ref{a:ex39}.}
The factorization is
\[
 x^4+x^2+1=(x^2+x+1)(x^2-x+1).
\]
Also \((x^2-1)(x^4+x^2+1)=x^6-1\).
The four roots are therefore the sixth roots of unity other than \(1,-1\):
\(e^{i\pi/3},e^{-i\pi/3},e^{2i\pi/3},e^{-2i\pi/3}\).
Each quadratic has negative discriminant, so the displayed factorization
is complete over the reals.

\paragraph{Problem~\ref{a:ex40}.}
Let \(\zeta=e^{2\pi i/n}\). Taking real parts in
\(e^{i\theta}\sum_{k=0}^{n-1}\zeta^k=0\) proves the first formula.
For the second, use
\[
 \sum_{k=0}^{n-1}\cos^2(\theta+2\pi k/n)
 =\frac n2+\frac12\Re\left(e^{2i\theta}
                 \sum_{k=0}^{n-1}\zeta^{2k}\right).
\]
For \(n\ge3\), \(\zeta^2\ne1\), so the geometric sum is zero.
For \(n=2\), the sum is \(2\cos^2\theta\), which need not equal \(1\).

\paragraph{Problem~\ref{a:ex41}.}
Let \(\zeta=e^{2\pi i/n}\).
The identity
\(\prod_{k=1}^{n-1}(1-\zeta^k)=n\) follows by factoring
\(1+x+\cdots+x^{n-1}\) and setting \(x=1\).
Now \(|1-\zeta^k|=2\sin(k\pi/n)\), since \(0<k\pi/n<\pi\).
Take absolute values of the product and divide by \(2^{n-1}\).

\paragraph{Problem~\ref{a:ex42}.}
Let \(Q(x)=1+x+\cdots+x^{n-1}=\prod_{k=1}^{n-1}(x-\zeta^k)\),
where \(\zeta=e^{2\pi i/n}\).
Product differentiation gives
\[
 A:=\sum_{k=1}^{n-1}\frac1{1-\zeta^k}
   =\frac{Q'(1)}{Q(1)}=\frac{n-1}{2}.
\]
Differentiate \(Q'/Q=\sum1/(x-\zeta^k)\) once more to obtain
\[
\begin{aligned}
 B:=\sum_{k=1}^{n-1}\frac1{(1-\zeta^k)^2}
 &=\left(\frac{Q'(1)}{Q(1)}\right)^2-\frac{Q''(1)}{Q(1)}\\
 &=\frac{(n-1)^2}{4}-\frac{(n-1)(n-2)}3
 =\frac{(n-1)(5-n)}{12}.
\end{aligned}
\]
The derivative values follow from
\(Q(1)=n\), \(Q'(1)=n(n-1)/2\), and
\(Q''(1)=n(n-1)(n-2)/3\).
Since \(\overline{\zeta^k}=\zeta^{-k}\),
\[
 \frac1{|1-\zeta^k|^2}
 =-\frac{\zeta^k}{(1-\zeta^k)^2}
 =\frac1{1-\zeta^k}-\frac1{(1-\zeta^k)^2}.
\]
Sum to obtain \(A-B=(n^2-1)/12\).
Finally \(|1-\zeta^k|^2=4\sin^2(k\pi/n)\), so multiplication by \(4\)
gives the requested \((n^2-1)/3\).
Although intermediate sums are complex, the final identity is real.

\subsection*{Basic inequalities}

\paragraph{Problem~\ref{a:ex43}.}
Cauchy's fraction form and the sum-of-squares inequality give
\[
 \sum_{\cyc}\frac a{b+c}
 =\sum_{\cyc}\frac{a^2}{a(b+c)}
 \ge\frac{(a+b+c)^2}{2(ab+bc+ca)}
 \ge\frac32.
\]
The last step is equivalent to
\(a^2+b^2+c^2\ge ab+bc+ca\).
It is equality only at \(a=b=c\), which also gives equality in Cauchy.

\paragraph{Problem~\ref{a:ex44}.}
Cauchy gives \(1=(x+y+z)^2\le3(x^2+y^2+z^2)\).
Thus the minimum is \(1/3\), attained only at \(x=y=z=1/3\).
Positivity was not assumed or needed.

\paragraph{Problem~\ref{a:ex45}.}
When two paired lists have an inversion, swapping that pair changes
the sum by a product of their nonnegative differences.
Thus equal ordering maximizes the sum, and opposite ordering minimizes it.
The maximum is \(1\cdot4+2\cdot5+3\cdot6=32\);
the minimum is \(1\cdot6+2\cdot5+3\cdot4=28\).
Strict ordering makes these respective arrangements unique.

\paragraph{Problem~\ref{a:ex46}.}
The identity
\[
 a^2+b^2+c^2-ab-bc-ca
 =\frac12[(a-b)^2+(b-c)^2+(c-a)^2]
\]
proves the first inequality. AM--GM for \(ab,bc,ca\) gives
\(1=(ab+bc+ca)/3\ge(abc)^{2/3}\), hence \(abc\le1\).
In either equality case all three positive variables are equal;
the condition then gives \(a=b=c=1\).

\paragraph{Problem~\ref{a:ex47}.}
Apply \(u^2+v^2+w^2\ge uv+vw+wu\) to
\(u=ab,v=bc,w=ca\).
The right side becomes \(abc(a+b+c)\).
Equality holds exactly when \(ab=bc=ca\): either all positive variables
are equal, or at least two variables vanish.

\paragraph{Problem~\ref{a:ex48}.}
Each denominator is positive because \(a,b,c>0\) and their sum is \(1\).
Their sum is \(2\), so Cauchy gives
\[
 \sum_{\cyc}\frac{a^2}{1-a}
 \ge\frac{(a+b+c)^2}{(1-a)+(1-b)+(1-c)}=\frac12.
\]
Equality requires \(a/(1-a)=b/(1-b)=c/(1-c)\), hence
\(a=b=c=1/3\), which works.

\paragraph{Problem~\ref{a:ex49}.}
Write \(s=a+b+c\). The product constraint and AM--GM imply \(s\ge3\).
Then Cauchy yields \(\sum a^2\ge s^2/3\ge s\).
Equality forces \(s=3\) and equal variables, so \(a=b=c=1\).
Without the constraint the assertion is false, for example at
\(a=b=c=1/2\). The product condition supplies the crucial lower bound on \(s\).

\subsection*{Convexity and sharper methods}

\paragraph{Problem~\ref{a:ex50}.}
Weighted AM--GM gives
\(\lambda x+(1-\lambda)y\ge x^\lambda y^{1-\lambda}>0\).
The increasing logarithm preserves the inequality and turns the product
into the desired weighted sum of logarithms.
For \(0<\lambda<1\), equality holds exactly when \(x=y\);
at the endpoint weights it is always equality.

\paragraph{Problem~\ref{a:ex51}.}
Hölder with exponents \(3,3/2\) gives
\[
 \sum a_i
 =\sum \frac{a_i}{b_i^{2/3}}b_i^{2/3}
 \le\left(\sum\frac{a_i^3}{b_i^2}\right)^{1/3}
      \left(\sum b_i\right)^{2/3}.
\]
Cube and divide by the positive denominator.
The normalized weighted AM--GM proof of Hölder shows equality exactly
when \(a_i/b_i\) is a common nonnegative constant.
This includes all \(a_i=0\); otherwise the constant is positive.

\paragraph{Problem~\ref{a:ex52}.}
Put \(A=\sqrt{x^2+y^2}\), \(B=\sqrt{u^2+v^2}\).
Cauchy gives \(xu+yv\le AB\). Consequently
\[
 (x+u)^2+(y+v)^2=A^2+B^2+2(xu+yv)\le(A+B)^2.
\]
Taking nonnegative square roots proves the claim.
Equality holds when one vector is zero or the vectors are positive
scalar multiples of one another. Proportional vectors with opposite
directions do not generally give equality, because the dot product
must be \(+AB\).

\paragraph{Problem~\ref{a:ex53}.}
Apply weighted AM--GM with weights \(1/r,(r-1)/r\):
\[
 \frac{x^r}{r}+\frac{r-1}{r}a^r
 \ge(x^r)^{1/r}(a^r)^{(r-1)/r}=xa^{r-1}.
\]
Multiply by \(r\) and rearrange. Both weights are positive, so equality
holds exactly when \(x^r=a^r\), equivalently \(x=a\).

\paragraph{Problem~\ref{a:ex54}.}
For \(p=a+b+c=0\), all variables vanish. Otherwise the difference
is \(p^3-4pq+9r\), so Lemma~\ref{a:pqr} applies on fixed-\(p,q\) fibers.
If \(c=0\), the difference is
\(p[(a+b)^2-4ab]=p(a-b)^2\ge0\).
If \(a=b=t,c=s\), direct expansion gives
\[
 (2t+s)^3+9t^2s-4(2t+s)(t^2+2ts)=s(t-s)^2\ge0.
\]
Thus the inequality holds everywhere. The ordered proof of Schur,
or inspection of its nonnegative terms, shows equality precisely when
all three are equal or one is zero and the other two are equal.

\paragraph{Problem~\ref{a:ex55}.}
Set \(q=ab+bc+ca\). Since \(abc=1\),
\(\sum1/a=q\). Cauchy applied in the deliberately chosen form gives
\[
 \sum_{\cyc}\frac1{a^3(b+c)}
 =\sum_{\cyc}\frac{(1/a)^2}{a(b+c)}
 \ge\frac{(1/a+1/b+1/c)^2}{2q}
 =\frac q2\ge\frac32.
\]
The last step is AM--GM on \(ab,bc,ca\), whose product is \(1\).
It forces equality only at \(a=b=c=1\); that triple gives equality
in every step.

\paragraph{Problem~\ref{a:ex56}.}
Order \(x\ge y\ge z\ge0\). The sum is
\[
 (x-y)\bigl[x^t(x-z)-y^t(y-z)\bigr]
       +z^t(z-x)(z-y).
\]
Both terms are nonnegative. If \(x>y\), then
\(x^t(x-z)>y^t(y-z)\), so the first term is strictly positive.
Thus equality requires \(x=y\). The remaining term is
\(z^t(x-z)^2\), which vanishes exactly when \(z=0\) or \(z=x\).
Reordering gives precisely the all-equal triples and the triples
with one zero and two equal entries.

\subsection*{Functional equations}

\paragraph{Problem~\ref{a:ex57}.}
An additive function has \(f(q)=qf(1)=2q\) for rational \(q\).
The values at \(0,1\) show that a monotone \(f\) cannot be decreasing.
For real \(x\), take rationals \(r_j\uparrow x\), \(s_j\downarrow x\).
Then \(2r_j\le f(x)\le2s_j\), so \(f(x)=2x\).
The resulting function satisfies all assumptions.

\paragraph{Problem~\ref{a:ex58}.}
Define \(A(n)=f(n)-n^2\). Expanding \((m+n)^2\) turns the equation
into \(A(m+n)=A(m)+A(n)\). On the integers this gives
\(A(n)=nA(1)=cn\), \(c\in\ZZ\).
Thus all solutions are \(f(n)=n^2+cn\), and direct substitution works
for every integer \(c\).

\paragraph{Problem~\ref{a:ex59}.}
Put \(A(x)=f(x)+1\). The equation is exactly
\(A(x+y)=A(x)+A(y)\). Therefore all solutions are \(f=A-1\),
where \(A\) is an arbitrary real additive function.
If \(f\) is continuous, \(A\) is continuous too, and
Theorem~\ref{a:cauchy} gives \(f(x)=cx-1\), \(c\in\RR\).
Both descriptions are sufficient by substitution.

\paragraph{Problem~\ref{a:ex60}.}
Let \(c=f(0)\) and \(g(x)=f(x)-x^2-c\).
The given equation becomes
\[
 g(x+y)+g(x-y)=2g(x),\qquad g(0)=0.
\]
Taking \(y=x\) gives \(g(2x)=2g(x)\).
Now put \(x=(u+v)/2,y=(u-v)/2\) to obtain
\[
 g(u)+g(v)=2g((u+v)/2)=g(u+v).
\]
Thus \(g=A\) is additive. Conversely every function
\[
 f(x)=x^2+A(x)+c,
\]
where \(A\) is additive and \(c\in\RR\), satisfies the equation:
the quadratic parts contribute \(2x^2+2y^2\), the additive parts
contribute \(2A(x)\), and the constants contribute \(2c\).
No continuity assumption permits replacing \(A\) by a linear function.

\paragraph{Problem~\ref{a:ex61}.}
Since values are positive, \(f(1)=f(1)^2\) gives \(f(1)=1\).
With \(f(2)=4\), monotonicity must be increasing.
Define \(A(t)=\log f(e^t)\). It is additive because
\(f(e^{s+t})=f(e^s)f(e^t)\), and increasing because each component
of the composition is increasing. Therefore \(A(t)=ct\).
At \(t=\log2\), \(c\log2=\log4=2\log2\), so \(c=2\).
For \(x>0\), \(f(x)=e^{A(\log x)}=x^2\), which verifies the equation.

\paragraph{Problem~\ref{a:ex62}.}
Let \(c=f(0)\). Taking \(x=0\) gives \(f(f(y))=c+y\).
Thus \(f\) is surjective, and if \(f(u)=f(v)\), applying \(f\)
shows \(u=v\), so it is injective.
Taking \(y=0\) gives \(f(x+c)=f(x)\).
Injectivity forces \(x+c=x\), hence \(c=0\).
We now have \(f(f(y))=y\). Replace \(y\) by \(f(t)\) in the original
equation to get
\[
 f(x+t)=f(x)+f(t).
\]
Thus \(f\) is an additive involution. Conversely any additive
involution satisfies \(f(x+f(y))=f(x)+f(f(y))=f(x)+y\).
This is the complete description without regularity.
If continuous, \(f(x)=kx\), and the involution property gives \(k^2=1\).
The continuous solutions are \(f(x)=x\) and \(f(x)=-x\).

\paragraph{Problem~\ref{a:ex63}.}
Taking \(y=0\) gives \(f(x)=f(x)f(0)\).
Either \(f=0\), which works, or \(f(0)=1\).
In the second case \(f(x)f(-x)=1\), so no value is zero;
also \(f(x)=f(x/2)^2>0\).
Hence \(A(x)=\log f(x)\) is continuous and additive, so \(A(x)=cx\).
The nonzero solutions are \(f(x)=e^{cx}\), \(c\in\RR\), each verified
by the exponential multiplication law.

\paragraph{Problem~\ref{a:ex64}.}
Evaluating the third iterate in the two possible ways gives
\[
 f(n+1)=f(f(f(n)))=f(n)+1.
\]
It follows for all integers, forward and backward, that \(f(n)=n+c\),
where \(c=f(0)\in\ZZ\). Then \(f(f(n))=n+2c=n+1\), impossible
because no integer satisfies \(2c=1\). Thus there is no such function.

\paragraph{Problem~\ref{a:ex65}.}
If \(f(x)\ne x\), the involution property gives the distinct pair
\(\{x,f(x)\}\), whose two elements map to each other. Different pairs
are disjoint because an element has exactly one image.
If there are \(k\) pairs and \(r\) fixed points, then \(N=2k+r\).
Therefore \(r\equiv N\pmod2\).

\paragraph{Problem~\ref{a:ex66}.}
The identically zero function fails the equation, for example at \(x=y=1\).
Taking one input zero therefore gives \(f(0)=1\).
Compute \(f(x+y+z)\) using \((x+y)+z\) and \(x+(y+z)\), then cancel
the common product \(f(x)f(y)f(z)\). The remaining equality is
\[
 xy(f(z)-1)=yz(f(x)-1).
\]
Set \(y=z=1\). There is a real constant \(c=f(1)-1\) with
\(f(x)=1+cx\) for every \(x\).
Substitution gives \(c^2xy=xy\) for all \(x,y\), so \(c^2=1\).
Thus precisely \(f(x)=1+x\) and \(f(x)=1-x\) work.

\paragraph{Problem~\ref{a:ex67}.}
Constant solutions satisfy \(c=c^2\), giving \(0,1\).
If \(\deg P=d\ge1\), degrees in the equation give \(d^2=2d\),
hence \(d=2\). Write \(P=ax^2+bx+c\), \(a\ne0\).
Leading coefficients give \(a^3=a^2\), so \(a=1\).
Then \(P(P(x))-P(x)^2=bP(x)+c=0\).
Its degree forces \(b=0\), and then \(c=0\).
Thus the only nonconstant solution is \(x^2\); all three candidates
verify directly.

\paragraph{Problem~\ref{a:ex68}.}
Denote the given identity by \(E(x,y)\), and set \(c=f(0)\).
From \(E(0,y)\),
\[
 f(f(y))=y+c^2.
\]
It follows that \(f\) is surjective and injective.
We first derive a translation identity by substituting
\(y=u^2+f(v)\) into \(E(x,y)\). Since
\(f(u^2+f(v))=v+f(u)^2\), this gives
\[
 f(x^2+v+f(u)^2)=u^2+f(v)+f(x)^2.
\]
Taking \(x=0\) yields
\[
 f(v+f(u)^2)=f(v)+u^2+c^2. \tag{*}
\]
Surjectivity provides \(t\) with \(f(t)=0\).
At \(u=t\), equation (*) gives \(t^2+c^2=0\).
Thus \(t=0,c=0\), and in particular \(f(f(y))=y\).

Every \(h\ge0\) can be written as \(f(u)^2\), because \(f\) is onto.
At \(v=0\), (*) then says \(f(h)=u^2\ge0\).
For arbitrary real \(v\), the same equation becomes
\[
 f(v+h)=f(v)+f(h)\qquad(h\ge0).
\]
Taking \(v=-h\) gives oddness at those \(h\).
For a negative shift \(h=-k\), \(k>0\), the positive-shift equation
applied at \(v-k\) gives \(f(v-k)=f(v)-f(k)=f(v)+f(-k)\).
Therefore \(f\) is additive on all reals.

Since \(f\) is nonnegative on nonnegative inputs, additivity makes it
increasing. Theorem~\ref{a:cauchy} gives \(f(x)=kx\), \(k\ge0\).
The identity \(f(f(y))=y\) forces \(k^2=1\), so \(k=1\).
Finally \(f(x)=x\) satisfies the original equation. It is the unique solution.

\paragraph{Problem~\ref{a:ex69}.}
Let \(c=f(0)\). The case \(x=0\) gives
\(f(f(y))=c^2+y\), so \(f\) is onto and injective.
Choose \(t\) with \(f(t)=0\). The original equation at \((t,0)\)
gives \(f(c)=0\), while its value at \((0,0)\) gives \(f(c)=c^2\).
Hence \(c=0\), and \(f(f(y))=y\).

Replace \(y\) by \(f(z)\) in the original equation:
\[
 f(xf(x)+z)=f(x)^2+f(z). \tag{*}
\]
Replace \(x\) in this identity by \(f(x)\). The shift \(xf(x)\)
does not change, since \(f(f(x))=x\); the right side becomes
\(x^2+f(z)\). Comparing gives \(f(x)^2=x^2\) for every real \(x\).
So each nonzero value is either \(x\) or \(-x\); we must still show
that these signs cannot be mixed.

For a real \(x\), put \(t_x=xf(x)\). Then
\(t_x\) is either \(x^2\) or \(-x^2\), and (*) at \(z=0\) gives
\(f(t_x)=x^2\). For any nonzero \(x,y\), (*) also gives
\[
 f(t_x+t_y)=x^2+f(t_y)=x^2+y^2.
\]
If \(t_x,t_y\) had opposite signs, the equality
\(|f(u)|=|u|\) at \(u=t_x+t_y\) would imply
\[
 x^2+y^2=|t_x+t_y|=|x^2-y^2|<x^2+y^2,
\]
a contradiction. All nonzero \(t_x\) therefore have the same sign.
Consequently \(f(x)=x\) for all \(x\), or \(f(x)=-x\) for all \(x\).
For the first candidate both sides of the original identity are \(x^2+y\).
For the second the left side is \(f(-x^2-y)=x^2+y\), also equal to
the right side. Both work.

\paragraph{Problem~\ref{a:ex70}.}
At \(m=n=0\), the equation gives \(f(0)=0\).
At \(m=0\), it gives \(f(-n)=f(n)\).
Let \(c=f(1)\in\ZZ\). Setting \(n=1\) yields
\[
 f(m+1)-2f(m)+f(m-1)=2c.
\]
Subtract the particular solution \(cm^2\). The remaining function
\(g(m)=f(m)-cm^2\) has zero second difference and \(g(0)=g(1)=0\).
The recurrence determines \(g=0\) for all nonnegative integers,
and evenness supplies the negative integers.
Thus \(f(n)=cn^2\), for any integer \(c\), and substitution verifies all.

\subsection*{Synthesis}

\paragraph{Problem~\ref{a:ex71}.}
Permute the coordinates so \(x\le y\le z\).
As a quadratic in the last coordinate, the equation is
\[
 F(T)=T^2-3xyT+x^2+y^2=0.
\]
Its other root is \(z'=3xy-z=(x^2+y^2)/z\), a positive integer.
If \(x=y=1\), the possible \(z\) are \(1,2\); the latter changes
to \(1\), and the former is the seed triple.

In every other case,
\[
 F(y)=x^2-(3x-2)y^2<0.
\]
Indeed for \(x=1,y>1\) this is \(1-y^2<0\);
for \(x\ge2\), it is at most \(x^2-4y^2<0\).
The upward-opening quadratic is negative strictly between its roots.
Thus \(0<z'<y<z\).
Replacing \(z\) by \(z'\) preserves positive integer solutions and
strictly decreases their sum. After reordering, repeat.
The positive integer sum cannot decrease indefinitely, so the process
reaches \((1,1,1)\).

Each replacement is reversible, since
\(3xy-(3xy-z)=z\). Reversing the descent gives the claimed generation
of every solution. Conversely every listed operation preserves the
equation and, whenever applied to a positive solution, the alternate
root is positive as proved above. This establishes the full classification.

\paragraph{Problem~\ref{a:ex72}.}
A nonzero polynomial has an eventual constant sign and is asymptotic
to its nonzero leading term. Direct division by the leading power of
\(n\) shows
\[
 \frac{P(n+1)}{P(n)}\longrightarrow1.
\]
For all sufficiently large \(n\), this ratio is an integer by the
divisibility assumption, is positive by the eventual sign, and lies
strictly between \(1/2\) and \(3/2\). Hence it equals \(1\).
The polynomial \(P(x+1)-P(x)\) therefore has infinitely many roots
and is identically zero. Degree reduction shows \(P\) is constant.
Every nonzero integer constant satisfies the assumptions.

\paragraph{Problem~\ref{a:ex73}.}
For the first inequality use the identity
\[
 \frac1{x^2}-x^2+4(x-1)
 =\frac{(x-1)^2(-x^2+2x+1)}{x^2}.
\]
The factor \(-x^2+2x+1\) is positive for \(0<x\le7/3\),
as is seen from its concavity and the endpoint values \(1,2/9\)
on \([0,7/3]\).
If \(a,b,c\le7/3\), sum the identity and use \(a+b+c=3\).
Equality then requires each variable to be \(1\).
If, say, \(c>7/3\), then one of \(a,b\) is below \(1/3\).
Its reciprocal square exceeds \(9\), whereas
\(a^2+b^2+c^2<9\), so the inequality is strict.

For the second inequality its difference from \(4\) is
\(5-2q+r\), with \(q=ab+bc+ca,r=abc\). Lemma~\ref{a:pqr} applies.
At \(c=0\) the sum of squares is at least \(9/2\).
At \(a=b=t,c=3-2t\), \(0\le t\le3/2\), the difference is
\[
 (t-1)^2(5-2t)\ge0.
\]
For equality classification without relying just on this reduction,
suppose the minimum \(0\) occurred at three distinct positive entries.
The root-perturbation argument of Lemma~\ref{a:pqr} could decrease \(r\)
while preserving \(p,q\), making \(5-2q+r<0\), impossible.
Thus equality requires a repeated variable or a zero, and the two
checked cases force \(a=b=c=1\).

\paragraph{Problem~\ref{a:ex74}.}
Let \(R(x)=\prod_{j=1}^n(x-a_j)\), \(P=R-1\).
If \(P\) factored over \(\QQ\), Gauss's lemma and monicity would give
monic integer factors \(P=GH\) of positive degrees. Choose the smaller
degree \(d=\deg G\le n/2\).
At each \(a_j\), \(G(a_j)H(a_j)=-1\), so \(G(a_j)=\pm1\).
Hence \(G^2-1\) has the \(n\) distinct roots \(a_j\).

If \(2d<n\), the root bound makes \(G^2-1=0\), impossible for
positive-degree \(G\).
If \(2d=n\), leading coefficients and roots give \(G^2-1=R\).
Then \(P=G^2-2\), so \(G\mid P\) implies \(G\mid2\), again impossible
for a positive-degree polynomial. This proves irreducibility.
The case \(n=1\) is also immediate from the degree.

For the opposite sign, group the factors using \(u=x^2+3x\):
\[
 x(x+1)(x+2)(x+3)+1=u(u+2)+1=(x^2+3x+1)^2.
\]
Thus the sign cannot simply be changed.

\paragraph{Problem~\ref{a:ex75}.}
Put \(U=x^2+y^2+z^2\). Under \(x+y+z=0\),
\(xy+yz+zx=-U/2\), and
\[
 x^2y^2+y^2z^2+z^2x^2
 =(xy+yz+zx)^2-2xyz(x+y+z)=U^2/4.
\]
It follows that
\(\sum x^4=U^2-2(U^2/4)=U^2/2\).
The given fourth-power sum therefore yields \(U=2\).

After permuting and, if necessary, reversing all signs, assume
\(x,y\ge0\), \(z=-(x+y)\).
Then \(x^2+xy+y^2=1\).
Set \(s=x+y,t=xy\). We have \(s^2-t=1\) and \(0\le t\le s^2/4\).
Consequently \(1\le s\le2/\sqrt3\), and
\[
 |xyz|=st=s(s^2-1).
\]
This expression is strictly increasing on \(s\ge1\):
for \(v>u\ge1\), its increase is
\((v-u)(v^2+uv+u^2-1)>0\).
The maximum is therefore \(2/(3\sqrt3)\).
Equality requires \(s=2/\sqrt3,t=s^2/4\), so
\(x=y=1/\sqrt3,z=-2/\sqrt3\), together with all permutations and
simultaneous sign reversals.

\paragraph{Problem~\ref{a:ex76}.}
Let \(s=x+y+z>0\). If \(s>3\), put \(t=3/s\in(0,1)\).
The scaled triple \(a=tx,b=ty,c=tz\) has sum \(3\).
Since \(x^2+y^2+z^2>0\),
\[
 a^2+b^2+c^2+abc
 =t^2(x^2+y^2+z^2)+t^3xyz
 <x^2+y^2+z^2+xyz=4.
\]
This contradicts the inequality of Problem~\ref{a:ex73}, whose proof
also allows zero variables. Thus \(s\le3\).
When \(s=3\), the equality case there gives \(x=y=z=1\);
this triple satisfies the constraint.

\paragraph{Problem~\ref{a:ex77}.}
For a constant \(P=c\), the equation is \(c^2=c+2\), giving \(c=2,-1\).
For a nonconstant solution of degree \(n\ge1\) and leading coefficient
\(a\), comparing leading coefficients gives \(a^2=a\), hence \(a=1\).

Define \(C_n(x)=2T_n(x/2)\) for \(n\ge1\). It is monic of degree \(n\).
At \(x=2\cos\theta\),
\[
 C_n(x)^2=4\cos^2(n\theta)
 =2\cos(2n\theta)+2=C_n(x^2-2)+2.
\]
Both sides are polynomials and agree at infinitely many real \(x\),
so the identity holds everywhere.

Compare an arbitrary degree-\(n\) solution \(P\) with \(C_n\).
If \(D=P-C_n\ne0\), its degree \(k\) is less than \(n\), since
both leading coefficients are \(1\). Subtracting their equations gives
\[
 D(x)\bigl(P(x)+C_n(x)\bigr)=D(x^2-2).
\]
The left side has degree \(n+k\), whereas the right side has degree \(2k\).
These degrees cannot agree when \(k<n\). Thus \(D=0\).
All solutions are
\[
 P=-1,\qquad P(x)=2T_n(x/2)\quad(n\ge0),
\]
where \(T_0=1\) includes the constant solution \(2\).
This final degree comparison proves completeness of the classification.



% ---- solutions/combinatorics ----
\section{Combinatorics}

\subsection*{Counting and double counting}
\begin{description}[style=nextline]
\item[\ref{c:ex01}]
The five images can be chosen independently, giving \(3^5=243\) maps.
No injection exists from a five-element set into a three-element set:
five distinct images would be needed. An injection \([3]\to[5]\)
has \(5\cdot4\cdot3=60\) choices for its ordered images.

\item[\ref{c:ex02}]
The first count is \(9!/(4!3!2!)=1260\). Once the first and last
positions contain \(C\), the seven middle positions contain four \(A\)'s
and three \(B\)'s, giving \(7!/(4!3!)=35\).
The identical \(C\)'s introduce no extra factor of two.

\item[\ref{c:ex03}]
The first answer is \(\binom{14}{2}=91\).
Write \(x=2+x', y=3+y', z=1+z'\); the new nonnegative total is six,
so the second answer is \(\binom82=28\).

\item[\ref{c:ex04}]
For \(k=0\) there is one subset. For \(k\ge1\), write the chosen
positions \(a_1<\cdots<a_k\), with \(a_{i+1}\ge a_i+2\).
Then \(b_i=a_i-(i-1)\) satisfy
\(1\le b_1<\cdots<b_k\le n-k+1\).
Conversely, any such \(b_i\) give legal \(a_i=b_i+i-1\).
Thus the count is \(\binom{n-k+1}{k}\) when \(n-k+1\ge k\);
it is zero when \(2k>n+1\). Negative \(k\) is not a subset size.

\item[\ref{c:ex05}]
There are \(n+1\) gaps before, between, and after the positions.
For each run of ones, choose the gap just before it starts and the gap
just after it ends. In increasing order the \(2k\) chosen gaps alternate
start, end, start, end; this reconstructs a unique word. Hence the answer
is \(\binom{n+1}{2k}\) if \(0\le2k\le n+1\), with the interpretation
one all-zero word for \(k=0\); otherwise it is zero.
Adjacent chosen boundary gaps create a run of length one and cause no
problem.

\item[\ref{c:ex06}]
Fix one person, giving \((n-1)!\) circular orders.
For \(n\ge3\), merge the two specified people into one block.
The \(n-1\) distinct units have \((n-2)!\) circular orders, and the
block has two internal orders, giving \(2(n-2)!\).
When \(n=2\), there is one circular order, in which the pair is adjacent;
the block formula would count the same order twice and must not be used.

\item[\ref{c:ex07}]
Every element has exactly three states: outside \(B\), in \(B\setminus A\),
or in \(A\). Thus there are \(3^n\) pairs. If \(|B|=k\), there are
\(\binom nk\) ways to choose \(B\) and \(2^k\) ways to choose
\(A\subseteq B\). Summing over the disjoint possibilities for \(k\)
proves the identity.

\item[\ref{c:ex08}]
Under the union condition, each element is in \(A\setminus B\),
\(B\setminus A\), or \(A\cap B\). There are \(3^n\) words in these
three states. Words avoiding the third state number \(2^n\), so the
answer is \(3^n-2^n\). Alternatively, choose the nonempty intersection
of size \(j\), then assign each remaining element to exactly one set:
\(\sum_{j=1}^n\binom nj2^{n-j}=3^n-2^n\).

\item[\ref{c:ex09}]
A subset with a marked element is specified by choosing that element
in \(n\) ways and any subset of the other \(n-1\) elements.
Thus the first sum is \(n2^{n-1}\) for \(n\ge1\), and zero for \(n=0\).
For an ordered pair of different marked elements, the answer is
\(n(n-1)2^{n-2}\) for \(n\ge2\), and zero for \(n=0,1\).
Stating these small cases avoids an unnecessary negative exponent.

\item[\ref{c:ex10}]
Inclusion--exclusion gives
\(3^5-3\cdot2^5+3\cdot1^5=150\).
The recurrence gives
\(S(3,2)=3,\ S(3,3)=1,\ S(4,2)=7,\ S(4,3)=6\), hence
\(S(5,3)=S(4,2)+3S(4,3)=25\). Label the three blocks in \(3!=6\)
ways to obtain the same \(150\).

\item[\ref{c:ex11}]
Without upper bounds there are \(\binom{13}{3}=286\) solutions.
A specified coordinate at least five leaves total five after subtraction,
giving \(\binom83=56\) solutions; a specified pair at least five leaves
total zero, giving one. Three bad coordinates are impossible.
The answer is \(286-4\cdot56+\binom42=68\).

\item[\ref{c:ex12}]
Selecting \(n\) elements from two \(n\)-element rows gives
\(\sum_k\binom nk\binom n{n-k}=\binom{2n}{n}\).
For the weighted sum, mark one chosen element in the first row:
choose it in \(n\) ways and choose the other \(n-1\) selected elements
from the remaining \(2n-1\) positions. Thus the sum equals
\(n\binom{2n-1}{n-1}\), also \(\frac n2\binom{2n}{n}\).

\item[\ref{c:ex13}]
In a derangement, let \(j=\pi(1)\ne1\). If \(\pi(j)=1\), remove the
two-cycle \((1,j)\), leaving a derangement of \(n-2\) elements.
Otherwise remove \(1\) from its cycle: replace the arrow leading to \(1\)
by an arrow to \(j\). This gives a derangement of the other \(n-1\)
elements, and inserting \(1\) immediately before \(j\) reverses it.
Therefore \(D_n=(n-1)(D_{n-1}+D_{n-2})\).
Starting with \(D_0=1,D_1=0\), induction gives even \(D_n\) for odd
\(n\): the multiplier \(n-1\) is even. For even \(n\), the two previous
values have parities even and odd, so the result is odd.

\item[\ref{c:ex14}]
Every triple contains three pairs. Since no pair is repeated,
the \(3|\mathcal F|\) resulting pairs are distinct members of the
\(\binom n2\) possible pairs. This proves the bound.

\item[\ref{c:ex15}]
If \(0\le k\le n\), choose the \(k\) represented couples and then one
spouse from each: \(\binom nk2^k\). If \(k>n\), the answer is zero.
For inclusion--exclusion, begin with all \(\binom{2n}{k}\) committees.
Prescribing \(j\) complete couples uses \(2j\) members, and leaves
\(\binom{2n-2j}{k-2j}\) ways to choose the others. Hence
\[
 \sum_{j=0}^{\min(n,\lfloor k/2\rfloor)}
 (-1)^j\binom nj\binom{2n-2j}{k-2j}
 =
 \begin{cases}\binom nk2^k,&0\le k\le n,\\0,&n<k\le2n.\end{cases}
 \]
Both sides count exactly the committees with no complete couple.

\item[\ref{c:ex16}]
Choose \(n\) positions from two rows of \(n\) positions and assign sign
\((-1)^k\), where \(k\) is the number chosen from the first row.
The total signed count is the proposed sum. If a column contains exactly
one chosen position, swap that choice between its rows in the first
such column. The first such column remains the same, so this is an
involution; its sign changes. Unpaired selections use either both
positions or neither in each column. There are none when \(n\) is odd.
When \(n=2m\), choose the \(m\) full columns in \(\binom{2m}{m}\)
ways; all have sign \((-1)^m\). Thus the answer is zero for odd \(n\)
and \((-1)^{n/2}\binom n{n/2}\) for even \(n\).
\end{description}

\subsection*{Recurrences and pigeonholes}
\begin{description}[style=nextline]
\item[\ref{c:ex17}]
Let \(f_n\) be the count and set \(f_0=1\) for the empty composition.
For \(n\ge2\), the first part is one or two, leaving a composition of
\(n-1\) or \(n-2\). Hence \(f_n=f_{n-1}+f_{n-2}\), with \(f_1=1\).
These are the strip numbers \(t_n\); for example \(f_5=8\).

\item[\ref{c:ex18}]
The empty word gives \(u_0=1\), and \(u_1=2,u_2=4\).
For \(n\ge3\), classify by the final run of ones. It has length zero,
one, or two, so the suffix is \(0\), \(01\), or \(011\), following
an arbitrary legal word of length \(n-1,n-2,n-3\), respectively.
Thus \(u_n=u_{n-1}+u_{n-2}+u_{n-3}\).
Words consisting entirely of ones occur only among the initial cases.

\item[\ref{c:ex19}]
Replace an opening parenthesis by \(U\) and a closing one by \(D\).
At each prefix the height is the number of unmatched opening parentheses.
If it is negative, a closing parenthesis cannot be matched.
If every height is nonnegative and the final height is zero, match
each closing parenthesis with the most recent still-unmatched opening
parenthesis; this process is well defined and matches everything.
Thus the map is a bijection between balanced words and Dyck paths,
and the count is \(C_n\).

\item[\ref{c:ex20}]
The coefficient of \(x^{n-1}\) in \(C(x)^2\) is
\(\sum_{i=0}^{n-1}C_iC_{n-1-i}=C_n\).
The constant coefficient is \(C_0=1\), proving the equation.
Reading coefficients successively gives
\(C_1=1,C_2=2,C_3=5,C_4=14,C_5=42\).
The equation has a unique formal solution with constant term one,
since each next coefficient depends only on earlier ones.

\item[\ref{c:ex21}]
The series is \((1-x)^{-1}(1-x^2)^{-1}(1-x^3)^{-1}\).
For \(c=0,1,2\), the remaining equation is \(a+2b=8-3c\),
with respectively \(5,3,2\) choices for \(b\). Thus the coefficient is
ten. Values \(c\ge3\) are impossible.

\item[\ref{c:ex22}]
There are \(\binom{r+s}{s}\) unrestricted words. Reflecting the prefix
through the first height \(-1\) changes a bad word into one with
\(r+1\) up-steps and \(s-1\) down-steps. Its endpoint has height
\(r-s+2\ge2\), so a first visit to height one exists and gives the
inverse reflection. Hence the answer is
\[
 \binom{r+s}{s}-\binom{r+s}{s-1}
 =\frac{r-s+1}{r+1}\binom{r+s}{s}.
 \]
For \(s=0\), the subtracted binomial coefficient is zero.

\item[\ref{c:ex23}]
For \(k=0\) there is one all-zero coloring. For \(k\ge1\), mark one
of the ones. Choose its labelled position in \(n\) ways, then list the
numbers of zeros after successive ones, reading clockwise from the mark.
These are \(k\) positive integers with sum \(n-k\). There are
\(\binom{n-k-1}{k-1}\) lists when \(n\ge2k\), and none otherwise.
Each unmarked coloring occurs \(k\) times, so the answer is
\[
 \frac nk\binom{n-k-1}{k-1}
 =\frac n{n-k}\binom{n-k}{k}\quad(1\le k\le\lfloor n/2\rfloor).
 \]
The gap description reconstructs the coloring uniquely; that fact
justifies the multiplication and division.

\item[\ref{c:ex24}]
Formally,
\[
 \prod_{j\ge1}(1+x^j)
 =\prod_{j\ge1}\frac{1-x^{2j}}{1-x^j}
 =\prod_{\substack{j\ge1\\j\ {\rm odd}}}\frac1{1-x^j}.
 \]
For any fixed coefficient, all manipulations can be truncated to
finitely many relevant factors, so cancellation is justified.
For a bijection, take a partition into odd parts. If odd part \(u\)
occurs \(m=\sum_t\epsilon_t2^t\) times, replace its copies by the
distinct parts \(2^tu\) with \(\epsilon_t=1\).
Every positive integer has a unique odd part and power of two,
so the resulting parts are distinct, and splitting them back reverses
the construction. Sums are preserved.

\item[\ref{c:ex25}]
Two of six integers occupy the same one of five residue classes.
Their difference is divisible by five. The integers \(0,1,2,3,4\)
show that five do not suffice.

\item[\ref{c:ex26}]
There are four parity pairs \((0,0),(0,1),(1,0),(1,1)\).
Two points share a pair, so both sums of corresponding coordinates
are even. Their midpoint is integral. Distinctness of points is not
needed; coincident points already have an integral midpoint.

\item[\ref{c:ex27}]
Use \(s_0,\ldots,s_n\) modulo \(n\) as in the chapter.
An equal pair yields a nonempty consecutive block, which is in
particular a nonempty subset of the indexed terms. Repeated values
among the original integers cause no difficulty.

\item[\ref{c:ex28}]
The \(n\) boxes \(\{1,2\},\{3,4\},\ldots,\{2n-1,2n\}\)
partition \([2n]\). A choice of \(n+1\) numbers includes both members
of one box. Consecutive integers are coprime: a common divisor would
divide their difference one.

\item[\ref{c:ex29}]
Use
\[
4,3,2,1,\quad8,7,6,5,\quad12,11,10,9,\quad16,15,14,13.
 \]
An increasing subsequence has at most one term per block; a decreasing
one remains inside a block. Both lengths are at most four.
Theorem~\ref{c:es}, with \(r=s=5\), forces a length-five monotone
subsequence among \(4\cdot4+1=17\) distinct numbers.

\item[\ref{c:ex30}]
Let \(s_j=a_1+\cdots+a_j\), with \(s_0=s_n=0\), and choose
\(i\in\{0,\ldots,n-1\}\) where the prefix sum is minimal.
Rotate to \(a_{i+1},\ldots,a_n,a_1,\ldots,a_i\).
Before wrapping, a partial sum is \(s_j-s_i\ge0\); after wrapping,
it is \((s_n-s_i)+s_j=s_j-s_i\ge0\). Thus every partial sum is
nonnegative. A strict minimum is unnecessary.

\item[\ref{c:ex31}]
For \(j=0,\ldots,M^d\), form the vector
\((\{j\alpha_1\},\ldots,\{j\alpha_d\})\) in the half-open unit cube.
Divide each coordinate into \(M\) half-open intervals, giving \(M^d\)
boxes. Two vectors, with indices \(j<k\), share a box. Put \(q=k-j\).
For each coordinate their difference has absolute value less than
\(1/M\) and equals \(q\alpha_i-p_i\) for an integer \(p_i\).
Also \(1\le q\le M^d\), as required.

\item[\ref{c:ex32}]
For every \(N\ge1\), pigeonhole gives \(1\le q\le N\) and an integer
\(p\) with \(|q\alpha-p|<1/N\). Therefore
\(|\alpha-p/q|<1/(qN)\le1/q^2\).
If a fraction is not reduced, reduce it first: this keeps \(q\le N\)
and can only decrease \(|q\alpha-p|\).
If only finitely many rational numbers arose, their positive distances
from irrational \(\alpha\), multiplied by their chosen positive
denominators after reduction, would have a positive lower bound.
More directly, bounded denominators \(q\le Q\) imply the positive
minimum \(\delta=\min_{1\le q\le Q}\operatorname{dist}(q\alpha,\ZZ)>0\).
Taking \(N>1/\delta\) forces a denominator above \(Q\).
Thus denominators are unbounded and infinitely many distinct rationals
satisfy the claimed inequality. Reducing a fraction only decreases its
denominator and preserves the inequality \(<1/q_{\rm reduced}^2\).
\end{description}
\subsection*{Invariants and extremal choices}
\begin{description}[style=nextline]
\item[\ref{c:ex33}]
Any pair of squares of the same color is ruled out: removing them
creates unequal black and white counts. Removing opposite colors keeps
the counts equal, so the invariant no longer rules out a tiling.
It gives no construction and cannot prove sufficiency for an arbitrary
region. For example, two isolated, nonadjacent unit squares of opposite
colors have balanced counts but cannot be tiled by dominoes.

\item[\ref{c:ex34}]
The reachable configurations are exactly those with an even number of
heads. Necessity follows because every move changes the head count
by \(0\) or \(2\) in absolute value. Conversely, take an even-head
configuration and correct it to all tails: if position one is a head,
flip positions one and two; then correct position two using two and three,
and continue. Previously corrected positions stay fixed. At the end the
last position is a tail because parity is even and every other coin is
a tail. Reversing these moves constructs the configuration from all tails.
For \(n=1\), the same classification leaves only all tails.

\item[\ref{c:ex35}]
The invariant is
\(\prod_{j=1}^n(1+j)=2\cdot3\cdots(n+1)=(n+1)!\).
If the last entry is \(z\), then \(1+z=(n+1)!\), so
\(z=(n+1)!-1\).

\item[\ref{c:ex36}]
For integers \(a,b,c\),
\(\gcd(a,b,c)=\gcd(a\pm b,b,c)\), with gcd taken nonnegative.
A common divisor on one side divides all entries on the other by
subtraction, proving equality. Thus every allowed move preserves the
gcd. It starts at one and would have to end at two, an impossibility.

\item[\ref{c:ex37}]
The coordinate sum changes by an odd number at every knight move,
so its parity reverses each time. After an odd number of moves the
knight is on the opposite color and cannot be at its initial square.

\item[\ref{c:ex38}]
The sorted permutation has zero inversions. Since one adjacent swap
can reduce the count by at most one, at least \(I\) swaps are needed
from a permutation with \(I\) inversions. While the permutation is
unsorted, some adjacent pair is inverted: otherwise all adjacent
inequalities, and therefore the whole ordering, would be increasing.
Swapping such a pair reduces \(I\) by exactly one. Repeating reaches
the sorted order in exactly \(I\) swaps, proving both bounds.

\item[\ref{c:ex39}]
The counts of residues \(0,1,2\) among the coordinate values
\(0,\ldots,7\) are \(3,3,2\). Convolving these counts, the numbers of
squares with \(i+j\equiv0,1,2\pmod3\) are \(21,22,21\).
Removing \((0,0)\) changes them to \(20,22,21\).
Every horizontal or vertical straight triomino has one square of each
color, so a tiling would require equal counts. The unequal counts
give a contradiction even though the remaining area, 63, is divisible
by three.

\item[\ref{c:ex40}]
Let \(\omega\) be a primitive complex \(k\)th root of unity, meaning
\(\omega^k=1\) and \(\omega^j\ne1\) for \(1\le j<k\).
Assign square \((i,j)\), \(0\le i<m,\ 0\le j<n\), weight
\(\omega^{i+j}\). Each allowed tile has total weight a power of
\(\omega\) times \(1+\omega+\cdots+\omega^{k-1}=0\).
Hence tileability forces
\[
 0=\left(\sum_{i=0}^{m-1}\omega^i\right)
   \left(\sum_{j=0}^{n-1}\omega^j\right).
 \]
In the complex numbers a product is zero only if a factor is zero.
Since \(\omega\ne1\), the first geometric sum is zero exactly when
\(\omega^m=1\), equivalently \(k\mid m\); similarly for \(n\).
Conversely, if \(k\) divides one side length, tile each parallel strip
in blocks of length \(k\). This proves necessity and sufficiency.

\item[\ref{c:ex41}]
Choose a partition maximizing the number of crossing edges.
If a vertex has \(a\) neighbors on its side and \(b<a\) across,
moving it to the other side changes the crossing count by \(a-b>0\).
This contradicts maximality. The maximum exists because there are
finitely many vertex partitions.

\item[\ref{c:ex42}]
Let \(v\) have maximum outdegree. If a vertex \(w\) is neither beaten
by \(v\) nor reachable through an outneighbor of \(v\), then \(w\to v\)
and \(w\) beats every outneighbor of \(v\). Its outdegree is therefore
at least \(d^+(v)+1\), a contradiction.
Thus \(v\) reaches every other vertex in one or two steps.

\item[\ref{c:ex43}]
Among all pairs \((P,\ell)\), where \(\ell\) is a line through two
points of \(S\) and \(P\in S\setminus\ell\), choose one of smallest
positive distance \(h\). Such pairs exist because \(S\) is not collinear,
and there are finitely many. Let \(H\) be the perpendicular foot from
\(P\) to \(\ell\). If \(\ell\) contained at least three points of \(S\),
choose two, \(A,B\), on the same closed ray of \(\ell\) starting at
\(H\), with \(A\) nearer \(H\). Then
\(|AB|\le|HB|<|PB|\).
Computing the area of triangle \(PAB\) with bases \(AB\) and \(PB\),
the distance from \(A\) to line \(PB\) is
\(|AB|h/|PB|<h\). It is positive, because \(P\notin\ell\).
This contradicts the chosen minimum. Thus \(\ell\) contains exactly
two points of \(S\).

\item[\ref{c:ex44}]
Let the special edge be \(e=uv\), and suppose a minimum tree \(T\)
omits it. The unique \(u\)-to-\(v\) path in \(T\), together with
\(e\), forms a cycle. Because \(u,v\) are on opposite sides of the
specified partition, that path contains another crossing edge \(f\).
Removing \(f\) after adding \(e\) restores a spanning tree. Its
weight is smaller, since \(w(e)<w(f)\), contradicting minimality.
The uniqueness of a tree path follows because two different paths
would contain a cycle.

\item[\ref{c:ex45}]
The intervals \([0,10]\), \([1,2]\), and \([8,9]\) suffice.
The longest-first rule selects only \([0,10]\), but the two shorter
intervals are disjoint. Replacing an optimal first interval by the
longest one can create intersections with later intervals; it does
not preserve the solution class. The earliest-finish proof relied
precisely on the replacement ending no later.

\item[\ref{c:ex46}]
Among the finitely many cyclic orders of the points, choose one of
minimum perimeter. If nonadjacent sides \(AB\) and \(CD\) cross at
an interior point \(X\), take their order along the cycle to be
\(A,B,\ldots,C,D,\ldots,A\).
Delete these sides, leaving paths from \(B\) to \(C\) and from \(D\)
to \(A\). Reconnect using \(AC\) and \(BD\); this makes one cycle
through all points, with the first path traversed in reverse.
Strict triangle inequalities give
\(|AC|+|BD|<|AB|+|CD|\), a contradiction.
No vertex lies on another side, because no three points are collinear.
Hence the minimizing order forms the required simple polygon.

\item[\ref{c:ex47}]
Let the occupancies be \(a_1,\ldots,a_m\), all nonnegative integers,
with sum \(N\). Moving from a nonempty box of occupancy \(a\)
to one of occupancy \(b\ge a\) changes the sum of squares by
\[
 (a-1)^2+(b+1)^2-a^2-b^2=2(b-a)+2>0.
 \]
That sum is an integer at most \(N^2\), since
\((\sum_i a_i)^2\ge\sum_i a_i^2\). Therefore every process terminates.
If two boxes remained nonempty, a move from the less occupied to the
more occupied would still be allowed. Thus a terminal state has all
tokens in one box.

\item[\ref{c:ex48}]
Strong connectivity supplies a directed cycle; a shortest one has
length three, since a chord of a longer cycle produces a shorter cycle
or a triangle. Choose a longest directed cycle
\(C=(v_1,\ldots,v_k)\).
If an outside vertex \(x\) has arrows in both directions between it
and \(C\), some consecutive pair satisfies
\(v_i\to x\to v_{i+1}\); inserting \(x\) makes a longer cycle.
Consequently every outside vertex is of one of two types:
\(U\), beaten by every vertex of \(C\), or \(W\), beating every
vertex of \(C\).
If any outside vertices exist, strong connectivity requires both
types and some arrow \(u\to w\) with \(u\in U,w\in W\).
Indeed, without such an arrow a vertex in \(U\) could never leave
\(U\) to reach \(C\); if \(U\) or \(W\) were empty, one direction
between \(C\) and the outside would likewise be impossible.
Replace an edge \(v_i\to v_{i+1}\) by
\(v_i\to u\to w\to v_{i+1}\), obtaining a longer cycle.
This contradiction proves that the longest cycle uses all vertices.
\end{description}

\subsection*{Graphs, matching, and order}
\begin{description}[style=nextline]
\item[\ref{c:ex49}]
If a vertex has degree zero, no vertex has degree \(n-1\); otherwise
degree zero is absent. Thus the \(n\) degrees occupy at most \(n-1\)
of the possible values \(0,\ldots,n-1\). Two must coincide.
The graph being simple is essential for this degree range.

\item[\ref{c:ex50}]
Both endpoints of a longest path have no neighbor off the path,
or it could be extended. Neither has another neighbor on the path,
or a cycle would result. Hence both have degree one.
A path graph has exactly its two endpoints as leaves, so two is possible.

\item[\ref{c:ex51}]
Add a new vertex \(w\) joined to each of the \(2k\) odd-degree vertices.
All old degrees become even and \(d(w)=2k\) is even. The augmented
graph is connected, so it has an Euler circuit. Splitting this circuit
at its visits to \(w\), and removing the \(2k\) added edges, leaves
\(k\) trails partitioning all original edges. No resulting trail is
empty: each old odd vertex has only one edge to \(w\), so arrival from
\(w\) cannot be followed immediately by a different edge back to \(w\).

\item[\ref{c:ex52}]
Take a partition maximizing crossing edges. As in
Problem~\ref{c:ex41}, each vertex has at least half its neighbors across.
Summing crossing degrees gives at least half the total degree, so the
crossing edge count is at least \(e/2\). Keep only those edges to obtain
the desired bipartite subgraph.

\item[\ref{c:ex53}]
In every nonempty subgraph, at least one degree is strictly below
\(2d\), hence at most \(2d-1\). Repeated removal and reverse greedy
coloring use \(2d\) colors.
The complete graph \(K_{2d}\) satisfies the average-degree assumption
for every subgraph, but requires \(2d\) colors. For \(d\ge2\),
this exceeds \(d+1\). Thus the stronger conclusion is false in general.

\item[\ref{c:ex54}]
For \(n\ge3\), the average degree is
\(2e/n\le6-12/n<6\), so some vertex has degree at most five.
For fewer vertices the statement is immediate. Delete such a vertex,
six-color the remaining planar graph inductively, and give the vertex
a color not used by its at most five neighbors.
This proves the six-color bound; it does not establish the stronger
four-color theorem.

\item[\ref{c:ex55}]
The graph is connected: any component has at least \(n/2+1\) vertices,
so two components cannot fit. Take a longest path \(v_1,\ldots,v_k\).
Every neighbor of either endpoint lies on it. Define
\[
 A=\{i: v_1v_{i+1}\in E,\ 1\le i<k\},\qquad
 B=\{i: v_iv_k\in E,\ 1\le i<k\}.
 \]
Their sizes sum to \(d(v_1)+d(v_k)\ge n\ge k\), exceeding the
\(k-1\) available indices. Choose \(i\in A\cap B\).
Then
\[
 v_1,v_2,\ldots,v_i,v_k,v_{k-1},\ldots,v_{i+1},v_1
 \]
is a cycle through all \(k\) path vertices. If \(k<n\), connectivity
gives an edge from an outside vertex to this cycle (take the first
crossing along a connecting path). Start at that outside vertex and
then traverse the cycle without closing it, producing a path with
\(k+1\) vertices. This contradicts maximality. Hence \(k=n\).

\item[\ref{c:ex56}]
Every three-vertex tournament is either a directed triangle or transitive.
In a transitive triple exactly one vertex beats the other two.
Thus transitive triples are counted once by
\(\sum_i\binom{d_i}{2}\), giving the formula.
The outdegrees sum to \(n(n-1)/2\), and
\[
 \sum_i\left(d_i-\frac{n-1}{2}\right)^2\ge0
 \]
shows that \(\sum_i\binom{d_i}{2}\) is minimized when the degrees
are equal, if such a tournament exists. For odd \(n\), label vertices
by residues modulo \(n\) and direct \(i\to j\) when
\(j-i\pmod n\in\{1,\ldots,(n-1)/2\}\).
This gives all outdegrees \((n-1)/2\), and the largest triangle count is
\[
 \binom n3-n\binom{(n-1)/2}{2}=\frac{n(n^2-1)}{24}.
 \]

\item[\ref{c:ex57}]
Assign jobs \(1,2,3\) in that order to the three applicants.
If the third list is changed to \(\{1,2\}\), all three applicants
together have only two neighboring jobs, violating Hall.
Each single applicant still has two options; checking only individuals
would miss the obstruction.

\item[\ref{c:ex58}]
The edge count and Hall argument in the chapter give a perfect matching.
Delete its edges. Every vertex now has degree \(d-1\), in the same
two equally sized parts. If \(d=1\), nothing remains; otherwise induct.
The removed matchings are disjoint and exhaust all edges.

\item[\ref{c:ex59}]
For a nonempty \(S\subseteq L\) of size at most \(n/2\), a vertex
of \(S\) already has at least \(n/2\) neighbors, so
\(|N(S)|\ge|S|\).
If \(|S|>n/2\) and Hall failed, then \(|N(S)|<|S|\le n\).
Choose \(y\in R\setminus N(S)\). All its neighbors lie in
\(L\setminus S\), a set of size \(n-|S|<n/2\), contradicting the
minimum degree. Thus Hall holds for all \(S\).

\item[\ref{c:ex60}]
Construct a bipartite graph whose left vertices are rows, right vertices
are columns, and edges are allowed squares. A nonattacking rook set is
exactly a matching: no two use a row or column.
A set of rows and columns covering every allowed square is exactly a
vertex cover. Konig's theorem equates their optimum sizes.
Rows or columns with no allowed squares may remain as isolated vertices.

\item[\ref{c:ex61}]
The \(n\) integers \(n+1,\ldots,2n\) form an antichain: a larger
multiple of one would exceed \(2n\). Partition \([2n]\) by odd part:
for each odd \(q\le2n\), take \(q,2q,4q,\ldots\) up to \(2n\).
There are \(n\) chains. Every antichain meets each in at most one
element, so its size is at most \(n\). Hence the width is \(n\).

\item[\ref{c:ex62}]
Let \(a_k\) be the number of size-\(k\) sets in the family.
Counting permutation-chain incidences as in LYM now gives
\(\sum_k a_k/\binom nk\le r\), since each chain meets the family
at most \(r\) times. Also \(0\le a_k\le\binom nk\).
Reorder the \(n+1\) binomial coefficients as
\(b_1\ge\cdots\ge b_{n+1}>0\), and let \(x_i\) denote the
corresponding level counts. With \(b=b_r\),
\[
 \begin{split}
 \sum_i x_i
 &\le \sum_{i\le r}x_i+b\sum_{i>r}\frac{x_i}{b_i}\\
 &\le rb+\sum_{i\le r}x_i(1-b/b_i)
 \le rb+\sum_{i\le r}b_i(1-b/b_i)
 =\sum_{i\le r}b_i.
 \end{split}
 \]
The first inequality uses \(b_i\le b\) for \(i>r\).
For attainment, take all subsets in \(r\) size levels with largest
binomial coefficients. A strict inclusion chain uses different sizes,
so it contains at most \(r\) members.

\item[\ref{c:ex63}]
Give \(x\) the label \(h(x)\), the maximum length of a chain ending
at \(x\). If \(x<y\), appending \(y\) gives \(h(y)>h(x)\).
Thus each label class is an antichain. If the longest chain has length
\(h\), these classes partition the poset into \(h\) antichains.
Conversely, any antichain partition needs at least \(h\) classes,
since a chain meets each class at most once.
Apply this to positions ordered by increasing index and increasing
value. Its chains are increasing subsequences, and its antichains
are decreasing subsequences, proving the stated deduction.

\item[\ref{c:ex64}]
Call a pair allowable if it occurs in at least one stable matching.
Suppose the algorithm first rejects an allowable pair \((a,B)\).
At that time college \(B\) holds an applicant \(a'\) whom it prefers
to \(a\). Fix a stable matching \(S\) containing \((a,B)\).
Applicant \(a'\) has already been rejected by every college preferred
to \(B\). By the choice of the first rejected allowable pair, none of
those earlier rejected pairs can occur in \(S\).
Also \(a'\) cannot have \(B\) in \(S\), since \(B\) has \(a\).
Therefore \(a'\) prefers \(B\) to the partner assigned in \(S\).
College \(B\) prefers \(a'\) to its partner \(a\), forming a blocking
pair for \(S\), a contradiction.
Thus no allowable pair is ever rejected. An applicant proposes in
preference order and ends at the first college that never rejects them;
every preferred college has rejected them and is therefore unavailable
in any stable matching. The final partner is their favorite stable partner.
\end{description}

\subsection*{Extremal and probabilistic methods}
\begin{description}[style=nextline]
\item[\ref{c:ex65}]
Mantel gives at most \(\lfloor49/4\rfloor=12\) edges.
The complete bipartite graph with parts of sizes three and four has
twelve edges and no triangle, so the maximum is twelve.

\item[\ref{c:ex66}]
The bound with \(r=3,n=12\) is
\(\frac{2}{6}\cdot12^2=48\).
Partition the vertices into three parts of size four and join exactly
the pairs in different parts. It has \(3\cdot4\cdot4=48\) edges.
A clique takes at most one vertex from each part, so no \(K_4\) exists.

\item[\ref{c:ex67}]
The recurrence gives \(R(3,4)\le R(2,4)+R(3,3)\le4+6=10\).
For the lower bound, label the eight vertices modulo eight.
Red differences are \(\pm1\) and \(4\).
There is no red triangle: the neighbors of a vertex,
its predecessor, successor, and opposite vertex, have no red edge
between them. A blue \(K_4\) would be an independent four-set in the
red graph. In the eight-cycle alone, four independent vertices must
alternate: the four cyclic gaps each need at least one unchosen vertex,
and there are only four such vertices. Each alternating set contains
opposite-vertex red edges, so neither is independent in the full red
graph. Hence the coloring has neither forbidden configuration,
showing \(R(3,4)>8\). These arguments give the requested bounds,
without determining which of nine or ten is the exact value.

\item[\ref{c:ex68}]
Choose a uniform random permutation of the vertices.
Keep \(v\) if it precedes all its neighbors.
Two adjacent vertices cannot both be kept, so the result is independent.
Among \(v\) and its \(d(v)\) neighbors, each is equally likely to come
first, so \(P(v\text{ kept})=1/(d(v)+1)\).
Linearity and existence above the average give the first bound.
Cauchy applied to the positive numbers \(d(v)+1\) yields
\[
 \left(\sum_v\frac1{d(v)+1}\right)
 \left(\sum_v(d(v)+1)\right)\ge n^2.
 \]
The second factor is \(2e+n=n(\bar d+1)\), giving
\(\sum_v1/(d(v)+1)\ge n/(\bar d+1)\).

\item[\ref{c:ex69}]
Let column \(j\) contain \(d_j\) ones. A pair of rows can contain
ones in at most one common column, or there would be a forbidden
rectangle. Thus
\[
 \sum_{j=1}^n\binom{d_j}{2}\le\binom n2.
 \]
Since \(\sum_jd_j=m\) and \(\sum_jd_j^2\ge m^2/n\),
we get \(m^2/n-m\le n(n-1)\).
The nonnegative solution range of this quadratic is
\(m\le\frac n2(1+\sqrt{4n-3})\), as claimed.
The hypotheses concern different rows and columns, exactly matching
the counted pairs.

\item[\ref{c:ex70}]
Put \(N=\lfloor2^{k/2}\rfloor\), and independently color each edge
red or blue with equal probability. A fixed \(k\)-set is monochromatic
with probability \(2^{1-\binom k2}\).
Let \(X\) be the number of monochromatic \(k\)-sets. Then
\[
 \mathbb E X=\binom Nk2^{1-\binom k2}
 \le\frac{N^k}{k!}2^{1-k(k-1)/2}
 \le\frac{2^{1+k/2}}{k!}<1.
 \]
For \(k=3\), the last inequality is \(4\sqrt2<6\);
increasing \(k\) multiplies the ratio by \(\sqrt2/(k+1)<1\).
If every coloring had a monochromatic \(k\)-set, the integer \(X\)
would always be at least one, contradicting the expectation.
Hence some coloring has \(X=0\).

\item[\ref{c:ex71}]
Color each ground-set element independently and uniformly.
For each ten-set \(A_i\), let \(E_i\) be its being monochromatic.
Then \(P(E_i)=2/2^{10}=1/512\).
It is independent of every joint condition on events from disjoint
sets; at most one hundred others intersect it.
Choose \(x=1/101\). The binomial expansion gives
\[
 (1+1/100)^{100}
 \le\sum_{j=0}^{100}\frac1{j!}<3.
 \]
The last bound follows from \(j!\ge2^{j-1}\) for \(j\ge1\), with
strict inequality from \(j=3\), by comparison with a geometric series.
Consequently
\[
 x(1-x)^{100}
 =\frac1{101}(100/101)^{100}
 >\frac1{303}>\frac1{512}.
 \]
The local lemma gives positive probability that no bad event occurs,
which proves the desired coloring exists.

\item[\ref{c:ex72}]
Assign every variable independently true or false with probability
one half. The three literals of a clause involve different variables,
so the probability all three are false is \(1/8\), regardless of their
negations. The expected number of satisfied clauses is \(7m/8\),
so some assignment attains at least that many.
For an algorithm, fix variables one at a time. Conditional on the
current choices, a clause already true has success probability one;
one whose all assigned literals are false and has \(r>0\) unassigned
literals has success probability \(1-2^{-r}\); a fully assigned false
clause has probability zero. Sum these probabilities.
For each next variable try both truth values and keep one for which
the sum does not decrease. The current sum is the mean of those two
options, so such a choice always exists. After all assignments the sum
is the actual number satisfied and remains at least \(7m/8\).
\end{description}

\subsection*{Symmetry and games}
\begin{description}[style=nextline]
\item[\ref{c:ex73}]
Before identification there are \(2^p\) colorings.
The identity fixes all of them. Every nonidentity rotation is a single
cycle on the positions because \(p\) is prime, so fixes only the
two constant colorings. Burnside gives
\((2^p+2(p-1))/p=2+(2^p-2)/p\) necklaces.
This is an integer, hence \(p\mid(2^p-2)\).
For \(p=2\) the same formula and reasoning remain valid.

\item[\ref{c:ex74}]
The square's eight symmetries have position-cycle counts
\(4,1,1,2,3,3,2,2\).
Thus Burnside gives
\[
 \frac{3^4+2\cdot3+3^2+2\cdot3^3+2\cdot3^2}{8}
 =\frac{168}{8}=21.
 \]
The colors have names, so exchanging all red and blue vertices is not
an additional allowed symmetry.

\item[\ref{c:ex75}]
The six rotations contribute \(84\) fixed colorings in total.
Of the six reflections, three pass through two opposite beads and
have four position cycles, fixing \(2^4=16\) each.
The other three pass through opposite gaps and have three cycles,
fixing \(2^3=8\) each. Hence the number is
\((84+3\cdot16+3\cdot8)/12=13\).

\item[\ref{c:ex76}]
Using only legal subtractions gives
\[
 (g_0,g_1,\ldots,g_{10})=(0,1,0,1,2,3,2,0,1,0,1).
 \]
For \(n\ge4\),
\(g_n=\operatorname{mex}\{g_{n-1},g_{n-3},g_{n-4}\}\).
The computed values establish \(g_{n+7}=g_n\) for \(0\le n\le3\).
For \(n\ge4\), assume it for all smaller indices; the three values
defining \(g_{n+7}\) equal those defining \(g_n\), so their mexes
are equal. Thus the period-seven pattern is
\(0,1,0,1,2,3,2\). Losing heap sizes are residues zero or two
modulo seven. Computing a few terms alone would not prove periodicity;
the finite-memory recurrence supplies the induction.

\item[\ref{c:ex77}]
If all nonempty heaps have size one, every move removes exactly one
heap, so an odd number of such heaps is losing and an even number is
winning. If some heap has size at least two, the losing positions are
exactly those with xor zero.
To prove this, a zero-xor position with a large heap has at least two
large heaps: one large heap cannot xor with only zeros and ones to zero.
Any move leaves at least one large heap and changes xor to nonzero.
Conversely, from nonzero xor with at least two large heaps, the ordinary
Nim move reaches xor zero; at least one large heap remains, so it is
in the proposed losing class.
If there is exactly one large heap and \(t\) heaps of size one, reduce
the large heap to zero when \(t\) is odd and to one when \(t\) is even.
The result has an odd number of ones, a losing position.
Finally, moving from an odd number of ones leaves an even number;
when it leaves none, the mover has taken the last token and loses
immediately. These cases verify the classification and explain the
endgame exception to ordinary Nim.

\item[\ref{c:ex78}]
The game terminates because every move removes a square.
Suppose the initial rectangle were losing for the first player.
Remove just its top-right square, a legal move because the board has
at least two squares. The resulting position must be winning for the
next player, so it has a reply \(M\) to a losing position \(P\).
But any legal move \(M\) from the original rectangle also removes its
top-right square: that corner is weakly above and to the right of
every selected square. Playing \(M\) directly from the original board
therefore leaves exactly \(P\). This would be a move from the supposed
losing initial position to a losing position, a contradiction.
Thus a winning first move exists. The proof does not say that removing
the corner is itself winning, or provide a general explicit strategy.

\item[\ref{c:ex79}]
For one heap, \(g_0=0,g_1=1,g_2=2\).
For \(n\ge2\), the followers have values \(g_{n-1},g_{n-2}\);
induction gives \(g_n=n\bmod3\), since the two other residues occur
and the current residue does not. The components form a disjoint sum
of finite impartial normal-play games, so the xor criterion for their
Grundy values gives exactly the assertion.

\item[\ref{c:ex80}]
Write \(a_k=\lfloor k\phi\rfloor\) and
\(b_k=\lfloor k\phi^2\rfloor=a_k+k\).
First prove complementarity. If irrational \(r,s>1\) satisfy
\(1/r+1/s=1\), the number of terms of the two sequences
\(\lfloor jr\rfloor,\lfloor js\rfloor\), \(j\ge1\), at most an integer
\(m\ge0\), counted with multiplicity, is
\[
 \left\lfloor\frac{m+1}{r}\right\rfloor+
 \left\lfloor\frac{m+1}{s}\right\rfloor=m.
 \]
Indeed, each fraction is noninteger and they sum to \(m+1\).
Subtracting the count at \(m-1\) shows that each positive integer
occurs exactly once across the two sequences. Apply this with
\(r=\phi,s=\phi^2\).

Let \(\mathcal P=\{(a_k,b_k):k\ge0\}\), with sorted coordinates.
No move joins two different members. A diagonal move preserves the
difference, which is \(k\), so it cannot reach a different indexed
pair. A one-heap move would preserve one coordinate, but every positive
coordinate belongs to exactly one member of \(\mathcal P\); zero
occurs only in \((0,0)\).

Now take a sorted position \((a,b)\notin\mathcal P\) and put \(d=b-a\).
If \(a>a_d\), remove \(a-a_d\) from both heaps to reach
\((a_d,b_d)\). Equality \(a=a_d\) would already give that pair.
Suppose \(a<a_d\). If \(a=0\), remove all of \(b\).
Otherwise complementarity says either \(a=a_k\) or \(a=b_k\)
for a unique \(k\ge1\).
In the first case \(k<d\), because \(a_k<a_d\), so
\(b_k=a+k<a+d=b\); reduce the larger heap to \(b_k\).
In the second case reduce the larger heap to \(a_k<a\), obtaining
the same pair after sorting. Thus every nonmember has a move into
\(\mathcal P\). The terminal pair is \((0,0)\in\mathcal P\).
Since the total number of tokens decreases, the two losing-position
conditions prove the result by induction.
\end{description}

\subsection*{Algebraic methods}
\begin{description}[style=nextline]
\item[\ref{c:ex81}]
Oddtown gives at most \(n\) members. All singleton subsets form a
family of \(n\) odd sets, with pairwise intersection size zero.
Thus the maximum is \(n\), including the empty-family answer zero
when \(n=0\).

\item[\ref{c:ex82}]
Eventown gives at most \(2^{\lfloor11/2\rfloor}=32\).
Pair ten of the elements and take all unions of these five pairs,
never using the eleventh. This gives exactly 32 sets, each with all
required intersections even.

\item[\ref{c:ex83}]
In residues modulo eight, take \(A=B=\{0,4\}\).
Then \(A+B=\{0,4\}\), of size two, below \(2+2-1=3\).
Residues modulo eight are not a field: for example, the nonzero
residue four has no inverse. Lagrange denominators involving
differences of distinct residues need not be invertible, and factorial
coefficients can vanish or be nonunits. The polynomial grid argument
therefore loses its crucial field property.

\item[\ref{c:ex84}]
The case \(k=1\) is immediate.
By Cauchy--Davenport,
\(|kA|\ge\min\{p,|(k-1)A|+|A|-1\}\).
If \((k-1)A=\FF_p\), adding any nonempty \(A\) still gives
\(\FF_p\). Otherwise use the inductive lower bound
\(|(k-1)A|\ge(k-1)(|A|-1)+1\), unless that bound already reaches
\(p\), in which case the previous case applies.
This gives the claimed minimum for every \(k\).

\item[\ref{c:ex85}]
The unrestricted path counts are
\[
 M=\begin{pmatrix}
 n+1 & \binom{n+2}{2}\\
 1 & n+1
 \end{pmatrix}.
 \]
Paths using the swapped source--sink pairing must meet.
To see this directly, at each integer horizontal coordinate \(j\),
each right-and-up path occupies a vertical interval of vertices.
If the intervals are disjoint, one is entirely below the other.
Their order cannot reverse between consecutive columns without the
horizontal steps sharing a vertex or crossing a vertical interval.
The swapped endpoints force such a reversal.
Therefore LGV counts the desired pairs without a negative contribution:
\[
 \det M=(n+1)^2-\frac{(n+2)(n+1)}2=\frac{n(n+1)}2.
 \]

\item[\ref{c:ex86}]
The set-pair inequality gives \(m/\binom{2k}{k}\le1\).
Take all \(k\)-subsets \(A\) of a fixed \(2k\)-element set and put
\(B\) equal to the complement. Different \(k\)-sets are not contained
in one another, so \(A_i\cap B_j=A_i\setminus A_j\ne\varnothing\)
for \(i\ne j\). The maximum is \(\binom{2k}{k}\).
For \(k=0\), the hypotheses allow at most one pair, also matching the
formula; an empty own pair is harmless.

\item[\ref{c:ex87}]
Consider any four of the points. If one lies inside the triangle of
the other three, no half-plane can select the inside point while
excluding all three vertices: a linear function negative at all three
vertices is negative at their positive convex combination.
Otherwise the four points form a convex quadrilateral. No half-plane
can select one opposite pair while excluding the other opposite pair:
their joining diagonals intersect, whereas a linear function would
have to be nonnegative at the intersection as a combination of the
selected pair, and negative there as a combination of the excluded
pair. The no-three-collinear hypothesis ensures these two cases
are exhaustive.
Thus no four-point set is shattered by the family of half-plane traces.
Apply Theorem~\ref{c:sauer} with \(d=3\) to obtain
\(1+n+\binom n2+\binom n3\).
For \(n<4\), the same expression is \(2^n\) under the zero convention.

\item[\ref{c:ex88}]
Let \(S=\{x+y:x,y\in A,\ x\ne y\}\).
First suppose \(2a-3\le p\). If \(|S|\le2a-4\), enlarge it to
a set \(C\subseteq\FF_p\) of size \(2a-4\).
Define
\[
 P(x,y)=(x-y)\prod_{c\in C}(x+y-c).
 \]
Choose any \(B\subseteq A\) of size \(a-1\).
Polynomial \(P\) vanishes on \(A\times B\): for equal coordinates
the first factor is zero, and for unequal coordinates the sum belongs
to \(S\subseteq C\).
Its total degree is \(2a-3=(a-1)+(a-2)\).
Its coefficient of \(x^{a-1}y^{a-2}\) is
\[
 \binom{2a-4}{a-2}-\binom{2a-4}{a-1}
 =\frac1{a-1}\binom{2a-4}{a-2}.
 \]
This is nonzero modulo \(p\): the upper factorial has argument
\(2a-4<p\), and \(1\le a-1<p\).
The coefficient test on the grid of sizes \(a,a-1\) contradicts the
vanishing. Hence \(|S|\ge2a-3\).

If \(2a-3>p\), then \(2a-p>3\).
For any \(z\in\FF_p\), the sets \(A\) and \(z-A\) have intersection
size at least \(2a-p>3\). Among these elements at most one solves
\(x=z-x\) when \(p\) is odd, so choose a different one and write
\(z=x+(z-x)\) with distinct elements of \(A\).
The large-\(a\) case cannot occur for \(p=2\).
Thus every \(z\) lies in \(S\), proving \(|S|=p\) in this case.
Both ranges give the claimed minimum.
\end{description}




% ---- solutions/convex-lattice ----
\subsection*{Combinatorics in the Plane}
\paragraph{Problem \ref{l:ex01}.}
The displacement is $(12,8)=4(3,2)$ with primitive vector $(3,2)$.
All lattice points are $(2,-1)+j(3,2)$ for $j=0,1,2,3,4$,
so there are five. If a point has displacement $t(3,2)$ with
both coordinates integers, the identity $3-2=1$ makes
$t=3t-2t$ an integer, proving no point is missing.

\paragraph{Problem \ref{l:ex02}.}
There are $(m-1)(n-1)$ interior points and $2m+2n$ boundary points.
Pick's expression is
$(m-1)(n-1)+(2m+2n)/2-1=mn$, the rectangle's area.
The formulas include $m=1$ or $n=1$, when there are no interior points.

\paragraph{Problem \ref{l:ex03}.}
Classify a lattice point by $(x\bmod2,y\bmod2)$, giving four classes.
Two of five points share a class, so their coordinate sums are even
and their midpoint is integral. For the polygon, choose any five
vertices. The midpoint is in the polygon by convexity, and cannot
be another vertex: a vertex is extreme and cannot lie strictly
between two distinct polygon points. Thus it is a nonvertex lattice
point on or inside the polygon.

\paragraph{Problem \ref{l:ex04}.}
The area is $27$, and
$B=9+6+\gcd(9,6)=18$.
Thus $I=27-18/2+1=19$.
The sloping edge contributes gcd $3$, not Euclidean length.

\paragraph{Problem \ref{l:ex05}.}
Write the intervals as $[a_i,b_i]$. Set $a=\max a_i$ and
$b=\min b_i$. Pairwise intersection of an interval attaining $a$
and one attaining $b$ gives $a\le b$. Every interval therefore
contains $[a,b]$, in particular $a$.
The proof uses the order of endpoints on one line. In the plane
there is no single largest left endpoint controlling all directions;
Helly's dimension-two threshold is three sets.

\paragraph{Problem \ref{l:ex06}.}
Take a convex representation with the fewest positive coefficients.
If there are more than three, solve the three homogeneous affine
dependence equations to obtain a nonzero $t$ with both positive
and negative entries. Let
$c=\min_{t_i>0}\lambda_i/t_i$.
For positive $t_i$, $\lambda_i-ct_i\ge0$ by the choice of $c$.
For negative $t_i$ the coefficient increases, and for zero $t_i$
it stays unchanged. The sum of coefficients and the represented
point are unchanged, while at least one positive coefficient vanishes.
This contradicts minimality. The remaining at most three points
form the required triangle or degenerate triangle.

\paragraph{Problem \ref{l:ex07}.}
Primitive edges contain only their endpoints, so $B=n$.
Pick gives $A=n/2-1=(n-2)/2$.
For even $n$ this is integral; for odd $n$ it lies halfway between
two consecutive integers. This is a necessary area formula, not
an assertion that such a polygon exists for every $n$.

\paragraph{Problem \ref{l:ex08}.}
First cover a triple. A collinear triple fits in the disk with
diameter its farthest pair, radius at most $D/2$.
For an obtuse or right triangle, the disk with its longest side
as diameter contains the third vertex, again radius at most $D/2$.
For an acute triangle, take its circumdisk. Its largest angle
$A$ satisfies $60^\circ\le A<90^\circ$ and its opposite side
$a\le D$, so the sine law gives
$R=a/(2\sin A)\le D/\sqrt3$.
Thus every triple is coverable by a disk of that radius.
Helly applied to possible-center disks covers the entire finite
set. Sets with one or two points are immediate.

For sharpness take an equilateral triangle of side $D$ and centroid $G$.
For any center $P$, the sum of its three squared vertex distances is
$3PG^2+D^2\ge D^2$, by expanding after translating $G$ to the origin.
At least one squared distance is at least $D^2/3$, so no disk of
radius smaller than $D/\sqrt3$ can cover the triangle.



% ---- solutions/geometry ----
\section{Geometry}

\subsection{Foundations, similarity, and circles}
\begin{description}[style=nextline,leftmargin=2.2em]
\item[\ref{g:ex01}]
If $AB=AC$, triangles $ABC,ACB$ are congruent by SSS, so
$\angle ABC=\angle ACB$. Conversely, if those angles are equal,
the two triangles have two corresponding angles equal and share the
corresponding side $BC$. ASA gives $AB=AC$.
The proof uses a correspondence, not the appearance of symmetry.

\item[\ref{g:ex02}]
For parallelogram $ABCD$, diagonal $AC$ creates triangles $ABC,CDA$.
Parallelism gives $\angle BAC=\angle DCA$ and
$\angle BCA=\angle DAC$, while $AC$ is common. ASA gives
$AB=CD$, $BC=DA$. If diagonals meet at $P$, the triangles
$APB,CPD$ have equal alternate angles and $AB=CD$; ASA gives
$AP=CP$, $BP=DP$.

\item[\ref{g:ex03}]
Through midpoint $M$ of $AB$ draw the parallel to $BC$, meeting
$AC$ at $N'$. AA gives $AN'/AC=AM/AB=1/2$, so $N'$ is the
midpoint $N$, and $MN/BC=1/2$. Line $BM$ is simply line $AB$.
Its intersection with the parallel through $A$ to $BC$ is the already
known point $A$; no new triangle is obtained from that attempted
construction.

\item[\ref{g:ex04}]
Let $\alpha=\angle BAM=\angle ABM$ and
$\beta=\angle MAC=\angle ACM$, using the two isosceles triangles.
Since $M$ is between $B,C$, the angles of $ABC$ sum to
$2\alpha+2\beta=180^\circ$. Thus
$\angle BAC=\alpha+\beta=90^\circ$.
If $A$ lies on $BC$, the distance equalities would force $A=B$ or
$C$, excluded by a nondegenerate triangle.

\item[\ref{g:ex05}]
For quadrilateral $ABCD$, let $E,F,G,H$ be successive side midpoints.
The midpoint theorem in $ABC,CDA$ gives $EF\parallel HG\parallel AC$
and $EF=HG=AC/2$. Similarly $FG\parallel HE\parallel BD$.
Thus $EFGH$ is a parallelogram. Its adjacent sides are perpendicular
exactly when $AC\perp BD$. This characterizes the rectangle case.
If the original quadrilateral has degeneracies making the midpoint
figure collapse, the rectangle assertion is interpreted only for
a nondegenerate midpoint parallelogram.

\item[\ref{g:ex06}]
Reflect $B$ in $\ell$ to $B'$. Then
$BX+XC=B'X+XC\ge B'C$, with equality precisely when $X$ lies
on segment $B'C$. Because $B',C$ are strictly on opposite sides
of $\ell$, this segment has exactly one intersection with $\ell$.
If the projections coincide, $B'C$ is perpendicular to $\ell$,
and that unique intersection is the common projection.

\item[\ref{g:ex07}]
Corresponding sides all scale by $k$, so their sums do too.
Corresponding altitudes also scale by $k$, either by right-triangle
similarity or by a similarity transformation. The area formula
base times altitude divided by $2$ therefore scales by $k^2$.

\item[\ref{g:ex08}]
The two triangles have a common altitude from $A$, giving ratio
$BD/DC$. Their areas from the two side-angle formulas have ratio
$AB\sin\angle BAD/(AC\sin\angle DAC)=AB/AC$.
Both denominators are positive because the bisector is internal
in a nondegenerate triangle.

\item[\ref{g:ex09}]
Since $P$ lies on both chords, $\angle PAC=\angle PDB$ and
$\angle PCA=\angle PBD$ by inscribed angles.
Thus triangles $PAC,PDB$ are similar, with
$PA/PD=PC/PB$. Rearrangement gives the result.
Equivalently use the signed secant theorem; both signed products
are negative for an internal $P$, so their unsigned magnitudes agree.

\item[\ref{g:ex10}]
The feet satisfy $BE\perp EC$ and $CF\perp FB$, so both belong
to the circle with diameter $BC$. This remains valid when a foot
is a vertex. If, for instance, $\angle A=90^\circ$, then $E=F=A$:
the statement means that the distinct points among $B,C,E,F$
belong to that circle. It does not claim four distinct points or
use an angle with a repeated arm.

\item[\ref{g:ex11}]
Write $D,E,F$ for the feet on $BC,CA,AB$ and $H$ for the
orthocenter. The quadrilaterals $ABDE$ and $ACDF$ are cyclic,
with diameters $AB,AC$. Hence
$\angle EDH=\angle EDA=\angle EBA=90^\circ-A$ and
$\angle HDF=\angle ADF=\angle ACF=90^\circ-A$.
The altitude bisects the angle at $D$. The other two statements
follow cyclically. Acuteness ensures these are internal angles
with the asserted order of rays.

\item[\ref{g:ex12}]
For the acute case, $\angle BHC=180^\circ-A$.
Reflection reverses the orientation, so the reflected point $H'$
sees $BC$ under the angle supplementary to $\angle BAC$ and is
on the circumcircle. To remove order assumptions, put
$B=(0,0)$, $C=(1,0)$, $A=(u,v)$, $v\ne0$.
Then $H=(u,u(1-u)/v)$ and
$H'=(u,-u(1-u)/v)$. Substitution into
$x^2+y^2-x+(u-u^2-v^2)y/v=0$ gives zero. Thus the general
statement includes obtuse and right configurations.
\end{description}

\subsection{Ratios and centers}
\begin{description}[style=nextline,leftmargin=2.2em]
\item[\ref{g:ex13}]
The angle-bisector ratios multiply to
$(AB/AC)(BC/BA)(CA/CB)=1$, so Ceva proves concurrence.
Alternatively, the first two internal bisectors meet inside the
triangle at a point equidistant from all three side lines.
It is therefore on the third internal bisector by the equal-distance
criterion inside an angle.

\item[\ref{g:ex14}]
Ceva gives $2\cdot3\cdot AF/FB=1$, so $AF/FB=1/6$.
There is exactly one point of the open side $AB$ with this ratio;
the converse of Ceva proves that it gives concurrence.

\item[\ref{g:ex15}]
The midpoint theorem gives $DE\parallel AB$, so no finite point
$F=DE\cap AB$ exists. The finite form of Menelaus cannot be applied
with a nonexistent denominator. In a projective extension there is
an ideal intersection, but one must use that formulation explicitly.

\item[\ref{g:ex16}]
With the signed convention, Menelaus gives
$2\cdot3\cdot\overline{AF}/\overline{FB}=-1$.
Thus the last ratio is $-1/6$: $F$ is external, and its unsigned
ratio is $1/6$. More precisely it lies beyond $A$ rather than $B$,
since a point beyond $B$ has $AF/FB>1$.

\item[\ref{g:ex17}]
Affine maps preserve ratios on lines and multiply every area by the
same nonzero factor. Use normalized barycentrics relative to $ABC$.
Put $L=t^2+t+1$. The intersections of successive cevians are
\[
 \begin{aligned}
 P=AD\cap BE&=(t^2:1:t),\\
 Q=BE\cap CF&=(t:t^2:1),\\
 R=CF\cap AD&=(1:t:t^2).
 \end{aligned}
\]
Indeed $AD$ has equation $z=ty$, $BE$ has equation $x=tz$,
and $CF$ has equation $y=tx$. Substitution verifies each pair.
The absolute determinant of their normalized rows equals
\[
 \frac{t^6+1-2t^3}{L^3}
 =\frac{(t^3-1)^2}{L^3}
 =\frac{(t-1)^2}{L}.
\]
This is the required area ratio. At $t=1$ all three intersections
coincide, correctly giving zero.

\item[\ref{g:ex18}]
Normalize the positive barycentric coordinates of $P$ to
$x+y+z=1$. The side point $D$ is $(0:y:z)$, so
$P=xA+(1-x)D$. Along $AD$, this gives $PD/AD=x$.
Similarly $PE/BE=y$, $PF/CF=z$; their sum is $1$.
The alternative sum is
$(1-x)/x+(1-y)/y+(1-z)/z$, which varies. For instance it is $6$
at the centroid and larger for unequal positive coordinates.

\item[\ref{g:ex19}]
Let $D,E$ be midpoints of $BC,CA$, and let $P=AD\cap BE$.
The midpoint theorem makes triangles $PDE,PAB$ similar with factor
$1/2$, so $PA=2PD$ and $PB=2PE$. Ceva places the third median
through the same $P$. Equal bases on each side give equal pairs
of small areas; since $[PBC]=[PCA]=[PAB]=[ABC]/3$, all six
small triangles have area $[ABC]/6$.

\item[\ref{g:ex20}]
Perpendicular distances from $I$ to all sides equal $r$; summing
the three areas gives $\Delta=r(a+b+c)/2=rs$.
For contact lengths, write the equal tangents from $A,B,C$ as
$x,y,z$. Solve $x+y=c$, $y+z=a$, $z+x=b$, obtaining
$x=s-a$, $y=s-b$, $z=s-c$. Positivity follows from strict
triangle inequalities.

\item[\ref{g:ex21}]
Let the $A$-excircle touch extensions of $AB,AC$ at $F,E$.
Equal tangents give $AF=AE=s$. Hence $BF=s-c$ and $CE=s-b$.
The tangents from $B,C$ to contact $D_A$ on $BC$ have the same
lengths, so $BD_A=s-c$, $CD_A=s-b$.
Cyclically, the three cevian ratios multiply to
$(s-c)/(s-b)\cdot(s-a)/(s-c)\cdot(s-b)/(s-a)=1$.
Ceva proves concurrence, the Nagel point.

\item[\ref{g:ex22}]
The side midpoints form a nondegenerate triangle. They lie on the
nine-point circle with center $(O+H)/2$ and radius $R/2$.
A nondegenerate triangle has a unique circumcircle, so this is its
circumcircle. Nothing in the argument requires $O\ne H$; for an
equilateral triangle the common center is still correct.

\item[\ref{g:ex23}]
With $O$ as origin, $h=a+b+c$ and the midpoint of $BC$ is
$m=(b+c)/2$. Reflection in that point sends $h$ to
$2m-h=-a$, the antipode of $A$ on the circumcircle.
This differs from reflection in the entire line $BC$, whose image
is generally another circumcircle point.

\item[\ref{g:ex24}]
Problem~\ref{g:ex11} shows that $DH,EH,FH$ are the internal
bisectors of the orthic triangle in the acute case. Thus $H$ is
its incenter. If $A$ is obtuse, the same directed-angle equalities
show that the line $DH$ is the internal bisector at $D$, while
$EH,FH$ are external bisectors at $E,F$. Thus $H$ is the
$D$-excenter of the orthic triangle. The original vertex $A$
is its incenter. Right triangles produce a degenerate orthic
triangle and are outside this incenter/excenter assertion.
\end{description}

\subsection{Metric geometry and power}
\begin{description}[style=nextline,leftmargin=2.2em]
\item[\ref{g:ex25}]
The cosine law gives
$\cos A=(b^2+c^2-a^2)/(2bc)$. The denominator is positive;
the sign of the numerator is positive, zero, or negative exactly
when $A$ is acute, right, or obtuse.

\item[\ref{g:ex26}]
The angle-bisector theorem gives
$BD=ac/(b+c)$ and $DC=ab/(b+c)$.
Stewart gives
$b^2BD+c^2DC=a(\ell^2+BD\cdot DC)$.
The left side is $abc$, so
$\ell^2=bc-a^2bc/(b+c)^2$, as required.

\item[\ref{g:ex27}]
The three formulas are
$m_a^2=(2b^2+2c^2-a^2)/4$ and its cyclic versions.
Every side square occurs with total coefficient $2+2-1=3$.
Adding yields the claimed $\tfrac34$ factor.

\item[\ref{g:ex28}]
Let $B,D$ be opposite angles between sides $(a,b)$ and $(c,d)$.
They are supplementary. The cosine law on the common diagonal gives
\[
 a^2+b^2-c^2-d^2=2(ab+cd)\cos B.
\]
Also $K=(ab+cd)\sin B/2$. Hence
\[
 16K^2=4(ab+cd)^2-(a^2+b^2-c^2-d^2)^2
 =16(s-a)(s-b)(s-c)(s-d).
\]
The last equality follows by factoring the difference of squares.
All factors are positive for a nondegenerate cyclic quadrilateral.

\item[\ref{g:ex29}]
Let its side be $s$, and choose
$A=(0,0)$, $B=(s,0)$, $C=(s/2,\sqrt3s/2)$, $P=(x,y)$.
Subtracting $PA^2=9$ from the other two distance equations yields
\[
 x=\frac{s^2-7}{2s},\qquad
 y=\frac{s^2-25}{2\sqrt3s}.
\]
Substitution into $x^2+y^2=9$ gives
$s^4-50s^2+193=0$. Thus $s^2=25\pm12\sqrt3$.
Since $P$ is interior, $y>0$, requiring $s^2>25$.
Therefore $s=\sqrt{25+12\sqrt3}$.
For this value $x=(9+6\sqrt3)/s$, $y=6/s$.
Both inequalities $y<\sqrt3x$ and
$y<\sqrt3(s-x)$ hold, so the point is indeed inside; direct
substitution checks all three distances.
The $60^\circ$ rotation in the hint explains why the three given
distances form a right-triangle configuration, while these coordinates
resolve the orientation without a diagram assumption.

\item[\ref{g:ex30}]
For any convex quadrilateral with sides $a,b,c,d$ and opposite
angles $B,D$, the two triangle area formulas and cosine laws give
\[
 K^2=(s-a)(s-b)(s-c)(s-d)
       -abcd\cos^2\frac{B+D}{2}.
\]
For completeness, put $h=a^2+b^2-c^2-d^2
=2ab\cos B-2cd\cos D$ and
$2K=ab\sin B+cd\sin D$.
Then
$16K^2+h^2=4(a^2b^2+c^2d^2-2abcd\cos(B+D))$.
Factoring as in the previous problem yields the displayed identity.
The correction is nonnegative, and vanishes precisely when
$B+D=180^\circ$, the cyclic case.

A cyclic quadrilateral with these side lengths exists.
To justify this point, let $m$ be the largest side and take
$R\ge m/2$. Each length $\ell$ gives a minor central angle
$2\arcsin(\ell/(2R))$. If their sum at $R=m/2$ is at least
$2\pi$, continuity and decrease to zero give a radius where
the four angles sum to $2\pi$. Otherwise use the major central
angle for the largest side. Its total starts below $2\pi$ and
is above $2\pi$ for sufficiently large $R$, since
$m$ is strictly smaller than the sum of the other sides.
Again continuity gives the required circle. In either case
place the four consecutive chords on it.
The maximum area is therefore
$\sqrt{(s-a)(s-b)(s-c)(s-d)}$.

\item[\ref{g:ex31}]
For tangency points $T,U$, radii $OT,OU$ are perpendicular to
$PT,PU$. Thus both tangent squares equal $PO^2-R^2$.
Taking positive square roots gives $PT=PU$.

\item[\ref{g:ex32}]
The power formula gives both positive products
$PA\cdot PB=PC\cdot PD=R^2-PO^2$.
Put $P=(0,0)$ and the two perpendicular chord lines on the axes;
let $O=(u,v)$. The signed endpoint coordinates on the first
axis have sum $2u$ and product $u^2+v^2-R^2$.
Their squared sum is $4u^2-2(u^2+v^2-R^2)$.
On the second axis replace $u$ by $v$. Adding cancels the center
coordinates and gives $4R^2$.

\item[\ref{g:ex33}]
Choose an auxiliary circle $\gamma$ meeting both given circles
at pairs $A,B$ and $C,D$, with its center off the line of the
given centers. Such a circle can be found by taking a sufficiently
distant center on their perpendicular bisector and choosing its
radius equal to the common distance to their centers.
The two common-chord lines $AB,CD$ are not parallel, since their
normal center differences are not parallel. Let them meet at $P$.
Transitivity gives equal powers of $P$ to the original circles.
The line through $P$ perpendicular to their center line is the
desired radical axis. One auxiliary circle suffices when the centers
are available.

\item[\ref{g:ex34}]
At an orthogonal intersection $T$, the radii $O_1T,O_2T$ are
perpendicular because their tangents are. Pythagoras yields
$O_1O_2^2=r_1^2+r_2^2$.
Conversely, put $d=\sqrt{r_1^2+r_2^2}$. Since $r_1,r_2>0$,
$|r_1-r_2|<d<r_1+r_2$, so the circles intersect in two distinct
points. The cosine law at either intersection gives a right angle
between the radii and therefore between the tangents.

\item[\ref{g:ex35}]
Each common-chord line is a radical axis. If two are nonparallel,
their intersection has equal power to all three circles and lies
on the third. If two are parallel, their normals, the corresponding
center differences, are parallel, so the three centers are collinear.
The third radical axis is then also perpendicular to that center
line and hence parallel, or coincides with an earlier one.
Coincident circles and coincident axes require the exclusions in
the statement; they do not determine three independent lines.

\item[\ref{g:ex36}]
Arc equality gives $LB=LC$. Since $A,I,L,I_A$ are collinear,
an angle chase in $BIL$ gives
$\angle LBI=\angle BIL=90^\circ-C/2$, hence $LB=LI$.
The external bisector at $B$ similarly gives
$\angle LBI_A=\angle BI_AL=C/2$, hence $LB=LI_A$.
All claims are unsigned lengths. The position of $L$ between
$I$ and $I_A$ is consistent with the equal distances, rather than
an extra assumption needed in the proof.
\end{description}

\subsection{Transformations and isogonality}
\begin{description}[style=nextline,leftmargin=2.2em]
\item[\ref{g:ex37}]
Choose orthonormal axes parallel and perpendicular to the reflecting
line. A displacement $(u,v)$ becomes $(u,-v)$. Its squared length
$u^2+v^2$ is unchanged, so every pairwise distance is unchanged.
Translating the coordinate origin onto the line completes the argument
for any reflecting line.

\item[\ref{g:ex38}]
In the stated order the composition is
\[
 x\longmapsto k\ell x+q,\qquad
 q=\ell(1-k)S+(1-\ell)T.
\]
If $k\ell\ne1$, it is the homothety of ratio $k\ell$ centered
at $q/(1-k\ell)$. If $k\ell=1$ and $q\ne0$, it is translation
by $q$. If $k\ell=1$ and $q=0$, it is the identity.
This classification includes ratio $1$ maps without inventing a
unique center for them.

\item[\ref{g:ex39}]
Write $f(z)=uz+v$. The first image gives $v=1+i$; the second
gives $u=1+i$. Thus
$s=v/(1-u)=(1+i)/(-i)=-1+i$ and the scale factor is $|u|=\sqrt2$.
Its rotation angle is $45^\circ$.

\item[\ref{g:ex40}]
The centers satisfy $TO_1=r$, $TO_2=s$ on opposite rays, so
$O_2-T=-(s/r)(O_1-T)$. A homothety of ratio $-s/r$ sends the
first center to the second and its radius $r$ to $s$; hence it
sends the whole circle to the other one. Every line through $T$
is fixed as a set. Points on it generally move to the opposite ray,
and the contact point itself remains fixed.

\item[\ref{g:ex41}]
Orient $ABC$ positively and rotate $P$ by $60^\circ$ about $A$
to $Q$, so $B$ maps to $C$. Then
$PA+PB=PQ+QC\ge PC$.
Equality means $Q\in PC$. Apart from $P=A,B$, collinearity
$P,Q,C$ gives the appropriate $60^\circ$ directed angle, placing
$P$ on $(ABC)$. Tracking the segment order, or applying Ptolemy
on each of its three arcs, selects exactly the closed arc $AB$
not containing $C$. Every point of that arc satisfies equality
by Ptolemy. Thus this closed arc, including $A,B$, is the full
equality locus.

\item[\ref{g:ex42}]
Assume $ABC$ is positively oriented; a reflection handles the
opposite orientation. Put $\tau=i/\sqrt3$ and
$\omega=(1+i\sqrt3)/2$. The centers of the external equilateral
triangles on $BC,CA,AB$ are respectively
\[
 U=\frac{(1+\tau)b+(1-\tau)c}{2},\quad
 V=\frac{(1+\tau)c+(1-\tau)a}{2},\quad
 W=\frac{(1+\tau)a+(1-\tau)b}{2}.
\]
The sign of the imaginary displacement puts each new triangle on
the right of its directed side, hence outside $ABC$.
Using $\omega(1-\tau)=1+\tau$ and
$\omega(1+\tau)=2\tau$ gives
$W-U=\omega(V-U)$. Thus the two sides at $U$ have equal
lengths and an angle of $60^\circ$.
They are nonzero: $U=V$ would imply
$a-b=\bar\omega(c-b)$, giving the opposite orientation for $ABC$.
Therefore $UVW$ is an equilateral triangle.

\item[\ref{g:ex43}]
For median $AM$,
$1=(c/b)\sin\angle BAM/\sin\angle MAC$, so its sine ratio is
$b/c$. The reflected cevian has sine ratio $c/b$.
Its side ratio is therefore $(c/b)(c/b)=c^2/b^2$.

\item[\ref{g:ex44}]
The centroid lies on all three medians. Their isogonal reflections
are precisely the symmedians. The concurrence theorem for reflected
interior cevians identifies the conjugate as their intersection,
the symmedian point.

\item[\ref{g:ex45}]
If all four lines coincide, the median is an angle bisector.
Its side ratio is both $1$ and $AB/AC$, so $AB=AC$.
Conversely, in an isosceles triangle the congruent halves show
that the median is both altitude and angle bisector. Reflection
fixes this line, so it is also the symmedian.

\item[\ref{g:ex46}]
For $P=(x:y:z)$, line $AP$ meets $BC$ with ratio $z/y$.
The candidate
$Q=(a^2/x:b^2/y:c^2/z)$ is a finite interior point because all
coordinates are positive. Its corresponding side ratio is
$(c^2/z)/(b^2/y)=(c^2/b^2)(y/z)$, the isogonal partner required
by the ratio identity. The same holds at $B,C$, so the three
reflected lines pass through $Q$.

\item[\ref{g:ex47}]
The product of the three sine ratios for the original cevians is
$1$ by trigonometric Ceva. Isogonal reflection replaces each ratio
by its reciprocal, so their product remains $1$. The reflected
cevians are interior and therefore concurrent at a finite point.

\item[\ref{g:ex48}]
The tangent construction gives $BK/KC=c^2/b^2$.
By the sine area formula,
$\sin\angle BAD/\sin\angle DAC=c/b$.
The inscribed-angle theorem converts this sine ratio into $DB/DC$.
Squaring yields
$BK/KC=(DB/DC)^2=c^2/b^2$.
All divisions are by positive side or chord lengths; the excluded
diameter case is the one with parallel tangents.
\end{description}

\subsection{Coordinates and inversion}
\begin{description}[style=nextline,leftmargin=2.2em]
\item[\ref{g:ex49}]
With $O=0$, all three vertex vectors have length $R$. Thus
$(b+c)\cdot(b-c)=|b|^2-|c|^2=0$.
The point $h=a+b+c$ belongs to the altitude from $A$, and the
cyclic arguments prove it is $H$. The nine-point center is
$n=h/2$: a side midpoint $m=(b+c)/2$ satisfies
$m-n=-a/2$, and a midpoint of $AH$ satisfies $u-n=a/2$.
Both distances are $R/2$. Orthocenter reflection supplies the
altitude feet, as proved in the text.

\item[\ref{g:ex50}]
For unit $a,b$, both satisfy $z+ab\bar z=a+b$.
This is a real line equation, since its conjugate is the original
equation divided by $ab$. Let $r=a+b-ab\bar p$ and
$f=(p+r)/2$. Substituting gives
$f+ab\bar f=a+b$, so $f$ is on the chord line.
Moreover
\[
 \overline{\frac{r-p}{b-a}}=-\frac{r-p}{b-a},
\]
using $\bar a=1/a$, $\bar b=1/b$.
The quotient is purely imaginary, so $r-p$ is perpendicular to
$b-a$. Hence $f$ is the perpendicular projection.
If $p$ lies on the line, $r=p=f$ and the same formula remains valid.

\item[\ref{g:ex51}]
After normalization $x+y+z=1$, set
$D=(yB+zC)/(y+z)$.
Then $P=xA+(1-x)D$, so $P\in AD$ and $PD/AD=x$.
The side division satisfies $BD/DC=z/y$.
Undoing normalization gives $PD/AD=x/(x+y+z)$.
Interiority ensures $x,y,z>0$, so no denominator vanishes.

\item[\ref{g:ex52}]
For three prescribed perpendicular lines through $X_A,X_B,X_C$
to the sides $BC,CA,AB$ of a nondegenerate triangle, concurrence
is equivalent to
\[
 X_A\cdot(B-C)+X_B\cdot(C-A)+X_C\cdot(A-B)=0.
\]
Necessity follows by replacing each $X$ with the common point $Q$
inside its dot product; the three side differences sum to zero.
For sufficiency, the first two nonparallel perpendicular lines
have a unique intersection $Q$. Subtract their two zero terms
from the displayed identity to obtain the equation for the third.

Apply this criterion with $(X_A,X_B,X_C)=(A_1,B_1,C_1)$ and
side triangle $(A_2,B_2,C_2)$. The resulting cyclic dot-product sum
is the negative of the sum with subscripts $1,2$ interchanged,
because the dot product is symmetric. Thus the reverse sum is
zero too, proving the second concurrence.

\item[\ref{g:ex53}]
Normalize $(ABC)$ to the unit circle and write
$a=e^{i\alpha}$, $b=e^{i\beta}$, $c=e^{i\gamma}$,
$p=e^{i\rho}$. Two of the foot formulas differ by
$(a-b)(1-c/p)/2$.
For unit numbers, the direction of
$e^{iu}-e^{iv}$ is $(u+v)/2+\pi/2$ modulo $\pi$.
Therefore the Simson-line direction is
\[
 \frac{\alpha+\beta+\gamma-\rho}{2}\pmod{\pi}.
\]
Replacing $p$ by its antipode changes $\rho$ by $\pi$ and the
line direction by $-\pi/2$. If the chosen two feet coincide,
choose a different pair. All three cannot coincide in a nondegenerate
triangle, since their common point would belong to all three
side lines. The same direction formula is obtained from that pair.

\item[\ref{g:ex54}]
If the feet are $X,Y,Z$, the reflections have vectors
$2X-P,2Y-P,2Z-P$. A homothety centered at $P$ with ratio $2$
takes the feet to these reflected points, so collinearity is
equivalent in the two triples. The Simson theorem supplies the
equivalence with $P\in(ABC)$, including repeated feet.

\item[\ref{g:ex55}]
The line $x=2$ becomes $x'=2(x'^2+y'^2)$, or
$(x'-1/4)^2+y'^2=1/16$, with the origin deleted.
The circle has equation $x^2+y^2-4x+3=0$ and becomes
$1-4x'+3(x'^2+y'^2)=0$, or
$(x'-2/3)^2+y'^2=1/9$.
Neither the original nor this image circle contains the origin,
so no point on the latter circle needs to be removed.

\item[\ref{g:ex56}]
Write a circle not through the inversion center as
$|x|^2-2c\cdot x+d=0$ with $d\ne0$.
Its image under radius-$k$ inversion is
$d|x'|^2-2k^2c\cdot x'+k^4=0$.
If $c\ne0$ and the two circles agree, comparison after normalizing
the quadratic coefficient gives $d=k^2$. Since
$d=|c|^2-\rho^2$, this is exactly
$OC^2=k^2+\rho^2$, the orthogonality condition.
The same coefficient calculation proves the forward direction.
If $c=0$, a circle of radius $\rho$ becomes one of radius
$k^2/\rho$; it is fixed precisely for $\rho=k$.
Thus the inversion circle itself is an additional fixed circle.

\item[\ref{g:ex57}]
Draw the line through $I$ perpendicular to $AI$, meeting the
two side rays at $E,F$. Through $E,F$ draw perpendiculars to the
respective sides. They meet at a point $J$ on $AI$ by symmetry.
The circle centered at $J$ through $E$ is constructible and tangent
to both sides. To prove it internally tangent to the circumcircle,
use the coordinates of the midpoint lemma:
$I=(d,0)$, $J=(t,0)$, $\theta=A/2$.
Then $t=d/\cos^2\theta$, $\rho=t\sin\theta=r/\cos^2\theta$,
and Euler's formula gives
$(t-u)^2+v^2=(R-\rho)^2$.

It remains to justify the internal, rather than wrong squared,
branch. The standard half-angle identity
$r=4R\sin(A/2)\sin(B/2)\sin(C/2)$ follows from
$r=\Delta/s$, $a=2R\sin A$, and the half-angle formulas.
Also
$\sin(B/2)\sin(C/2)\le(1-\sin\theta)/2$,
by the product-to-sum identity and $\cos((B-C)/2)\le1$.
Consequently
$r\le2R\sin\theta(1-\sin\theta)<R\cos^2\theta$.
Thus $\rho<R$ and $OJ=R-\rho$, proving existence and internal
tangency. The midpoint construction also proves uniqueness
among nonzero circles in the angle satisfying the tangencies.

\item[\ref{g:ex58}]
Invert at $A$. Equality is exactly equality in
$B'D'\le B'C'+C'D'$, hence $B',C',D'$ lie on a line with
$C'$ between the other two. That line cannot pass through the
inversion center, since then $A,B,C,D$ would be collinear.
Its inverse is a circle through $A$.
Along that circle, deleting $A$ leaves a linearly ordered arc, and
inversion sends it continuously and one-to-one onto the image
line. The between condition is precisely the order $A,B,C,D$
or its reverse. Conversely Ptolemy gives equality for that cyclic
order. The assumption of no three collinear excludes additional
degenerate equality cases.

\item[\ref{g:ex59}]
Put $T=(0,0)$ and the two centers at $(R,0)$ and $(r,0)$.
Their equations are $x^2+y^2=2Rx$ and
$x^2+y^2=2rx$. Inversion of radius $k$ gives
$x'=k^2/(2R)$ and $x'=k^2/(2r)$.
The lines are parallel and distinct, and their distance is
$k^2(R-r)/(2Rr)$. The deleted tangency point has no finite image;
it is not an additional intersection of the image lines.

\item[\ref{g:ex60}]
Place the outer center at the origin and put $P_i=Ru_i$,
$|u_i|=1$. The inner centers are $C_i=(R-r_i)u_i$.
For common external tangents, the right triangle between the centers
gives
$t_{ij}^2=|C_i-C_j|^2-(r_i-r_j)^2$.
Expanding yields
\[
 t_{ij}^2=(R-r_i)(R-r_j)|u_i-u_j|^2.
\]
This is positive because contacts are distinct and $r_i<R$,
so these external tangents exist and
$t_{ij}=\sqrt{(R-r_i)(R-r_j)}\,P_iP_j/R$.
Ptolemy for the contact quadrilateral, multiplied by
$\sqrt{\prod_i(R-r_i)}/R^2$, is exactly the claimed identity.
The argument proves this necessary relation in the given
configuration; it does not furnish an unrestricted converse.
\end{description}

\subsection{Projective tools and synthesis}
\begin{description}[style=nextline,leftmargin=2.2em]
\item[\ref{g:ex61}]
For $f(t)=(\alpha t+\beta)/(\gamma t+\delta)$,
\[
 f(x)-f(y)=
 \frac{(\alpha\delta-\beta\gamma)(x-y)}
      {(\gamma x+\delta)(\gamma y+\delta)}.
\]
Insert four such differences in the cross-ratio expression.
Every determinant factor and every denominator cancels, leaving
the original cross ratio. If a point goes to infinity, cancel
the common factor first and then evaluate there, or use
homogeneous pairs $(\alpha t+\beta:\gamma t+\delta)$.
Nonzero determinant ensures distinct source points stay distinct.

\item[\ref{g:ex62}]
The polar is $px+qy=R^2$. The pole of $ux+vy=w$ with
$w\ne0$ is $(R^2u/w,R^2v/w)$.
If $w=0$, the line passes through the center and has no finite
pole; projectively its pole is the ideal point $(u:v:0)$.
The center has no finite polar line; its projective polar is
the line at infinity.

\item[\ref{g:ex63}]
Use barycentrics $D=(0:d:e)$ with $de\ne0$.
Since $E\in AD$ and $E\ne A$, scale to $E=(u:d:e)$.
The stated nondegeneracies force $u\ne0$.
Then $F=(u:0:e)$ and $G=(u:d:0)$.
Their joining line has equation
$-de\,x+ue\,y+ud\,z=0$, so
$X=(0:d:-e)$.
The signed ratios are $BD/DC=e/d$ and $BX/XC=-e/d$.
Their quotient is the cross ratio $(B,C;D,X)$, which is $-1$.
If $d=e$, $X$ is ideal; that case is excluded by the finite
statement but is consistent with harmonic division projectively.

\item[\ref{g:ex64}]
For six points listed cyclically or in any specified hexagon order,
Pascal applies because a circle is a nondegenerate conic.
The three required opposite-side intersections are exactly its
$X,Y,Z$. If one pair of opposite sides is parallel, its
intersection is their ideal direction point. Collinearity then
means that the line through the other two finite intersections is
parallel to those sides. If additional intersections are ideal or
coincident, state the corresponding projective formulation rather
than assigning a fictitious finite point.

\item[\ref{g:ex65}]
Let $\ell_i$ be the six tangents at successive contact points $T_i$,
and let $V_i=\ell_i\cap\ell_{i+1}$, with indices modulo $6$.
Pascal says that
$T_1T_2\cap T_4T_5$,
$T_2T_3\cap T_5T_6$, and
$T_3T_4\cap T_6T_1$ are collinear.
The polar of the first intersection joins the poles of
$T_1T_2$ and $T_4T_5$, namely $V_1,V_4$;
this is an application of reciprocity.
The other two polars are $V_2V_5,V_3V_6$.
The polars of collinear points concur at the pole of their line,
so these three main diagonals concur.
If that pole is ideal, the three diagonals are parallel.

\item[\ref{g:ex66}]
The invertible linear map $(x,y)\mapsto(x/a,y/b)$ sends the
ellipse to the unit circle. Its contact point becomes
$(x_0/a,y_0/b)$, whose tangent is
$Xx_0/a+Yy_0/b=1$. Substitution gives
$xx_0/a^2+yy_0/b^2=1$.
One can verify tangency without invoking a tangent-map theorem:
the weighted Cauchy inequality shows that all ellipse points are
on one side of this line, with equality only at the contact.
When $x_0y_0\ne0$, the intercepts are $(a^2/x_0,0)$ and
$(0,b^2/y_0)$. If $x_0=0$ the tangent is horizontal; if $y_0=0$
it is vertical, so the missing finite intercept is not to be divided by.

\item[\ref{g:ex67}]
Heron's algebraic form gives
$16\Delta^2=2(a^2b^2+b^2c^2+c^2a^2)-(a^4+b^4+c^4)$.
Therefore
\[
 (a^2+b^2+c^2)^2-48\Delta^2
 =2\bigl((a^2-b^2)^2+(b^2-c^2)^2+(c^2-a^2)^2\bigr)\ge0.
\]
Both sides of the desired inequality are nonnegative, so taking
square roots is valid. Equality forces $a=b=c$.

\item[\ref{g:ex68}]
Use a unit circumcircle and complex coordinates. The reflection of
$p$ in $BC$ is $r_1=b+c-bc/p$; that in $CA$ is
$r_2=c+a-ca/p$, and $h=a+b+c$.
When $r_1\ne r_2$,
\[
 \frac{h-r_1}{r_2-r_1}
 =\frac{ap+bc}{(a-b)(p-c)}
\]
is real: its conjugate is the same expression after multiplying
numerator and denominator by $abcp$.
Thus $h$ lies on the line through these two reflections.
If this pair coincides, choose another pair; all three cannot
coincide, as explained in Problem~\ref{g:ex53}.
Problem~\ref{g:ex54} places the three reflections on one line.
The homothety of center $P$ and ratio $1/2$ sends their line to
the Simson line and sends $H$ to the midpoint of $PH$.

\item[\ref{g:ex69}]
Scale to $M=(0,0)$, $P=(-1,0)$, $Q=(1,0)$.
The circle has equation $x^2+y^2-2ky=1$ for some real $k$.
Write $A=(a_x,a_y)$, $C=(c_x,c_y)$ and
$q_A=|A|^2$, $q_C=|C|^2$.
The chord power at $M$ is $-1$, so
$B=-A/q_A$ and $D=-C/q_C$.
Put $\delta=a_xc_y-a_yc_x$.
The line-intersection formulas give the signed abscissas
\[
 x_X=\frac{\delta}{c_y+q_Ca_y},\qquad
 x_Y=-\frac{\delta}{q_Ac_y+a_y}.
\]
The denominators are nonzero by the finite-intersection assumptions.
The circle equation gives
$q_A-1=2ka_y$ and $q_C-1=2kc_y$, hence
$(q_C-1)a_y=(q_A-1)c_y$. The two denominators agree.
Thus $x_X=-x_Y$, proving $MX=MY$, including the case $k=0$.

\item[\ref{g:ex70}]
An exterior point cannot improve the sum. Let $Q$ be its nearest
point in the closed triangle, which exists by compactness.
For every vertex $V$, minimality along segment $QV$ implies
$(P-Q)\cdot(V-Q)\le0$: expand the squared distance to
$Q+t(V-Q)$ for $t>0$ and let $t$ decrease to zero.
Then
$PV^2=PQ^2+QV^2-2(P-Q)\cdot(V-Q)>QV^2$.
All three distances decrease at $Q$.

Inside the triangle, the existence and directional-change argument
in Chapter~\ref{g:synthesis} gives the interior $120^\circ$ point
when all vertex angles are below $120^\circ$. The unit-vector
inequality proves its global minimality, and the minimum is
\[
 \sqrt{\frac{a^2+b^2+c^2+4\sqrt3\,\Delta}{2}}.
\]
If $\angle A\ge120^\circ$, let $p=P-A$ and unit side directions
$e_B,e_C$. Then
\[
 PA+PB+PC\ge AB+AC+|p|-p\cdot(e_B+e_C)
 \ge AB+AC,
\]
since $|e_B+e_C|\le1$. Equality is attained at $P=A$.
The cases for $B,C$ are cyclic, and at most one angle can be
at least $120^\circ$.

\item[\ref{g:ex71}]
Orthogonal projection onto $\ell$ has form $T(x)+q$, where $T$
is the symmetric matrix $ee^{\mathsf T}$ for a unit direction
vector $e$, and $q$ is a fixed vector perpendicular to $\ell$.
By the concurrency criterion in Problem~\ref{g:ex52}, it is enough
to show
\[
 (TA+q)\cdot(B-C)+(TB+q)\cdot(C-A)
 +(TC+q)\cdot(A-B)=0.
\]
The $q$ terms sum to zero. The remaining terms cancel in pairs
because $TA\cdot B=A\cdot TB$.
Thus the three lines concur at the orthopole of $\ell$.
If $\ell$ is shifted by a perpendicular vector $v$, each projection
is shifted by $v$. Shifting the concurrence point by $v$ preserves
all three perpendicular conditions, and uniqueness proves the claim.

\item[\ref{g:ex72}]
For any $X\in BC$, reflect $X$ in $AB,AC$ to $X_1,X_2$.
For $Z\in AB$, $Y\in AC$,
\[
 XY+YZ+ZX=X_2Y+YZ+ZX_1\ge X_1X_2=2AX\sin A.
\]
The last identity uses $AX_1=AX_2=AX$ and the angle between
their rays, $2A$. Since $AX\ge AD$,
the perimeter is at least $2AD\sin A$.
For $X=D$, the line from $D_1$ to $D_2$, its two reflections,
crosses $AB,AC$ at $F,E$ in that order. This follows by the
angle-bisector relations of the orthic triangle: reflection in the
side at $F$ straightens $DF$ and $FE$, and similarly at $E$.
Acuteness ensures both intersections lie in the open sides.
Thus
$DF+FE+ED=D_1F+FE+ED_2=D_1D_2=2AD\sin A$.
Equality requires $AX=AD$, so $X=D$; then straightness fixes
$Z=F$ and $Y=E$. These are the unique equality points.

\begin{center}
\begin{tikzpicture}[scale=.69]
\coordinate (A) at (0,3); \coordinate (B) at (-2,0);
\coordinate (C) at (3,0); \coordinate (D) at (0,0);
\coordinate (E) at (.5,2.5); \coordinate (F) at (-.461538,2.307692);
\coordinate (D1) at (-2.769231,1.846154); \coordinate (D2) at (3,3);
\draw (A)--(B)--(C)--cycle;
\draw[main] (D)--(E)--(F)--cycle;
\draw[aux] (D1)--(D2);
\draw[aux] (D1)--(D) (D)--(D2);
\foreach \X/\lab/\pos in {A/A/above,B/B/below left,C/C/below right,D/D/below,E/E/above right,F/F/above left,D1/D_1/left,D2/D_2/right}
 \fill (\X) circle(1.3pt) node[\pos]{$\lab$};
\end{tikzpicture}

{\small Reflection straightens the orthic perimeter into $D_1D_2$.}
\end{center}
\end{description}



% ---- solutions/number-theory ----
\section{Number Theory}
\subsection*{Divisibility, the Euclidean Algorithm, and Structure}
\paragraph{Problem \ref{n:ex01}.}
The remainders are $10,4,2,0$. Thus the gcd is $2$, and $2=10-2(144-14\cdot10)=29\cdot2026-408\cdot144$.

\paragraph{Problem \ref{n:ex02}.}
Every common positive divisor divides $(n+1)-n=1$, so it is $1$. The integer $1$ divides both, proving the gcd is $1$.

\paragraph{Problem \ref{n:ex03}.}
A common divisor divides $(n^2+n+1)-n(n+1)=1$, so it is $1$. This works even if $n+1=0$, since the other entry is then $1$.

\paragraph{Problem \ref{n:ex04}.}
Divide by $6$ to get $3x+5y=7$. The pair $(4,-1)$ works. All solutions are $x=4+5t,y=-1-3t$, $t\in\ZZ$. Substitution checks them; coprime cancellation after subtracting $(4,-1)$ shows no others.

\paragraph{Problem \ref{n:ex05}.}
Since $A=mk$ and $n\mid A$, cancellation gives $n\mid k$. Thus $A=mn\ell$ for an integer $\ell$. The coprimality is essential, as $m=n=2,A=2$ shows.

\paragraph{Problem \ref{n:ex06}.}
The product identity gives $ab=1080$, so $xy=30$ with coprime $x\le y$. The allowed factor pairs are $(1,30),(2,15),(3,10),(5,6)$. Thus $(a,b)$ is $(6,180),(12,90),(18,60),(30,36)$. Each pair has the required gcd and lcm.

\paragraph{Problem \ref{n:ex07}.}
Write $F=ab-a-b$. Both cannot be representable because adding their representations would represent $F$, contradicting Theorem~\ref{n:frobenius}. Choose $0\le y\le a-1$ with $N=ax+by$, $x\in\ZZ$. If $x\ge0$, $N$ is represented. Otherwise $x\le-1$ and $F-N=a(-x-1)+b(a-y-1)$ has nonnegative coefficients. Thus at least one, and hence exactly one, is representable.

\paragraph{Problem \ref{n:ex08}.}
Let $d$ be the gcd. Since $\gcd(a+b,b)=1$, also $\gcd(d,b)=1$. Reducing $a\equiv-b\pmod d$ gives $d\mid2b^2$, so coprime cancellation gives $d\mid2$. If $a,b$ are both odd, both expressions are even and $d=2$. Otherwise exactly one is even (both even would violate coprimality), so $a+b$ is odd and $d=1$.

\subsection*{Primes, Factorization, and Arithmetic Functions}
\paragraph{Problem \ref{n:ex09}.}
The formulas give $\tau(72)=4\cdot3=12$, $\sigma(72)=(1+2+4+8)(1+3+9)=195$, and $\varphi(72)=72(1-1/2)(1-1/3)=24$.

\paragraph{Problem \ref{n:ex10}.}
The product $\tau(n)=\prod(e_i+1)$ is odd exactly when every $e_i$ is even, exactly the square condition. For $n=1$ the empty product is $1$ and the assertion remains true.

\paragraph{Problem \ref{n:ex11}.}
Take $(k+1)!+2,\ldots,(k+1)!+(k+1)$. Each is larger than its divisor in $\{2,\ldots,k+1\}$, so it is composite; every term exceeds $k$.

\paragraph{Problem \ref{n:ex12}.}
The exponent patterns are $(5)$ or $(2,1)$, so $n=p^5$ or $n=p^2q$ with distinct primes $p,q$. Conversely the divisor counts are $6$ and $3\cdot2=6$.

\paragraph{Problem \ref{n:ex13}.}
If the list is $p_1,\ldots,p_k$, take $N=4p_1\cdots p_k-1$ (if the list were empty, $3$ already contradicts it). No $p_i$ divides $N$. It is odd and congruent to $3$ modulo $4$. A product of primes all congruent to $1$ modulo $4$, or with an even number of $3$ modulo $4$ factors counted with multiplicity, is $1$ modulo $4$. Therefore $N$ has a prime divisor congruent to $3$ modulo $4$, outside the list.

\paragraph{Problem \ref{n:ex14}.}
For $\gcd(m,n)=1$, disjoint prime supports make every divisor of $mn$ uniquely $ab$ with $a\mid m,b\mid n$. Hence $\sigma(mn)=\sum_{a\mid m}\sum_{b\mid n}ab=\sigma(m)\sigma(n)$. Without coprimality, $\sigma(4)=7\ne9=\sigma(2)^2$.

\paragraph{Problem \ref{n:ex15}.}
Inversion gives $\varphi(n)=\sum_{d\mid n}\mu(d)n/d$. Only squarefree $d$ contribute; expanding the choices of prime divisors yields $n\prod_{p\mid n}(1-1/p)$. The case $n=1$ is an empty product and gives $1$.

\paragraph{Problem \ref{n:ex16}.}
The map $d\mapsto n/d$ permutes the divisors, so if their product is $P>0$, then $P^2=\prod_{d\mid n}d(n/d)=n^{\tau(n)}$. Taking the positive square root gives the formula. If $\tau(n)$ is odd, $n$ is a square, so the result is still an integer.

\subsection*{Congruences, Inverses, and the Chinese Remainder Theorem}
\paragraph{Problem \ref{n:ex17}.}
The reduced congruence is $3x\equiv5\pmod7$. Since $3^{-1}\equiv5$, $x\equiv4\pmod7$. Thus the six classes modulo $42$ are $4,11,18,25,32,39$.

\paragraph{Problem \ref{n:ex18}.}
We need $12t\equiv6\pmod{18}$, or $2t\equiv1\pmod3$, hence $t\equiv2\pmod3$. Thus $x\equiv29\pmod{36}$. The compatibility condition holds since both residues are $5$ modulo $6$.

\paragraph{Problem \ref{n:ex19}.}
The first makes $x$ odd and the second makes it even. Equivalently $\gcd(4,6)=2$ does not divide $2-1=1$.

\paragraph{Problem \ref{n:ex20}.}
An odd integer $2k+1$ has square $4k(k+1)+1\equiv1\pmod8$. An even integer $2k$ has square $4k^2$, equal to $0$ modulo $8$ for even $k$ and $4$ for odd $k$. The examples $0,1,2$ attain $0,1,4$.

\paragraph{Problem \ref{n:ex21}.}
Since $7^4=2401\equiv1\pmod{100}$ and $2026=4\cdot506+2$, the residue is $7^2=49$. Thus the last two digits are $49$.

\paragraph{Problem \ref{n:ex22}.}
For $N=\sum d_i B^i$, congruence arithmetic gives $N\equiv\sum d_i\pmod{B-1}$ and $N\equiv\sum(-1)^i d_i\pmod{B+1}$. Divisibility holds precisely when the corresponding sum is zero modulo the indicated modulus. The modulus $B-1=1$ for $B=2$ gives a true but uninformative test.

\paragraph{Problem \ref{n:ex23}.}
Choose distinct primes $p_1,\ldots,p_k$ and solve $N\equiv-i\pmod{p_i^2}$ by CRT. Then $p_i^2\mid N+i$ for every $i$. Adding a sufficiently large multiple of $\prod_i p_i^2$ makes $N+1$ positive. The $k$ integers are consecutive and satisfy all conditions.

\paragraph{Problem \ref{n:ex24}.}
There are $m^2$ possible pairs, so a pair repeats. The transition is bijective, with inverse $(u,v)\mapsto(v-u,u)$. If the first repeated pair occurred after a nonempty initial tail, applying the inverse would give an earlier repetition. Thus the initial pair $(0,1)$ recurs at some positive index $T$. Hence the entire pair sequence is periodic and $a_T\equiv0\pmod m$. Also $a_T>0$ for $T\ge1$.

\subsection*{Powers, Orders, and Primitive Roots}
\paragraph{Problem \ref{n:ex25}.}
We have $3\not\equiv1$, $3^2=9\not\equiv1$, and $3^3=27\equiv1$, so the order is $3$. Since $2026\equiv1\pmod3$, the residue is $3$.

\paragraph{Problem \ref{n:ex26}.}
If $n$ is composite it has a prime divisor $q<n$. The factorial includes $q$, so it is $0$ modulo $q$. The assumed congruence then gives $0\equiv-1\pmod q$, impossible. Thus $n$ is prime.

\paragraph{Problem \ref{n:ex27}.}
The prime $q$ is odd. Its order is $p$, since order $1$ would imply $q\mid1$. Thus $p\mid q-1$, and also $2\mid q-1$. As $\gcd(2,p)=1$, $2p\mid q-1$.

\paragraph{Problem \ref{n:ex28}.}
A solution is nonzero. With primitive root $g$, write $x=g^t$. Then $(p-1)\mid kt$. This linear congruence has $\gcd(k,p-1)$ classes modulo $p-1$, each giving a different nonzero residue. Thus the answer is $\gcd(k,p-1)$.

\paragraph{Problem \ref{n:ex29}.}
If $p_1,\ldots,p_k$ listed all such primes, set $A=2p_1\cdots p_k$ (take $A=2$ if the list is empty). Any prime divisor $q$ of $A^2+1$ is odd and does not divide $A$. Modulo $q$, $A^2=-1\ne1$, so $A$ has order $4$. Hence $4\mid q-1$, but no listed $p_i$ divides $A^2+1$. Contradiction.

\paragraph{Problem \ref{n:ex30}.}
If $p\mid a$, both terms are divisible by $p$. Otherwise Fermat gives $a^{p-1}\equiv1$; multiplying by $a$ gives $a^p\equiv a$.

\paragraph{Problem \ref{n:ex31}.}
Let $g$ be primitive. The sum is $\sum_{j=0}^{p-2}(g^r)^j$. If $p-1\mid r$, every term is $1$, giving $p-1\equiv-1$. Otherwise $g^r\not\equiv1$, and multiplication by $g^r-1$ makes the sum $(g^r)^{p-1}-1\equiv0$. The multiplier is invertible, so the sum is $0$.

\paragraph{Problem \ref{n:ex32}.}
The divisor $q$ is odd and $2^{2^n}\equiv-1\not\equiv1\pmod q$. Squaring gives exponent $2^{n+1}$ with value $1$. The order is a divisor of this power of $2$, but cannot divide $2^n$, so it equals $2^{n+1}$. The order divides $q-1$, proving the claim.

\subsection*{Valuations, Factorials, and Lifting Exponents}
\paragraph{Problem \ref{n:ex33}.}
The exponent of $2$ is $50+25+12+6+3+1=97$. The exponent of $5$ is $20+4=24$, smaller than $97$, so there are $24$ trailing zeroes.

\paragraph{Problem \ref{n:ex34}.}
Since $3\nmid2-1$, the direct hypotheses fail. Instead $2^{60}-1=4^{30}-1$ gives $v_3= v_3(3)+v_3(30)=1+1=2$.

\paragraph{Problem \ref{n:ex35}.}
The two-adic formula gives $v_2(2)+v_2(4)+v_2(100)-1=1+2+2-1=4$.

\paragraph{Problem \ref{n:ex36}.}
Legendre gives $v_2((2^m)!)=2^m-1$ and $v_2((2^{m-1})!)=2^{m-1}-1$. Subtraction gives $(2^m-1)-2(2^{m-1}-1)=1$. The formula holds also for $m=1$.

\paragraph{Problem \ref{n:ex37}.}
Lucas gives $\binom21\binom00\binom21=4\pmod7$. All digits have been compared in matching positions.

\paragraph{Problem \ref{n:ex38}.}
The expansion is $2026=1024+512+256+128+64+32+8+2$, with eight one digits. Thus exactly $2^8=256$ coefficients are odd.

\paragraph{Problem \ref{n:ex39}.}
LTE gives the first formula. For a positive exponent $k$, the power of $(1+p)^{p^r}$ is $1$ modulo $p^s$ exactly when $1+r+v_p(k)\ge s$. The least such positive $k$ is $p^{s-r-1}$, proving the exact order.

\paragraph{Problem \ref{n:ex40}.}
Every base-$p$ digit of $p^r-1$ is $p-1$. Lucas makes $\binom{p^r-1}{k-1}$ a product of $\binom{p-1}{d}$, each nonzero modulo $p$ because none of the factorials through $p-1$ contains $p$. Its valuation is therefore zero. Taking valuations in the identity gives $r-v_p(k)$.

\subsection*{Diophantine Equations: Factor, Bound, and Parametrize}
\paragraph{Problem \ref{n:ex41}.}
Let $u=x-y,v=x+y$, so $uv=2025$, $1\le u\le v$, and both are odd. Since $2025=3^4 5^2$, its divisors at most $45$ are $1,3,5,9,15,25,27,45$. Using $x=(u+v)/2,y=(v-u)/2$ gives $(1013,1012),(339,336),(205,200),(117,108),(75,60),(53,28),(51,24),(45,0)$. Each pair checks in the original equation.

\paragraph{Problem \ref{n:ex42}.}
The equation is $(x-1)(y-1)=6$. Both factors are nonnegative, neither zero, so positive. The four ordered factor pairs give $(2,7),(3,4),(4,3),(7,2)$, all valid.

\paragraph{Problem \ref{n:ex43}.}
Factor as $(n^2-2n+2)(n^2+2n+2)$. For $n\ge2$, the first factor is $(n-1)^2+1\ge2$ and the second is larger. For $n=1$, the number is $5$, prime. Thus $n=1$ only.

\paragraph{Problem \ref{n:ex44}.}
Here $m\le8$ and $n<m$. Testing $n=1,\ldots,5$ (since $2n^2<65$) gives $(m,n)=(8,1),(7,4)$ only. They are coprime with opposite parity, yielding $(63,16,65)$ and $(33,56,65)$, up to interchanging the legs.

\paragraph{Problem \ref{n:ex45}.}
Modulo $3$, two squares sum to $0$ only when both are $0$. Thus $3\mid x,y$, and the equation then gives $3\mid z$. Dividing all three by $3$ produces a smaller positive solution. A least positive $z$ cannot exist, so there is no solution.

\paragraph{Problem \ref{n:ex46}.}
The smaller positive divisor of $36$ can be $1,2,3,4,6$. Thus $(x,y)=(7,42),(8,24),(9,18),(10,15),(12,12)$. The product construction proves both validity and completeness.

\paragraph{Problem \ref{n:ex47}.}
Let $t$ be the minimum of $v_2(x),v_2(y),v_2(z)$ and set $x=2^tX,y=2^tY,z=2^tZ$, with at least one odd. Then $X^2+Y^2+Z^2=2^{t+1}XYZ$. If exactly one is odd, the left side is $1$ modulo $4$ and the right side $0$. If exactly two are odd, the left side is $2$ and the even third factor makes the right side $0$. If all three are odd, the left side is $3$ and the right side is $0$ or $2$. All cases are impossible.

\paragraph{Problem \ref{n:ex48}.}
If $d\mid a,b$, the first given divisibility makes $d\mid b^2+1$, hence $d\mid1$. Thus the gcd is $1$. The expression $a^2+b^2+1$ is divisible by $a$ and by $b$, using the corresponding hypothesis, so their coprime product divides it. The next chapter develops the much stronger classification hidden here.

\subsection*{Descent and Vieta Jumping}
\paragraph{Problem \ref{n:ex49}.}
The substitutions give $3=3$, $6=6$, $30=30$, and $195=195$. The next entry is $3\cdot13-5=34$, so the next pair is $(13,34)$. The recurrence preserves the equation by Theorem~\ref{n:markov-chain}.

\paragraph{Problem \ref{n:ex50}.}
If $d$ divides $x,y$, the equation gives $d\mid1$, so $d=1$. This necessary condition is weaker than the full chain classification.

\paragraph{Problem \ref{n:ex51}.}
Modulo $3$, $3\mid x,y$, then the equation gives $3\mid z$. For any nonzero solution choose the least positive $|x|+|y|+|z|$; dividing all entries by $3$ preserves integrality and nonzeroness and reduces this measure. Contradiction. The all-zero triple is a solution.

\paragraph{Problem \ref{n:ex52}.}
Set $k=(x^2+y^2+1)/(xy)$ and choose a positive solution for this fixed $k$ with least sum, ordered $x\le y$. The other root $y'=kx-y=(x^2+1)/y$ is a positive integer. If $x\ge2$, then $y'\le x$ and equality would require $x\mid1$, impossible. Thus $y'<x$, contradicting minimality. Hence $x=1$, and the equation $1+y^2+1=ky$ implies $y\mid2$, so $y=1$ or $2$. Both give $k=3$.

\paragraph{Problem \ref{n:ex53}.}
If $u^4-v^4=w^2$, the identity $(u^4-v^4)^2+(2u^2v^2)^2=(u^4+v^4)^2$ gives a positive integer right triangle with legs $w^2,2u^2v^2$. Its area is $u^2v^2w^2=(uvw)^2$, contradicting Fermat's right-triangle theorem.

\paragraph{Problem \ref{n:ex54}.}
Problem~\ref{n:ex48} gives $\gcd(a,b)=1$ and $ab\mid a^2+b^2+1$. Problem~\ref{n:ex52} makes the quotient $3$. Thus the pairs are exactly the ordered consecutive terms of the chain in Theorem~\ref{n:markov-chain}, allowing interchanging the entries. Conversely the chain equation gives $b^2+1=a(3b-a)$ and $a^2+1=b(3a-b)$, so both required divisibilities hold.

\paragraph{Problem \ref{n:ex55}.}
For $t\ge1$, $(t^2+t^6)/(t^4+1)=t^2$, so every positive square occurs. The numerator and denominator are both positive, so the quotient cannot be $0$. The worked descent proves that no nonsquare quotient occurs.

\paragraph{Problem \ref{n:ex56}.}
A smaller object outside the domain need not contradict minimality among allowed objects. In the quotient-square proof, the cases $b'=0$ and $b'<0$ are treated separately before invoking minimality. Only after they are excluded is $(b',a)$ a positive solution. In a different problem, zero might be a valid terminal seed, but that would require an explicit change in the classification.

\subsection*{Quadratic Residues and Sums of Two Squares}
\paragraph{Problem \ref{n:ex57}.}
The squares are $1,4,9,5,3$ modulo $11$. They are distinct and account for all five nonzero square classes. Thus $5$ is a square, for instance $4^2\equiv5$.

\paragraph{Problem \ref{n:ex58}.}
Reciprocity gives $\left(\frac3p\right)=(-1)^{(p-1)/2}\left(\frac p3\right)$. A prime $p>3$ is $1,5,7,$ or $11$ modulo $12$. For these classes, the two factors give respectively $1,-1,-1,1$. Therefore the square classes of moduli are exactly $1,11$ modulo $12$.

\paragraph{Problem \ref{n:ex59}.}
Lemma~\ref{n:three-mod-four} gives $q\mid x,y$. Hence each square is divisible by $q^2$, so their sum is. The stronger conclusion that the exact exponent is even follows by dividing out common powers repeatedly.

\paragraph{Problem \ref{n:ex60}.}
Because $3\mid x^2+y^2$, both coordinates are divisible by $3$. Dividing by $9$ gives a sum of squares equal to $225$, so both are again divisible by $3$. Thus $x=9u,y=9v$ with $u^2+v^2=25$. For $u\ge v\ge0$, the possibilities are $(5,0),(4,3)$. Hence $(x,y)=(45,0),(36,27)$. Checking proves both work.

\paragraph{Problem \ref{n:ex61}.}
The symbols at $3$ and $5$ are both $-1$, so their product is $1$. But a solution modulo $15$ would be a solution modulo $3$, where squares are $0,1$, never $2$. Thus the product symbol conceals two local failures.

\paragraph{Problem \ref{n:ex62}.}
The remainder is $3-2i$. Since $5+i=(1+i)(3-2i)$, this remainder is a gcd up to a unit. Its norm is $3^2+(-2)^2=13$, yielding $13=9+4$.

\paragraph{Problem \ref{n:ex63}.}
Write $x=a/c,y=b/c$ with integers and $c>0$. Then $a^2+b^2=3c^2$, a nonzero integer solution forbidden by descent. Equivalently the prime $3$ has odd exponent on the right after removing common powers, whereas in a sum of two squares it must have even exponent.

\paragraph{Problem \ref{n:ex64}.}
Suppose $p=a^2+b^2=c^2+d^2$. The Gaussian integer $\pi=a+bi$ has norm $p$, so any factorization of $\pi$ has a unit factor; it is irreducible. In a Euclidean domain an irreducible divides a product only by dividing a factor: its gcd with the first factor is either an associate or a unit, and in the unit case Bezout cancels that factor. Since $\pi\mid(c+di)(c-di)=p$, it divides one factor. Equal norms then make the quotient a unit. Multiplying by $\pm1,\pm i$ and conjugating corresponds exactly to changing signs and interchanging coordinates, proving uniqueness.

\subsection*{Approximation, Continued Fractions, and Pell Equations}
\paragraph{Problem \ref{n:ex65}.}
The successive convergents are $1,3/2,7/5,17/12$, so the answer is $17/12$. Direct expansion gives the same reduced fraction.

\paragraph{Problem \ref{n:ex66}.}
The four pairs are $(3,2),(17,12),(99,70),(577,408)$. The classification theorem says every positive solution is one of these powers in sequence. Their positive $y$ coefficients strictly increase under the recurrence, so no solution is skipped.

\paragraph{Problem \ref{n:ex67}.}
The smallest positive solution is $(1,1)$, and every positive solution has $v=x+y\sqrt2\ge1+\sqrt2$. The quotient $v/(1+\sqrt2)=v(\sqrt2-1)$ has integer coefficients, norm $1$, and is at least $1$. If it is greater than $1$, its positive reciprocal conjugate makes both coefficients positive, so it is a power of $3+2\sqrt2$. If it is $1$, take the zeroth power. Thus $x+y\sqrt2=(1+\sqrt2)(3+2\sqrt2)^k$, $k\ge0$, equivalently $(1+\sqrt2)^{2k+1}$. Conversely every such expression has positive integer coefficients and norm $-1$.

\paragraph{Problem \ref{n:ex68}.}
For $1\le q\le Q$, irrationality makes the distance from $q\alpha$ to the nearest integer positive. The minimum of finitely many such distances is $\delta>0$. An inequality $|q\alpha-p|<1/N$ with $N>1/\delta$ cannot use these denominators. Thus the pigeonhole constructions eventually require $q>Q$, for every $Q$.

\paragraph{Problem \ref{n:ex69}.}
The smallest positive $y=1$ gives $x=2$, so the fundamental unit is $2+\sqrt3$. All positive solutions are its powers. The first three pairs are $(2,1),(7,4),(26,15)$.

\paragraph{Problem \ref{n:ex70}.}
A common positive divisor of $P_n,Q_n$ divides $P_nQ_{n-1}-P_{n-1}Q_n=\pm1$. Hence that divisor is $1$. This also covers $n=0$ with the initial values.

\paragraph{Problem \ref{n:ex71}.}
For $0<p/q\le2$, the difference is $|p^2-2q^2|/[q^2(p/q+\sqrt2)]>1/(4q^2)$, because the numerator is a nonzero integer and the denominator factor is less than $4$. If $p/q\le0$, the difference is at least $\sqrt2>1/100$. If $p/q>2$, it exceeds $2-\sqrt2>1/100$. Since $q\ge1$, neither outside case can satisfy the proposed bound.

\paragraph{Problem \ref{n:ex72}.}
It is equivalent to $(2t+1)^2-2h^2=-1$. The negative-Pell classification gives $(2t+1)+h\sqrt2=(1+\sqrt2)^{2k+1}$. The value $k=0$ gives $t=0$, excluded; for every $k\ge1$ the first coefficient is odd and exceeds $1$, so $t$ is a positive integer. The first three triangles are $(3,4,5),(20,21,29),(119,120,169)$. All satisfy the original equation, and the Pell classification proves completeness.

\subsection*{Prime Powers, Integer Polynomials, and Further Constructions}
\paragraph{Problem \ref{n:ex73}.}
Modulo $7$ the roots are $3,4=-3$. Their derivatives $2x$ are nonzero modulo $7$, so each lifts uniquely. The positive branch computed in the example is $108$; the other is its negative, $235$, modulo $343$. Every root reduces to one of these two initial classes, proving completeness.

\paragraph{Problem \ref{n:ex74}.}
Reducing modulo $p$ gives $p\mid x^2$, so $p\mid x$. Then $p^2\mid x^2$, contradicting $x^2\equiv p\pmod{p^2}$. Modulo $p$, the residue $0$ is indeed a root.

\paragraph{Problem \ref{n:ex75}.}
A solution is odd. Of $x-1,x+1$, one has exactly one factor $2$ and the other must have at least $k-1$, so $x\equiv1$ or $-1\pmod{2^{k-1}}$. This gives the four classes $1,-1,1+2^{k-1},-1+2^{k-1}$ modulo $2^k$. They are distinct for $k\ge3$, and direct expansion shows all square to $1$.

\paragraph{Problem \ref{n:ex76}.}
Modulo $7$, $3^3=-1$ and $3,3^2$ are not $1$, so the order is $6$. Also $3^6-1=728=7\cdot104$ has valuation $1$. For exponents $6t$, LTE gives $v_7(3^{6t}-1)=1+v_7(t)$. Hence the exact order modulo $7^k$ is $6\cdot7^{k-1}$.

\paragraph{Problem \ref{n:ex77}.}
Choose $a$ with $c=|P(a)|>1$. For integers $t$, $P(a+ct)$ is divisible by $c$. Since the polynomial is nonconstant, its absolute value exceeds $c$ for sufficiently large $|t|$. Such a value is divisible by $c$ and has larger absolute value, so its absolute value is composite.

\paragraph{Problem \ref{n:ex78}.}
If $n=ab$ is composite, $R_n=R_a(1+10^a+\cdots+10^{a(b-1)})$ is a product of two integers larger than $1$. For $n=1$, $R_1=1$ is not prime. Thus a prime repunit requires prime $n$. The converse fails at $n=3$, since $111=3\cdot37$.

\paragraph{Problem \ref{n:ex79}.}
If the listed primes have product $M$, then $P(Mt)\equiv P(0)=1\pmod{p}$ for each listed $p$. For sufficiently large $|t|$, $|P(Mt)|>1$, so it has a prime divisor outside the list. The value is nonzero by its absolute value, giving the contradiction.

\paragraph{Problem \ref{n:ex80}.}
The value $a^2-a+1$ is odd, so $q$ is odd, and $q\nmid a$. Modulo $q$, $a^3=-1$, hence the order divides $6$ but not $3$. It is therefore $2$ or $6$. Order $2$ would give $a=-1$ modulo $q$, making $a^2-a+1=3$ modulo $q$, contrary to $q\ne3$. Thus the order is $6$, which divides $q-1$.




% ---- solutions/capstones ----
\section{Mixed Problems}
\paragraph{Problem \ref{m:ex01}.}
If $P$ is constant $c$, the equation is $c=c^2$, giving $0,1$.
For degree $d\ge1$, composition has degree $d^2$ while the square
has degree $2d$, so $d=2$. Write $P(x)=ax^2+bx+c$, $a\ne0$.
The leading coefficients give $a^3=a^2$, hence $a=1$.
Now $P(P(x))=P(x)^2+bP(x)+c$, so the equation makes
$bP(x)+c$ the zero polynomial. Nonconstancy gives $b=0,c=0$.
Thus $P(x)=x^2$, and direct substitution verifies it and both
constants.

\paragraph{Problem \ref{m:ex02}.}
Set $x=1,y=0$ to obtain $f(0)=0$.
Then $y=0$ gives $f(x^2)=xf(x)$.
Substituting this into the original relation shows
$f(t+y)=f(t)+f(y)$ for every $t\ge0$ and every real $y$.
Taking $y=-t$ yields $f(-t)=-f(t)$.
For a negative increment $u=-t$, use
$f((y-t)+t)=f(y-t)+f(t)$ to get
$f(y+u)=f(y)+f(u)$.
Thus $f$ is additive on all $\RR$.

By additivity,
$f((x+1)^2)=f(x^2)+2f(x)+f(1)
=xf(x)+2f(x)+f(1)$.
The square rule also gives
$f((x+1)^2)=(x+1)(f(x)+f(1))$.
Subtracting yields $f(x)=x f(1)$.
Therefore $f(x)=cx$ for arbitrary real $c$, and substitution verifies
every such function.
The choice $x=1,y=0$ establishes $f(0)=0$;
$x=y=0$ alone would have been a tautology.

\paragraph{Problem \ref{m:ex03}.}
The integer $n$ is odd, since $2^n+1$ is odd. Let $p$ be its
least prime divisor. Then $p$ is odd and $2^n\equiv-1\pmod p$.
The order $d$ of $2$ modulo $p$ divides $2n$ but not $n$,
so $d=2e$ with $e\mid n$.
Also $d\mid p-1$, hence $e\le(p-1)/2<p$.
If $e>1$, it has a prime divisor at least $p$, since it divides $n$,
contradicting $e<p$. Therefore $e=1$, and $2$ has order $2$
modulo $p$. The unique nontrivial square root of $1$ is $-1$,
so $2\equiv-1\pmod p$, giving $p=3$.

\paragraph{Problem \ref{m:ex04}.}
If a row has $b$ black and $6-b$ white entries, it contributes
$\binom b2+\binom{6-b}2=(b-3)^2+6\ge6$
same-color pairs of columns. Across six rows there are at least
$36$ such pairs with their row occurrences counted.
There are only $2\binom62=30$ possible records consisting of
two columns and one color. Some record occurs in two different
rows. Its four entries form the required monochromatic rectangle.

\paragraph{Problem \ref{m:ex05}.}
If the quadrilateral is cyclic, the intersecting-chords theorem gives
the product equality. Conversely draw the circle through $A,B,C$.
Since $P$ lies inside segment $AC$, it is inside this circle.
Ray $PD$ therefore meets it at a unique point $D'$.
Power gives $PB\cdot PD'=PA\cdot PC=PB\cdot PD$.
Since $PB>0$, $PD'=PD$, and the points lie on the same ray.
Hence $D'=D$, so all four vertices are cyclic.
Convexity supplies the inside and ray conditions that make the
unsigned products valid.

\paragraph{Problem \ref{m:ex06}.}
For rational $t$, the formula
$((1-t^2)/(1+t^2),\,2t/(1+t^2))$
is rational and has squared-coordinate sum $1$.
Conversely a rational point with $x\ne-1$ has rational slope
$t=y/(x+1)$ from $(-1,0)$; solving the line-circle equation
gives exactly that formula. The remaining point is $(-1,0)$.
No finite slope parameter produces it.

\paragraph{Problem \ref{m:ex07}.}
Since $0<a,b,c<1$, all denominators are positive.
Cauchy's inequality gives
\[
\sum_{\mathrm{cyc}}\frac{a^2}{1-a}
\ge\frac{(a+b+c)^2}{(1-a)+(1-b)+(1-c)}=\frac12.
\]
Equality requires $a/(1-a)=b/(1-b)=c/(1-c)$, hence
$a=b=c=1/3$, which does attain the bound.

\paragraph{Problem \ref{m:ex08}.}
A nonzero divisor of $n$ has absolute value at most $n$.
Thus $\deg P\ge2$ is impossible, because the leading term would make
$|P(n)|>n$ for all sufficiently large positive $n$.
If $P=c$ is constant, using $n=1$ gives $c=\pm1$.
If $P(X)=aX+b$, $a\ne0$, the divisibility implies
$an+b\mid(an+b)-an=b$.
For $b\ne0$, this contradicts the unbounded absolute value of
$an+b$. Therefore $b=0$, and $an\mid n$ forces $a=\pm1$.
All four polynomials $1,-1,X,-X$ satisfy the conditions.

\paragraph{Problem \ref{m:ex09}.}
If a vertex has at least four red neighbors, then either a red
edge within those neighbors makes a red triangle with the vertex,
or all their edges are blue, giving a blue $K_4$.
Otherwise every red degree is at most $3$. Nine vertices cannot
all have odd degree $3$, since the degree sum is twice the number
of red edges. Some vertex $v$ therefore has red degree at most $2$
and has at least six blue neighbors.
Among six vertices, a red or blue triangle always occurs:
from one vertex, three incident edges share a color; either an
edge of that color joins their other endpoints, or those endpoints
form a triangle of the other color.
A red triangle finishes directly; a blue triangle with $v$ is a blue $K_4$.

For sharpness, label eight vertices by residues modulo $8$.
Color an edge red when its endpoints differ by $1,-1$, or $4$,
and color the others blue. The red graph is an eight-cycle with
opposite chords. It has no triangle: two neighbors of any vertex
have none of the three permitted red differences.
A red-independent set of four vertices on an eight-cycle would
have to alternate, since all four circular gaps must be at least two
and sum to eight. Each alternating set contains opposite vertices,
joined by a red chord. Thus no four vertices are red-independent,
equivalently no blue $K_4$ exists.

\paragraph{Problem \ref{m:ex10}.}
Let $k=(a^2+b^2)/(ab-1)$. Since $a^2+b^2\ge2ab>2(ab-1)$,
$k$ is an integer at least $3$.
For this fixed $k$, choose a positive solution with $ab>1$ minimizing
$a+b$, ordered $a\le b$.
The other root is $b'=ka-b=(a^2+k)/b$, a positive integer.
If $a\ge2$, evaluate
$F(T)=T^2-kaT+a^2+k$ at $T=a$:
\[
F(a)=k-(k-2)a^2\le k-4(k-2)=8-3k<0.
\]
The upward-opening quadratic has roots $b,b'$ and is negative
at $a$, so $b'<a<b$. The new pair $(b',a)$ has product greater
than $1$ because $a\ge2$, satisfies the same equation, and has
smaller sum, a contradiction.
Thus $a=1$ and $b\ge2$. Divisibility becomes
$b-1\mid b^2+1$, or $b-1\mid2$, giving $b=2,3$.
Both yield $k=5$, which is therefore the only possible quotient.

\paragraph{Problem \ref{m:ex11}.}
Among all positive distances from a point of $S$ to a line determined
by two other points of $S$, choose the least, $h$, attained by
$P$ and line $\ell$. Such distances form a nonempty finite set
because $S$ is finite and noncollinear.
If $\ell$ contains exactly two points, we are done.
Otherwise let $H$ be the perpendicular foot from $P$.
Among three distinct points on $\ell$, two, called $A,B$, lie on
the same closed half-line from $H$, ordered with $A$ closer to $H$.
Then $AB\le HB<PB$.
The distance from $A$ to line $PB$ is, by comparing two area
formulas for triangle $PAB$, $AB\cdot h/PB<h$.
It is positive because $P,A,B$ are noncollinear.
This contradicts minimality. Hence $\ell$ contains exactly two points.

\begin{center}
\begin{tikzpicture}[scale=.9]
\coordinate (P) at (0,2);\coordinate (H) at (0,0);
\coordinate (A) at (1,0);\coordinate (B) at (4,0);
\coordinate (Q) at (1.6,1.2);
\draw (-1.4,0)--(4.5,0) node[right]{$\ell$};
\draw (P)--(A)--(B)--cycle;
\draw[aux] (P)--(H);
\draw[main] (A)--(Q);
\foreach \X/\pos in {P/above,H/below,A/below,B/below,Q/above right}
{\fill (\X) circle (1.3pt) node[\pos]{$\X$};}
\node[left] at (0,1){$h$};
\end{tikzpicture}
\end{center}
The shorter altitude is the segment $AQ$; the choice $AB<PB$
is what makes it shorter.

\paragraph{Problem \ref{m:ex12}.}
There are finitely many pairings, at most $n!$, each with a real
total length, so a least one exists. Crossing segments can be
uncrossed exactly as in the worked example: strict triangle
inequalities reduce the total while preserving the red--blue
bijection. Thus a minimum pairing has no crossings.
Existence of a minimum here follows from finiteness, not from
an unsupported optimization claim over an infinite family.

\paragraph{Problem \ref{m:ex13}.}
Choose a positive integer $D$ clearing all three rational side
denominators. Scaling gives an integer right triangle with area
$D^2$. This is a positive integer square, forbidden by
Theorem~\ref{n:fermat-area}. Thus no original triangle exists.

\paragraph{Problem \ref{m:ex14}.}
Let the internal angle bisector $AI$ meet $(ABC)$ again at $L$,
on the arc $BC$ not containing $A$. This is the arc midpoint.
Since $I$ is interior, it lies between $A$ and $L$.
The standard incenter angle is $\angle BIA=90^\circ+C/2$,
so $\angle BIL=90^\circ-C/2$.
Also $\angle IBL=B/2+A/2=90^\circ-C/2$.
Hence $IL=LB$.
The sine law on the circumcircle gives $LB=2R\sin(A/2)$.
Dropping a perpendicular from $I$ to a side through $A$ gives
$AI=r/\sin(A/2)$, so $AI\cdot IL=2Rr$.
Power of $I$ gives
$OI^2-R^2=-AI\cdot IL=-2Rr$,
which is the identity.
Since $OI^2\ge0$ and $R>0$, $R\ge2r$.

Equality means $O=I$. In that case the circumcenter has equal
perpendicular distances $r$ from the three sides, so the corresponding
chords in the circumcircle have equal lengths. Thus $a=b=c$.
Conversely an equilateral triangle has $R=2r$.

\begin{center}
\begin{tikzpicture}[scale=.65]
\pgfmathsetmacro{\radc}{5*sqrt(5)/4}
\pgfmathsetmacro{\ix}{(-10+6*sqrt(5))/(10+2*sqrt(5))}
\pgfmathsetmacro{\iy}{20/(10+2*sqrt(5))}
\coordinate (A) at (0,4);\coordinate (B) at (-2,0);
\coordinate (C) at (3,0);\coordinate (O) at (.5,1.25);
\coordinate (I) at (\ix,\iy);
\coordinate (L) at (.5,{1.25-\radc});
\coordinate (D) at (\ix,0);
\draw[main] (O) circle (\radc);
\draw (A)--(B)--(C)--cycle;
\draw (A)--(L) (B)--(L);
\draw[aux] (I)--(D) (O)--(I);
\foreach \X/\pos in {A/above,B/below left,C/below right,L/below,D/below}
{\fill (\X) circle (1.5pt) node[\pos]{$\X$};}
\fill (I) circle (1.5pt);
\draw (I)--++(-.8,.5) node[left]{$I$};
\fill (O) circle (1.5pt);
\draw (O)--++(.8,.5) node[right]{$O$};
\node[left] at (\ix,{.5*\iy}){$r$};
\end{tikzpicture}
\end{center}
The diagram depicts a scalene triangle; the proof uses no acute-only
property or assumed order of the centers.

\paragraph{Problem \ref{m:ex15}.}
Take $n-1$ zeroes and $n-1$ ones.
Every $n$-element submultiset has between $1$ and $n-1$ ones,
so its sum is not divisible by $n$. Thus one fewer entry cannot
always work, while Theorem~\ref{m:egz} proves $2n-1$ suffices.
The partial-sums argument gives a nonempty consecutive block
with a divisible sum, but its length is the difference of two
colliding indices and is not controlled to equal $n$.

\paragraph{Problem \ref{m:ex16}.}
Four chosen vertices, taken in cyclic order, form a convex
quadrilateral whose two diagonals meet once inside the polygon.
Conversely any interior crossing involves two diagonals with four
distinct endpoints; these endpoints determine exactly that
quadrilateral and its crossing pair.
The no-three-concurrent condition ensures distinct four-vertex
choices do not produce the same intersection point through a
third diagonal. Hence the correspondence between four-element
vertex subsets and intersection points is bijective, giving $\binom n4$.


\backmatter

% ---- front/references ----
\chapter{Sources and Further Problems}
Classical results and problems in this book belong to the mathematical
tradition. The explanations, organization, diagrams, and many exercise
variants were prepared for this course. A familiar problem presented
without a contest label should not be interpreted as a claim of
new discovery.

The following public sources were consulted for scope or specific
configurations. Links were checked on 3 November 2026.
\begingroup\raggedright
\begin{enumerate}
\item Evan Chen, \emph{Unofficial syllabus for math olympiads},
31 January 2023.
\url{https://web.evanchen.cc/handouts/Syllabus/Syllabus.pdf}.
This provides an informal scope comparison, not an official or
exhaustive curriculum.
\item HBCSE, Mathematical Olympiad portion of \emph{Syllabus}.
\url{https://olympiads.hbcse.tifr.res.in/syllabus/}.
This offers a national-program comparison for the elementary foundations.
\item International Mathematical Olympiad, official problem archive.
\url{https://www.imo-official.org/problems/}.
The quotient-square worked example in the descent chapter is IMO
1988, Problem 6. Use the year selector in the archive for the
original contest paper. The book's argument is written out
in its own exposition.
\item Evan Chen, \emph{A Guessing Game: Mixtilinear Incircles}.
\url{https://web.evanchen.cc/handouts/Mixt-GeoGuessr/Mixt-GeoGuessr.pdf}.
Used to cross-check the mixtilinear contact-chord configuration.
\item Evan Chen, \emph{Euclidean Geometry in Mathematical Olympiads},
author's overview.
\url{https://web.evanchen.cc/geombook.html}.
This was a reference for the user's desired style of geometry
instruction. The commercial text has not been copied into this book.
\end{enumerate}
\endgroup

\section*{Continuing the problem bank}
For further independent problems, use the official archive and record
the year and problem number when adding one to the source. Choose a
mixture of subjects and difficulties, and keep hints separate from
full solutions. A new problem earns its place by teaching a decision
or testing transfer, not merely by carrying a difficult contest label.


\end{document}
